% De Scratch a Python: cuaderno de trabajo con Google Colab para un niño de 9 años (edición en español).
% Compilar con:  tectonic de-scratch-a-python.tex   (desde la carpeta es/)
% El aspecto compartido está en preamble.tex; cada parte está en chapters/.
% Es la traducción de ../scratch-to-python-workbook.tex (edición ampliada en inglés).
\documentclass[12pt,a4paper]{article}
% Preámbulo compartido de la edición en español (libro y arnés de prueba de cada capítulo).
\usepackage[a4paper,margin=2cm,top=2cm,bottom=2.2cm]{geometry}
\usepackage{fontspec}
\usepackage{polyglossia}
\setdefaultlanguage{spanish}
\usepackage{FiraSans}
\setmonofont{FiraMono}[Extension=.otf, UprightFont=*-Regular, BoldFont=*-Bold, Ligatures=NoCommon]
\renewcommand{\familydefault}{\sfdefault}
\usepackage{xcolor}
\usepackage{amssymb}
\usepackage{array,longtable,booktabs}
\usepackage{enumitem}
\usepackage{needspace}
\newcolumntype{P}[1]{>{\raggedright\arraybackslash}p{#1}}
\usepackage{listings}
\usepackage[most]{tcolorbox}
\usepackage{tikz}
\usetikzlibrary{arrows.meta,calc,shapes.geometric}
\usepackage{scratch3}
\usepackage{fancyhdr}
\usepackage[hidelinks]{hyperref}

\linespread{1.22}
\setlength{\emergencystretch}{3em}
\setlength{\parindent}{0pt}
\setlength{\parskip}{7pt}
\setlist{itemsep=3pt,topsep=4pt}
\setscratch{scale=1, else word=si no}

% ---------- Scratch colours ----------
\definecolor{motion}{HTML}{4C97FF}
\definecolor{looks}{HTML}{9966FF}
\definecolor{sound}{HTML}{CF63CF}
\definecolor{events}{HTML}{FFBF00}
\definecolor{control}{HTML}{FFAB19}
\definecolor{sensing}{HTML}{5CB1D6}
\definecolor{operators}{HTML}{59C059}
\definecolor{variables}{HTML}{FF8C1A}
\definecolor{lists}{HTML}{FF661A}
\definecolor{myblocks}{HTML}{FF6680}
\definecolor{pen}{HTML}{0FBD8C}
\definecolor{ink}{HTML}{1B2550}
\definecolor{muted}{HTML}{56618A}
\definecolor{codebg}{HTML}{F3F6FD}
\definecolor{codeframe}{HTML}{9AA8D4}
\definecolor{kw}{HTML}{1D4FC4}
\definecolor{str}{HTML}{1B7F3B}
\definecolor{com}{HTML}{7A84A8}
\definecolor{warnbg}{HTML}{FFF3DB}
\definecolor{warnframe}{HTML}{E89A12}
\definecolor{chalbg}{HTML}{E6F7EC}
\definecolor{chalframe}{HTML}{2E9C57}
\definecolor{thinkbg}{HTML}{E8F0FF}
\definecolor{thinkframe}{HTML}{4C7FE0}
\definecolor{errbg}{HTML}{FFECEC}
\definecolor{errframe}{HTML}{D23B3B}\colorlet{errhead}{errframe!35}
\color{ink}

% ---------- Python code ----------
\lstdefinestyle{py}{
  language=Python, basicstyle=\ttfamily\normalsize, keepspaces=true, columns=fixed,
  upquote=true, showstringspaces=false, keywordstyle=\color{kw}\bfseries,
  stringstyle=\color{str}, commentstyle=\color{com}\itshape,
  morekeywords={True,False,None}, deletekeywords={print,input,range,int,str,len,abs,round},
  numbers=left, numberstyle=\scriptsize\color{muted}, numbersep=9pt, xleftmargin=0pt,
  tabsize=4, breaklines=true
}
\lstdefinestyle{inl}{style=py, numbers=none, basicstyle=\ttfamily}
\newcommand{\py}{\lstinline[style=inl]}
\DeclareTextCommand{\textvisiblespace}{TU}{\makebox[0.6em][c]{\color{muted}\rule{0.5pt}{0.35ex}\rule{0.45em}{0.5pt}\rule{0.5pt}{0.35ex}}}
\makeatletter\def\lst@visiblespace{\textvisiblespace}\makeatother

\newtcblisting{typethis}[1][Escribe esto en una celda de código]{
  listing only, enhanced, colback=codebg, colframe=codeframe, boxrule=0.8pt, arc=6pt,
  left=22pt, right=8pt, top=4pt, bottom=4pt,
  title={\textbf{#1}}, coltitle=white, colbacktitle=ink, fonttitle=\small,
  listing options={style=py}}
\newtcblisting{code}{
  listing only, enhanced, colback=codebg, colframe=codeframe, boxrule=0.6pt, arc=5pt,
  left=22pt, right=6pt, top=2pt, bottom=2pt, listing options={style=py}}
\newtcblisting{codenonum}{
  listing only, enhanced, colback=codebg, colframe=codeframe, boxrule=0.6pt, arc=5pt,
  left=8pt, right=6pt, top=2pt, bottom=2pt, listing options={style=py,numbers=none}}
\newtcblisting{spaces}{
  listing only, enhanced, colback=codebg, colframe=codeframe, boxrule=0.6pt, arc=5pt,
  left=22pt, right=6pt, top=2pt, bottom=2pt, listing options={style=py,showspaces=true}}
\newtcolorbox{outputbox}[1][Deberías ver]{
  enhanced, colback=white, colframe=ink!40, boxrule=0.6pt, arc=5pt, left=8pt,
  title={\textbf{#1}}, coltitle=ink, colbacktitle=ink!10, fonttitle=\small,
  fontupper=\ttfamily}

% ---------- Helper boxes ----------
\newtcolorbox{watchout}[1][¡Cuidado!]{enhanced, breakable, colback=warnbg, colframe=warnframe,
  boxrule=0.8pt, arc=6pt, title={\textbf{#1}}, coltitle=ink, colbacktitle=warnframe!35, left=8pt, right=8pt}
\newtcolorbox{challenge}[1][Hazlo tuyo]{enhanced, breakable, colback=chalbg, colframe=chalframe,
  boxrule=0.8pt, arc=6pt, title={\textbf{#1}}, coltitle=white, colbacktitle=chalframe, left=8pt, right=8pt}
\newtcolorbox{think}[1][¿Te diste cuenta?]{enhanced, breakable, colback=thinkbg, colframe=thinkframe,
  boxrule=0.8pt, arc=6pt, title={\textbf{#1}}, coltitle=white, colbacktitle=thinkframe, left=8pt, right=8pt}
\newtcolorbox{errorbox}{enhanced, colback=errbg, colframe=errframe, boxrule=0.8pt, arc=5pt,
  fontupper=\ttfamily\small, left=8pt}
\newtcbox{\chip}[1][motion]{on line, colback=#1!45!white, colframe=#1!80!black, boxrule=0.5pt,
  arc=7pt, left=4pt, right=4pt, top=0.5pt, bottom=0.5pt, boxsep=0pt, fontupper=\sffamily\small\bfseries}
\newcommand{\tick}{$\square$\ }
\newcommand{\step}[1]{\par\needspace{9\baselineskip}\medskip{\large\bfseries\color{ink}#1}\par}

% Botón "play" de Colab dibujado en línea
\newcommand{\playbtn}{\tikz[baseline=-0.6ex]{\fill[ink] (0,0) circle (1.7ex); \fill[white] (-0.5ex,-0.8ex) -- (-0.5ex,0.8ex) -- (0.8ex,0) -- cycle;}}
\newcommand{\stopbtn}{\tikz[baseline=-0.6ex]{\draw[ink,line width=0.9pt] (0,0) circle (1.7ex); \fill[ink] (-0.6ex,-0.6ex) rectangle (0.6ex,0.6ex);}}

% Section banners
\newcommand{\banner}[3]{% colour, small label, title
  \clearpage
  \begin{tcolorbox}[enhanced, colback=#1, colframe=#1!80!black, arc=10pt, boxrule=0pt,
    left=14pt, top=10pt, bottom=10pt]
    {\small\bfseries\color{ink}#2}\par\vspace{2pt}
    {\fontsize{26}{30}\selectfont\bfseries\color{ink}#3}
  \end{tcolorbox}
  \addcontentsline{toc}{section}{#3}
  \vspace{4pt}}
\newcommand{\missionhead}[3]{% number, title, you will learn
  \clearpage
  \begin{tcolorbox}[enhanced, colback=ink, colframe=ink, arc=10pt, left=14pt, top=10pt, bottom=10pt]
    {\small\bfseries\color{events}MISIÓN #1}\par\vspace{2pt}
    {\fontsize{26}{30}\selectfont\bfseries\color{white}#2}\par\vspace{4pt}
    {\color{white}\textbf{Aprenderás:} #3}
  \end{tcolorbox}
  \addcontentsline{toc}{section}{Misión #1: #2}
  \vspace{2pt}}

\pagestyle{fancy}
\fancyhf{}
\renewcommand{\headrulewidth}{0pt}
\fancyfoot[C]{\small\color{muted}\thepage}
\fancyfoot[R]{\small\color{muted}De Scratch a Python}


% ---------- Auto-numbered missions (use in chapters) ----------
\newcounter{mission}
\newcommand{\mission}[2]{% title, you will learn
  \stepcounter{mission}\missionhead{\arabic{mission}}{#1}{#2}}
% A part opener: \partopener{colour}{PART n}{Title}{one-line promise}
\newcommand{\partopener}[4]{\banner{#1}{#2}{#3}{\large #4\par}\vspace{6pt}}
% Block dictionary helpers
\newcommand{\dictrow}{\\[0pt]\midrule}
\newcommand{\cat}[2]{\multicolumn{3}{l}{\rule{0pt}{18pt}\tcbox[on line, colback=#1, colframe=#1!80!black, arc=6pt, boxrule=0pt, left=6pt, right=6pt, top=1pt, bottom=1pt]{\bfseries\color{ink}#2}} \\*[4pt]}
% 2-line puzzle box at the end of each part
\newtcolorbox{puzzle}[1][Acertijo de dos líneas]{enhanced, breakable, colback=looks!12, colframe=looks!80!black,
  boxrule=0.8pt, arc=6pt, title={\textbf{#1}}, coltitle=white, colbacktitle=looks!80!black, left=8pt, right=8pt}


\begin{document}

% =====================================================================
% 00-front.tex : portada, ¡Hola, programador!, Conoce Colab, palabras de
% Scratch, hábitos, en qué se diferencia Colab.  (Sin archivo de respuestas.)
% =====================================================================

% ---------- PORTADA ----------
\thispagestyle{empty}
\begin{tikzpicture}[remember picture, overlay]
  \fill[codebg] (current page.north west) rectangle (current page.south east);
  % una estrella dibujada por la tortuga
  \begin{scope}[shift={($(current page.center)+(0,-1.2)$)}, line width=3pt, line join=round]
    \draw[looks] (-3.2,0.6) -- ++(0:6.4) -- ++(-144:6.4) -- ++(72:6.4) -- ++(-72:6.4) -- cycle;
  \end{scope}
  \foreach \c/\x in {motion/0, looks/1, events/2, control/3, operators/4, variables/5, myblocks/6, pen/7}
    \fill[\c] ($(current page.south west)+(\x*2.625,0)$) rectangle ++(2.625,1.1);
\end{tikzpicture}

\vspace*{1.2cm}
{\color{muted}\Large\bfseries Un libro de trabajo para jóvenes programadores}\par\vspace{6pt}
{\fontsize{46}{52}\selectfont\bfseries De Scratch\\[2pt] a Python}\par\vspace{10pt}
{\Large Con Google Colab}\par
\vspace{11.5cm}
\begin{tcolorbox}[colback=white, colframe=ink, arc=8pt, width=11cm, boxrule=1pt]
\Large Este libro pertenece a:\par\vspace{9mm}\rule{9.5cm}{0.6pt}\par\vspace{2mm}
\end{tcolorbox}

% =====================================================================
\banner{events}{EMPIEZA AQUÍ}{¡Hola, programador!}

Tú ya sabes programar. Cada vez que encajas bloques de Scratch, estás escribiendo un programa. Le dices a la computadora qué hacer, paso a paso.

En este libro vas a aprender \textbf{Python}. Python es un lenguaje de programación de verdad. Los programadores adultos lo usan para crear páginas web, juegos, robots ¡y hasta misiones espaciales! La gran noticia es que Python usa las \textbf{mismas ideas} que Scratch:

\begin{itemize}
  \item un bloque \chip[motion]{mover 100 pasos} se convierte en \py|forward(100)|
  \item un bloque \chip[control]{repetir 4} se convierte en \py|for i in range(4):|
  \item un bloque \chip[looks]{decir ¡Hola!} se convierte en \py|print("¡Hola!")|
\end{itemize}

La diferencia es que en Python \textbf{escribes} las palabras en lugar de arrastrar bloques. Esa es la única gran novedad. Todo lo demás, como los bucles, las variables y crear tus propios bloques, ya lo entiendes.

\begin{think}[Python habla inglés]
Las palabras propias de Python (\texttt{for}, \texttt{print}, \texttt{forward}, \texttt{range}...) siempre van en inglés, porque así las entiende Python. Por ejemplo, \texttt{forward} significa «adelante» y \texttt{print} significa «imprimir» o «mostrar». Pero los nombres que tú inventas y los textos entre comillas pueden ir en español: \py|print("¡Hola!")|.
\end{think}

\subsection*{Cómo usar este libro}

Recorre el libro en orden. Cada \textbf{misión} tiene cinco pasos:

\begin{enumerate}[label=\textbf{\arabic*.}]
  \item \textbf{Constrúyelo en Scratch.} Primero haz el programa en Scratch, para saber qué debe hacer.
  \item \textbf{Escríbelo en Colab.} Copia el código de Python en Colab, letra por letra.
  \item \textbf{Léelo línea por línea.} Descubre a qué bloque de Scratch corresponde cada línea.
  \item \textbf{Ejecútalo.} Presiona el botón de reproducir y comprueba que funcionó.
  \item \textbf{Hazlo tuyo.} Cambia el programa y mira qué pasa.
\end{enumerate}

\needspace{12\baselineskip}
\subsection*{Busca estos recuadros}

\noindent
\begin{minipage}[t]{0.485\linewidth}\vspace{0pt}
\begin{typethis}[Escribe esto]
print("¡Hola!")
\end{typethis}
\begin{think}
Algo ingenioso que descubrir en el código.
\end{think}
\begin{puzzle}
Un acertijo pequeñito al final de cada Parte. La respuesta son dos líneas de Python.
\end{puzzle}
\end{minipage}\hfill
\begin{minipage}[t]{0.485\linewidth}\vspace{0pt}
\begin{watchout}
Un error que comete mucha gente, para que tú lo puedas evitar.
\end{watchout}
\begin{challenge}
Ideas para cambiar el programa. ¡Marca cada una que termines!
\end{challenge}
\end{minipage}

Las respuestas de los acertijos y de las páginas «Tradúcelo tú» están al final del libro. ¡Primero inténtalo y después revisa!

\needspace{17\baselineskip}
\subsection*{Mapa del libro}

El libro tiene siete \textbf{Partes}. Cada Parte toma una familia de bloques de Scratch y te muestra su Python. Después de las Partes vienen cuatro secciones especiales al final.

\begin{center}
\renewcommand{\arraystretch}{1.3}
\begin{tabular}{l P{5.4cm} P{8.8cm}}
\toprule
 & \textbf{Parte} & \textbf{Lo que harás} \\
\midrule
1 & \chip[motion]{Movimiento y Lápiz} & Mover la tortuga y dibujar figuras. \\
2 & \chip[control]{Bucles} & Repetir cosas sin volver a escribirlas. \\
3 & \chip[variables]{Variables y Listas} & Guardar puntos, nombres y listas de cosas. \\
4 & \chip[operators]{Operadores} & Hacer cuentas, elegir números al azar, unir palabras. \\
5 & \chip[sensing]{Sensores y decisiones} & Hacer preguntas y lograr que el programa elija. \\
6 & \chip[myblocks]{Mis Bloques (funciones)} & Crear tus propios superbloques. \\
7 & \chip[looks]{Apariencia, Sonido y Eventos} & Figuras, colores, sonidos y reaccionar. \\
\midrule
 & \textbf{Grandes proyectos} & Programas más grandes que usan todo. \\
 & \textbf{Zoológico de Bichos} & Los mensajes de error de Python y cómo atrapar cada bicho. \\
 & \textbf{Diccionario de bloques} & Busca cualquier bloque de Scratch y encuentra su Python. \\
 & \textbf{Para los adultos} & Notas para la persona que te ayuda. \\
\bottomrule
\end{tabular}
\end{center}

% =====================================================================
\banner{sensing}{PREPARÁNDONOS}{Conoce Google Colab}

\textbf{Google Colab} es una página web donde puedes escribir Python y ejecutarlo enseguida. No tienes que instalar nada. Colab guarda tu trabajo en Google Drive, así que estará ahí la próxima vez.

\step{Paso 1: Abre Colab}
Pídele a un adulto que abra esta página web y que inicie sesión con una cuenta de Google:
\begin{center}\Large\texttt{colab.research.google.com}\end{center}

\step{Paso 2: Crea un cuaderno nuevo}
Haz clic en \textbf{Nuevo cuaderno} (\textit{New notebook}). Un cuaderno es como un proyecto de Scratch: es donde vive tu código.

Arriba dice \texttt{Untitled0.ipynb}. Haz clic en ese nombre y cámbialo por \textbf{Mis aventuras en Python}.

\step{Paso 3: Encuentra la celda de código}
Verás un recuadro gris con un botón redondo a la izquierda. Este recuadro se llama \textbf{celda de código}. Ahí escribes Python.

\begin{center}
\begin{tikzpicture}
  \draw[codeframe, fill=codebg, rounded corners=6pt] (0,0) rectangle (12,1.4);
  \fill[ink] (0.75,0.7) circle (0.38);
  \fill[white] (0.62,0.5) -- (0.62,0.9) -- (0.98,0.7) -- cycle;
  \node[anchor=west, font=\ttfamily\large] at (1.5,0.7) {print("¡Hola!")};
  \draw[-{Stealth[length=3mm]}, line width=1.2pt, operators!70!black] (0.75,-1.1) -- (0.75,0.2);
  \node[anchor=north, font=\bfseries, text=operators!60!black, align=center] at (0.75,-1.1) {El botón de\\ reproducir};
  \node[anchor=north west, align=left, font=\small] at (3.3,-0.6) {¡Esta es tu bandera verde!\\ Haz clic en él para ejecutar el código de la celda.};
\end{tikzpicture}
\end{center}

\step{Paso 4: Saluda}
Haz clic dentro de la celda y escribe esto exactamente:

\begin{typethis}
print("¡Hola!")
\end{typethis}

Ahora haz clic en el botón de reproducir \playbtn. También puedes mantener presionada la tecla \textbf{Mayús} (\textit{Shift}) y presionar \textbf{Enter}. La primera vez puede tardar unos segundos en arrancar. Luego verás esto debajo:

\begin{outputbox}
¡Hola!
\end{outputbox}

\textbf{¡Acabas de ejecutar tu primer programa de Python!}

\step{Paso 5: Agrega más celdas}
Haz clic en \textbf{+ Código} (\textit{+ Code}), arriba, para obtener una celda nueva y vacía. Usa una celda nueva para cada programa. Puedes volver a ejecutar cualquier celda cuando quieras.

\needspace{16\baselineskip}
\subsection*{La celda «En el principio»: cómo conseguir la tortuga}

En Scratch, consigues los bloques de lápiz haciendo clic en el botón \textbf{Agregar Extensión}, en la esquina de abajo. Python funciona igual: agregamos las herramientas de dibujo que necesitamos. En Python, el objeto que dibuja se llama \textbf{tortuga}. Camina por el escenario y deja una línea detrás, igual que un objeto de Scratch después de \chip[pen]{bajar lápiz}.

Cada vez que abras tu cuaderno, ejecuta primero esta celda. La llamamos la \textbf{celda «En el principio»}. La Biblia empieza con «En el principio creó Dios los cielos y la tierra» (Génesis 1:1, Reina-Valera 1960), y cada sesión de Colab empieza con esta celda:

\begin{typethis}[La celda «En el principio»: ejecútala una vez cada vez que abras Colab]
!pip install ColabTurtlePlus
\end{typethis}

El \texttt{!} del principio significa: «esto no es Python. Es una instrucción para que Colab descargue las herramientas de la tortuga». Va a pasar mucho texto por la pantalla. No pasa nada. Espera a que se detenga.

\begin{watchout}[¡Cuidado: Plus!]
Debe decir \textbf{ColabTurtlePlus}, con \textbf{Plus} al final. Existe otra biblioteca, más vieja, llamada \texttt{ColabTurtle}. Tiene un escenario negro, otro tamaño y el (0, 0) en la esquina equivocada, así que ninguno de los programas de este libro funcionaría con ella.
\end{watchout}

\needspace{20\baselineskip}
\subsection*{Las tres líneas de inicio}

Todo programa que dibuja empieza con estas tres líneas. Hacen el mismo trabajo que este bloque de Scratch:\par\chip[events]{al presionar bandera verde}

\begin{typethis}[Empieza cada dibujo con estas líneas]
from ColabTurtlePlus.Turtle import *
clearscreen()
initializeTurtle((480, 360))
\end{typethis}

\begin{longtable}{>{\centering\arraybackslash}p{1.1cm} P{14.2cm}}
\toprule
\textbf{Línea} & \textbf{Qué hace} \\
\midrule
1 & \textbf{Trae las herramientas de la tortuga.} Es como hacer clic en Agregar Extensión y elegir Lápiz. El \texttt{*} significa «tráelas todas». \\
2 & \textbf{Limpia el escenario.} Borra el dibujo viejo y las tortugas viejas, para que empieces de cero cada vez que ejecutes la celda. \\
3 & \textbf{Hace un escenario de 480 de ancho y 360 de alto, y prepara la tortuga.} ¡Es exactamente el mismo tamaño que el escenario de Scratch! Los dos números van dentro de su propio par extra de paréntesis, \texttt{(480, 360)}, porque el ancho y el alto viajan juntos como una sola cosa. \\
\bottomrule
\end{longtable}

Después de estas líneas, simplemente le dices a la tortuga qué hacer. La tortuga es como el \textbf{objeto seleccionado} en Scratch: cuando haces clic en un objeto, los bloques que encajas le pertenecen a él. En Python, solo escribes \py|forward(100)|, y la tortuga sabe que le hablas a ella.

\needspace{14\baselineskip}
¡Pruébalo! Escribe las tres líneas de inicio y luego una más:

\begin{typethis}
from ColabTurtlePlus.Turtle import *
clearscreen()
initializeTurtle((480, 360))
forward(100)
\end{typethis}

La tortuga empieza en el centro del escenario, en (0, 0), mirando a la \textbf{derecha}, igual que un objeto nuevo de Scratch. Cuando ejecutas la celda, camina 100 pasos a la derecha y dibuja una línea.

Cuando ejecutes las líneas de inicio, puede que Colab también muestre una notita sobre \py|clearscreen()|. No pasa nada. Es solo un recordatorio, y puedes ignorarlo.

\begin{watchout}[Si Colab se atora o se olvida]
\begin{itemize}
  \item \textbf{¿Una celda se queda ejecutándose para siempre?} Haz clic en el botón \stopbtn{}, donde estaba el botón de reproducir. Es igual que la señal roja de detener en Scratch.
  \item \textbf{¿Ves \texttt{No module named 'ColabTurtlePlus'}?} Colab lo olvida todo cuando lo cierras o lo dejas solo un rato. Vuelve a ejecutar la celda «En el principio».
  \item \textbf{¿Ves un \texttt{NameError}?} Python todavía no conoce una palabra, porque una celda de arriba no se ha ejecutado. Ejecuta primero las celdas de arriba, empezando desde la primera. O haz clic en \textbf{Entorno de ejecución} (\textit{Runtime}) y luego en \textbf{Ejecutar todas} (\textit{Run all}), para ejecutar todas las celdas en orden.
\end{itemize}
Python siempre escribe sus mensajes de error en inglés. La palabra importante es el nombre del error, como \texttt{NameError}.
\end{watchout}

\begin{think}[Truco: activa los números de línea]
Las tablas línea por línea de este libro hablan de «la línea 3» y «la línea 4». Colab también puede numerar las líneas de tus celdas. Haz clic en \textbf{Herramientas} (\textit{Tools}), luego en \textbf{Configuración} (\textit{Settings}), luego en \textbf{Editor}, y marca \textbf{Mostrar números de línea} (\textit{Show line numbers}).
\end{think}

% =====================================================================
\banner{motion}{EL PANORAMA COMPLETO}{Palabras de Scratch y de Python}

Aquí están las cosas de Scratch que ya conoces y cómo las llama Python.

\begin{center}
\renewcommand{\arraystretch}{1.35}
\begin{tabular}{P{3.9cm} P{4.6cm} P{7.0cm}}
\toprule
\textbf{En Scratch} & \textbf{En Python con Colab} & \textbf{Ejemplo} \\
\midrule
El objeto seleccionado & La tortuga & \py|forward(100)| \\
Otro objeto & Una tortuga nueva con su propio nombre & \py|gato = Turtle()| \\
El escenario & El área de dibujo & \lstinline[style=inl,basicstyle=\ttfamily\small]|initializeTurtle((480, 360))| \\
La bandera verde & El botón de reproducir \playbtn & o Mayús + Enter \\
La señal de detener & El botón de detener \stopbtn & \\
Un bloque & Una línea de código & \py|forward(10)| \\
Botón Agregar Extensión & \texttt{import} & \py|import random| \\
Un hueco blanco para escribir & Paréntesis \texttt{( )} & \py|right(90)| \\
Palabras en un hueco & Palabras entre comillas & \py|"¡Hola!"| \\
Crear una variable & Solo dale un valor & \py|puntos = 0| \\
Mis Bloques & Una función & \py|def saltar():| \\
\bottomrule
\end{tabular}
\end{center}

\begin{think}[Más de un objeto]
La mayoría de los programas solo necesitan una tortuga, así que solo escribes \py|forward(100)|. Cuando quieres un segundo objeto, creas una tortuga nueva y le pones nombre: \py|gato = Turtle()|. Después le dices a esa qué hacer escribiendo primero su nombre y un punto: \py|gato.forward(10)| significa «gato, muévete 10 pasos».
\end{think}

\needspace{20\baselineskip}
\subsection*{El escenario es el mismo}

Buena noticia: el escenario de la tortuga tiene las mismas coordenadas que el escenario de Scratch. El centro es (0, 0). \textbf{x} va de −240 a la izquierda a 240 a la derecha. \textbf{y} va de −180 abajo a 180 arriba.

\begin{center}
\begin{tikzpicture}[scale=0.016]
  \draw[ink, line width=1pt, fill=white] (-240,-180) rectangle (240,180);
  \draw[muted!50, -{Stealth}] (-240,0) -- (255,0) node[right, font=\small, text=ink] {x};
  \draw[muted!50, -{Stealth}] (0,-180) -- (0,195) node[above, font=\small, text=ink] {y};
  \fill[motion] (0,0) circle (7);
  \node[font=\small, anchor=south west] at (5,5) {(0, 0)};
  \node[font=\small, anchor=east] at (-245,0) {−240};
  \node[font=\small, anchor=west] at (262,-18) {240};
  \node[font=\small, anchor=north] at (0,-185) {−180};
  \node[font=\small, anchor=south west] at (6,182) {180};
  \fill[looks] (100,80) circle (7) node[right, font=\small, text=ink] {\ (100, 80)};
\end{tikzpicture}
\end{center}

Así que \chip[motion]{ir a x: 100 y: 80} es \py|goto(100, 80)|, y la tortuga cae en el punto morado, exactamente en el mismo lugar que en Scratch.

\needspace{12\baselineskip}
\subsection*{¡Pero la brújula está girada!}

\begin{watchout}[¡Cuidado: la brújula está girada!]
En Scratch, \chip[motion]{apuntar en dirección 90} apunta a la \textbf{derecha}, y 0 apunta \textbf{arriba}.

En Python, \py|setheading(0)| apunta a la \textbf{derecha}, y 90 apunta \textbf{arriba}.
\end{watchout}

\begin{center}
\begin{tikzpicture}[scale=1.1]
  \foreach \cx/\title/\up/\right/\down/\left in {0/Scratch/0/90/180/-90, 7/Python/90/0/270/180}{
    \begin{scope}[shift={(\cx,0)}]
      \draw[ink, line width=1pt] (0,0) circle (1.3);
      \draw[-{Stealth}, line width=1.2pt, motion] (0,0) -- (0,1.2);
      \draw[-{Stealth}, line width=1.2pt, motion] (0,0) -- (1.2,0);
      \draw[-{Stealth}, line width=1.2pt, motion] (0,0) -- (0,-1.2);
      \draw[-{Stealth}, line width=1.2pt, motion] (0,0) -- (-1.2,0);
      \node[above, font=\bfseries] at (0,1.35) {\up};
      \node[right, font=\bfseries] at (1.35,0) {\right};
      \node[below, font=\bfseries] at (0,-1.35) {\down};
      \node[left, font=\bfseries] at (-1.35,0) {\left};
      \node[font=\Large\bfseries] at (0,-2.4) {\title};
    \end{scope}}
\end{tikzpicture}
\end{center}

\begin{think}[La regla secreta]
Para convertir cualquier dirección de Scratch en una dirección de Python, réstala de 90:
\begin{center}\textbf{dirección de Python = 90 − dirección de Scratch}\end{center}
¡Ni siquiera tienes que hacer la resta tú! ¡Python la puede hacer por ti!
\end{think}

\begin{tcbraster}[raster columns=2, raster equal height=rows, raster column skip=12pt, raster valign=center]
\begin{tcolorbox}[colback=white, colframe=white, boxrule=0pt, valign=center]
\begin{scratch}
\blockmove{apuntar en dirección \ovalnum{180}}
\end{scratch}
\end{tcolorbox}
\begin{code}
setheading(90 - 180)
\end{code}
\end{tcbraster}

90 − 180 es −90. Python también entiende −90: apunta \textbf{abajo}, igual que el 180 de Scratch.

\begin{center}
\renewcommand{\arraystretch}{1.25}
\begin{tabular}{lll}
\toprule
\textbf{Dirección en Scratch} & \textbf{Python} & \textbf{Apunta} \\
\midrule
0 & \py|setheading(90 - 0)| & arriba \\
90 & \py|setheading(90 - 90)| & derecha \\
180 & \py|setheading(90 - 180)| & abajo \\
−90 & \py|setheading(90 - -90)| & izquierda \\
\bottomrule
\end{tabular}
\end{center}

La buena noticia es que \textbf{girar} funciona igual. \chip[motion]{girar a la derecha 90 grados} es \py|right(90)|, y \chip[motion]{girar a la izquierda 90 grados} es \py|left(90)|. Casi siempre vas a usar giros, no direcciones.

% =====================================================================
\banner{control}{EL PANORAMA COMPLETO}{Cuatro hábitos nuevos}

Python entiende las mismas ideas que Scratch. Pero como escribes en lugar de encajar bloques, hay cuatro hábitos nuevos que aprender.

\subsection*{Hábito 1: Los espacios muestran qué va adentro}

En Scratch, un bloque queda \textbf{dentro} de un bucle cuando lo sueltas en la boca en forma de C.

En Python, dos cosas hacen ese trabajo:
\begin{enumerate}
  \item unos \textbf{dos puntos} \texttt{:} al final de la línea abren la boca.
  \item \textbf{cuatro espacios} al principio de una línea la meten adentro.
\end{enumerate}

\begin{tcbraster}[raster columns=2, raster equal height=rows, raster column skip=12pt, raster valign=center]
\begin{tcolorbox}[colback=white, colframe=white, boxrule=0pt, valign=center]
\begin{scratch}
\blockrepeat{repetir \ovalnum{4}}{
  \blockmove{mover \ovalnum{100} pasos}
  \blockmove{girar \turnright{} \ovalnum{90} grados}
}
\blockpen{subir lápiz}
\end{scratch}
\end{tcolorbox}
\begin{spaces}
for i in range(4):
    forward(100)
    right(90)
penup()
\end{spaces}
\end{tcbraster}

Las marquitas \texttt{\textvisiblespace} muestran los espacios, que normalmente son invisibles. Las líneas 2 y 3 tienen cuatro espacios, así que están \textbf{dentro} del bucle y se repiten 4 veces. La línea 4 no tiene espacios, así que está \textbf{fuera} del bucle. Pasa una sola vez, cuando el bucle ya terminó.

\begin{think}[Colab te ayuda]
Cuando escribes unos dos puntos \texttt{:} y presionas Enter, Colab pone los cuatro espacios por ti. Para salir del bucle, presiona \textbf{Retroceso} (\textit{Backspace}) una vez.
\end{think}

\needspace{14\baselineskip}
\subsection*{Hábito 2: Escríbelo exactamente igual}

Los bloques ya vienen escritos. En Python, tú escribes, y Python es muy quisquilloso. Las mayúsculas, los paréntesis y las comillas tienen que estar perfectos.

\begin{tcbraster}[raster columns=2, raster equal height=rows, raster column skip=12pt, raster valign=center]
\begin{tcolorbox}[colback=white, colframe=white, boxrule=0pt, valign=center]
\begin{scratch}
\blockmove{mover \ovalnum{100} pasos}
\end{scratch}
\end{tcolorbox}
\begin{code}
forward(100)
\end{code}
\end{tcbraster}

\begin{center}
\renewcommand{\arraystretch}{1.3}
\begin{tabular}{P{4.2cm} P{11cm}}
\toprule
\py|forward(100)| & \textbf{¡Funciona!} \\
\py|Forward(100)| & \textbf{Error:} F mayúscula. Para Python, \texttt{Forward} y \texttt{forward} son palabras diferentes. \\
\py|foward(100)| & \textbf{Error:} falta la r. Python no puede adivinar lo que quisiste decir. \\
\py|forward 100| & \textbf{Error:} no hay paréntesis alrededor del 100. \\
\py|forward(100| & \textbf{Error:} falta el paréntesis de cierre. \\
\py|print("Hola)| & \textbf{Error:} faltan las comillas de cierre. \\
\bottomrule
\end{tabular}
\end{center}

\begin{think}[La regla de los pares]
¡Los paréntesis y las comillas siempre van en pares, como los calcetines! Cada \texttt{(} necesita un \texttt{)}, y cada \texttt{"} necesita otro \texttt{"}.
\end{think}

\needspace{20\baselineskip}
\subsection*{Hábito 3: Una línea a la vez, desde arriba}

En Scratch, puedes tener muchos programas que empiezan cuando haces clic en la bandera verde. Todos se ejecutan al mismo tiempo.

Python es diferente. Lee tu celda como un libro: la línea 1, luego la línea 2, luego la línea 3, hasta llegar abajo. Solo salta de un lado a otro cuando un bucle o una de tus propias funciones se lo pide.

\begin{tcbraster}[raster columns=2, raster equal height=rows, raster column skip=12pt, raster valign=center]
\begin{tcolorbox}[colback=white, colframe=white, boxrule=0pt, valign=center]
\begin{scratch}
\blockinit{al presionar \greenflag}
\blocklook{decir \ovalnum{primero}}
\blocklook{decir \ovalnum{segundo}}
\blocklook{decir \ovalnum{tercero}}
\end{scratch}
\end{tcolorbox}
\begin{code}
print("primero")
print("segundo")
print("tercero")
\end{code}
\end{tcbraster}

\needspace{24\baselineskip}
\subsection*{Hábito 4: Un = guarda, dos == comparan}

Scratch tiene dos bloques diferentes: uno para \textbf{fijar} una variable y otro para \textbf{comprobar} si es igual a algo. Python usa signos de igual para las dos cosas, así que necesita una forma de distinguirlas:

\begin{itemize}
  \item \textbf{Un} signo de igual \texttt{=} significa «pon este valor en la caja». Es como \chip[variables]{fijar puntos a 0}.
  \item \textbf{Dos} signos de igual \texttt{==} significan «¿son iguales?». Es como el bloque verde \chip[operators]{\ = \ }.
\end{itemize}

\begin{tcbraster}[raster columns=2, raster equal height=rows, raster column skip=12pt, raster valign=center]
\begin{tcolorbox}[colback=white, colframe=white, boxrule=0pt, valign=center]
\begin{scratch}
\blockvariable{fijar \selectmenu{puntos} a \ovalnum{5}}
\blockif{si \booloperator{\ovalvariable{puntos} = \ovalnum{5}} entonces}{
  \blocklook{decir \ovalnum{¡Ganaste!}}
}
\end{scratch}
\end{tcolorbox}
\begin{code}
puntos = 5
if puntos == 5:
    print("¡Ganaste!")
\end{code}
\end{tcbraster}

% =====================================================================
\banner{pen}{EL PANORAMA COMPLETO}{En qué se diferencia Colab de Scratch}
\enlargethispage{3\baselineskip}

\subsection*{Colab pregunta en lugar de escuchar}

Colab no oye las teclas ni los clics del ratón. Por eso, en vez de \chip[events]{al presionar tecla espacio}, tu programa te \textbf{pregunta} con \py|input()|, y tú escribes tu respuesta en un recuadro debajo de la celda.

\begin{typethis}
respuesta = input("¿Te gustan los gatos? ")
if respuesta == "si":
    print("¡A mí también!")
\end{typethis}

Escribe \texttt{si} (sin acento), presiona Enter, y Python te contesta. Más adelante en el libro conocerás un programa que sí puede oír las teclas, para hacer juegos de verdad con teclas.

\subsection*{Un solo bucle mueve todos los objetos}

\noindent
\begin{minipage}[c]{0.64\linewidth}
\begin{typethis}[Escribe esto]
from ColabTurtlePlus.Turtle import *
clearscreen()
initializeTurtle((480, 360))
gato = Turtle()
perro = Turtle()
perro.left(90)
for i in range(10):
    gato.forward(10)
    perro.forward(10)
\end{typethis}
\end{minipage}\hfill
\begin{minipage}[c]{0.33\linewidth}\raggedright
Python hace una línea a la vez, así que \textbf{un solo bucle mueve todos los objetos} un poquito en cada vuelta.

Aquí el gato camina a la derecha y el perro camina hacia arriba. En Scratch, esto serían dos programas.
\end{minipage}

\subsection*{Figuras y colores en lugar de disfraces}

En Colab no hay dibujos de disfraces, pero la tortuga puede ponerse muchas \textbf{figuras}, como \texttt{"turtle"} (tortuga), \texttt{"arrow"} (flecha) o \texttt{"circle"} (círculo), del color que más te guste: \py|shape("turtle")| y \py|color("gold")|. Las figuras y los colores se escriben en inglés y en minúsculas, porque Python solo los entiende así: \texttt{"red"} es rojo, \texttt{"blue"} azul, \texttt{"gold"} dorado... (\texttt{color} es como Python dice «color».) Además, puede dibujar todo lo que te imagines, línea por línea. \textbf{Ya estás listo. ¡Vamos!}

% Generado por chapters/_make_mission_chart.py -- no editar a mano.
\banner{events}{SIGUE TU PROGRESO}{Mis misiones}
Colorea una estrella cada vez que termines una misión: una estrella por hacer que funcione, dos por hacer un reto, ¡y tres por hacerlos todos!
\begin{center}\renewcommand{\arraystretch}{1.12}
\begin{longtable}{c P{10.5cm} l}
\toprule
\textbf{Misión} & \textbf{Lo que haces} & \textbf{¿Lista?} \\
\midrule
\endhead
\multicolumn{3}{l}{\rule{0pt}{14pt}\textbf{Parte 1: Movimiento y Lápiz}} \\
1 & Un cuadrado por el camino largo & $\star$\ \ $\star$\ \ $\star$ \\
2 & Dibuja una casa & $\star$\ \ $\star$\ \ $\star$ \\
3 & Cara de robot & $\star$\ \ $\star$\ \ $\star$ \\
\multicolumn{3}{l}{\rule{0pt}{14pt}\textbf{Parte 2: Bucles}} \\
4 & Dibuja figuras por el camino corto & $\star$\ \ $\star$\ \ $\star$ \\
5 & Estrellas & $\star$\ \ $\star$\ \ $\star$ \\
6 & Espiral cuadrada & $\star$\ \ $\star$\ \ $\star$ \\
7 & Flor de cuadrados & $\star$\ \ $\star$\ \ $\star$ \\
\multicolumn{3}{l}{\rule{0pt}{14pt}\textbf{Parte 3: Variables y listas}} \\
8 & Contador de puntos & $\star$\ \ $\star$\ \ $\star$ \\
9 & Espiral que crece & $\star$\ \ $\star$\ \ $\star$ \\
10 & Arcoíris con una lista de colores & $\star$\ \ $\star$\ \ $\star$ \\
11 & Inventor de historias al azar & $\star$\ \ $\star$\ \ $\star$ \\
\multicolumn{3}{l}{\rule{0pt}{14pt}\textbf{Parte 4: Operadores}} \\
12 & Lanzador de dados & $\star$\ \ $\star$\ \ $\star$ \\
13 & Impresora de tablas & $\star$\ \ $\star$\ \ $\star$ \\
14 & Arte de puntos al azar & $\star$\ \ $\star$\ \ $\star$ \\
15 & Máquina de nombres & $\star$\ \ $\star$\ \ $\star$ \\
\multicolumn{3}{l}{\rule{0pt}{14pt}\textbf{Parte 5: Sensores y decisiones}} \\
16 & Di hola & $\star$\ \ $\star$\ \ $\star$ \\
17 & Adivina mi número & $\star$\ \ $\star$\ \ $\star$ \\
18 & Concurso con puntos & $\star$\ \ $\star$\ \ $\star$ \\
19 & Prueba de reflejos & $\star$\ \ $\star$\ \ $\star$ \\
\multicolumn{3}{l}{\rule{0pt}{14pt}\textbf{Parte 6: Mis Bloques}} \\
20 & Haz tu propio bloque & $\star$\ \ $\star$\ \ $\star$ \\
21 & Máquina de figuras & $\star$\ \ $\star$\ \ $\star$ \\
22 & Máquina de respuestas & $\star$\ \ $\star$\ \ $\star$ \\
23 & Ejército de objetos & $\star$\ \ $\star$\ \ $\star$ \\
\multicolumn{3}{l}{\rule{0pt}{14pt}\textbf{Parte 7: Apariencia, sonido y eventos}} \\
24 & Desfile de disfraces & $\star$\ \ $\star$\ \ $\star$ \\
25 & Hacer música & $\star$\ \ $\star$\ \ $\star$ \\
26 & Carrera de tortugas & $\star$\ \ $\star$\ \ $\star$ \\
27 & Cuento con mensajes & $\star$\ \ $\star$\ \ $\star$ \\
\multicolumn{3}{l}{\rule{0pt}{14pt}\textbf{Grandes proyectos}} \\
28 & Anima tu nombre & $\star$\ \ $\star$\ \ $\star$ \\
29 & Flor de lápiz & $\star$\ \ $\star$\ \ $\star$ \\
30 & Concurso de tablas & $\star$\ \ $\star$\ \ $\star$ \\
31 & Adivina el número de lujo & $\star$\ \ $\star$\ \ $\star$ \\
32 & Elige tu propia historia & $\star$\ \ $\star$\ \ $\star$ \\
33 & Mascota virtual & $\star$\ \ $\star$\ \ $\star$ \\
34 & Juego de persecución & $\star$\ \ $\star$\ \ $\star$ \\
35 & Juego de atrapar & $\star$\ \ $\star$\ \ $\star$ \\
\bottomrule
\end{longtable}
\end{center}


% =====================================================================
% PARTE 1: Movimiento y Lápiz
% =====================================================================
\partopener{motion}{PARTE 1}{Movimiento y Lápiz}{Haz que la tortuga camine, gire y dibuje, igual que un objeto con la extensión Lápiz.}

En Scratch, los bloques azules de \textbf{Movimiento} mueven tu objeto por el escenario, y los bloques verdes de \textbf{Lápiz} le permiten dibujar. En Python, la \textbf{tortuga} hace los dos trabajos. Es un objeto que siempre lleva un lápiz.

En esta parte aprenderás el Python de bloques como estos:

{\raggedright \chip[motion]{mover 10 pasos} \chip[motion]{girar $\circlearrowright$ 15 grados} \chip[motion]{ir a x: 0 y: 0} \chip[motion]{desplazar en 1 segs a x: 0 y: 0} \chip[motion]{apuntar en dirección 90} \chip[motion]{cambiar x en 10} \chip[motion]{posición en x} \chip[pen]{bajar lápiz} \chip[pen]{subir lápiz} \chip[pen]{fijar color de lápiz a} \chip[pen]{fijar tamaño de lápiz a 1} \chip[pen]{borrar todo} \chip[pen]{sellar}\par}

Y aprenderás algunos trucos de Python que Scratch \textbf{no puede} hacer, como círculos, puntos, figuras rellenas de color y escribir palabras en el escenario.

\subsection*{Antes de empezar}

Todos los programas de dibujo de esta parte empiezan con \textbf{las tres líneas de inicio}. Son como \chip[events]{al presionar bandera verde}. Traen a la tortuga, limpian el escenario y lo hacen del mismo tamaño que el escenario de Scratch.

\begin{typethis}[Las tres líneas de inicio]
from ColabTurtlePlus.Turtle import *
clearscreen()
initializeTurtle((480, 360))
\end{typethis}

En las páginas de conceptos que vienen, los recuadros muestran solo las líneas nuevas. Escribe primero las tres líneas de inicio, luego las líneas nuevas debajo, y presiona \playbtn.

\begin{think}[¿Dónde empieza la tortuga?]
\begin{itemize}
  \item En el centro del escenario, en $(0, 0)$, igual que un objeto nuevo.
  \item Mirando a la \textbf{derecha}, igual que un objeto nuevo (dirección 90 en Scratch).
  \item Con el lápiz \textbf{abajo}. ¡Esto es diferente! Un objeto de Scratch solo dibuja después de \chip[pen]{bajar lápiz}, pero la tortuga dibuja una línea cada vez que se mueve.
\end{itemize}
\end{think}

% ---------------------------------------------------------------------
\needspace{10\baselineskip}
\subsection*{Mover y girar}

\begin{tcbraster}[raster columns=2, raster equal height=rows, raster column skip=12pt, raster valign=center]
\begin{tcolorbox}[colback=white, colframe=white, boxrule=0pt, valign=center]
\begin{scratch}
\blockmove{mover \ovalnum{100} pasos}
\blockmove{girar \turnright{} \ovalnum{90} grados}
\blockmove{mover \ovalnum{100} pasos}
\blockmove{girar \turnleft{} \ovalnum{45} grados}
\blockmove{mover \ovalnum{-50} pasos}
\end{scratch}
\end{tcolorbox}
\begin{code}
forward(100)
right(90)
forward(100)
left(45)
backward(50)
\end{code}
\end{tcbraster}

\textbf{Cómo funciona:} \py|forward(100)| es \chip[motion]{mover 100 pasos}. El número va dentro de los paréntesis \texttt{( )}, igual que va dentro de la burbuja blanca del bloque. Un paso es un píxel en los dos, así que tus números de Scratch también sirven en Python. Para caminar hacia atrás, Scratch usa un número negativo, pero Python tiene su propia palabra, \py|backward| («hacia atrás»).

\py|right(90)| gira en el sentido de las agujas del reloj, como la flecha $\circlearrowright$, y \py|left(45)| gira al revés, como la flecha $\circlearrowleft$. La buena noticia: girar funciona \textbf{exactamente} igual en Scratch y en Python.

Python habla inglés: sus palabras (\py|forward|, \py|right|, \py|left|, \py|print|...) siempre van en inglés. \py|forward| quiere decir «adelante», \py|right| «derecha» y \py|left| «izquierda».

\begin{watchout}[¡Cuidado: solo minúsculas!]
A Python le importan mucho las mayúsculas. \py|forward(100)| funciona, pero \py|Forward(100)| da un error. Todos los comandos de esta parte se escriben con letras minúsculas. (Python siempre escribe sus mensajes de error en inglés. Lo importante es el nombre del error, como \texttt{NameError}.)
\end{watchout}

% ---------------------------------------------------------------------
\needspace{14\baselineskip}
\subsection*{Ir a un lugar}

\begin{tcbraster}[raster columns=2, raster equal height=rows, raster column skip=12pt, raster valign=center]
\begin{tcolorbox}[colback=white, colframe=white, boxrule=0pt, valign=center]
\begin{scratch}
\blockmove{ir a x: \ovalnum{100} y: \ovalnum{50}}
\blockmove{ir a x: \ovalnum{-100} y: \ovalnum{50}}
\blockmove{ir a x: \ovalnum{0} y: \ovalnum{0}}
\end{scratch}
\end{tcolorbox}
\begin{code}
goto(100, 50)
goto(-100, 50)
goto(0, 0)
\end{code}
\end{tcbraster}

\textbf{Cómo funciona:} \py|goto(100, 50)| quiere decir «ir a x: 100, y: 50». Siempre va primero el número x, luego una coma y luego el número y. El escenario es igual que en Scratch: $(0, 0)$ es el centro, x va de $-240$ a $240$ y y va de $-180$ a $180$. Como el lápiz está abajo, ¡este programa dibuja un triángulo!

% ---------------------------------------------------------------------
\needspace{13\baselineskip}
\subsection*{Desplazarse despacio}

\begin{tcbraster}[raster columns=2, raster equal height=rows, raster column skip=12pt, raster valign=center]
\begin{tcolorbox}[colback=white, colframe=white, boxrule=0pt, valign=center]
\begin{scratch}[scale=0.89]
\blockmove{desplazar en \ovalnum{3} segs a x: \ovalnum{150} y: \ovalnum{100}}
\blockmove{desplazar en \ovalnum{0.2} segs a x: \ovalnum{-150} y: \ovalnum{100}}
\end{scratch}
\end{tcolorbox}
\begin{code}
speed(1)
goto(150, 100)
speed(13)
goto(-150, 100)
\end{code}
\end{tcbraster}

\textbf{Cómo funciona:} La tortuga siempre se desplaza cuando se mueve. Tú eliges qué tan rápido con \py|speed()| («velocidad»). Lleva un número del \textbf{1} (lo más lento, como un caracol) al \textbf{13} (lo más rápido, como un cohete). Python no usa segundos, así que eliges la velocidad que se vea bien. La velocidad no cambia hasta que tú la cambies, así que pon \py|speed(13)| cerca del principio cuando quieras que un dibujo grande termine rápido.

\begin{watchout}[¡Cuidado: no uses la velocidad 0!]
Podrías pensar que \py|speed(0)| quiere decir «quédate quieta». ¡Pues no! En Colab quiere decir «no muestres nada hasta el final», así que tu escenario se queda vacío. Usa solo del 1 al 13.
\end{watchout}

% ---------------------------------------------------------------------
\needspace{13\baselineskip}
\subsection*{Apuntar en una dirección}

\begin{tcbraster}[raster columns=2, raster equal height=rows, raster column skip=12pt, raster valign=center]
\begin{tcolorbox}[colback=white, colframe=white, boxrule=0pt, valign=center]
\begin{scratch}
\blockmove{apuntar en dirección \ovalnum{0}}
\blockmove{mover \ovalnum{100} pasos}
\blockmove{apuntar en dirección \ovalnum{180}}
\blockmove{mover \ovalnum{50} pasos}
\end{scratch}
\end{tcolorbox}
\begin{code}
setheading(90 - 0)
forward(100)
setheading(90 - 180)
forward(50)
\end{code}
\end{tcbraster}

\textbf{Cómo funciona:} \py|setheading| quiere decir «apuntar en dirección». ¡Pero la brújula de Python está girada! En Scratch, el 0 apunta \textbf{arriba} y los números van en el sentido de las agujas del reloj. En Python, el 0 apunta a la \textbf{derecha} y los números van al revés. Python sabe hacer cuentas dentro de los paréntesis, así que dejamos que haga la traducción por nosotros.

\begin{center}
\begin{tikzpicture}[scale=0.7]
  \foreach \cx/\title/\up/\right/\down/\left in {0/Scratch/0/90/180/-90, 7/Python/90/0/-90/180}{
    \begin{scope}[shift={(\cx,0)}]
      \draw[ink, line width=1pt] (0,0) circle (1.3);
      \foreach \a in {0,90,180,270} \draw[-{Stealth}, line width=1.2pt, motion] (0,0) -- (\a:1.2);
      \node[above, font=\bfseries] at (0,1.35) {\up};
      \node[right, font=\bfseries] at (1.35,0) {\right};
      \node[below, font=\bfseries] at (0,-1.35) {\down};
      \node[left, font=\bfseries] at (-1.35,0) {\left};
      \node[font=\bfseries] at (0,-2.7) {\title};
    \end{scope}}
\end{tikzpicture}
\end{center}

\begin{think}[La regla secreta]
\begin{center}\textbf{dirección en Python} $=$ \textbf{90} $-$ \textbf{dirección en Scratch}\end{center}
Así, \chip[motion]{apuntar en dirección 30} se convierte en \py|setheading(90 - 30)|. Escribe el número de Scratch después del \texttt{90 -} y Python calcula el resto. ¡Nunca tienes que hacer la cuenta de cabeza!
\end{think}

\needspace{12\baselineskip}
\textbf{Completa la tabla.} Calcula cada dirección en Python. Luego pruébala en Colab con \py|forward(100)| después, y revisa que la tortuga vaya hacia donde esperabas.

\begin{center}
\renewcommand{\arraystretch}{1.5}
\begin{tabular}{P{2.8cm} P{6.0cm} P{3.0cm} P{2.8cm}}
\toprule
\textbf{Dirección en Scratch} & \textbf{Código de Python} & \textbf{Dirección en Python} & \textbf{¿Hacia dónde?} \\
\midrule
90 & \py|setheading(90 - 90)| & 0 & derecha \\
0 & \py|setheading(90 - 0)| & 90 & arriba \\
-90 & \py|setheading(90 - -90)| & \rule{2cm}{0.5pt} & \rule{2cm}{0.5pt} \\
45 & \py|setheading(90 - 45)| & \rule{2cm}{0.5pt} & \rule{2cm}{0.5pt} \\
135 & \py|setheading(90 - 135)| & \rule{2cm}{0.5pt} & \rule{2cm}{0.5pt} \\
\bottomrule
\end{tabular}
\end{center}

% ---------------------------------------------------------------------
\needspace{13\baselineskip}
\subsection*{Apuntar hacia algo}

\begin{tcbraster}[raster columns=5, raster equal height=rows, raster column skip=12pt, raster valign=center]
\begin{tcolorbox}[raster multicolumn=2, colback=white, colframe=white, boxrule=0pt, valign=center]
\begin{scratch}
\blockmove{apuntar hacia \selectmenu{Estrella}}
\blockmove{mover \ovalnum{100} pasos}
\end{scratch}
\end{tcolorbox}
\begin{tcolorbox}[raster multicolumn=3, blank]
\begin{code}
setheading(towards(100, 100))
forward(100)
\end{code}
\end{tcolorbox}
\end{tcbraster}

\textbf{Cómo funciona:} En Scratch apuntas hacia otro objeto, como una Estrella que está en x: 100, y: 100. En Python apuntas hacia un \textbf{lugar}. \py|towards(100, 100)| calcula en qué dirección está ese lugar, y \py|setheading( )| gira a la tortuga para que lo mire. Un comando va dentro del otro, como un bloque redondo que metes en una burbuja blanca.

\begin{think}[Esta vez no hace falta el 90 menos]
\py|towards| ya da la respuesta en direcciones de Python, así que aquí no necesitas la regla secreta. Solo necesitas el \texttt{90 -} cuando \textbf{tú} escribes un número de dirección de Scratch.
\end{think}

% ---------------------------------------------------------------------
\needspace{13\baselineskip}
\subsection*{Cambiar x y y, fijar x y y}

\begin{tcbraster}[raster columns=2, raster equal height=rows, raster column skip=12pt, raster valign=center]
\begin{tcolorbox}[colback=white, colframe=white, boxrule=0pt, valign=center]
\begin{scratch}
\blockmove{cambiar x en \ovalnum{50}}
\blockmove{cambiar y en \ovalnum{50}}
\blockmove{fijar x a \ovalnum{-100}}
\blockmove{fijar y a \ovalnum{0}}
\end{scratch}
\end{tcolorbox}
\begin{code}
setx(xcor() + 50)
sety(ycor() + 50)
setx(-100)
sety(0)
\end{code}
\end{tcbraster}

\textbf{Cómo funciona:} \py|setx(-100)| es \chip[motion]{fijar x a -100}, y \py|sety(0)| es \chip[motion]{fijar y a 0}. Python no tiene un comando «cambiar x», así que lo construimos. \py|xcor()| quiere decir «mi posición en x». Entonces \py|setx(xcor() + 50)| dice: «fija x a (mi posición en x + 50)». Eso mueve a la tortuga 50 pasos a la derecha, ¡que es exactamente lo que hace \chip[motion]{cambiar x en 50}!

\begin{watchout}
\py|xcor()| necesita sus paréntesis vacíos. Le dicen a Python «ve y trae la respuesta ahora». Sin ellos, Python se confunde.
\end{watchout}

% ---------------------------------------------------------------------
\needspace{13\baselineskip}
\subsection*{¿Dónde estoy? ¿Hacia dónde miro?}

\begin{tcbraster}[raster columns=2, raster equal height=rows, raster column skip=12pt, raster valign=center]
\begin{tcolorbox}[colback=white, colframe=white, boxrule=0pt, valign=center]
\begin{scratch}[scale=0.85]
\blockmove{mover \ovalnum{100} pasos}
\blockmove{girar \turnleft{} \ovalnum{90} grados}
\blocklook{decir \ovalmove{posición en x} durante \ovalnum{2} segundos}
\blocklook{decir \ovalmove{posición en y} durante \ovalnum{2} segundos}
\blocklook{decir \ovalmove{dirección} durante \ovalnum{2} segundos}
\end{scratch}
\end{tcolorbox}
\begin{code}
forward(100)
left(90)
print(xcor())
print(ycor())
print(heading())
\end{code}
\end{tcbraster}

\textbf{Cómo funciona:} Los bloques redondos \chip[motion]{posición en x}, \chip[motion]{posición en y} y \chip[motion]{dirección} se convierten en \py|xcor()|, \py|ycor()| y \py|heading()|. \py|print( )| muestra una respuesta debajo de la celda, un poco como cuando un objeto la dice.

\begin{outputbox}
100.0\\
0.0\\
90
\end{outputbox}

\begin{think}[Dos cosas que descubrir]
\begin{itemize}
  \item \texttt{100.0} es solo 100 escrito con punto decimal. A Python le gusta mostrar las posiciones así.
  \item Scratch diría que la dirección es \textbf{0}, porque el objeto apunta hacia arriba. Python dice \textbf{90}, porque \py|heading()| responde en direcciones de \textbf{Python}. Para obtener el número de Scratch, usa otra vez la regla secreta: \py|print(90 - heading())|.
\end{itemize}
\end{think}

% ---------------------------------------------------------------------
\needspace{14\baselineskip}
\subsection*{Un adelanto: ir a un lugar al azar}

\begin{tcbraster}[raster columns=5, raster equal height=rows, raster column skip=12pt, raster valign=center]
\begin{tcolorbox}[raster multicolumn=2, colback=white, colframe=white, boxrule=0pt, valign=center]
\begin{scratch}
\blockmove{ir a \selectmenu{posición aleatoria}}
\end{scratch}
\end{tcolorbox}
\begin{tcolorbox}[raster multicolumn=3, blank]
\begin{code}
import random
goto(random.randint(-240, 240),
     random.randint(-180, 180))
\end{code}
\end{tcolorbox}
\end{tcbraster}

\textbf{Cómo funciona:} \py|import random| trae los dados de Python. \py|random.randint(-240, 240)| elige un número entero de $-240$ a $240$, así que es una x sorpresa. El segundo elige una y sorpresa. La línea era demasiado larga, así que sigue en la línea de abajo. Python lo permite mientras haya un paréntesis abierto. Conocerás \py|random| a fondo en la Parte 4.

% ---------------------------------------------------------------------
\needspace{10\baselineskip}
\subsection*{Si toca un borde, rebotar}

\begin{tcbraster}[raster columns=2, raster equal height=rows, raster column skip=12pt, raster valign=center]
\begin{tcolorbox}[colback=white, colframe=white, boxrule=0pt, valign=center]
\begin{scratch}
\blockmove{si toca un borde, rebotar}
\end{scratch}
\end{tcolorbox}
\begin{tcolorbox}[colback=warnbg, colframe=warnframe, boxrule=0.6pt, arc=5pt]
Python no tiene un comando para este.
\end{tcolorbox}
\end{tcbraster}

El escenario de la tortuga no tiene paredes. Prueba \py|forward(1000)|: ¡la tortuga se sale por el borde y sigue caminando! En la \textbf{Parte 5} construirás tu propio «si toca un borde, rebotar» con un \py|if|. Entonces conocerás el secreto de lo que hace el bloque de Scratch.

% ---------------------------------------------------------------------
\needspace{13\baselineskip}
\subsection*{Subir lápiz, bajar lápiz}

\begin{tcbraster}[raster columns=2, raster equal height=rows, raster column skip=12pt, raster valign=center]
\begin{tcolorbox}[colback=white, colframe=white, boxrule=0pt, valign=center]
\begin{scratch}
\blockpen{bajar lápiz}
\blockmove{mover \ovalnum{50} pasos}
\blockpen{subir lápiz}
\blockmove{mover \ovalnum{50} pasos}
\blockpen{bajar lápiz}
\blockmove{mover \ovalnum{50} pasos}
\end{scratch}
\end{tcolorbox}
\begin{code}
forward(50)
penup()
forward(50)
pendown()
forward(50)
\end{code}
\end{tcbraster}

\textbf{Cómo funciona:} \py|penup()| levanta el lápiz del papel, así que la tortuga se mueve sin dibujar. \py|pendown()| lo vuelve a bajar. Los paréntesis están vacíos porque estos comandos no necesitan número, pero Python igual quiere los paréntesis. Quieren decir «¡hazlo ahora!». La tortuga empieza con el lápiz abajo, así que el primer \chip[pen]{bajar lápiz} no necesita ninguna línea de Python.

\begin{think}[Moverse sin dejar línea]
Para saltar a otro lugar sin dibujar, usa tres líneas: \py|penup()|, luego \py|goto(x, y)| y luego \py|pendown()|. Vas a usar mucho este truco en las misiones.
\end{think}

% ---------------------------------------------------------------------
\needspace{14\baselineskip}
\subsection*{Elegir un color de lápiz}

\begin{tcbraster}[raster columns=2, raster equal height=rows, raster column skip=12pt, raster valign=center]
\begin{tcolorbox}[colback=white, colframe=white, boxrule=0pt, valign=center]
\begin{scratch}
\blockpen{fijar color de lápiz a \ovalnum{rojo}}
\blockmove{mover \ovalnum{100} pasos}
\blockpen{fijar color de lápiz a \ovalnum{naranja}}
\blockmove{mover \ovalnum{100} pasos}
\end{scratch}
\end{tcolorbox}
\begin{code}
pencolor("red")
forward(100)
pencolor("#ff8800")
forward(100)
\end{code}
\end{tcbraster}

\textbf{Cómo funciona:} En Scratch haces clic en un color. En Python escribes su nombre entre \textbf{comillas}, todo en minúsculas: \py|"red"|, \py|"gold"|, \py|"skyblue"|, \py|"hotpink"|. Las comillas le dicen a Python «esto es una palabra, no un comando». ¡Hay 140 nombres de colores para probar!

\textbf{Los colores se escriben en inglés}, porque Python solo entiende sus nombres en inglés: \py|"red"| es rojo, \py|"blue"| azul, \py|"gold"| dorado, \py|"skyblue"| azul cielo y \py|"hotpink"| rosa fuerte. Si escribes \py|"rojo"|, Python no lo reconoce.

Para cualquier color que quieras, usa un \textbf{código hexadecimal}, como \py|"#ff8800"| (naranja). Es un \texttt{\#} y seis números o letras. Los dos primeros dicen cuánto rojo, los dos del medio cuánto verde y los dos últimos cuánto azul. \texttt{00} es nada y \texttt{ff} es lo máximo.

\begin{watchout}[¡Cuidado: \texttt{pencolor} va sin «u»!]
En inglés británico se escribe «colour», pero Python lo escribe a la manera estadounidense: \py|pencolor|, sin \texttt{u}. Por suerte, es igual que en español. Si escribes \py|pencolour|, Python no conoce la palabra.
\end{watchout}

% ---------------------------------------------------------------------
\needspace{10\baselineskip}
\subsection*{Tamaño del lápiz y borrar todo}

\begin{tcbraster}[raster columns=2, raster equal height=rows, raster column skip=12pt, raster valign=center]
\begin{tcolorbox}[colback=white, colframe=white, boxrule=0pt, valign=center]
\begin{scratch}
\blockpen{fijar tamaño de lápiz a \ovalnum{10}}
\blockmove{mover \ovalnum{100} pasos}
\blockpen{borrar todo}
\end{scratch}
\end{tcolorbox}
\begin{code}
pensize(10)
forward(100)
clear()
\end{code}
\end{tcbraster}

\textbf{Cómo funciona:} Un número más grande en \py|pensize| da una línea más gruesa. La tortuga empieza con tamaño 1, igual que en Scratch. El tamaño tiene que ser mayor que 0: \py|pensize(0)| da un error.

\py|clear()| borra todo lo que la tortuga ha dibujado. La tortuga se queda donde está, y conserva su color y su tamaño, igual que con \chip[pen]{borrar todo}. (La línea de inicio \py|clearscreen()| es el borrador más grande: limpia todo el escenario y empieza de nuevo.)

% ---------------------------------------------------------------------
\needspace{13\baselineskip}
\subsection*{Sellar}

\begin{tcbraster}[raster columns=2, raster equal height=rows, raster column skip=12pt, raster valign=center]
\begin{tcolorbox}[colback=white, colframe=white, boxrule=0pt, valign=center]
\begin{scratch}
\blockpen{sellar}
\blockmove{mover \ovalnum{100} pasos}
\blockpen{sellar}
\end{scratch}
\end{tcolorbox}
\begin{code}
shape("turtle")
stamp()
forward(100)
stamp()
\end{code}
\end{tcbraster}

\textbf{Cómo funciona:} \py|stamp()| imprime una copia de la tortuga en el escenario, como un sello de goma. En Scratch el sello se ve como tu disfraz. En Python se ve como la \textbf{forma} de la tortuga. \py|shape("turtle")| hace que se vea como una tortuga de verdad. También puedes probar las formas \texttt{arrow} (flecha), \texttt{circle} (círculo), \texttt{square} (cuadrado), \texttt{triangle} (triángulo) y \texttt{classic} (la flecha con la que empieza). Los nombres de las formas también van en inglés. ¡Recuerda las comillas!

% ---------------------------------------------------------------------
\needspace{18\baselineskip}
\subsection*{Superpoderes de Python que Scratch no tiene}

El Lápiz de Scratch solo puede dibujar líneas. ¡La tortuga puede hacer más!

\begin{center}
\renewcommand{\arraystretch}{1.4}
\begin{tabular}{P{5.4cm} P{10.4cm}}
\toprule
\textbf{Python} & \textbf{Qué hace} \\
\midrule
\py|circle(50)| & Dibuja un círculo completo con un radio de 50 (es la distancia del centro al borde). El centro del círculo queda a la \textbf{izquierda} de la tortuga. Un número negativo, \py|circle(-50)|, lo pone a la derecha. \\
\py|dot(30, "gold")| & Dibuja un punto relleno, de 30 de ancho, justo donde está la tortuga. Funciona aunque el lápiz esté arriba. \\
\py|fillcolor("skyblue")|\newline \py|begin_fill()|\newline \ldots\newline \py|end_fill()| & Rellena una figura de color. Elige el color, di \py|begin_fill()|, dibuja la figura y luego di \py|end_fill()|. El interior se pinta. \\
\py|write("¡Hola!")| & Escribe palabras en el escenario donde está la tortuga. \\
\py|shape("turtle")| & Cambia cómo se ve la tortuga. Es un poco como elegir un disfraz. \\
\bottomrule
\end{tabular}
\end{center}

\needspace{12\baselineskip}
Pruébalos todos juntos:

\begin{tcbraster}[raster columns=2, raster equal height=rows, raster column skip=12pt, raster valign=center]
\begin{code}
shape("turtle")
circle(50)
dot(30, "gold")
fillcolor("skyblue")
begin_fill()
circle(-40)
end_fill()
write("¡Hola!")
\end{code}
\begin{tcolorbox}[colback=white, colframe=white, boxrule=0pt, valign=center]
Deberías ver un círculo vacío arriba de la tortuga, un círculo azul relleno debajo, un punto dorado donde se tocan y la palabra \texttt{¡Hola!} encima. ¡Parece un muñeco de nieve!
\end{tcolorbox}
\end{tcbraster}

% =====================================================================
\mission{Un cuadrado por el camino largo}{\texttt{forward}, \texttt{right}, \texttt{pencolor}, \texttt{pensize} y \texttt{shape}, todo en un programa}

Es hora de juntar los bloques. Vas a dibujar un cuadrado azul, un lado a la vez.

\step{Paso 1: Constrúyelo en Scratch}
Agrega la extensión \textbf{Lápiz} y luego construye:

\begin{center}
\begin{tabular}{c@{\hspace{1cm}}c}
\begin{scratch}
\blockinit{al presionar \greenflag}
\blockpen{borrar todo}
\blockpen{bajar lápiz}
\blockpen{fijar color de lápiz a \ovalnum{azul}}
\blockpen{fijar tamaño de lápiz a \ovalnum{3}}
\blockmove{mover \ovalnum{100} pasos}
\blockmove{girar \turnright{} \ovalnum{90} grados}
\end{scratch}
&
\begin{scratch}
\blockmove{mover \ovalnum{100} pasos}
\blockmove{girar \turnright{} \ovalnum{90} grados}
\blockmove{mover \ovalnum{100} pasos}
\blockmove{girar \turnright{} \ovalnum{90} grados}
\blockmove{mover \ovalnum{100} pasos}
\blockmove{girar \turnright{} \ovalnum{90} grados}
\end{scratch}
\end{tabular}
\end{center}
Los bloques de la derecha se pegan debajo de los de la izquierda.

\step{Paso 2: Escríbelo en Colab}

\begin{typethis}
from ColabTurtlePlus.Turtle import *
clearscreen()
initializeTurtle((480, 360))
shape("turtle")
pencolor("blue")
pensize(3)
forward(100)
right(90)
forward(100)
right(90)
forward(100)
right(90)
forward(100)
right(90)
\end{typethis}

\step{Paso 3: Léelo línea por línea}

\begin{longtable}{>{\centering\arraybackslash}p{1.1cm} P{5.3cm} P{9.1cm}}
\toprule
\textbf{Línea} & \textbf{Bloque de Scratch} & \textbf{Qué significa} \\
\midrule
1--3 & \chip[events]{al presionar bandera verde} \chip[pen]{borrar todo} & Las tres líneas de inicio. \py|clearscreen()| limpia el escenario, así que también hace el trabajo de \chip[pen]{borrar todo}. \\
--- & \chip[pen]{bajar lápiz} & No hace falta ninguna línea. El lápiz de la tortuga empieza abajo. \\
4 & (elegir un disfraz) & \py|shape("turtle")| hace que la tortuga se vea como una tortuga. \\
5 & \chip[pen]{fijar color de lápiz a azul} & El nombre del color va entre comillas, y en inglés: \py|"blue"| es azul. \\
6 & \chip[pen]{fijar tamaño de lápiz a 3} & Una línea un poco más gruesa. \\
7--8 & \chip[motion]{mover 100 pasos} \chip[motion]{girar $\circlearrowright$ 90 grados} & El primer lado, y luego un giro en la esquina. \\
9--10 & lo mismo otra vez & El segundo lado. \\
11--12 & lo mismo otra vez & El tercer lado. \\
13--14 & lo mismo otra vez & El cuarto lado. El último giro deja a la tortuga mirando hacia donde empezó. \\
\bottomrule
\end{longtable}

\step{Paso 4: Ejecútalo}
Presiona \playbtn. La tortuga camina a la derecha, gira, camina hacia abajo, y así sigue. Deberías ver un cuadrado azul debajo del centro del escenario, con la tortuga otra vez donde empezó, mirando a la derecha.

\begin{think}[¡Uf, qué largo!]
Las líneas 7 y 8 se escriben \textbf{cuatro veces}. En Scratch usarías \chip[control]{repetir 4} para construirlas solo una vez. Python tiene un atajo igualito, y la \textbf{Parte 2} te enseña cómo. Por ahora, escribirlo por el camino largo te ayuda a ver cada paso.
\end{think}

\step{Paso 5: Hazlo tuyo}

\begin{challenge}
\tick Cambia \py|"blue"| por tu color favorito, y el tamaño del lápiz a 8.\par
\tick Haz un cuadrado más grande: cambia los cuatro \texttt{100} por \texttt{150}.\par
\tick Haz cada lado de un color diferente. Agrega una línea \py|pencolor(...)| antes de cada \py|forward|.\par
\tick Cambia cada \py|right| por \py|left|. ¿Hacia dónde va ahora el cuadrado? Primero adivina y luego ejecútalo.
\end{challenge}

% =====================================================================
\mission{Dibuja una casa}{a moverte sin dibujar con \texttt{penup} y \texttt{goto}, y a rellenar figuras de color}

Ahora vas a dibujar una casa con techo, puerta y una ventana redonda. Usarás \py|goto| para cada esquina, y \py|penup()| para moverte entre las partes sin dejar una línea.

\step{Paso 1: Constrúyelo en Scratch}
Scratch no puede rellenar figuras de color, así que la casa de Scratch es solo el contorno. Aquí están las paredes y el techo:

\begin{center}
\begin{tabular}{c@{\hspace{1cm}}c}
\begin{scratch}
\blockinit{al presionar \greenflag}
\blockpen{borrar todo}
\blockpen{fijar tamaño de lápiz a \ovalnum{3}}
\blockpen{subir lápiz}
\blockmove{ir a x: \ovalnum{-100} y: \ovalnum{-150}}
\blockpen{bajar lápiz}
\blockmove{ir a x: \ovalnum{100} y: \ovalnum{-150}}
\blockmove{ir a x: \ovalnum{100} y: \ovalnum{50}}
\blockmove{ir a x: \ovalnum{-100} y: \ovalnum{50}}
\blockmove{ir a x: \ovalnum{-100} y: \ovalnum{-150}}
\end{scratch}
&
\begin{scratch}
\blockpen{subir lápiz}
\blockmove{ir a x: \ovalnum{-120} y: \ovalnum{50}}
\blockpen{bajar lápiz}
\blockmove{ir a x: \ovalnum{120} y: \ovalnum{50}}
\blockmove{ir a x: \ovalnum{0} y: \ovalnum{150}}
\blockmove{ir a x: \ovalnum{-120} y: \ovalnum{50}}
\end{scratch}
\end{tabular}
\end{center}

Los bloques de la derecha se pegan debajo de los de la izquierda. La puerta se hace de la misma manera.

\step{Paso 2: Escríbelo en Colab}
Las líneas vacías dividen el programa en partes: paredes, techo, puerta y ventana. A Python no le molestan las líneas vacías, y hacen que sea más fácil de leer.

\begin{typethis}
from ColabTurtlePlus.Turtle import *
clearscreen()
initializeTurtle((480, 360))
speed(8)
pensize(3)
\end{typethis}
\begin{tcbraster}[raster columns=2, raster column skip=10pt, raster valign=top]
\begin{tcblisting}{listing only, enhanced, colback=codebg, colframe=codeframe, boxrule=0.8pt, arc=6pt, left=22pt, right=8pt, top=4pt, bottom=4pt, title={\textbf{\ldots y sigue en la misma celda}}, coltitle=white, colbacktitle=ink, fonttitle=\small, listing options={style=py, firstnumber=7}}
penup()
goto(-100, -150)
pendown()
fillcolor("tan")
begin_fill()
goto(100, -150)
goto(100, 50)
goto(-100, 50)
goto(-100, -150)
end_fill()

penup()
goto(-120, 50)
pendown()
fillcolor("firebrick")
begin_fill()
goto(120, 50)
goto(0, 150)
goto(-120, 50)
end_fill()
\end{tcblisting}
\begin{tcblisting}{listing only, enhanced, colback=codebg, colframe=codeframe, boxrule=0.8pt, arc=6pt, left=22pt, right=8pt, top=4pt, bottom=4pt, title={\textbf{\ldots y sigue en la misma celda}}, coltitle=white, colbacktitle=ink, fonttitle=\small, listing options={style=py, firstnumber=28}}
penup()
goto(-25, -150)
pendown()
fillcolor("saddlebrown")
begin_fill()
goto(-25, -70)
goto(25, -70)
goto(25, -150)
end_fill()

penup()
goto(55, -10)
dot(40, "skyblue")
hideturtle()
\end{tcblisting}
\end{tcbraster}

\step{Paso 3: Léelo línea por línea}

\begin{longtable}{>{\centering\arraybackslash}p{1.1cm} P{5.3cm} P{9.1cm}}
\toprule
\textbf{Línea} & \textbf{Bloque de Scratch} & \textbf{Qué significa} \\
\midrule
1--3 & \chip[events]{al presionar bandera verde} & Las tres líneas de inicio. \\
4 & (ninguno) & \py|speed(8)| hace que la tortuga sea más rápida, para que no esperes demasiado. \\
5 & \chip[pen]{fijar tamaño de lápiz a 3} & Líneas más gruesas para el contorno. \\
7--9 & \chip[pen]{subir lápiz} \chip[motion]{ir a x: -100 y: -150} \chip[pen]{bajar lápiz} & El truco de «moverse sin dejar línea». La tortuga salta a la esquina de abajo a la izquierda de las paredes. \\
10--16 & \chip[motion]{ir a x: 100 y: -150} y las otras esquinas & Elige el color de pintura y di \py|begin_fill()|: «empieza a recordar esta figura». Recorre las cuatro esquinas y vuelve al inicio. Luego \py|end_fill()| pinta el interior de color \py|"tan"| (café claro). ¡Scratch no puede rellenar! \\
18--26 & el techo & Otra vez el mismo patrón: salta, elige un color, dibuja tres esquinas, rellena. ¡Un triángulo solo necesita tres esquinas! \\
28--36 & la puerta & El mismo patrón por tercera vez. La puerta empieza en la parte de abajo de las paredes. \\
38--41 & \chip[looks]{esconder} & Salta al lugar de la ventana y deja caer un punto azul grande. Luego \py|hideturtle()| esconde a la tortuga para que puedas admirar tu casa. \\
\bottomrule
\end{longtable}

\step{Paso 4: Ejecútalo}
Deberías ver una casa café claro con techo rojo, una puerta café y una ventana azul redonda.

\begin{watchout}[Si una línea cruza tu casa]
Seguramente se te olvidó un \py|penup()|, y la tortuga dibujó una línea al moverse. Revisa las líneas 7, 18, 28 y 38.
\end{watchout}

\step{Paso 5: Hazlo tuyo}

\enlargethispage{\baselineskip}
\begin{challenge}\small
\tick Pinta tu casa con tus propios colores.\par
\tick Agrega una segunda ventana a la izquierda. Copia las líneas 38 a 40 al final y cambia el número x.\par
\tick Agrega un sol en la esquina de arriba: salta a \py|goto(-190, 140)| y usa un \mbox{\texttt{dot(60, "gold")}} grande.\par
\tick Escribe el número de tu casa en la puerta. Salta a la puerta y luego usa \py|write("12")|.
\end{challenge}

% =====================================================================
\mission{Cara de robot}{\texttt{circle}, \texttt{dot}, \texttt{write} y \texttt{stamp}: los superpoderes de Python}

¡Te presento a Robo! Vas a dibujar la cabeza de un robot con una cara redonda, ojos grandes, una boca y una antena. Después Robo tendrá una tortuguita de mascota.

\step{Paso 1: Constrúyelo en Scratch}
Scratch no tiene un bloque de círculo, así que la cabeza y los ojos son difíciles de hacer en Scratch. Esta es la parte que Scratch \textbf{sí} puede hacer: la boca, la antena y la mascota.

\begin{center}
\begin{tabular}{c@{\hspace{1cm}}c}
\begin{scratch}
\blockinit{al presionar \greenflag}
\blockpen{borrar todo}
\blockpen{fijar tamaño de lápiz a \ovalnum{4}}
\blockpen{subir lápiz}
\blockmove{ir a x: \ovalnum{-50} y: \ovalnum{-40}}
\blockpen{bajar lápiz}
\blockpen{fijar color de lápiz a \ovalnum{rojo}}
\blockmove{ir a x: \ovalnum{50} y: \ovalnum{-40}}
\end{scratch}
&
\begin{scratch}
\blockpen{subir lápiz}
\blockmove{ir a x: \ovalnum{0} y: \ovalnum{100}}
\blockpen{bajar lápiz}
\blockpen{fijar color de lápiz a \ovalnum{gris}}
\blockmove{ir a x: \ovalnum{0} y: \ovalnum{140}}
\blockpen{subir lápiz}
\blockmove{ir a x: \ovalnum{-150} y: \ovalnum{-120}}
\blockpen{sellar}
\blocklook{esconder}
\end{scratch}
\end{tabular}
\end{center}

\step{Paso 2: Escríbelo en Colab}

\begin{typethis}
from ColabTurtlePlus.Turtle import *
clearscreen()
initializeTurtle((480, 360))
speed(10)
pensize(4)
pencolor("slategray")
\end{typethis}
\begin{tcbraster}[raster columns=2, raster column skip=10pt, raster valign=top]
\begin{tcblisting}{listing only, enhanced, colback=codebg, colframe=codeframe, boxrule=0.8pt, arc=6pt, left=22pt, right=8pt, top=4pt, bottom=4pt, title={\textbf{\ldots y sigue en la misma celda}}, coltitle=white, colbacktitle=ink, fonttitle=\small, listing options={style=py, firstnumber=8}}
penup()
goto(0, -100)
pendown()
fillcolor("lightgray")
begin_fill()
circle(100)
end_fill()

penup()
goto(-40, 30)
dot(40, "white")
dot(16, "black")
goto(40, 30)
dot(40, "white")
dot(16, "black")
\end{tcblisting}
\begin{tcblisting}{listing only, enhanced, colback=codebg, colframe=codeframe, boxrule=0.8pt, arc=6pt, left=22pt, right=8pt, top=4pt, bottom=4pt, title={\textbf{\ldots y sigue en la misma celda}}, coltitle=white, colbacktitle=ink, fonttitle=\small, listing options={style=py, firstnumber=24}}
goto(-50, -40)
pendown()
pencolor("red")
goto(50, -40)

penup()
goto(0, 100)
pendown()
pencolor("slategray")
goto(0, 140)
dot(20, "gold")

penup()
goto(120, -150)
write("¡Soy Robo!")
shape("turtle")
color("seagreen")
goto(-150, -120)
stamp()
hideturtle()
\end{tcblisting}
\end{tcbraster}

\step{Paso 3: Léelo línea por línea}

\begin{longtable}{>{\centering\arraybackslash}p{1.1cm} P{5.3cm} P{9.1cm}}
\toprule
\textbf{Línea} & \textbf{Bloque de Scratch} & \textbf{Qué significa} \\
\midrule
1--3 & \chip[events]{al presionar bandera verde} & Las tres líneas de inicio. \\
4--6 & \chip[pen]{fijar tamaño de lápiz a 4} \chip[pen]{fijar color de lápiz a} & Dibujo rápido, una línea gruesa y un color gris azulado que se llama \py|"slategray"|. \\
8--10 & \chip[pen]{subir lápiz} \chip[motion]{ir a x: 0 y: -100} \chip[pen]{bajar lápiz} & Salta a la parte de abajo de donde va a estar la cara. \\
11--14 & (Scratch no puede hacer esto) & \py|circle(100)| dibuja la cara. El centro del círculo queda 100 pasos a la izquierda de la tortuga. La tortuga mira a la derecha, así que su izquierda es \textbf{arriba}, y el centro cae en $(0, 0)$. El relleno la pinta de gris claro. \\
16--22 & \chip[pen]{subir lápiz} \chip[motion]{ir a x: -40 y: 30} & Salta al ojo izquierdo. Un punto blanco grande con un punto negro pequeño encima forma un ojo. Luego lo mismo para el ojo derecho. Los puntos funcionan con el lápiz arriba. \\
24--27 & \chip[motion]{ir a x: -50 y: -40} \ldots & La boca: lápiz abajo, rojo, y una línea recta de lado a lado. \\
29--34 & \chip[motion]{ir a x: 0 y: 100} \ldots & La antena: una línea corta hacia arriba desde lo alto de la cabeza, y un punto dorado en la punta. \\
36--38 & (Scratch no puede hacer esto) & Salta debajo de la cara y escribe un mensaje con \py|write|. Las palabras van entre comillas. \\
39--40 & (elegir un disfraz) & Se convierte en forma de tortuga. \py|color("seagreen")| fija el color del lápiz \textbf{y} el color de relleno al mismo tiempo, así que la tortuga se vuelve verde. \\
41--42 & \chip[motion]{ir a x: -150 y: -120} \chip[pen]{sellar} & Camina hasta abajo a la izquierda y sella una copia de la tortuga: ¡la mascota de Robo! \\
43 & \chip[looks]{esconder} & Esconde a la tortuga de verdad. El sello se queda. \\
\bottomrule
\end{longtable}

\step{Paso 4: Ejecútalo}
Deberías ver una cara de robot redonda y gris con dos ojos grandes, una boca roja recta y una antena con una bolita dorada en la punta. «¡Soy Robo!» está escrito abajo a la derecha, y un sello de tortuga verde está abajo a la izquierda.


\step{Paso 5: Hazlo tuyo}

\begin{challenge}
\tick Dale a Robo ojos de otros colores, como \py|"lime"| (verde lima) o \py|"hotpink"| (rosa fuerte).\par
\tick Cambia el mensaje por el nombre de tu propio robot.\par
\tick Haz una sonrisa: cambia la línea 27 por \py|goto(0, -60)| y agrega \py|goto(50, -40)| después. Ahora la boca tiene forma de V.\par
\tick Sella toda una familia de mascotas. Después de la línea 42, agrega más parejas de \py|goto(...)| y \py|stamp()|. Prueba cambiar \py|color(...)| entre un sello y otro.
\end{challenge}

% =====================================================================
\needspace{20\baselineskip}
\subsection*{Tradúcelo tú}

Aquí hay un programa de Scratch que no has visto. Escribe su Python en las líneas de abajo. Empieza con las tres líneas de inicio. Luego escríbelo en Colab para comprobarlo. Pistas: recuerda la regla secreta para \chip[motion]{apuntar en dirección}, y recuerda que el lápiz de la tortuga empieza abajo. (La respuesta está al final del libro.)

\begin{tcbraster}[raster columns=2, raster equal height=rows, raster column skip=12pt, raster valign=center]
\begin{tcolorbox}[colback=white, colframe=white, boxrule=0pt, valign=center]
\begin{scratch}
\blockinit{al presionar \greenflag}
\blockpen{subir lápiz}
\blockmove{ir a x: \ovalnum{-100} y: \ovalnum{0}}
\blockpen{bajar lápiz}
\blockpen{fijar color de lápiz a \ovalnum{rojo}}
\blockpen{fijar tamaño de lápiz a \ovalnum{4}}
\blockmove{apuntar en dirección \ovalnum{0}}
\blockmove{mover \ovalnum{80} pasos}
\blockmove{girar \turnright{} \ovalnum{90} grados}
\blockmove{mover \ovalnum{80} pasos}
\blockpen{sellar}
\end{scratch}
\end{tcolorbox}
\begin{tcolorbox}[colback=white, colframe=codeframe, arc=5pt, title=\textbf{Mi Python}, coltitle=ink, colbacktitle=codebg]
\vspace{4pt}
\foreach \n in {1,...,13}{\rule{\linewidth}{0.4pt}\par\vspace{9pt}}
\end{tcolorbox}
\end{tcbraster}


% =====================================================================
\needspace{12\baselineskip}
\begin{puzzle}
La tortuga está en el centro, mirando a la derecha, como siempre. Haz que apunte \textbf{directo hacia abajo} y camine \textbf{100 pasos}. Solo puedes usar \textbf{dos líneas} de Python después de las tres líneas de inicio. (Pista: en Scratch, la dirección hacia abajo es 180.)

\vspace{6pt}
\rule{\linewidth}{0.4pt}\par\vspace{12pt}
\rule{\linewidth}{0.4pt}
\end{puzzle}

% =====================================================================
\needspace{14\baselineskip}
\subsection*{Lista de la parte 1}

Marca cada una que ya sabes hacer.

\tick Ya sé mover a la tortuga con \py|forward| y \py|backward|, y girarla con \py|right| y \py|left|.\par
\tick Ya sé mandar a la tortuga a un lugar con \py|goto(x, y)|, y elegir qué tan rápido va con \py|speed|.\par
\tick Ya sé convertir una dirección de Scratch a Python con \py|setheading(90 - d)|.\par
\tick Ya sé moverme sin dibujar usando \py|penup()| y \py|pendown()|.\par
\tick Ya sé elegir el color y el tamaño del lápiz, borrar el dibujo con \py|clear()| y sellar la tortuga con \py|stamp()|.\par
\tick Puedo usar un superpoder de Python: \py|circle|, \py|dot|, \py|begin_fill| o \py|write|.

% =====================================================================
% PARTE 2: Bucles
% =====================================================================
\partopener{control}{PARTE 2}{Bucles}{Haz que la computadora repita las cosas por ti.}

En la Parte 1 dibujaste figuras orden por orden. Para dibujar un cuadrado escribiste \py|forward(100)| y \py|right(90)| cuatro veces cada una. ¡Eso es mucho escribir! En Scratch nunca harías eso. Tomarías un bloque \chip[control]{repetir 4} y dejarías que él hiciera el trabajo.

En esta parte vas a aprender el Python de los bloques naranjas de \textbf{Control} que repiten y esperan:

\begin{center}
\chip[control]{repetir 10}\quad \chip[control]{por siempre}\quad \chip[control]{esperar 1 segundos}\quad \chip[control]{detener este programa}
\end{center}

Al final vas a dibujar estrellas, espirales y flores con solo unas pocas líneas de código.

% ---------------------------------------------------------------------
\needspace{14\baselineskip}
\subsection*{Repetir se convierte en \texttt{for}}

Aquí está el cuadrado de la Parte 1, pero con un bucle. Mira los dos lados y une cada bloque con su línea. (En esta parte, los dibujos lado a lado no muestran las tres líneas de inicio para ahorrar espacio. Cuando escribas un programa de dibujo en Colab, empieza siempre con ellas).

\begin{tcbraster}[raster columns=2, raster equal height=rows, raster column skip=12pt, raster valign=center]
\begin{tcolorbox}[colback=white, colframe=white, boxrule=0pt, valign=center]
\begin{scratch}[scale=0.85]
\blockinit{al presionar \greenflag}
\blockrepeat{repetir \ovalnum{4}}{
  \blockmove{mover \ovalnum{100} pasos}
  \blockmove{girar \turnright{} \ovalnum{90} grados}
}
\end{scratch}
\end{tcolorbox}
\begin{code}
for i in range(4):
    forward(100)
    right(90)
\end{code}
\end{tcbraster}

\textbf{Cómo funciona.} La línea 1, \py|for i in range(4):|, es el bloque \chip[control]{repetir 4}. Puedes leerla como «por cada vuelta, 4 veces». La parte \py|range(4)| quiere decir «cuenta 4 veces». Las líneas 2 y 3 son los bloques que están dentro del bucle, así que pasan 4 veces.

Python habla inglés: sus palabras (\py|for|, \py|print|, \py|forward|...) siempre van en inglés. Lo que tú inventas, como los nombres y los textos entre comillas, sí puede ir en español.

% ---------------------------------------------------------------------
\needspace{14\baselineskip}
\subsection*{La boca de la C: dos puntos y cuatro espacios}

En Scratch, un bloque va \textbf{dentro} de un bucle cuando lo sueltas en la boca con forma de C. Python no tiene forma de C, así que usa dos cosas en su lugar:

\begin{enumerate}
  \item Los \textbf{dos puntos} \texttt{:} al final de la línea del \py|for|. Los dos puntos \textbf{abren la boca}. Dicen «lo de adentro empieza en la siguiente línea».
  \item \textbf{Cuatro espacios} al principio de una línea. Esas líneas están \textbf{dentro} de la boca. A esos espacios se les llama \textbf{sangría}.
\end{enumerate}

La primera línea que vuelve al borde izquierdo, sin espacios, es el bloque pegado \textbf{debajo} de la C. Se ejecuta una sola vez, cuando el bucle ya terminó.

\begin{tcbraster}[raster columns=2, raster equal height=rows, raster column skip=12pt, raster valign=center]
\begin{tcolorbox}[colback=white, colframe=white, boxrule=0pt, valign=center]
\begin{scratch}[scale=0.85]
\blockrepeat{repetir \ovalnum{4}}{
  \blockmove{mover \ovalnum{100} pasos}
  \blockmove{girar \turnright{} \ovalnum{90} grados}
}
\blocklook{decir \ovalnum{¡Cuadrado listo!}}
\end{scratch}
\end{tcolorbox}
\begin{spaces}
for i in range(4):
    forward(100)
    right(90)
print("¡Cuadrado listo!")
\end{spaces}
\end{tcbraster}

Las marquitas \texttt{\textvisiblespace} muestran los espacios, que normalmente son invisibles. Las líneas 2 y 3 tienen cuatro espacios, así que se repiten 4 veces. La línea 4 no tiene espacios, así que está fuera del bucle. Ves \texttt{¡Cuadrado listo!} una sola vez, cuando el cuadrado ya está terminado.

\begin{think}[Colab te ayuda]
Cuando escribes los dos puntos y presionas \textbf{Enter}, Colab pone los cuatro espacios por ti. Para salir del bucle, presiona \textbf{Retroceso} (\textit{Backspace}) una vez al principio de la línea.
\end{think}

% ---------------------------------------------------------------------
\needspace{12\baselineskip}
\subsection*{¿Qué cuenta la \texttt{i}?}

La \texttt{i} de \py|for i in range(4):| es un \textbf{contador}. Cada vez que el bucle da una vuelta, Python cambia la \texttt{i} él solito. El bloque \chip[control]{repetir} de Scratch no tiene un contador que puedas ver. En Scratch tendrías que hacerte uno tú.

Vamos a espiar la \texttt{i}. Este programa no dibuja, así que no necesita las tres líneas de inicio.

\begin{tcbraster}[raster columns=2, raster equal height=rows, raster column skip=12pt, raster valign=center]
\begin{code}
for i in range(4):
    print(i)
\end{code}
\begin{outputbox}
0\\
1\\
2\\
3
\end{outputbox}
\end{tcbraster}

Python empieza a contar en \textbf{0}, así que \texttt{i} vale 0, luego 1, luego 2, luego 3. ¡Siguen siendo 4 vueltas! Fíjate en que \py|range(4)| se detiene \textbf{antes} de llegar a 4. No tienes que usar la \texttt{i}, pero siempre está ahí si la quieres.

% ---------------------------------------------------------------------
\needspace{12\baselineskip}
\subsection*{Usar la \texttt{i} dentro del bucle}

Como la \texttt{i} cambia en cada vuelta, puedes usarla para que cada vuelta del bucle sea un poquito distinta. La estrella \texttt{*} quiere decir \textbf{por} (multiplicar) en Python (en el teclado no hay tecla $\times$).

\begin{tcbraster}[raster columns=2, raster equal height=rows, raster column skip=12pt, raster valign=center]
\begin{code}
for i in range(5):
    print(i * 10)
\end{code}
\begin{outputbox}
0\\
10\\
20\\
30\\
40
\end{outputbox}
\end{tcbraster}

Ahora imagina \py|forward(i * 5)| en lugar de \py|print(i * 10)|. La tortuga caminaría un poco más lejos en cada vuelta. ¡Así vas a dibujar una espiral en una misión muy pronto! Un nombre como \texttt{i} que guarda un número se llama \textbf{variable}. Vas a crear tus propias variables en la Parte 3.

% ---------------------------------------------------------------------
\needspace{14\baselineskip}
\subsection*{Cuando los espacios fallan}

Python es muy exigente con los dos puntos y los espacios. Aquí están los dos mensajes rojos que vas a ver más seguido. Léelos como la nota con pistas de un detective: la última línea te dice qué está mal. Python siempre escribe sus errores en inglés, aunque tu Colab esté en español. La palabra importante es el nombre del error (\texttt{SyntaxError}, \texttt{IndentationError}...).

\textbf{¿Olvidaste los dos puntos?}

\begin{codenonum}
for i in range(3)
    print("hola")
\end{codenonum}
\begin{errorbox}
SyntaxError: expected ':'
\end{errorbox}

Python quiere los dos puntos que abren la boca (\texttt{expected ':'} quiere decir «esperaba \texttt{:}»). Agrega \texttt{:} al final de la línea del \py|for|.

\needspace{10\baselineskip}
\textbf{¿Olvidaste los espacios?}

\begin{codenonum}
for i in range(3):
print("hola")
\end{codenonum}
\begin{errorbox}
IndentationError: expected an indented block after 'for' statement on line 1
\end{errorbox}

Esto quiere decir «la C tiene la boca vacía». No hay nada dentro del bucle. Pon cuatro espacios delante de \py|print("hola")|.

\needspace{16\baselineskip}
\begin{watchout}[¡Cuidado! El bicho tramposo sin letras rojas]
¿Y si solo \textbf{una} línea pierde sus espacios?
\begin{codenonum}
for i in range(4):
    forward(100)
right(90)
\end{codenonum}
Python está contento, así que \textbf{no hay error}. Pero ahora \py|right(90)| está fuera del bucle, como un bloque soltado \textbf{debajo} de la C en lugar de adentro. La tortuga camina hacia adelante 4 veces seguidas (una línea larguísima que se sale del escenario) y gira una sola vez, al final de todo. Si tu dibujo se ve mal, ¡revisa primero los espacios!
\end{watchout}

% ---------------------------------------------------------------------
\needspace{14\baselineskip}
\subsection*{Un bucle dentro de otro bucle}

En Scratch puedes poner un bloque \chip[control]{repetir} dentro de otro bloque \chip[control]{repetir}. Python puede hacer lo mismo. El bucle de adentro lleva \textbf{más} espacios.

\begin{tcbraster}[raster columns=2, raster equal height=rows, raster column skip=12pt, raster valign=center]
\begin{tcolorbox}[colback=white, colframe=white, boxrule=0pt, valign=center]
\begin{scratch}[scale=0.85]
\blockrepeat{repetir \ovalnum{6}}{
  \blockrepeat{repetir \ovalnum{3}}{
    \blockmove{mover \ovalnum{80} pasos}
    \blockmove{girar \turnright{} \ovalnum{120} grados}
  }
  \blockmove{girar \turnright{} \ovalnum{60} grados}
}
\end{scratch}
\end{tcolorbox}
\begin{spaces}
for i in range(6):
    for j in range(3):
        forward(80)
        right(120)
    right(60)
\end{spaces}
\end{tcbraster}

\textbf{Cómo funciona.} Cuenta los espacios al principio de cada línea:
\begin{itemize}
  \item \textbf{0 espacios} (línea 1): el bucle de afuera. Da 6 vueltas.
  \item \textbf{4 espacios} (líneas 2 y 5): dentro del bucle de afuera.
  \item \textbf{8 espacios} (líneas 3 y 4): dentro de \textbf{los dos} bucles. Dibujan un triángulo.
\end{itemize}
La línea 5 tiene 4 espacios, así que no es parte del triángulo. Gira un poquito la tortuga después de cada triángulo. El bucle de adentro usa otra letra para su contador, \texttt{j}, para que los dos contadores no se revuelvan. El resultado son 6 triángulos girando alrededor del centro, 18 líneas en total.

% ---------------------------------------------------------------------
\needspace{14\baselineskip}
\subsection*{Por siempre se convierte en \texttt{while True}}

El bloque \chip[control]{por siempre} repite sin fin. En Python es \py|while True:|. Léelo como «sigue mientras verdadero sea verdadero» (\texttt{while} es «mientras» y \texttt{True} es «verdadero»), ¡y verdadero siempre es verdadero!

\begin{tcbraster}[raster columns=2, raster equal height=rows, raster column skip=12pt, raster valign=center]
\begin{tcolorbox}[colback=white, colframe=white, boxrule=0pt, valign=center]
\begin{scratch}[scale=0.85]
\blockinit{al presionar \greenflag}
\blockinfloop{por siempre}{
  \blockmove{mover \ovalnum{100} pasos}
  \blockmove{girar \turnright{} \ovalnum{90} grados}
}
\end{scratch}
\end{tcolorbox}
\begin{code}
while True:
    forward(100)
    right(90)
\end{code}
\end{tcbraster}

\textbf{Cómo funciona.} Los dos puntos y los cuatro espacios funcionan igual que en un bucle \py|for|. La diferencia es que este bucle \textbf{nunca se detiene solo}. La tortuga dibuja el mismo cuadrado una y otra vez, y la celda sigue ejecutándose. Mientras una celda se está ejecutando, su botón de reproducir se convierte en un círculo que gira. Haz clic en él (\stopbtn) para detener la celda, igual que la señal roja de alto en Scratch. También puedes usar el menú: \textbf{Entorno de ejecución} (\textit{Runtime}) y luego \textbf{Interrumpir ejecución} (\textit{Interrupt execution}).

\begin{watchout}[¡Cuidado! Por siempre quiere decir por siempre]
Cuando detienes la celda, Colab muestra unas letras rojas que terminan con \texttt{KeyboardInterrupt}. Eso no es un error tuyo. Solo quiere decir «presionaste detener». Mientras un bucle por siempre se está ejecutando, ninguna otra celda puede ejecutarse, así que acuérdate de detenerlo. Para los dibujos, casi siempre es mejor un bucle \py|for| con un número: termina solo, y de todos modos ves el dibujo completo.
\end{watchout}

% ---------------------------------------------------------------------
\needspace{14\baselineskip}
\subsection*{Esperar se convierte en \texttt{time.sleep}}

El bloque \chip[control]{esperar 1 segundos} pausa el programa. Python guarda su reloj en una caja de herramientas llamada \py|time| («tiempo»). Abres la caja de herramientas una vez, arriba de todo, con \py|import time|. Después, \py|time.sleep(2)| quiere decir «duerme 2 segundos».

\begin{tcbraster}[raster columns=2, raster equal height=rows, raster column skip=12pt, raster valign=center]
\begin{tcolorbox}[colback=white, colframe=white, boxrule=0pt, valign=center]
\begin{scratch}[scale=0.85]
\blockinit{al presionar \greenflag}
\blocklook{decir \ovalnum{En sus marcas...}}
\blockcontrol{esperar \ovalnum{2} segundos}
\blocklook{decir \ovalnum{¡Fuera!}}
\end{scratch}
\end{tcolorbox}
\begin{code}
import time
print("En sus marcas...")
time.sleep(2)
print("¡Fuera!")
\end{code}
\end{tcbraster}

El punto en \py|time.sleep| quiere decir «la herramienta \texttt{sleep} (dormir) de la caja \texttt{time}». Si olvidas la línea 1, Python da un \texttt{NameError}, porque todavía no sabe qué es \texttt{time}.

\needspace{12\baselineskip}
\textbf{Cohete con cuenta regresiva.} Aquí un bucle y una espera trabajan juntos. En cada vuelta, \py|5 - i| calcula el siguiente número: $5-0=5$, luego $5-1=4$, y así sucesivamente.

\begin{tcbraster}[raster columns=2, raster equal height=rows, raster column skip=12pt, raster valign=center]
\begin{code}
import time
for i in range(5):
    print(5 - i)
    time.sleep(1)
print("¡Despegue!")
\end{code}
\begin{outputbox}
5\\
4\\
3\\
2\\
1\\
¡Despegue!
\end{outputbox}
\end{tcbraster}

\begin{think}[Adivina primero]
¿Qué pasaría si \py|print("¡Despegue!")| tuviera cuatro espacios delante? Di lo que crees y luego pruébalo. (¡Acuérdate del bicho tramposo!)
\end{think}

% ---------------------------------------------------------------------
\needspace{14\baselineskip}
\subsection*{Detener este programa se convierte en \texttt{break}}

El bloque \chip[control]{detener este programa} termina el programa. Dentro de un bucle de Python, la palabra \py|break| («romper») hace el mismo trabajo: \textbf{rompe} el bucle y se sale de inmediato.

\begin{tcbraster}[raster columns=2, raster equal height=rows, raster column skip=12pt, raster valign=center]
\begin{tcolorbox}[colback=white, colframe=white, boxrule=0pt, valign=center]
\begin{scratch}[scale=0.85]
\blockinit{al presionar \greenflag}
\blockinfloop{por siempre}{
  \blocklook{decir \ovalnum{¡Hola!}}
  \blockstop{detener \selectmenu{este programa}}
}
\end{scratch}
\end{tcolorbox}
\begin{code}
while True:
    print("¡Hola!")
    break
print("Fuera del bucle")
\end{code}
\end{tcbraster}

Esto escribe \texttt{¡Hola!} una sola vez, y luego \texttt{Fuera del bucle}. Aquí hay algo que Scratch no puede hacer: en Scratch no se puede pegar nada debajo de un \chip[control]{por siempre}, pero en Python la línea que va después del bucle se ejecuta en cuanto \py|break| se sale. Por ahora \py|break| parece un poco tonto, porque siempre detiene el bucle en la primera vuelta. En la Parte 5 vas a aprender \py|if|, y entonces \py|break| se vuelve muy útil: «si adivinaste, detente».

\begin{think}[¿Y repetir hasta que?]
Scratch también tiene \chip[control]{repetir hasta que}. En Python se convierte en \py|while not ...:|, donde los puntos son una pregunta como «¿ya llegó el gato al borde?». Necesitas unas cuantas herramientas más para hacer esas preguntas, así que \chip[control]{repetir hasta que} espera hasta la Parte 5.
\end{think}

% =====================================================================
\mission{Dibuja figuras por el camino corto}{bucles \texttt{for}, los dos puntos y los cuatro espacios, y el truco del 360}

Ahora puedes dibujar cualquier figura con solo tres líneas de bucle. Empecemos con un triángulo. (En Scratch, primero agrega la extensión \textbf{Lápiz} con el botón \textbf{Agregar Extensión}).

\begin{tcbraster}[raster columns=2, raster column skip=12pt, raster valign=top, raster force size=false,
  raster column 1/.style={add to width=-2.0cm}, raster column 2/.style={add to width=2.0cm}]
\begin{tcolorbox}[blankest]{\normalsize\bfseries\color{ink}Paso 1: Constrúyelo en Scratch}\end{tcolorbox}
\begin{tcolorbox}[blankest]{\normalsize\bfseries\color{ink}Paso 2: Escríbelo en Colab}\end{tcolorbox}
\begin{tcolorbox}[blankest]
\begin{scratch}[scale=0.75]
\blockinit{al presionar \greenflag}
\blockpen{fijar color de lápiz a \ovalnum{azul}}
\blockpen{fijar tamaño de lápiz a \ovalnum{3}}
\blockrepeat{repetir \ovalnum{3}}{
  \blockmove{mover \ovalnum{120} pasos}
  \blockmove{girar \turnright{} \ovalnum{120} grados}
}
\end{scratch}
\end{tcolorbox}
\begin{typethis}
from ColabTurtlePlus.Turtle import *
clearscreen()
initializeTurtle((480, 360))
pencolor("blue")
pensize(3)
for i in range(3):
    forward(120)
    right(120)
\end{typethis}
\end{tcbraster}

\step{Paso 3: Léelo línea por línea}

\begin{longtable}{>{\centering\arraybackslash}p{1.1cm} P{5.3cm} P{9.1cm}}
\toprule
\textbf{Línea} & \textbf{Bloque de Scratch} & \textbf{Qué significa} \\
\midrule
1--3 & \chip[events]{al presionar bandera verde} & Las tres líneas de inicio. Preparan el escenario y la tortuga. El lápiz empieza abajo. \\
4 & \chip[pen]{fijar color de lápiz a azul} & El color es una palabra entre comillas, en minúsculas. Los colores se escriben en inglés: \py|"blue"| es azul, \py|"red"| rojo, \py|"gold"| dorado... \\
5 & \chip[pen]{fijar tamaño de lápiz a 3} & Una línea más gruesa. \\
6 & \chip[control]{repetir 3} & \py|for i in range(3):| quiere decir «haz esto 3 veces». Los dos puntos abren la boca. \\
7 & \chip[motion]{mover 120 pasos} & Lleva cuatro espacios delante, así que está dentro del bucle. \\
8 & \chip[motion]{girar a la derecha 120 grados} & También lleva cuatro espacios, así que también está dentro del bucle. \\
\bottomrule
\end{longtable}

\step{Paso 4: Ejecútalo}
Presiona \playbtn. La tortuga dibuja un triángulo azul y se detiene otra vez donde empezó, apuntando hacia el mismo lado que al principio.

\needspace{14\baselineskip}
\step{Paso 5: Hazlo tuyo}

Aquí va un truco: \textbf{cuando la tortuga da toda la vuelta a una figura, sus giros suman 360 grados}, un círculo completo. Así que, para encontrar el giro, divide 360 entre el número de lados. Llena la tabla y luego dibuja cada figura cambiando los números de las líneas 6, 7 y 8. Las figuras con más lados necesitan un avance más corto, o se salen del escenario.

\begin{center}
\renewcommand{\arraystretch}{1.3}
\begin{tabular}{P{3.2cm} P{2.6cm} P{2.0cm} P{3.0cm} P{2.3cm}}
\toprule
\textbf{Figura} & \textbf{Lados (range)} & \textbf{Avance} & \textbf{Giro} & \textbf{¿Dibujada?} \\
\midrule
Triángulo & 3 & 120 & $360 \div 3 = 120$ & $\square$ \\
Cuadrado & 4 & 120 & $360 \div 4 = 90$ & $\square$ \\
Pentágono & 5 & 100 & \rule{2.2cm}{0.5pt} & $\square$ \\
Hexágono & 6 & 80 & \rule{2.2cm}{0.5pt} & $\square$ \\
Octágono & 8 & 60 & \rule{2.2cm}{0.5pt} & $\square$ \\
Decágono & 10 & 50 & \rule{2.2cm}{0.5pt} & $\square$ \\
Casi un círculo & 36 & 10 & \rule{2.2cm}{0.5pt} & $\square$ \\
\bottomrule
\end{tabular}
\end{center}

\begin{challenge}
\tick Llena la tabla y dibuja todas las figuras. ¿Cuál se parece más a un círculo?\par
\tick Agrega \py|speed(13)| justo después de la línea 3 para que la tortuga dibuje más rápido.\par
\tick Pon \py|pencolor("red")| \textbf{dentro} del bucle, con cuatro espacios, justo antes de \py|forward(120)|. ¿Cambia algo? ¿Por qué?\par
\tick Agrega \py|print("¡Tarán!")| al final de todo, \textbf{sin} espacios. ¿Cuántas veces se escribe? Ahora ponle cuatro espacios y ejecútalo otra vez.
\end{challenge}


% =====================================================================
\mission{Estrellas}{a dar más de una vuelta, y a moverte antes de dibujar}

Una estrella es una figura que gira \textbf{más} que un pentágono. Cada punta es un giro muy cerrado de 144 grados.

\begin{tcbraster}[raster columns=2, raster column skip=12pt, raster valign=top, raster force size=false,
  raster column 1/.style={add to width=-2.0cm}, raster column 2/.style={add to width=2.0cm}]
\begin{tcolorbox}[blankest]{\normalsize\bfseries\color{ink}Paso 1: Constrúyelo en Scratch}\end{tcolorbox}
\begin{tcolorbox}[blankest]{\normalsize\bfseries\color{ink}Paso 2: Escríbelo en Colab}\end{tcolorbox}
\begin{tcolorbox}[blankest]
\begin{scratch}[scale=0.75]
\blockinit{al presionar \greenflag}
\blockpen{fijar color de lápiz a \ovalnum{dorado}}
\blockpen{fijar tamaño de lápiz a \ovalnum{4}}
\blockpen{subir lápiz}
\blockmove{ir a x: \ovalnum{-75} y: \ovalnum{30}}
\blockpen{bajar lápiz}
\blockrepeat{repetir \ovalnum{5}}{
  \blockmove{mover \ovalnum{150} pasos}
  \blockmove{girar \turnright{} \ovalnum{144} grados}
}
\end{scratch}
\end{tcolorbox}
\begin{typethis}
from ColabTurtlePlus.Turtle import *
clearscreen()
initializeTurtle((480, 360))
pencolor("gold")
pensize(4)
penup()
goto(-75, 30)
pendown()
for i in range(5):
    forward(150)
    right(144)
\end{typethis}
\end{tcbraster}

\step{Paso 3: Léelo línea por línea}

\begin{longtable}{>{\centering\arraybackslash}p{1.1cm} P{5.3cm} P{9.1cm}}
\toprule
\textbf{Línea} & \textbf{Bloque de Scratch} & \textbf{Qué significa} \\
\midrule
1--3 & \chip[events]{al presionar bandera verde} & Las tres líneas de inicio. \\
4--5 & \chip[pen]{fijar color de lápiz a dorado} \chip[pen]{fijar tamaño de lápiz a 4} & Un lápiz dorado y grueso (\py|"gold"| es dorado). \\
6 & \chip[pen]{subir lápiz} & Levanta el lápiz para que el siguiente movimiento no deje línea. \\
7 & \chip[motion]{ir a x: -75 y: 30} & Salta un poco a la izquierda, para que la estrella quede en el centro del escenario. \\
8 & \chip[pen]{bajar lápiz} & Listo para dibujar otra vez. \\
9 & \chip[control]{repetir 5} & 5 puntas, así que 5 vueltas. \\
10--11 & \chip[motion]{mover 150 pasos} \chip[motion]{girar a la derecha 144 grados} & Dentro del bucle: una línea larga y un giro cerrado por cada punta. \\
\bottomrule
\end{longtable}

\step{Paso 4: Ejecútalo}
Presiona \playbtn. Aparece una gran estrella dorada en el centro del escenario, y la tortuga termina otra vez en el punto de partida.

\begin{think}[¡Dos vueltas!]
Un pentágono gira $5 \times 72 = 360$ grados: una vuelta. La estrella gira $5 \times 144 = 720$ grados. Eso es $360 + 360$, así que la tortuga da \textbf{dos} vueltas completas mientras dibuja la estrella. Por eso sus líneas se cruzan en el centro.
\end{think}

\needspace{12\baselineskip}
\step{Paso 5: Hazlo tuyo}

Para una estrella con un número \textbf{impar} de puntas (5, 7, 9, \ldots), aquí tienes una receta para el giro:
\[ \text{giro} = 180 - (180 \div \text{puntas}) \]
Compruébala con 5 puntas: $180 \div 5 = 36$, y $180 - 36 = 144$. ¡Funciona! Llena la tabla y dibuja cada estrella. Usa también el número de puntas dentro de \py|range( )|.

\begin{center}
\renewcommand{\arraystretch}{1.3}
\begin{tabular}{P{2.6cm} P{3.8cm} P{3.4cm} P{2.3cm}}
\toprule
\textbf{Puntas} & \textbf{$180 \div$ puntas} & \textbf{Giro} & \textbf{¿Dibujada?} \\
\midrule
5 & 36 & 144 & $\square$ \\
9 & \rule{2.2cm}{0.5pt} & \rule{2.2cm}{0.5pt} & $\square$ \\
15 & \rule{2.2cm}{0.5pt} & \rule{2.2cm}{0.5pt} & $\square$ \\
45 & \rule{2.2cm}{0.5pt} & \rule{2.2cm}{0.5pt} & $\square$ \\
\bottomrule
\end{tabular}
\end{center}

\begin{challenge}
\tick Llena la tabla de estrellas. ¡La estrella de 45 puntas parece un sol peludo! Para esa, agrega \py|speed(13)| después de la línea 3.\par
\tick Las \textbf{estrellas pares} necesitan otro giro. Prueba \py|range(8)| con \py|right(135)|, y \py|range(12)| con \py|right(150)|.\par
\tick Cambia el color y el tamaño del lápiz. Prueba \py|"skyblue"| (azul cielo) o \py|"hotpink"| (rosa fuerte).\par
\tick Cambia la línea 7 para poner la estrella en una esquina del escenario. ¿Puedes hacer que quede toda dentro del escenario?
\end{challenge}

% =====================================================================
\mission{Espiral cuadrada}{a usar el contador \texttt{i} para cambiar cada paso}

En esta misión cada línea es un poco más larga que la anterior, así que la tortuga hace una espiral hacia afuera. En Scratch necesitas una variable para esto. ¡En Python el contador \texttt{i} hace el trabajo gratis!

En Scratch, primero crea una variable llamada \textbf{medida} con \textbf{Crear una variable}.

\begin{tcbraster}[raster columns=2, raster column skip=12pt, raster valign=top, raster force size=false,
  raster column 1/.style={add to width=-2.0cm}, raster column 2/.style={add to width=2.0cm}]
\begin{tcolorbox}[blankest]{\normalsize\bfseries\color{ink}Paso 1: Constrúyelo en Scratch}\end{tcolorbox}
\begin{tcolorbox}[blankest]{\normalsize\bfseries\color{ink}Paso 2: Escríbelo en Colab}\end{tcolorbox}
\begin{tcolorbox}[blankest]
\begin{scratch}[scale=0.75]
\blockinit{al presionar \greenflag}
\blockpen{fijar color de lápiz a \ovalnum{morado}}
\blockvariable{fijar \selectmenu{medida} a \ovalnum{0}}
\blockrepeat{repetir \ovalnum{40}}{
  \blockmove{mover \ovalvariable{medida} pasos}
  \blockmove{girar \turnright{} \ovalnum{90} grados}
  \blockvariable{cambiar \selectmenu{medida} en \ovalnum{5}}
}
\end{scratch}
\end{tcolorbox}
\begin{typethis}
from ColabTurtlePlus.Turtle import *
clearscreen()
initializeTurtle((480, 360))
pencolor("purple")
speed(13)
for i in range(40):
    forward(i * 5)
    right(90)
\end{typethis}
\end{tcbraster}

\step{Paso 3: Léelo línea por línea}

\begin{longtable}{>{\centering\arraybackslash}p{1.1cm} P{5.3cm} P{9.1cm}}
\toprule
\textbf{Línea} & \textbf{Bloque de Scratch} & \textbf{Qué significa} \\
\midrule
1--3 & \chip[events]{al presionar bandera verde} & Las tres líneas de inicio. \\
4 & \chip[pen]{fijar color de lápiz a morado} & Un lápiz morado (\py|"purple"|). \\
5 & (ninguno) & Dibujo lo más rápido posible, para que no esperes mucho por 40 líneas. \\
--- & \chip[variables]{fijar medida a 0} & ¡No hace falta ninguna línea! El contador \texttt{i} empieza en 0 él solito. \\
6 & \chip[control]{repetir 40} & 40 vueltas. \texttt{i} cuenta 0, 1, 2, \ldots{} hasta 39. \\
7 & \chip[motion]{mover medida pasos} & \py|i * 5| quiere decir «\texttt{i} por 5». La primera línea mide $0 \times 5 = 0$ pasos, luego 5, luego 10, y así sucesivamente, hasta $39 \times 5 = 195$ pasos. \\
8 & \chip[motion]{girar a la derecha 90 grados} & Una esquina cuadrada cada vez. \\
--- & \chip[variables]{cambiar medida en 5} & ¡No hace falta ninguna línea! \texttt{i} sube de 1 en 1 en cada vuelta, así que \py|i * 5| sube de 5 en 5. \\
\bottomrule
\end{longtable}

\step{Paso 4: Ejecútalo}
Presiona \playbtn. Una espiral cuadrada morada crece desde el centro del escenario. Las líneas se hacen cada vez más largas, de 5 en 5 pasos.

\begin{think}[Un acertijo para ti]
El bucle da 40 vueltas, pero si cuentas con cuidado, la tortuga dibuja solo 39 líneas. ¿Adónde se fue una? (Fíjate en la primera vuelta. La respuesta está al final del libro).
\end{think}

\needspace{12\baselineskip}
\step{Paso 5: Hazlo tuyo}

\begin{challenge}
\tick Cambia \py|right(90)| por \py|right(91)|. ¡Con solo un grado más, la espiral se tuerce como un remolino!\par
\tick Cambia el giro a \py|right(120)| para una espiral triangular, o a \py|right(60)| para una espiral hexagonal.\par
\tick Prueba \py|range(60)| con \py|forward(i * 3)|. ¿Por qué la espiral tiene más líneas pero queda más o menos del mismo tamaño?\par
\tick \textbf{Súper reto:} ¿qué número en la línea 6 haría que la espiral se saliera del escenario? Primero adivina y luego pruébalo.
\end{challenge}

% =====================================================================
\mission{Flor de cuadrados}{bucles dentro de bucles, con 8 espacios}

Un cuadrado solo es aburrido. ¡Doce cuadrados, cada uno un poquito girado, forman una flor!

\begin{tcbraster}[raster columns=2, raster column skip=12pt, raster valign=top, raster force size=false,
  raster column 1/.style={add to width=-2.0cm}, raster column 2/.style={add to width=2.0cm}]
\begin{tcolorbox}[blankest]{\normalsize\bfseries\color{ink}Paso 1: Constrúyelo en Scratch}\end{tcolorbox}
\begin{tcolorbox}[blankest]{\normalsize\bfseries\color{ink}Paso 2: Escríbelo en Colab}\end{tcolorbox}
\begin{tcolorbox}[blankest]
\begin{scratch}[scale=0.75]
\blockinit{al presionar \greenflag}
\blockpen{fijar color de lápiz a \ovalnum{rosa}}
\blockrepeat{repetir \ovalnum{12}}{
  \blockrepeat{repetir \ovalnum{4}}{
    \blockmove{mover \ovalnum{80} pasos}
    \blockmove{girar \turnright{} \ovalnum{90} grados}
  }
  \blockmove{girar \turnright{} \ovalnum{30} grados}
}
\end{scratch}
\end{tcolorbox}
\begin{typethis}
from ColabTurtlePlus.Turtle import *
clearscreen()
initializeTurtle((480, 360))
pencolor("deeppink")
speed(13)
for i in range(12):
    for j in range(4):
        forward(80)
        right(90)
    right(30)
\end{typethis}
\end{tcbraster}

\step{Paso 3: Léelo línea por línea}

\begin{longtable}{>{\centering\arraybackslash}p{1.1cm} P{5.3cm} P{9.1cm}}
\toprule
\textbf{Línea} & \textbf{Bloque de Scratch} & \textbf{Qué significa} \\
\midrule
1--3 & \chip[events]{al presionar bandera verde} & Las tres líneas de inicio. \\
4--5 & \chip[pen]{fijar color de lápiz a rosa} & Un lápiz rosa intenso (\py|"deeppink"|) que dibuja rápido. \\
6 & \chip[control]{repetir 12} & El bucle de afuera: 12 cuadrados. Sin espacios delante. \\
7 & \chip[control]{repetir 4} & El bucle de adentro: un cuadrado. Lleva 4 espacios, así que está dentro del bucle de afuera. Tiene su propio contador, \texttt{j}. \\
8--9 & \chip[motion]{mover 80 pasos} \chip[motion]{girar a la derecha 90 grados} & 8 espacios: dentro de \textbf{los dos} bucles. Estas dos líneas se ejecutan $12 \times 4 = 48$ veces. \\
10 & \chip[motion]{girar a la derecha 30 grados} & Otra vez 4 espacios: después de cada cuadrado, gira un poquito. Se ejecuta 12 veces. \\
\bottomrule
\end{longtable}

\step{Paso 4: Ejecútalo}
Presiona \playbtn. La tortuga dibuja un cuadrado, gira un poquito, dibuja otro cuadrado, y así sucesivamente. Después de 12 cuadrados tienes una flor rosa. Los giros suman $12 \times 30 = 360$, así que la flor da exactamente una vuelta.

\needspace{14\baselineskip}
\step{Paso 5: Hazlo tuyo}

\begin{challenge}
\tick Cambia la línea 6 a \py|range(36)| y la línea 10 a \py|right(10)|. ¿Por qué esos dos números van juntos? (Pista: el truco del 360).\par
\tick Haz una flor de triángulos: cambia el bucle de adentro a \py|range(3)| y \py|right(120)|.\par
\tick Haz una flor de estrellas: bucle de adentro \py|range(5)| con \py|right(144)|, bucle de afuera \py|range(8)| con \py|right(45)|.\par
\tick Agrega \py|hideturtle()| al final de todo para que solo se vea la flor. ¿Cuántos espacios debe llevar?
\end{challenge}

\begin{watchout}
Si te sale un cuadrado y luego un desorden, revisa la línea 10. Necesita exactamente \textbf{4} espacios. Con 8 espacios queda dentro del bucle del cuadrado, y cada esquina gira 30 grados de más.
\end{watchout}

% =====================================================================
\clearpage
\subsection*{Tradúcelo tú}

Aquí hay un programa de Scratch que no has visto antes. Escribe el Python en las líneas. Empieza con las tres líneas de inicio.

\begin{enumerate}
  \item \textbf{Traza una línea} desde cada bloque hasta la palabra de Python que hace el mismo trabajo.
  \item \textbf{Escribe el Python} en las líneas. ¡Cuenta tus espacios con cuidado!
  \item \textbf{Escríbelo en Colab} y ejecútalo. ¿La tortuga dibujó lo mismo que el objeto de Scratch?
\end{enumerate}

\begin{tcbraster}[raster columns=2, raster equal height=rows, raster column skip=12pt, raster valign=top]
\begin{tcolorbox}[colback=white, colframe=white, boxrule=0pt]
\begin{scratch}[scale=0.85]
\blockinit{al presionar \greenflag}
\blockpen{fijar tamaño de lápiz a \ovalnum{4}}
\blockpen{fijar color de lápiz a \ovalnum{morado}}
\blockrepeat{repetir \ovalnum{8}}{
  \blockrepeat{repetir \ovalnum{6}}{
    \blockmove{mover \ovalnum{40} pasos}
    \blockmove{girar \turnright{} \ovalnum{60} grados}
  }
  \blockmove{girar \turnright{} \ovalnum{45} grados}
}
\blockpen{subir lápiz}
\blocklook{decir \ovalnum{¡Listo!}}
\end{scratch}
\end{tcolorbox}
\begin{tcolorbox}[colback=white, colframe=codeframe, arc=5pt, title=\textbf{Mi Python}, coltitle=ink, colbacktitle=codebg]
\vspace{4pt}
\foreach \n in {1,...,12}{\rule{\linewidth}{0.4pt}\par\vspace{9pt}}
\end{tcolorbox}
\end{tcbraster}

\begin{think}[Pistas]
¿Cuántos espacios necesita \chip[motion]{mover 40 pasos}? Está dentro de dos bloques en forma de C. ¿Y dónde va \chip[pen]{subir lápiz}: dentro de un bucle, o después?
\end{think}

% ---------------------------------------------------------------------
\needspace{9\baselineskip}
\begin{puzzle}
Escribe un programa que muestre los números \textbf{1, 2, 3, 4, 5}, uno en cada línea. Solo puedes usar \textbf{dos líneas} de Python, y una de ellas tiene que ser un bucle \py|for|.

\textbf{Pista:} \texttt{i} empieza en 0, pero tú quieres empezar en 1. ¿Qué le puedes hacer a la \texttt{i} dentro del bucle?

\vspace{4pt}
\foreach \n in {1,2}{\rule{\linewidth}{0.4pt}\par\vspace{9pt}}
\end{puzzle}

% ---------------------------------------------------------------------
\needspace{16\baselineskip}
\subsection*{Lista de la Parte 2}
Marca cada una cuando puedas hacerla sin mirar atrás.

\tick Puedo convertir \chip[control]{repetir 10} en \py|for i in range(10):| con dos puntos y cuatro espacios.\par
\tick Ya sé que \texttt{i} cuenta 0, 1, 2, \ldots{} y puedo usarla dentro de un bucle, como en \py|forward(i * 5)|.\par
\tick Puedo poner una línea \textbf{después} de un bucle quitándole los espacios.\par
\tick Puedo poner un bucle dentro de otro bucle, con 8 espacios para las líneas de más adentro.\par
\tick Puedo usar \py|while True:| y detenerlo con el botón de detener \stopbtn.\par
\tick Puedo hacer que Python espere con \py|import time| y \py|time.sleep(1)|.

% =====================================================================
% PARTE 3: Variables y listas
% =====================================================================
\partopener{variables}{PARTE 3}{Variables y listas}{Dale memoria a tus programas: cajas que recuerdan una cosa y mochilas que cargan muchas.}

En Scratch, los bloques naranjas de \textbf{Variables} y los bloques rojo-naranja de \textbf{Listas} le dan memoria a tu proyecto. Una variable recuerda \textbf{una} cosa, como los puntos de un juego. Una lista recuerda \textbf{muchas} cosas en orden, como todo lo que llevas en la mochila.

En esta parte vas a conocer el Python de bloques como estos:

\chip[variables]{fijar puntos a 0} \quad \chip[variables]{cambiar puntos en 1} \quad \chip[variables]{mostrar variable puntos}

\chip[lists]{añadir manzana a mochila} \quad \chip[lists]{elemento 1 de mochila}

\chip[lists]{longitud de mochila} \quad \chip[lists]{¿espada está en mochila?}

Buenas noticias: en Python, las variables y las listas \textbf{no} necesitan líneas de inicio. Los programas que no dibujan pueden empezar enseguida.

% ---------------------------------------------------------------------
\needspace{12\baselineskip}
\subsection*{Una variable es una caja con etiqueta}

En Scratch, haces clic en \textbf{Crear una variable}, escribes un nombre y Scratch te hace una caja. En Python no hay botón. Haces una caja con solo meter algo en ella.

\begin{tcbraster}[raster columns=2, raster equal height=rows, raster column skip=12pt, raster valign=center]
\begin{tcolorbox}[colback=white, colframe=white, boxrule=0pt, valign=center]
\begin{scratch}
\blockinit{al presionar \greenflag}
\blockvariable{fijar \selectmenu{puntos} a \ovalnum{0}}
\blockvariable{fijar \selectmenu{nombre} a \ovalnum{Ana}}
\end{scratch}
\end{tcolorbox}
\begin{code}
puntos = 0
nombre = "Ana"
\end{code}
\end{tcbraster}

\begin{center}
\begin{tikzpicture}[font=\sffamily]
  \foreach \x/\lab/\val in {0/puntos/0, 5/nombre/"Ana"}{
    \fill[variables!25] (\x,0) -- (\x+3,0) -- (\x+3,1.6) -- (\x,1.6) -- cycle;
    \draw[variables!80!black, line width=1pt] (\x,0) rectangle (\x+3,1.6);
    \draw[variables!80!black, line width=1pt, fill=variables!45] (\x-0.15,1.6) -- (\x+3.15,1.6) -- (\x+3.15,1.95) -- (\x-0.15,1.95) -- cycle;
    \node[fill=white, draw=variables!80!black, rounded corners=2pt, font=\sffamily\bfseries] at (\x+1.5,0.55) {\lab};
    \node[font=\ttfamily\Large] at (\x+1.5,2.45) {\val};
    \draw[-{Stealth[length=7pt]}, line width=1pt, ink] (\x+1.5,2.2) -- (\x+1.5,1.35);
  }
\end{tikzpicture}
\end{center}

\textbf{Cómo funciona:} \py|puntos = 0| quiere decir «mete 0 en la caja con la etiqueta \texttt{puntos}». En Python, un \texttt{=} \textbf{no} significa «es igual a». Significa \textbf{mete en la caja}. Léelo de derecha a izquierda: el valor de la derecha entra en la caja de la izquierda. Una caja guarda una sola cosa a la vez, así que si metes algo nuevo, lo viejo sale volando.

\begin{think}[¿Por qué \texttt{"Ana"} lleva comillas?]
Las comillas le dicen a Python «esto es texto». Sin ellas, \py|nombre = Ana| hace que Python busque una caja llamada \texttt{Ana}, y no la encuentra. Los números como \texttt{0} no necesitan comillas.
\end{think}

\begin{think}[Python habla inglés]
Las palabras de Python (\texttt{print}, \texttt{for}, \texttt{append}...) siempre van en inglés. Pero los nombres que inventas tú, como \texttt{puntos} o \texttt{nombre}, pueden ir en español. En este libro los escribimos sin acentos y sin ñ, porque así son más fáciles de teclear.
\end{think}

% ---------------------------------------------------------------------
\needspace{12\baselineskip}
\subsection*{Cambia los puntos}

\begin{tcbraster}[raster columns=2, raster equal height=rows, raster column skip=12pt, raster valign=center]
\begin{tcolorbox}[colback=white, colframe=white, boxrule=0pt, valign=center]
\begin{scratch}
\blockvariable{fijar \selectmenu{puntos} a \ovalnum{0}}
\blockvariable{cambiar \selectmenu{puntos} en \ovalnum{1}}
\blockvariable{cambiar \selectmenu{puntos} en \ovalnum{1}}
\end{scratch}
\end{tcolorbox}
\begin{code}
puntos = 0
puntos = puntos + 1
puntos += 1
\end{code}
\end{tcbraster}

\textbf{Cómo funciona:} La línea 2 se ve rara si piensas que \texttt{=} significa «es igual a». Pero recuerda: significa \textbf{mete en la caja}. Python calcula primero el lado derecho: «lo que hay en \texttt{puntos}, más 1». Eso da 1. Luego vuelve a meter el 1 en la caja. La línea 3 es un atajo que hace exactamente lo mismo: \py|puntos += 1| quiere decir «cambiar puntos en 1». Después de las dos líneas, \texttt{puntos} guarda 2.

\begin{think}[Hacia abajo]
\chip[variables]{cambiar puntos en -1} es \py|puntos -= 1| (o \py|puntos = puntos - 1|). Puedes usar \texttt{+=}, \texttt{-=}, \texttt{*=} y \texttt{/=}.
\end{think}

% ---------------------------------------------------------------------
\needspace{12\baselineskip}
\subsection*{Muestra la variable}

En Scratch, si marcas la casilla junto a una variable, aparece un \textbf{monitor} en el escenario que se actualiza solito. Python no tiene monitor. En su lugar, haces \texttt{print} de la caja cada vez que quieras ver lo que tiene dentro.

\begin{tcbraster}[raster columns=2, raster equal height=rows, raster column skip=12pt, raster valign=center]
\begin{tcolorbox}[colback=white, colframe=white, boxrule=0pt, valign=center]
\begin{scratch}
\blockvariable{fijar \selectmenu{puntos} a \ovalnum{7}}
\blockvariable{mostrar variable \selectmenu{puntos}}
\blocklook{decir \ovaloperator{unir \ovalnum{Puntos: } \ovalvariable{puntos}}}
\end{scratch}
\end{tcolorbox}
\begin{code}
puntos = 7
print(puntos)
print("Puntos:", puntos)
\end{code}
\end{tcbraster}

\begin{outputbox}
7\\
Puntos: 7
\end{outputbox}

\textbf{Cómo funciona:} \py|print(puntos)| no lleva comillas alrededor de \texttt{puntos}, así que Python muestra lo que hay \textbf{dentro} de la caja, no la palabra «puntos». En la línea 3, la coma te deja imprimir dos cosas seguidas, y Python pone un espacio entre ellas por ti. ¡Es más fácil que el bloque \chip[operators]{unir} de Scratch!

% ---------------------------------------------------------------------
\needspace{12\baselineskip}
\subsection*{Ponles nombre a tus cajas}

Scratch te deja llamar a una variable casi como quieras, incluso \texttt{mis puntos} con un espacio. Python es más exigente. Un nombre de Python tiene que seguir estas reglas:

\begin{itemize}
  \item \textbf{Sin espacios.} Usa un guion bajo \texttt{\_} en su lugar: \texttt{mis\_puntos}.
  \item \textbf{No empieces con un dígito.} \texttt{puntos2} está bien, \texttt{2puntos} no.
  \item \textbf{Las mayúsculas importan.} \texttt{Puntos} y \texttt{puntos} son dos cajas distintas.
  \item \textbf{Que no sea una palabra de Python.} No puedes llamar a una caja \texttt{for} o \texttt{import}, porque Python ya usa esas palabras.
\end{itemize}

\begin{center}
\renewcommand{\arraystretch}{1.35}
\begin{tabular}{P{3.6cm} P{1.8cm} P{8.2cm}}
\toprule
\textbf{Nombre} & \textbf{¿Sirve?} & \textbf{Por qué} \\
\midrule
\texttt{puntos} & \textcolor{chalframe}{\textbf{Sí}} & Las minúsculas son perfectas. \\
\texttt{mis\_puntos} & \textcolor{chalframe}{\textbf{Sí}} & El guion bajo une las palabras. \\
\texttt{puntos2} & \textcolor{chalframe}{\textbf{Sí}} & Los dígitos están bien después de la primera letra. \\
\texttt{mis puntos} & \textcolor{errframe}{\textbf{No}} & Un espacio. Python cree que son dos nombres. \\
\texttt{2puntos} & \textcolor{errframe}{\textbf{No}} & Empieza con un dígito. \\
\texttt{for} & \textcolor{errframe}{\textbf{No}} & Python ya usa esta palabra para los bucles. \\
\bottomrule
\end{tabular}
\end{center}

\needspace{12\baselineskip}
\begin{watchout}[¡Cuidado! Las mayúsculas hacen otra caja]
\begin{codenonum}
Puntos = 10
print(puntos)
\end{codenonum}
\begin{errorbox}
NameError: name 'puntos' is not defined
\end{errorbox}
Python solo conoce las cajas que tú has hecho. Lee la última línea del error rojo: nombra la caja que Python no encuentra. Mira con atención: la caja que hiciste es \texttt{Puntos}, con P mayúscula. Usa minúsculas en tus nombres y casi nunca te toparás con este bicho.

Python siempre escribe sus errores en inglés (no los traduce). No te preocupes: la palabra importante es el nombre del error, aquí \texttt{NameError} (error de nombre), y después viene la pista. \texttt{name 'puntos' is not defined} quiere decir «el nombre \texttt{puntos} no está definido».
\end{watchout}

\needspace{10\baselineskip}
\begin{challenge}[Examen rápido: ¿caja o fracaso?]
Marca los nombres que Python va a aceptar. (Las respuestas están al final.)\par
\begin{tabular}{@{}p{4.6cm}p{4.6cm}p{4.6cm}@{}}
\tick \texttt{vidas} & \tick \texttt{mejor marca} & \tick \texttt{jugador\_1} \\
\tick \texttt{1er\_lugar} & \tick \texttt{mejorTiempo} & \tick \texttt{import} \\
\end{tabular}
\end{challenge}

% ---------------------------------------------------------------------
\needspace{12\baselineskip}
\subsection*{¿Texto o número? ¡Las comillas importan!}

Scratch adivina en silencio si algo es un número o un texto. Python nunca adivina. \texttt{5} es un número. \texttt{"5"} entre comillas es texto (Python llama al texto \textbf{string}, que en inglés quiere decir «cadena», como una cadena de letras).

\begin{tcbraster}[raster columns=2, raster equal height=rows, raster column skip=12pt, raster valign=center]
\begin{code}
print(5 + 5)
print("5" + "5")
\end{code}
\begin{outputbox}
10\\
55
\end{outputbox}
\end{tcbraster}

\textbf{Cómo funciona:} Sumar dos números hace matemáticas: 10. Sumar dos textos los pega, como el bloque \chip[operators]{unir}: \texttt{55}. En Scratch, los dos darían 10. Y si intentas sumar un texto y un número, Python se detiene y se queja:

\begin{codenonum}
puntos = 5
print("Puntos: " + puntos)
\end{codenonum}
\begin{errorbox}
TypeError: can only concatenate str (not "int") to str
\end{errorbox}

Esto quiere decir «solo puedo unir texto con texto, y tú me diste un número». (\texttt{str} es texto e \texttt{int} es un número entero.) El arreglo fácil es la coma que ya conoces: \py|print("Puntos:", puntos)|.

% ---------------------------------------------------------------------
\needspace{8\baselineskip}
\begin{think}[¿Solo para este objeto o para todos los objetos?]
Cuando creas una variable en Scratch, Scratch pregunta: ¿\textbf{Para todos los objetos} o \textbf{Sólo para este objeto}? En esta parte, todas tus cajas de Python son como «para todos los objetos»: cualquier línea de tu programa puede usarlas. En una parte posterior aprenderás a darle a cada objeto sus propias cajas, como un gato con su propia \texttt{salud} y un perro con su propia \texttt{salud}.
\end{think}

% ---------------------------------------------------------------------
\needspace{12\baselineskip}
\subsection*{Las variables también dibujan}

Una caja puede guardar la medida de una figura. Cambia la caja, y cambia también cada línea que la usa.

\noindent\begin{minipage}[t]{0.37\linewidth}\vspace{0pt}
\setscratch{scale=0.85}
\begin{scratch}
\blockinit{al presionar \greenflag}
\blockpen{bajar lápiz}
\blockvariable{fijar \selectmenu{medida} a \ovalnum{50}}
\blockrepeat{repetir \ovalnum{4}}{
  \blockmove{mover \ovalvariable{medida} pasos}
  \blockmove{girar \turnright{} \ovalnum{90} grados}
}
\blockvariable{cambiar \selectmenu{medida} en \ovalnum{50}}
\blockrepeat{repetir \ovalnum{4}}{
  \blockmove{mover \ovalvariable{medida} pasos}
  \blockmove{girar \turnright{} \ovalnum{90} grados}
}
\end{scratch}
\end{minipage}\hfill\begin{minipage}[t]{0.61\linewidth}\vspace{0pt}
\begin{code}
from ColabTurtlePlus.Turtle import *
clearscreen()
initializeTurtle((480, 360))
medida = 50
for i in range(4):
    forward(medida)
    right(90)
medida = medida + 50
for i in range(4):
    forward(medida)
    right(90)
\end{code}
\end{minipage}

\textbf{Cómo funciona:} \py|forward(medida)| no lleva comillas, así que la tortuga camina tantos pasos como haya en la caja \texttt{medida}. El primer cuadrado usa 50. La línea 8 hace que la caja guarde 100, así que el segundo cuadrado es el doble de grande. Los dos cuadrados empiezan en la misma esquina.

¿Y si pones \py|medida = medida + 50| \textbf{dentro} del bucle, para que la medida crezca después de cada línea? Lo descubrirás en la misión \textbf{Espiral que crece}\ldots

% ---------------------------------------------------------------------
\needspace{8\baselineskip}
\subsection*{Una lista es una mochila}

Una variable guarda una cosa. Una \textbf{lista} guarda muchas cosas, en orden, bajo un solo nombre. En Scratch haces clic en \textbf{Crear una lista}. En Python escribes corchetes.

\begin{tcbraster}[raster columns=2, raster equal height=rows, raster column skip=12pt, raster valign=center]
\begin{tcolorbox}[colback=white, colframe=white, boxrule=0pt, valign=center, halign=center]
\tcbox[on line, colback=white, colframe=muted, arc=4pt, boxrule=0.8pt]{\sffamily\small Crear una lista}
\par\medskip
Nombre para nueva lista: \textbf{mochila}
\end{tcolorbox}
\begin{code}
mochila = []
print(mochila)
\end{code}
\end{tcbraster}

\textbf{Cómo funciona:} Los corchetes \texttt{[ ]} son la mochila. Sin nada entre ellos, la mochila está vacía, y \py|print(mochila)| muestra \texttt{[]}. Imprimir una lista es el \chip[lists]{mostrar lista mochila} de Python.

% ---------------------------------------------------------------------
\needspace{12\baselineskip}
\subsection*{Mete cosas en la mochila}

\begin{tcbraster}[raster columns=2, raster equal height=rows, raster column skip=12pt, raster valign=center]
\begin{tcolorbox}[colback=white, colframe=white, boxrule=0pt, valign=center]
\begin{scratch}
\blocklist{Eliminar todos de \selectmenu{mochila}}
\blocklist{añadir \ovalnum{manzana} a \selectmenu{mochila}}
\blocklist{añadir \ovalnum{linterna} a \selectmenu{mochila}}
\blocklist{añadir \ovalnum{mapa} a \selectmenu{mochila}}
\end{scratch}
\end{tcolorbox}
\begin{code}
mochila = []
mochila.append("manzana")
mochila.append("linterna")
mochila.append("mapa")
print(mochila)
\end{code}
\end{tcbraster}

\begin{outputbox}
{[}'manzana', 'linterna', 'mapa']
\end{outputbox}

\textbf{Cómo funciona:} \py|append| quiere decir «añadir al final». El punto une el nombre de la lista con el trabajo que quieres que haga: \py|mochila.append("linterna")| quiere decir «mochila, añade linterna a tu final». También puedes llenar una mochila de una sola vez: \py|mochila = ["manzana", "linterna", "mapa"]|. Las comas van entre los elementos.

% ---------------------------------------------------------------------
\needspace{12\baselineskip}
\subsection*{¡Python cuenta desde 0!}

Piensa en una lista como un tren. Cada vagón lleva un elemento. Scratch numera los vagones como asientos: 1, 2, 3. Python cuenta \textbf{a cuántos vagones de distancia del frente} está cada uno, así que el primer vagón es el número \textbf{0}.

\begin{center}
\begin{tikzpicture}[font=\sffamily]
  \fill[ink] (-1.9,0) rectangle (-0.3,1.3);
  \fill[ink] (-1.6,1.3) rectangle (-1.2,1.75);
  \node[white, font=\sffamily\scriptsize\bfseries] at (-1.1,0.65) {locomotora};
  \foreach \i/\item in {0/manzana, 1/linterna, 2/mapa}{
    \draw[lists!80!black, line width=1pt, fill=lists!25, rounded corners=3pt] (\i*2.7,0) rectangle (\i*2.7+2.3,1.3);
    \node[font=\ttfamily\small] at (\i*2.7+1.15,0.65) {"\item"};
    \fill[ink] (\i*2.7+0.5,-0.15) circle (0.18);
    \fill[ink] (\i*2.7+1.8,-0.15) circle (0.18);
    \node[font=\sffamily\scriptsize, text=muted] at (\i*2.7+1.15,1.7) {Scratch: elemento \the\numexpr\i+1\relax};
    \node[font=\sffamily\bfseries, text=kw] at (\i*2.7+1.15,-0.75) {Python: [\i]};
  }
  \draw[ink, line width=1.5pt] (-0.3,0.4) -- (0,0.4);
  \draw[ink, line width=1.5pt] (2.3,0.4) -- (2.7,0.4);
  \draw[ink, line width=1.5pt] (5.0,0.4) -- (5.4,0.4);
\end{tikzpicture}
\end{center}

\begin{tcbraster}[raster columns=2, raster equal height=rows, raster column skip=12pt, raster valign=center]
\begin{tcolorbox}[colback=white, colframe=white, boxrule=0pt, valign=center]
\begin{scratch}[scale=0.96]
\blocklook{decir \ovallist{elemento \ovalnum{1} de \selectmenu{mochila}}}
\blocklook{decir \ovallist{elemento \ovalnum{3} de \selectmenu{mochila}}}
\end{scratch}
\end{tcolorbox}
\begin{code}
print(mochila[0])
print(mochila[2])
\end{code}
\end{tcbraster}

\begin{outputbox}
manzana\\
mapa
\end{outputbox}

\textbf{Cómo funciona:} Los corchetes después del nombre eligen un vagón. Esta es la regla que hay que recordar: \textbf{el elemento $n$ de Scratch es \texttt{[n - 1]} en Python}. El elemento 1 es \py|mochila[0]|, el elemento 3 es \py|mochila[2]|.

% ---------------------------------------------------------------------
\needspace{12\baselineskip}
\subsection*{¿Qué tan larga es la lista?}

\begin{tcbraster}[raster columns=2, raster equal height=rows, raster column skip=12pt, raster valign=center]
\begin{tcolorbox}[colback=white, colframe=white, boxrule=0pt, valign=center]
\begin{scratch}
\blocklook{decir \ovallist{longitud de \selectmenu{mochila}}}
\end{scratch}
\end{tcolorbox}
\begin{code}
print(len(mochila))
\end{code}
\end{tcbraster}

\textbf{Cómo funciona:} \py|len| es la abreviatura de \textbf{length}, que en inglés quiere decir «longitud». Cuenta los vagones: para nuestra mochila da \texttt{3}. Fíjate en que la longitud es 3, pero el último elemento es \py|mochila[2]|. Otra vez la regla de contar desde 0.

% ---------------------------------------------------------------------
\needspace{12\baselineskip}
\subsection*{Saca cosas}

\begin{tcbraster}[raster columns=2, raster equal height=rows, raster column skip=12pt, raster valign=center]
\begin{tcolorbox}[colback=white, colframe=white, boxrule=0pt, valign=center]
\begin{scratch}
\blocklist{eliminar \ovalnum{2} de \selectmenu{mochila}}
\blocklist{Eliminar todos de \selectmenu{mochila}}
\end{scratch}
\end{tcolorbox}
\begin{code}
del mochila[1]
mochila.clear()
\end{code}
\end{tcbraster}

\textbf{Cómo funciona:} \py|del| es la abreviatura de \textbf{delete}, «eliminar» en inglés. El elemento 2 de Scratch es el \texttt{[1]} de Python, así que \py|del mochila[1]| saca \texttt{"linterna"}, y el tren cierra el hueco: \texttt{['manzana', 'mapa']}. \py|mochila.clear()| vacía toda la mochila y deja \texttt{[]} (\textbf{clear} quiere decir «vaciar»).

% ---------------------------------------------------------------------
\needspace{12\baselineskip}
\subsection*{Mete a presión y cambia}

Empieza otra vez con \py|mochila = ["manzana", "linterna", "mapa"]|.

\begin{center}
\begin{scratch}
\blocklist{insertar \ovalnum{cuerda} en \ovalnum{2} de \selectmenu{mochila}}
\blocklist{reemplazar elemento \ovalnum{4} de \selectmenu{mochila} con \ovalnum{lámpara}}
\end{scratch}
\end{center}
\begin{code}
mochila.insert(1, "cuerda")
mochila[3] = "lámpara"
\end{code}

\begin{outputbox}[Después de la línea 1, y después de la línea 2]
{[}'manzana', 'cuerda', 'linterna', 'mapa']\\
{[}'manzana', 'cuerda', 'linterna', 'lámpara']
\end{outputbox}

\textbf{Cómo funciona:} \py|insert(1, "cuerda")| mete a presión un vagón nuevo en el lugar 1 (el lugar 2 de Scratch), y los demás se recorren hacia atrás. \py|mochila[3] = "lámpara"| usa el signo de «mete en la caja»: el vagón 3 (el elemento 4 de Scratch) recibe un elemento nuevo, y el viejo \texttt{"mapa"} sale volando. ¡Cada vagón funciona como una cajita!

% ---------------------------------------------------------------------
\needspace{12\baselineskip}
\subsection*{¿Está en la mochila?}

\begin{center}
\begin{scratch}
\blocklook{decir \boollist{¿\ovalnum{linterna} está en \selectmenu{mochila}?}}
\blocklook{decir \boollist{¿\ovalnum{espada} está en \selectmenu{mochila}?}}
\end{scratch}
\end{center}
\begin{code}
mochila = ["manzana", "linterna", "mapa"]
print("linterna" in mochila)
print("espada" in mochila)
\end{code}

\begin{outputbox}
True\\
False
\end{outputbox}

\textbf{Cómo funciona:} \py|"espada" in mochila| hace una pregunta de sí o no. Python contesta \texttt{True} (sí, verdadero) o \texttt{False} (no, falso), donde Scratch dice \texttt{true} o \texttt{false}. En la Parte 5 usarás preguntas como esta con \py|if|, para hacer algo solo cuando la espada está ahí.

\begin{watchout}[¡Cuidado! Otra vez las mayúsculas]
Scratch piensa que \texttt{Linterna} y \texttt{linterna} son iguales. Python no: \py|"Linterna" in mochila| es \texttt{False}. Escribe tus elementos exactamente igual.
\end{watchout}

% ---------------------------------------------------------------------
\needspace{24\baselineskip}
\subsection*{Un superpoder de Python: recorre la mochila}

En Scratch, para visitar cada elemento necesitas una variable contadora y estos dos bloques:\newline \chip[control]{repetir (longitud de mochila)} y \chip[lists]{elemento (i) de mochila}.\newline Python tiene una forma mucho más corta.

\begin{center}
\begin{scratch}
\blockvariable{fijar \selectmenu{i} a \ovalnum{1}}
\blockrepeat{repetir \ovallist{longitud de \selectmenu{mochila}}}{
  \blocklook{decir \ovallist{elemento \ovalvariable{i} de \selectmenu{mochila}} durante \ovalnum{1} segundos}
  \blockvariable{cambiar \selectmenu{i} en \ovalnum{1}}
}
\end{scratch}
\end{center}
\begin{code}
for cosa in mochila:
    print("En mi mochila:", cosa)
\end{code}

\begin{outputbox}
En mi mochila: manzana\\
En mi mochila: linterna\\
En mi mochila: mapa
\end{outputbox}

\textbf{Cómo funciona:} \py|for cosa in mochila:| quiere decir «saca cada elemento de la mochila, uno por uno, y llámalo \texttt{cosa}». La primera vuelta del bucle, \texttt{cosa} guarda \texttt{"manzana"}, luego \texttt{"linterna"} y luego \texttt{"mapa"}. Tú eliges el nombre: \py|for botana in botanas:| se lee muy bien. ¡Sin contador, sin \texttt{- 1}, sin errores!

% ---------------------------------------------------------------------
\needspace{12\baselineskip}
\subsection*{Elige uno al azar}

\begin{center}
\begin{scratch}
\blocklook{decir \ovallist{elemento \ovalnum{al azar} de \selectmenu{botanas}}}
\end{scratch}
\end{center}
\begin{code}
import random
botanas = ["papitas", "plátano", "galleta"]
print(random.choice(botanas))
\end{code}

\textbf{Cómo funciona:} ¿Sabías que puedes escribir \texttt{al azar} en el bloque \chip[lists]{elemento 1 de} de Scratch? La versión de Python es \py|random.choice(botanas)| (\textbf{choice} quiere decir «elección»): mete la mano en la lista sin mirar y saca un elemento. Igual que \py|import time| en la Parte 2, \py|import random| trae primero las herramientas de azar de Python. Ejecútalo varias veces: no siempre te va a tocar la misma botana.

% ---------------------------------------------------------------------
\needspace{12\baselineskip}
\begin{watchout}[¡Cuidado! Python te avisa cuando falta un vagón]
Pídele a Scratch el \chip[lists]{elemento 4 de mochila} cuando solo hay 3 elementos, y Scratch te da nada, calladito. Python se detiene y te avisa:
\begin{codenonum}
mochila = ["manzana", "linterna", "mapa"]
print(mochila[3])
\end{codenonum}
\begin{errorbox}
IndexError: list index out of range
\end{errorbox}
\textbf{Index} (índice) es la palabra de Python para el número de un vagón. «Out of range» quiere decir «fuera de rango»: «no hay ningún vagón con ese número». Una lista de 3 solo tiene los vagones 0, 1 y 2. Este error es una nota con una pista, no un regaño: te muestra exactamente qué línea arreglar.
\end{watchout}

% =====================================================================
\mission{Contador de puntos}{a crear una variable, cambiarla en un bucle e imprimirla}

Todo juego necesita puntos. Este programa te da 10 puntos cada segundo, cinco veces, y muestra los puntos cada vez, igual que un monitor de Scratch que va subiendo.

\step{Paso 1: Constrúyelo en Scratch}
\noindent\begin{minipage}[t]{0.58\linewidth}\vspace{0pt}
\begin{scratch}[scale=0.95]
\blockinit{al presionar \greenflag}
\blockvariable{fijar \selectmenu{puntos} a \ovalnum{0}}
\blockvariable{mostrar variable \selectmenu{puntos}}
\blockrepeat{repetir \ovalnum{5}}{
  \blockcontrol{esperar \ovalnum{1} segundos}
  \blockvariable{cambiar \selectmenu{puntos} en \ovalnum{10}}
}
\blocklook{decir \ovaloperator{unir \ovalnum{¡Fin del juego! Puntos finales: } \ovalvariable{puntos}}}
\end{scratch}
\end{minipage}\hfill\begin{minipage}[t]{0.33\linewidth}\vspace{0pt}
Haz clic en \textbf{Crear una variable} y llámala \textbf{puntos}. Luego construye:
\end{minipage}


\step{Paso 2: Escríbelo en Colab}
Este no dibuja, así que no necesitas las tres líneas de inicio.

\begin{typethis}
import time
puntos = 0
print("Puntos:", puntos)
for i in range(5):
    time.sleep(1)
    puntos = puntos + 10
    print("Puntos:", puntos)
print("¡Fin del juego! Puntos finales:", puntos)
\end{typethis}

\needspace{8\baselineskip}
\step{Paso 3: Léelo línea por línea}

\begin{longtable}{>{\centering\arraybackslash}p{1.1cm} P{5.3cm} P{9.1cm}}
\toprule
\textbf{Línea} & \textbf{Bloque de Scratch} & \textbf{Qué significa} \\
\midrule
--- & \chip[events]{al presionar bandera verde} & No hace falta código. El botón de reproducir \playbtn{} lo inicia. \\
1 & (esperar necesita esto) & Trae las herramientas de reloj de Python, para poder usar \py|time.sleep|. \\
2 & \chip[variables]{fijar puntos a 0} & Hace una caja llamada \texttt{puntos} y mete 0 en ella. \\
3 & \chip[variables]{mostrar variable puntos} & Muestra lo que hay en la caja: \texttt{Puntos: 0}. \\
4 & \chip[control]{repetir 5} & Haz las líneas con sangría 5 veces. \\
5 & \chip[control]{esperar 1 segundos} & Espera un segundo. \\
6 & \chip[variables]{cambiar puntos en 10} & Toma lo que hay en la caja, súmale 10 y vuelve a meterlo. \\
7 & (el monitor se actualiza) & El monitor de Scratch cambia solo. En Python, imprimimos otra vez para ver los puntos nuevos. \\
8 & \chip[looks]{decir unir ¡Fin del juego! \ldots} & No tiene espacios delante, así que se ejecuta una vez, después del bucle. \\
\bottomrule
\end{longtable}

\step{Paso 4: Ejecútalo}
Presiona \playbtn. Aparece una línea nueva cada segundo.

\noindent\begin{minipage}[t]{0.44\linewidth}\vspace{0pt}
\begin{outputbox}
Puntos: 0\\
Puntos: 10\\
Puntos: 20\\
Puntos: 30\\
Puntos: 40\\
Puntos: 50\\
¡Fin del juego! Puntos finales: 50
\end{outputbox}
\end{minipage}\hfill\begin{minipage}[t]{0.53\linewidth}\vspace{0pt}
\begin{think}[¿Y si olvidas la línea 2?]
Borra la línea 2 y haz clic en \textbf{Entorno de ejecución} (\textit{Runtime}), luego en \textbf{Reiniciar sesión} (\textit{Restart session}), y vuelve a ejecutar la celda. Python dice \texttt{NameError: name 'puntos' is not defined}. ¡No puedes sumarle 10 a una caja que todavía no existe! En Scratch, una variable nueva empieza en 0 solita. En Python, siempre llenas la caja primero.
\end{think}
\end{minipage}


\begin{challenge}
\tick Cambia la línea 6 por \py|puntos += 10|. ¿Cambia algo cuando lo ejecutas?\par
\tick Haz que cuente \textbf{hacia atrás}: empieza los puntos en 50 y quita 10 cada vez. (Pista: \texttt{-=}.)\par
\tick Agrega una segunda caja, \py|vidas = 3|, antes del bucle. Dentro del bucle, quita una vida e imprime las dos: \py|print("Puntos:", puntos, "Vidas:", vidas)|. ¿Qué les pasa a las vidas después de 5 vueltas?\par
\tick \textbf{Superreto:} haz que los puntos que ganas sean cada vez más. Agrega \py|premio = 10| antes del bucle. Dentro de él, suma \py|premio| a los puntos y luego haz \py|premio| 10 más grande. ¿Cuántos puntos tienes al final?
\end{challenge}

% =====================================================================
\mission{Espiral que crece}{a usar una variable para la medida y cambiarla dentro de un bucle}

En la página «Las variables también dibujan», la medida cambió una sola vez, entre dos cuadrados. Ahora la medida crece un poquito después de \textbf{cada} línea. Cada línea es más larga que la anterior, y la tortuga hace una espiral hacia afuera.

\step{Paso 1 y Paso 2: Constrúyelo en Scratch y luego escríbelo en Colab}
Crea una variable llamada \textbf{medida}. Agrega la extensión \textbf{Lápiz} y construye el programa de la izquierda. Luego escribe el Python de la derecha.

\noindent\begin{minipage}[t]{0.37\linewidth}\vspace{0pt}
\setscratch{scale=0.85}
\begin{scratch}
\blockinit{al presionar \greenflag}
\blockpen{borrar todo}
\blockpen{bajar lápiz}
\blockpen{fijar color de lápiz a \ovalnum{morado}}
\blockvariable{fijar \selectmenu{medida} a \ovalnum{5}}
\blockrepeat{repetir \ovalnum{40}}{
  \blockmove{mover \ovalvariable{medida} pasos}
  \blockmove{girar \turnright{} \ovalnum{90} grados}
  \blockvariable{cambiar \selectmenu{medida} en \ovalnum{5}}
}
\end{scratch}
\end{minipage}\hfill\begin{minipage}[t]{0.61\linewidth}\vspace{0pt}
\begin{typethis}
from ColabTurtlePlus.Turtle import *
clearscreen()
initializeTurtle((480, 360))
speed(13)
pencolor("purple")
medida = 5
for i in range(40):
    forward(medida)
    right(90)
    medida = medida + 5
\end{typethis}
\end{minipage}

\step{Paso 3: Léelo línea por línea}

\begin{longtable}{>{\centering\arraybackslash}p{1.1cm} P{5.3cm} P{9.1cm}}
\toprule
\textbf{Línea} & \textbf{Bloque de Scratch} & \textbf{Qué significa} \\
\midrule
1--3 & \chip[events]{al presionar bandera verde}\newline\chip[pen]{borrar todo} & Las tres líneas de inicio. \py|clearscreen()| limpia el escenario, como \chip[pen]{borrar todo}. El lápiz empieza abajo. \\
4 & & La velocidad más rápida, para que no esperes demasiado. \\
5 & \chip[pen]{fijar color de lápiz a}\newline\chip[pen]{morado} & Los nombres de colores van entre comillas, en minúsculas. Los colores se escriben en inglés: \texttt{"purple"} es morado, \texttt{"red"} rojo, \texttt{"blue"} azul, \texttt{"gold"} dorado\ldots \\
6 & \chip[variables]{fijar medida a 5} & La primera línea medirá 5 pasos. \\
7 & \chip[control]{repetir 40} & Haz las líneas con sangría 40 veces. \\
8 & \chip[motion]{mover (medida) pasos} & Camina hacia adelante lo que haya en la caja \texttt{medida} en este momento. \\
9 & \chip[motion]{girar a la derecha 90 grados} & Da vuelta en una esquina. \\
10 & \chip[variables]{cambiar medida en 5} & Hace la caja 5 más grande, \textbf{dentro} del bucle, para que la siguiente línea sea más larga. \\
\bottomrule
\end{longtable}

\step{Paso 4: Ejecútalo}
Una espiral cuadrada morada crece desde el centro del escenario. Las líneas miden 5, 10, 15, 20 \ldots{} pasos, y la última mide 200 pasos.

\begin{think}[Los espacios importan]
La línea 10 tiene cuatro espacios delante, así que está dentro del bucle. Quítale los espacios y ejecútalo. Ahora \py|medida| solo cambia \textbf{una vez}, cuando el bucle ya terminó, así que cada línea mide 5 pasos y la tortuga solo da vueltas a un cuadradito 40 veces. ¡No sale ningún error rojo, pero el dibujo está mal!
\end{think}

\enlargethispage{\baselineskip}
\begin{challenge}\small
\tick Cambia \py|right(90)| por \py|right(91)|. ¡La espiral se tuerce! Prueba también con 89, 120 y 144.\par
\tick Haz la espiral más apretada con \py|medida = medida + 2|, o más abierta con \py|medida = medida + 8|. Con \texttt{+ 8}, ¿todavía cabe en el escenario?\par
\tick Haz el lápiz más grueso también: agrega \py|grueso = 1| antes del bucle, y luego \py|pensize(grueso)| y \py|grueso = grueso + 0.2| dentro de él.\par
\tick \textbf{Superreto:} haz una espiral que se \textbf{encoja}. Empieza con \py|medida = 200| y quita 5 cada vez. (Pista: para que quepa en el escenario, salta a \py|goto(-100, 100)| con el lápiz arriba antes de empezar.)
\end{challenge}

% =====================================================================
\mission{Arcoíris con una lista de colores}{una lista de colores y un bucle \texttt{for} que la recorre}

¡Un arcoíris es una lista de colores! Este programa guarda los seis colores en una lista y luego dibuja una franja gruesa por cada color, de arriba abajo.

\step{Paso 1: Constrúyelo en Scratch}
\noindent\begin{minipage}[t]{0.61\linewidth}\vspace{0pt}
\begin{scratch}[scale=0.79]
\blockinit{al presionar \greenflag}
\blockpen{borrar todo}
\blockpen{fijar tamaño de lápiz a \ovalnum{15}}
\blockvariable{fijar \selectmenu{y} a \ovalnum{60}}
\blockvariable{fijar \selectmenu{i} a \ovalnum{1}}
\blockrepeat{repetir \ovallist{longitud de \selectmenu{colores}}}{
  \blockpen{fijar color de lápiz a \ovallist{elemento \ovalvariable{i} de \selectmenu{colores}}}
  \blockpen{subir lápiz}
  \blockmove{ir a x: \ovalnum{-150} y: \ovalvariable{y}}
  \blockpen{bajar lápiz}
  \blockmove{mover \ovalnum{300} pasos}
  \blockvariable{cambiar \selectmenu{y} en \ovalnum{-15}}
  \blockvariable{cambiar \selectmenu{i} en \ovalnum{1}}
}
\end{scratch}
\end{minipage}\hfill\begin{minipage}[t]{0.36\linewidth}\vspace{0pt}
Crea una lista llamada \textbf{colores} y dos variables llamadas \textbf{y} e \textbf{i}. El lápiz de Scratch no entiende nombres de colores, así que llena la lista en el escenario con seis códigos de color: \texttt{\#FF0000}, \texttt{\#FFA500}, \texttt{\#FFFF00}, \texttt{\#008000}, \texttt{\#0000FF}, \texttt{\#800080}. Luego construye:
\end{minipage}


\step{Paso 2: Escríbelo en Colab}

\begin{typethis}
from ColabTurtlePlus.Turtle import *
clearscreen()
initializeTurtle((480, 360))
speed(13)
pensize(15)
colores = ["red", "orange", "yellow",
           "green", "blue", "purple"]
y = 60
for color in colores:
    pencolor(color)
    penup()
    goto(-150, y)
    pendown()
    forward(300)
    y = y - 15
hideturtle()
\end{typethis}

\par\medskip{\large\bfseries\color{ink}Paso 3: Léelo línea por línea}\par\nopagebreak

\begin{longtable}{>{\centering\arraybackslash}p{1.1cm} P{5.3cm} P{9.1cm}}
\toprule
\textbf{Línea} & \textbf{Bloque de Scratch} & \textbf{Qué significa} \\
\midrule
1--3 & \chip[events]{al presionar bandera verde} & Las tres líneas de inicio. También hacen \chip[pen]{borrar todo}. \\
4--5 & \chip[pen]{fijar tamaño de lápiz a 15} & La velocidad más rápida, y franjas gruesas. \\
6--7 & (la lista en el escenario) & Seis nombres de colores, en inglés: rojo, naranja, amarillo, verde, azul y morado. Dentro de los corchetes puedes seguir en la línea siguiente. \\
8 & \chip[variables]{fijar y a 60} & La primera franja está en $y = 60$. \\
9 & \chip[control]{repetir longitud de colores} & Saca cada color, uno por uno, y llámalo \texttt{color}. ¡Sin contador \texttt{i}! \\
10 & \chip[pen]{fijar color de lápiz a}\newline\chip[lists]{elemento i de colores} & Usa el color que acabamos de sacar. \texttt{color} no lleva comillas: es una caja, no la palabra «color». Python entiende \textbf{nombres} de colores, así que no hacen falta códigos. \\
11--13 & \chip[pen]{subir lápiz}\newline\chip[motion]{ir a x: -150 y: (y)}\newline\chip[pen]{bajar lápiz} & Salta al extremo izquierdo de esta franja sin dibujar. \\
14 & \chip[motion]{mover 300 pasos} & Dibuja la franja. La tortuga mira a la derecha. \\
15 & \chip[variables]{cambiar y en -15} & Baja 15 para la siguiente franja. Luego la línea 16 esconde la tortuga. \\
\bottomrule
\end{longtable}

\step{Paso 4: Ejecútalo}
Aparecen seis franjas gruesas, una debajo de otra: roja arriba, luego naranja, amarilla, verde, azul y morada hasta abajo.

\begin{think}[\texttt{color} y \texttt{pencolor}]
¡Qué suerte! En español decimos \textbf{color}, y en inglés también se escribe \textbf{color}. Por eso Python usa esa misma palabra en sus nombres: \py|pencolor| quiere decir «color del lápiz». Pero ojo: \py|pencolor| es una palabra de Python, y \texttt{colores} y \texttt{color} son cajas que inventamos nosotros. Nuestros nombres de listas y cajas pueden ser los que queramos, siempre que los escribamos igual cada vez.
\end{think}

\begin{challenge}
\tick Agrega una séptima franja con \py|colores.append("pink")| (rosa) en una línea nueva después de la línea 7.\par
\tick Deja huecos entre las franjas: cambia el \texttt{15} de la línea 15 por \texttt{20}.\par
\tick Haz tu propia bandera. Cambia la lista por \py|["red", "white", "blue"]| (rojo, blanco, azul). ¿Qué países podrías hacer?\par
\tick \textbf{Superreto:} haz que cada franja sea más corta que la de arriba. Agrega \py|largo = 300| antes del bucle, usa \py|forward(largo)| y quítale 30 a \texttt{largo} en cada vuelta.
\end{challenge}

% =====================================================================
\mission{Inventor de historias al azar}{a llenar listas, añadirles cosas y elegir con \texttt{random.choice}}

Tres listas, una pizca de azar, y te sale una oración chistosa cada vez. ¿Se comerá un dragón tu tarea?

\step{Paso 1: Constrúyelo en Scratch}
\noindent\begin{minipage}[t]{0.70\linewidth}\vspace{0pt}
\setscratch{scale=0.92}
\begin{scratch}[scale=0.89]
\blockinit{al presionar \greenflag}
\blocklist{añadir \ovalnum{mi tarea} a \selectmenu{cosas}}
\blockrepeat{repetir \ovalnum{3}}{
  \blockvariable{fijar \selectmenu{heroe} a \ovallist{elemento \ovalnum{al azar} de \selectmenu{heroes}}}
  \blockvariable{fijar \selectmenu{accion} a \ovallist{elemento \ovalnum{al azar} de \selectmenu{acciones}}}
  \blockvariable{fijar \selectmenu{cosa} a \ovallist{elemento \ovalnum{al azar} de \selectmenu{cosas}}}
  \blocklook{decir \ovaloperator{unir \ovalvariable{heroe} \ovaloperator{unir \ovalvariable{accion} \ovalvariable{cosa}}} durante \ovalnum{2} segundos}
}
\end{scratch}
\end{minipage}\hfill\begin{minipage}[t]{0.27\linewidth}\vspace{0pt}
Crea tres listas, \textbf{heroes}, \textbf{acciones} y \textbf{cosas}, y llénalas en el escenario. Crea tres variables, \textbf{heroe}, \textbf{accion} y \textbf{cosa}. Luego construye:
\end{minipage}


\step{Paso 2: Escríbelo en Colab}

\begin{typethis}
import random
heroes = ["Un dragón", "Mi maestra", "Un robotito"]
acciones = ["se comió", "le cantó a", "persiguió a"]
cosas = ["un plátano", "la luna", "un calcetín"]
cosas.append("mi tarea")
for i in range(3):
    heroe = random.choice(heroes)
    accion = random.choice(acciones)
    cosa = random.choice(cosas)
    print(heroe, accion, cosa)
\end{typethis}

\needspace{8\baselineskip}
\step{Paso 3: Léelo línea por línea}

\begin{longtable}{>{\centering\arraybackslash}p{1.1cm} P{5.3cm} P{9.1cm}}
\toprule
\textbf{Línea} & \textbf{Bloque de Scratch} & \textbf{Qué significa} \\
\midrule
1 & & Trae las herramientas de azar de Python. \\
2--4 & (listas llenadas en el escenario) & Tres listas, cada una con tres elementos. Cada elemento es un texto, así que cada uno lleva comillas. Las comas van entre los elementos. \\
5 & \chip[lists]{añadir mi tarea a cosas} & Pone una cuarta cosa al final de la lista \texttt{cosas}. \\
6 & \chip[control]{repetir 3} & Haz tres oraciones. \\
7--9 & \chip[variables]{fijar heroe a}\newline\chip[lists]{elemento al azar de heroes} \ldots & Elige un elemento de cada lista y mételo en una caja. \\
10 & \chip[looks]{decir unir heroe \ldots} & Imprime las tres cajas seguidas. Las comas ponen espacios entre ellas, así que no necesitas tres bloques \chip[operators]{unir}. \\
\bottomrule
\end{longtable}

\step{Paso 4: Ejecútalo}
Te van a salir tres oraciones al azar. ¡Las tuyas serán diferentes!

\begin{outputbox}[Podrías ver]
Un dragón se comió un calcetín\\
Un robotito persiguió a mi tarea\\
Mi maestra le cantó a la luna
\end{outputbox}

\begin{think}[Haz clic dos veces en la bandera]
En Scratch, las listas recuerdan sus elementos aunque el proyecto se detenga. Haz clic dos veces en la bandera verde, ¡y \texttt{mi tarea} se añade \textbf{dos veces}! En Python, la línea 4 hace una lista \texttt{cosas} nuevecita cada vez que ejecutas la celda, así que siempre empieza con exactamente tres cosas.
\end{think}

\begin{challenge}
\tick Agrega tus propias palabras a cada lista. ¡No olvides las comillas y las comas!\par
\tick Haz 5 oraciones en vez de 3.\par
\tick Agrega una cuarta lista, \texttt{lugares}, como \py|["en la tina", "en el camión"]|, e imprime un lugar al final de cada oración.\par
\tick \textbf{Superreto:} después del bucle, imprime \py|len(cosas)| y la lista \texttt{cosas} completa. ¿Cuántas cosas hay, y a dónde fue a parar \texttt{"mi tarea"}?
\end{challenge}

% =====================================================================
\banner{variables}{TU TURNO}{Tradúcelo tú}

¡Ahora el traductor eres tú! Escribe el Python de este nuevo programa de Scratch y luego ejecútalo en Colab.

\noindent\begin{minipage}[t]{0.47\linewidth}\vspace{0pt}
\setscratch{scale=0.9}
\begin{scratch}[scale=0.63]
\blockinit{al presionar \greenflag}
\blocklist{Eliminar todos de \selectmenu{mascotas}}
\blocklist{añadir \ovalnum{gato} a \selectmenu{mascotas}}
\blocklist{añadir \ovalnum{perro} a \selectmenu{mascotas}}
\blocklist{añadir \ovalnum{pez} a \selectmenu{mascotas}}
\blocklook{decir \ovallist{longitud de \selectmenu{mascotas}} durante \ovalnum{2} segundos}
\blocklook{decir \ovallist{elemento \ovalnum{2} de \selectmenu{mascotas}} durante \ovalnum{2} segundos}
\end{scratch}
\end{minipage}\hfill\begin{minipage}[t]{0.5\linewidth}\vspace{0pt}
\begin{tcolorbox}[colback=white, colframe=codeframe, arc=5pt, title=\textbf{Mi Python}, coltitle=ink, colbacktitle=codebg]
\vspace{4pt}
\foreach \n in {1,...,6}{\rule{\linewidth}{0.4pt}\par\vspace{9pt}}
\end{tcolorbox}
\small\textbf{Pistas:} ¿Qué empieza una lista vacía? ¿Qué palabra añade al final? ¡El elemento 2 no es \texttt{[2]}!
\end{minipage}

\medskip
\begin{puzzle}
¿Puedes contar los elementos de una lista \textbf{sin} \py|len|? Escribe las \textbf{dos} líneas que faltan, con un bucle \py|for| y \texttt{+=}. Debe imprimir \texttt{4}.
\begin{codenonum}
palabras = ["tigre", "león", "cebra", "hipopótamo"]
cuenta = 0
\end{codenonum}
\vspace{10pt}
\rule{\linewidth}{0.4pt}\par\vspace{10pt}
\rule{\linewidth}{0.4pt}\par\vspace{2pt}
\hspace*{8pt}\texttt{print(cuenta)}
\end{puzzle}

\smallskip
\begin{tcolorbox}[colback=variables!10, colframe=variables!80!black, arc=6pt, boxrule=0.8pt, left=10pt, top=2pt, bottom=2pt, title={\textbf{Lista de la Parte 3}}, coltitle=white, colbacktitle=variables!80!black]
\tick Ya sé crear una variable con \py|=| y explicar que significa «mete en la caja».\par
\tick Puedo cambiar una variable con \py|puntos += 1| y mostrarla con \py|print|.\par
\tick Ya sé las reglas de los nombres, y que \texttt{5} y \texttt{"5"} son diferentes.\par
\tick Puedo crear una lista y usar \py|append|, \py|len|, \py|del|, \py|insert| y \py|clear|.\par
\tick Ya sé que el elemento 1 de Scratch es \texttt{[0]} en Python, y sé qué significa \texttt{IndexError}.\par
\tick Puedo recorrer una lista con \py|for| y elegir al azar con \py|random.choice|.
\end{tcolorbox}

% =====================================================================
% PARTE 4: Operadores
% =====================================================================
\partopener{operators}{PARTE 4}{Operadores}{Haz cuentas, tira dados, compara cosas y desarma palabras.}
\setscratch{scale=0.8}

Los bloques verdes de \textbf{Operadores} son los bloques que piensan en Scratch. Hacen cuentas, eligen números al azar, revisan si algo es verdadero y juegan con las palabras. En esta parte vas a conocer la versión en Python de cada uno de ellos:

\begin{itemize}
  \item \chip[operators]{7 + 2} se convierte en \py|7 + 2|
  \item \chip[operators]{elegir número al azar entre 1 y 6} se convierte en \py|random.randint(1, 6)|
  \item \chip[operators]{unir Hola Ana} se convierte en \py|"Hola " + "Ana"|
  \item \chip[operators]{letra 1 de banana} se convierte en \py|"banana"[0]|
\end{itemize}

Casi todos los programas de esta parte \textbf{no dibujan}, así que no necesitas las tres líneas de inicio. Solo crea una celda nueva con \textbf{+ Código} (\emph{+ Code}), escribe y presiona \playbtn.

\begin{think}[Un mesero amable y un chef robot estricto]
Scratch es como un mesero amable. Si le pides algo un poco revuelto, adivina lo que querías decir. Python es como un chef robot estricto. Hace \textbf{exactamente} lo que le dices, y si tu pedido no tiene sentido, se detiene y te dice por qué. Suena molesto, pero significa que Python nunca adivina mal sin avisarte. Busca los recuadros \textbf{¡Cuidado!} de esta parte: muestran los lugares donde el chef robot es más exigente que el mesero.
\end{think}

% ---------------------------------------------------------------------
\needspace{15\baselineskip}
\subsection*{Hacer cuentas: + -- * /}

En Scratch puedes meter un bloque verde de suma dentro de un bloque \chip[looks]{decir}. En Python pones la cuenta dentro de \py|print( )|.

\begin{tcbraster}[raster columns=2, raster equal height=rows, raster column skip=12pt, raster valign=center, raster force size=false, raster odd column/.style={add to width=-1.2cm}, raster even column/.style={add to width=1.2cm}]
\begin{tcolorbox}[colback=white, colframe=white, boxrule=0pt, valign=center]
\begin{scratch}
\blockinit{al presionar \greenflag}
\blocklook{decir \ovaloperator{\ovalnum{7} + \ovalnum{2}} durante \ovalnum{1} segundos}
\blocklook{decir \ovaloperator{\ovalnum{7} - \ovalnum{2}} durante \ovalnum{1} segundos}
\blocklook{decir \ovaloperator{\ovalnum{7} * \ovalnum{2}} durante \ovalnum{1} segundos}
\blocklook{decir \ovaloperator{\ovalnum{7} / \ovalnum{2}} durante \ovalnum{1} segundos}
\end{scratch}
\end{tcolorbox}
\begin{tcolorbox}[colback=white, colframe=white, boxrule=0pt, left=0pt, right=0pt, top=0pt, bottom=0pt, valign=center]
\begin{code}
print(7 + 2)
print(7 - 2)
print(7 * 2)
print(7 / 2)
\end{code}
\begin{outputbox}
9\\
5\\
14\\
3.5
\end{outputbox}
\end{tcolorbox}
\end{tcbraster}

\textbf{Cómo funciona.} Los signos son los mismos que en Scratch: \texttt{*} (una estrellita) significa por, y \texttt{/} (una diagonal) significa entre. No necesitas el óvalo verde: los números y el signo simplemente van uno junto al otro. Pon espacios alrededor del signo para que se lea fácil.

\needspace{10\baselineskip}
\begin{watchout}[¡Cuidado! Dividir siempre da un decimal]
En Python, \texttt{/} siempre da un número con punto decimal, aunque la división sea exacta. \py|print(6 / 2)| muestra \texttt{3.0}, no \texttt{3}. ¡Y sigue siendo la respuesta correcta!

\textbf{Extra:} dos diagonales, \texttt{//}, dividen y tiran a la basura lo que sobra. \py|print(7 // 2)| muestra \texttt{3}. Scratch no tiene un bloque para esto.
\end{watchout}

\needspace{14\baselineskip}
\begin{think}[Los paréntesis son los óvalos verdes]
En Scratch eliges qué cuenta se hace primero metiendo un bloque \textbf{dentro} de otro. En Python usas paréntesis. Sin paréntesis, Python hace las multiplicaciones y divisiones antes que las sumas y restas, igual que en la clase de matemáticas:
\begin{center}
\begin{tabular}{l@{\hspace{1.5cm}}l}
\ovaloperator{\ovalnum{2} + \ovaloperator{\ovalnum{3} * \ovalnum{4}}} & \py|2 + 3 * 4| da \texttt{14} \\[8pt]
\ovaloperator{\ovaloperator{\ovalnum{2} + \ovalnum{3}} * \ovalnum{4}} & \py|(2 + 3) * 4| da \texttt{20}
\end{tabular}
\end{center}
\end{think}

% ---------------------------------------------------------------------
\needspace{15\baselineskip}
\subsection*{Elegir número al azar}

\begin{tcbraster}[raster columns=1, raster valign=center, raster row skip=6pt]
\begin{tcolorbox}[colback=white, colframe=white, boxrule=0pt, valign=center]
\begin{scratch}
\blockinit{al presionar \greenflag}
\blockvariable{fijar \selectmenu{tirada} a \ovaloperator{elegir número al azar entre \ovalnum{1} y \ovalnum{6}}}
\blocklook{decir \ovalvariable{tirada}}
\end{scratch}
\end{tcolorbox}
\begin{code}
import random
tirada = random.randint(1, 6)
print(tirada)
\end{code}
\end{tcbraster}

\textbf{Cómo funciona.} \py|import random| es como hacer clic en el botón \textbf{Agregar Extensión}: trae las herramientas de números al azar de Python. Solo lo necesitas una vez, al principio de la celda. \py|random.randint(1, 6)| significa «de las herramientas de azar, elige un número entero (en inglés, un \textbf{int}eger) del 1 al 6». El punto significa «de adentro de». Igual que en Scratch, \textbf{los dos extremos cuentan}, así que puede salir 1 y puede salir 6. Ejecuta la celda varias veces: cada vez sale un número diferente.

Python habla inglés: sus palabras (\texttt{import}, \texttt{print}, \texttt{random}...) siempre van en inglés. Lo que tú inventas, como el nombre \py|tirada|, sí puede ir en español.

\begin{watchout}[¡Cuidado! El número más chico primero]
A Scratch no le importa \chip[operators]{elegir número al azar entre 6 y 1}: voltea los números por ti. Python se detiene con \texttt{ValueError: empty range} («rango vacío»). Pon siempre el número más chico primero.

Python siempre escribe sus errores en inglés (no los traduce). La palabra importante es el nombre del error, al principio: \texttt{ValueError}, \texttt{TypeError}, \texttt{IndexError}...
\end{watchout}

% ---------------------------------------------------------------------
\needspace{15\baselineskip}
\subsection*{Más grande, más chico, igual: > < ==}

Los bloques verdes puntiagudos responden una pregunta con \textbf{verdadero} o \textbf{falso}. Puedes meter uno en un bloque \chip[looks]{decir} para ver la respuesta. En Python, la respuesta es \py|True| (verdadero) o \py|False| (falso).

\begin{tcbraster}[raster columns=2, raster equal height=rows, raster column skip=12pt, raster valign=center, raster force size=false, raster odd column/.style={add to width=-1.2cm}, raster even column/.style={add to width=1.2cm}]
\begin{tcolorbox}[colback=white, colframe=white, boxrule=0pt, valign=center]
\begin{scratch}[scale=0.77]
\blockinit{al presionar \greenflag}
\blockvariable{fijar \selectmenu{puntos} a \ovalnum{10}}
\blocklook{decir \booloperator{\ovalnum{5} > \ovalnum{3}} durante \ovalnum{1} segundos}
\blocklook{decir \booloperator{\ovalnum{5} < \ovalnum{3}} durante \ovalnum{1} segundos}
\blocklook{decir \booloperator{\ovalvariable{puntos} = \ovalnum{10}}}
\end{scratch}
\end{tcolorbox}
\begin{tcolorbox}[colback=white, colframe=white, boxrule=0pt, left=0pt, right=0pt, top=0pt, bottom=0pt, valign=center]
\begin{code}
puntos = 10
print(5 > 3)
print(5 < 3)
print(puntos == 10)
\end{code}
\begin{outputbox}
True\\
False\\
True
\end{outputbox}
\end{tcolorbox}
\end{tcbraster}

\textbf{Cómo funciona.} \texttt{>} y \texttt{<} son iguales que en Scratch. Para revisar si dos cosas son iguales, Python usa \textbf{dos} signos de igual, \texttt{==}. Acuérdate de la Parte 3: \textbf{un} signo de igual guarda un valor en una caja (\py|puntos = 10|). Dos signos de igual hacen una pregunta: «¿\texttt{puntos} es igual a 10?». \py|True| y \py|False| siempre empiezan con mayúscula.

En la Parte 5 vas a meter estas preguntas dentro de \py|if|, para que el programa decida qué hacer. Por ahora, solo imprimimos las respuestas.

\begin{think}[Dos bloques que Scratch no tiene]
Python también puede preguntar «más grande \textbf{o igual}» con \texttt{>=}, y «más chico o igual» con \texttt{<=}. \py|print(10 >= 10)| muestra \texttt{True}. Y \texttt{!=} significa «\textbf{no} es igual a».
\end{think}

\needspace{12\baselineskip}
\begin{watchout}[¡Cuidado! ¡Las mayúsculas cuentan!]
En Scratch, \chip[operators]{Gato = gato} es \textbf{verdadero}. A Scratch no le importan las mayúsculas. A Python sí: \py|print("Gato" == "gato")| muestra \texttt{False}.

La solución es \py|.lower()|, que convierte el texto a minúsculas antes de comparar:\par
\py|print("Gato".lower() == "gato".lower())| muestra \texttt{True}.
\end{watchout}

% ---------------------------------------------------------------------
\needspace{15\baselineskip}
\subsection*{y, o, no: and, or, not}

\begin{tcbraster}[raster columns=2, raster equal height=rows, raster column skip=12pt, raster valign=center, raster force size=false, raster odd column/.style={add to width=-1.2cm}, raster even column/.style={add to width=1.2cm}]
\begin{tcolorbox}[colback=white, colframe=white, boxrule=0pt, valign=center]
\begin{scratch}[scale=0.74]
\blockinit{al presionar \greenflag}
\blockvariable{fijar \selectmenu{edad} a \ovalnum{9}}
\blocklook{decir \booloperator{\booloperator{\ovalvariable{edad} > \ovalnum{8}} y \booloperator{\ovalvariable{edad} < \ovalnum{12}}}}
\blocklook{decir \booloperator{\booloperator{\ovalvariable{edad} < \ovalnum{5}} o \booloperator{\ovalvariable{edad} > \ovalnum{10}}}}
\blocklook{decir \booloperator{no \booloperator{\ovalvariable{edad} > \ovalnum{8}}}}
\end{scratch}
\end{tcolorbox}
\begin{tcolorbox}[colback=white, colframe=white, boxrule=0pt, left=0pt, right=0pt, top=0pt, bottom=0pt, valign=center]
\begin{code}
edad = 9
print(edad > 8 and edad < 12)
print(edad < 5 or edad > 10)
print(not edad > 8)
\end{code}
\begin{outputbox}
True\\
False\\
False
\end{outputbox}
\end{tcolorbox}
\end{tcbraster}

\textbf{Cómo funciona.} Python usa las mismas palabras que Scratch, pero en inglés y en minúsculas: \py|and| es «y», \py|or| es «o» y \py|not| es «no». \py|and| es \py|True| solo cuando \textbf{los dos} lados son verdaderos. \py|or| es \py|True| cuando \textbf{al menos uno} de los lados es verdadero. \py|not| voltea la respuesta: verdadero se vuelve falso y falso se vuelve verdadero. No necesitas los hexágonos extra: Python resuelve primero las preguntas con \texttt{>} y \texttt{<}, y después el \py|and|.

% ---------------------------------------------------------------------
\needspace{15\baselineskip}
\subsection*{Unir: pegar palabras}

\begin{tcbraster}[raster columns=1, raster valign=center, raster row skip=6pt]
\begin{tcolorbox}[colback=white, colframe=white, boxrule=0pt, valign=center]
\begin{scratch}
\blockinit{al presionar \greenflag}
\blockvariable{fijar \selectmenu{nombre} a \ovalnum{Ana}}
\blockvariable{fijar \selectmenu{puntos} a \ovalnum{7}}
\blocklook{decir \ovaloperator{unir \ovalnum{Hola } \ovalvariable{nombre}} durante \ovalnum{1} segundos}
\blocklook{decir \ovaloperator{unir \ovalnum{Puntos: } \ovalvariable{puntos}}}
\end{scratch}
\end{tcolorbox}
\begin{tcolorbox}[colback=white, colframe=white, boxrule=0pt, left=0pt, right=0pt, top=0pt, bottom=0pt, valign=center]
\begin{code}
nombre = "Ana"
puntos = 7
print("Hola " + nombre)
print("Puntos: " + str(puntos))
\end{code}
\begin{outputbox}
Hola Ana\\
Puntos: 7
\end{outputbox}
\end{tcolorbox}
\end{tcbraster}

\textbf{Cómo funciona.} Cuando \texttt{+} tiene texto en los dos lados, los une, igual que \chip[operators]{unir}. Python llama a un texto \textbf{string}, que en inglés quiere decir «hilo» (imagina letras ensartadas en un hilo, como las cuentas de un collar). Fíjate en el espacio dentro de \py|"Hola "|: sin él saldría \texttt{HolaAna}.

\needspace{12\baselineskip}
\begin{watchout}[¡Cuidado! No puedes unir texto con un número]
\py|print("Puntos: " + puntos)| se detiene con esta pista:
\begin{errorbox}
TypeError: can only concatenate str (not "int") to str
\end{errorbox}
\emph{Concatenate} es «concatenar», una palabra larga para unir. Python está diciendo: «Solo puedo unir un string con otro string (\texttt{str}), y tú me diste un número (\texttt{int})». Scratch simplemente adivinaría. En Python, primero convierte el número en string con \py|str(puntos)|, como en la línea 4 de arriba.
\end{watchout}

\needspace{10\baselineskip}
\begin{think}[Extra: los f-strings, la forma fácil de unir]
Pon una \texttt{f} antes de las comillas, y todo lo que esté dentro de llaves \texttt{\{ \}} se rellena solito, también los números:\par
\py|print(f"Puntos: {puntos}")| muestra \texttt{Puntos: 7}. ¡Sin \py|str( )| y sin \texttt{+}!

Y un truco más: \py|print("ja" * 3)| muestra \texttt{jajaja}.
\end{think}

% ---------------------------------------------------------------------
\needspace{15\baselineskip}
\subsection*{Letra de, y longitud de}

\begin{tcbraster}[raster columns=1, raster valign=center, raster row skip=6pt]
\begin{tcolorbox}[colback=white, colframe=white, boxrule=0pt, valign=center]
\begin{scratch}
\blockinit{al presionar \greenflag}
\blockvariable{fijar \selectmenu{palabra} a \ovalnum{banana}}
\blocklook{decir \ovaloperator{letra \ovalnum{1} de \ovalvariable{palabra}} durante \ovalnum{1} segundos}
\blocklook{decir \ovaloperator{longitud de \ovalvariable{palabra}}}
\end{scratch}
\end{tcolorbox}
\begin{tcolorbox}[colback=white, colframe=white, boxrule=0pt, left=0pt, right=0pt, top=0pt, bottom=0pt, valign=center]
\begin{code}
palabra = "banana"
print(palabra[0])
print(len(palabra))
\end{code}
\begin{outputbox}
b\\
6
\end{outputbox}
\end{tcolorbox}
\end{tcbraster}

\textbf{Cómo funciona.} Los corchetes después de un texto sacan una letra, igual que sacar un elemento de una lista en la Parte 3. Y, igual que con las listas, Python \textbf{cuenta desde 0}. Así que la \chip[operators]{letra 1 de banana} de Scratch es \py|palabra[0]| en Python, la letra 2 es \py|palabra[1]|, y así. La regla es: \textbf{réstale 1} al número de Scratch. \py|len( )| es la forma corta de \emph{length}, «longitud» en inglés: cuenta las letras. También funciona con listas.

\begin{center}
\begin{tikzpicture}[x=1.15cm]
  \foreach \l [count=\i from 0] in {b,a,n,a,n,a}{
    \draw[codeframe, fill=codebg, rounded corners=3pt] (\i,0) rectangle ++(0.95,0.95);
    \node[font=\ttfamily\Large] at (\i+0.475,0.475) {\l};
    \node[font=\small\bfseries, text=operators!60!black] at (\i+0.475,-0.35) {\i};
    \pgfmathtruncatemacro{\s}{\i+1}
    \node[font=\small, text=muted] at (\i+0.475,1.3) {\s};
  }
  \node[anchor=east, font=\small, text=muted] at (-0.2,1.3) {Scratch};
  \node[anchor=east, font=\small\bfseries, text=operators!60!black] at (-0.2,-0.35) {Python};
\end{tikzpicture}
\end{center}

\begin{watchout}[¡Cuidado! Te pasaste del final]
\texttt{"banana"} tiene 6 letras, así que la última es \py|palabra[5]|. Pedir \py|palabra[6]| se detiene con \texttt{IndexError: string index out of range} («posición fuera del texto»). (Scratch solo da una respuesta vacía.) \textbf{Extra:} \py|palabra[-1]| es un atajo para la última letra.
\end{watchout}

% ---------------------------------------------------------------------
\needspace{15\baselineskip}
\subsection*{Contiene}

\begin{tcbraster}[raster columns=2, raster equal height=rows, raster column skip=12pt, raster valign=center, raster force size=false, raster odd column/.style={add to width=-1.2cm}, raster even column/.style={add to width=1.2cm}]
\begin{tcolorbox}[colback=white, colframe=white, boxrule=0pt, valign=center]
\begin{scratch}
\blockinit{al presionar \greenflag}
\blocklook{decir \booloperator{¿\ovalnum{banana} contiene \ovalnum{nan}?}}
\end{scratch}
\end{tcolorbox}
\begin{code}
print("nan" in "banana")
\end{code}
\end{tcbraster}

\textbf{Cómo funciona.} Muestra \texttt{True}. Léelo en voz alta: «¿\texttt{nan} está \textbf{en} (\py|in|) \texttt{banana}?». El orden está al revés que en Scratch: primero va el pedacito, luego \py|in| y luego la palabra grande. Igual que \texttt{==}, a \py|in| le importan las mayúsculas: \py|"B" in "banana"| es \py|False|. Usa \py|.lower()| otra vez si lo necesitas.

% ---------------------------------------------------------------------
\needspace{15\baselineskip}
\subsection*{Módulo: lo que sobra}

\begin{tcbraster}[raster columns=1, raster valign=center, raster row skip=6pt]
\begin{tcolorbox}[colback=white, colframe=white, boxrule=0pt, valign=center]
\begin{scratch}
\blockinit{al presionar \greenflag}
\blocklook{decir \ovaloperator{\ovalnum{7} módulo \ovalnum{3}} durante \ovalnum{1} segundos}
\blocklook{decir \booloperator{\ovaloperator{\ovalnum{10} módulo \ovalnum{2}} = \ovalnum{0}}}
\end{scratch}
\end{tcolorbox}
\begin{tcolorbox}[colback=white, colframe=white, boxrule=0pt, left=0pt, right=0pt, top=0pt, bottom=0pt, valign=center]
\begin{code}
print(7 % 3)
print(10 % 2 == 0)
\end{code}
\begin{outputbox}
1\\
True
\end{outputbox}
\end{tcolorbox}
\end{tcbraster}

\textbf{Cómo funciona.} En Python, el bloque \chip[operators]{módulo} es el signo de porcentaje, \texttt{\%}. Da lo que \textbf{sobra} después de dividir. Reparte 7 dulces entre 3 amigos: a cada uno le tocan 2, y sobra 1. Así que \py|7 % 3| es 1.

\begin{think}[¿Par o impar?]
Un número par dividido entre 2 no deja nada de sobra. Así que \py|n % 2 == 0| es \py|True| para los números pares y \py|False| para los impares. Los programadores usan este truco todo el tiempo.
\end{think}

% ---------------------------------------------------------------------
\needspace{15\baselineskip}
\subsection*{Valor absoluto, raíz cuadrada y sus amigos}

\begin{tcbraster}[raster columns=1, raster valign=center, raster row skip=6pt]
\begin{tcolorbox}[colback=white, colframe=white, boxrule=0pt, valign=center]
\begin{scratch}
\blockinit{al presionar \greenflag}
\blocklook{decir \ovaloperator{\selectmenu{valor absoluto} de \ovalnum{-5}} durante \ovalnum{1} segundos}
\blocklook{decir \ovaloperator{\selectmenu{raíz cuadrada} de \ovalnum{16}}}
\end{scratch}
\end{tcolorbox}
\begin{tcolorbox}[colback=white, colframe=white, boxrule=0pt, left=0pt, right=0pt, top=0pt, bottom=0pt, valign=center]
\begin{code}
import math
print(abs(-5))
print(math.sqrt(16))
\end{code}
\begin{outputbox}
5\\
4.0
\end{outputbox}
\end{tcolorbox}
\end{tcbraster}

\textbf{Cómo funciona.} \py|abs( )| (de \emph{absolute}, «absoluto») quita el signo de menos. Ya viene incluido en Python, así que puedes usarlo enseguida. La raíz cuadrada vive en las herramientas de matemáticas de Python, así que primero escribe \py|import math| y luego usa \py|math.sqrt( )| (de \emph{square root}, «raíz cuadrada»). La raíz cuadrada de 16 es 4, porque $4 \times 4 = 16$. Igual que dividir, da un decimal: \texttt{4.0}.

Las otras opciones de ese menú de Scratch también están en \texttt{math}: \chip[operators]{piso de} es \py|math.floor( )| (redondear hacia abajo) y \chip[operators]{techo de} es \py|math.ceil( )| (redondear hacia arriba). Las de ángulos, como \texttt{sen}, funcionan un poco diferente en Python, así que déjalas para después.

% ---------------------------------------------------------------------
\needspace{15\baselineskip}
\subsection*{Redondear}

\begin{tcbraster}[raster columns=1, raster valign=center, raster row skip=6pt]
\begin{tcolorbox}[colback=white, colframe=white, boxrule=0pt, valign=center]
\begin{scratch}
\blockinit{al presionar \greenflag}
\blocklook{decir \ovaloperator{redondear \ovalnum{2.7}} durante \ovalnum{1} segundos}
\blocklook{decir \ovaloperator{redondear \ovalnum{2.5}}}
\end{scratch}
\end{tcolorbox}
\begin{tcolorbox}[colback=white, colframe=white, boxrule=0pt, left=0pt, right=0pt, top=0pt, bottom=0pt, valign=center]
\begin{code}
print(round(2.7))
print(round(2.5))
\end{code}
\begin{outputbox}
3\\
2
\end{outputbox}
\end{tcolorbox}
\end{tcbraster}

\textbf{Cómo funciona.} \py|round( )| («redondear») va al número entero más cercano, así que 2.7 se vuelve 3. ¡Pero mira la segunda respuesta!

\begin{watchout}[¡Cuidado! Python redondea 2.5 a 2]
Scratch redondea 2.5 \textbf{hacia arriba}, a 3. Python redondea un número que termina en .5 al número \textbf{par} más cercano. Así que \py|round(2.5)| es 2, pero \py|round(3.5)| es 4. Los programadores adultos lo hacen así para que, con muchísimos números, redondear no empuje el total hacia arriba.

¿Quieres la forma de Scratch? Usa \py|math.floor(2.5 + 0.5)|, que da \texttt{3}.

\textbf{Extra:} \py|round(3.14159, 2)| deja 2 decimales y da \texttt{3.14}.
\end{watchout}

% ---------------------------------------------------------------------
\needspace{15\baselineskip}
\subsection*{Python es más exigente que Scratch}

Aquí están todos los momentos del chef robot de esta parte en una sola tabla. Regresa a ella cuando Python te muestre un error en rojo.

\begin{center}\small
\renewcommand{\arraystretch}{1.35}
\begin{tabular}{P{3.6cm} P{2.3cm} P{3.7cm} P{5.2cm}}
\toprule
\textbf{Tú escribes} & \textbf{Scratch dice} & \textbf{Python dice} & \textbf{Mejor haz esto} \\
\midrule
\texttt{"Puntos: " + 10} & Puntos: 10 & \texttt{TypeError} & \py|"Puntos: " + str(10)| \\
\texttt{"Gato" == "gato"} & verdadero & \texttt{False} & \py|.lower()| en los dos \\
letra 1 de \texttt{"gato"} & g & \py|"gato"[1]| es \texttt{a} & \py|"gato"[0]| \\
redondear 2.5 & 3 & \texttt{2} & \py|math.floor(2.5 + 0.5)| \\
\texttt{10 / 0} & Infinito & \texttt{ZeroDivisionError} & Nunca dividas entre 0 \\
elegir número al azar entre 6 y 1 & un número & \texttt{ValueError} & El número más chico primero \\
\bottomrule
\end{tabular}
\end{center}

\needspace{16\baselineskip}
\subsection*{¡Adivina qué sale!}
Antes de ejecutar cada línea, escribe lo que crees que va a mostrar Python. Luego escribe las seis líneas en una celda y ejecútala. ¿Cuántas le atinaste? (Las respuestas están al final del libro.)

\begin{center}
\renewcommand{\arraystretch}{1.6}
\begin{tabular}{l P{5cm}}
\toprule
\textbf{Código} & \textbf{Creo que muestra} \\
\midrule
\py|print(9 / 3)| & \rule{4cm}{0.4pt} \\
\py|print("4" + "4")| & \rule{4cm}{0.4pt} \\
\py|print("Perro" == "perro")| & \rule{4cm}{0.4pt} \\
\py|print("pizza"[1])| & \rule{4cm}{0.4pt} \\
\py|print(round(4.5))| & \rule{4cm}{0.4pt} \\
\py|print(17 % 5)| & \rule{4cm}{0.4pt} \\
\bottomrule
\end{tabular}
\end{center}

% =====================================================================
\setscratch{scale=0.7}
\mission{Lanzador de dados}{\texttt{random.randint}, ir sumando dentro de un bucle y unir números con \texttt{str( )}}

Tira un dado 10 veces, suma todas las tiradas y calcula el promedio. ¿El promedio se parece a lo que esperabas?

\step{Paso 1: Constrúyelo en Scratch}
Crea dos variables, \textbf{tirada} y \textbf{total}, y luego arma esto:

\begin{center}
\begin{scratch}
\blockinit{al presionar \greenflag}
\blockvariable{fijar \selectmenu{total} a \ovalnum{0}}
\blockrepeat{repetir \ovalnum{10}}{
  \blockvariable{fijar \selectmenu{tirada} a \ovaloperator{elegir número al azar entre \ovalnum{1} y \ovalnum{6}}}
  \blockvariable{cambiar \selectmenu{total} en \ovalvariable{tirada}}
  \blocklook{decir \ovaloperator{unir \ovalnum{Sacaste } \ovalvariable{tirada}} durante \ovalnum{1} segundos}
}
\blocklook{decir \ovaloperator{unir \ovalnum{Total: } \ovalvariable{total}} durante \ovalnum{2} segundos}
\blocklook{decir \ovaloperator{unir \ovalnum{Promedio: } \ovaloperator{\ovalvariable{total} / \ovalnum{10}}}}
\end{scratch}
\end{center}

\step{Paso 2: Escríbelo en Colab}

\begin{typethis}
import random
total = 0
for i in range(10):
    tirada = random.randint(1, 6)
    total = total + tirada
    print("Sacaste " + str(tirada))
print("Total: " + str(total))
print("Promedio: " + str(total / 10))
\end{typethis}

\step{Paso 3: Léelo línea por línea}

\begin{longtable}{>{\centering\arraybackslash}p{1.1cm} P{5.3cm} P{9.1cm}}
\toprule
\textbf{Línea} & \textbf{Bloque de Scratch} & \textbf{Qué significa} \\
\midrule
--- & \chip[events]{al presionar bandera verde} & No hace falta código: el botón de reproducir \playbtn{} lo arranca. \\
1 & (Agregar Extensión) & Trae las herramientas de azar. \\
2 & \chip[variables]{fijar total a 0} & Crea una caja llamada \py|total| con 0 adentro. \\
3 & \chip[control]{repetir 10} & Haz 10 veces las líneas con sangría. \\
4 & \chip[variables]{fijar tirada a}\newline\chip[operators]{elegir número al azar} & Tira el dado: un número entero del 1 al 6, con los dos extremos incluidos. \\
5 & \chip[variables]{cambiar total en tirada} & Toma lo que hay en \py|total|, súmale la tirada y guarda el resultado otra vez en \py|total|. \\
6 & \chip[looks]{decir unir Sacaste tirada} & \py|tirada| es un número, así que \py|str(tirada)| lo convierte en string antes de unirlo. \\
7 & \chip[looks]{decir unir Total: total} & No tiene espacios adelante, así que se ejecuta una sola vez, después del bucle. \\
8 & \chip[looks]{decir unir Promedio: \ldots} & Primero se calcula \py|total / 10| y luego se convierte en string. \\
\bottomrule
\end{longtable}

\step{Paso 4: Ejecútalo}
Presiona \playbtn. Tus tiradas van a ser diferentes, ¡porque son al azar! Esta es una ejecución real:

\begin{outputbox}[Podrías ver]
Sacaste 6\\
Sacaste 2\\
{\sffamily\itshape\color{muted} \ldots\ 8 tiradas más: 6, 2, 1, 3, 3, 1, 4, 6 \ldots}\\
Total: 34\\
Promedio: 3.4
\end{outputbox}

El promedio es un decimal porque \texttt{/} siempre da un decimal. Si tiras un dado justo muchísimas veces, el promedio se acerca a 3.5, exactamente a la mitad entre 1 y 6.

\begin{challenge}
\tick Tira 100 veces en vez de 10. Cambia la línea 3 \textbf{y} la línea 8. ¿Ahora el promedio está más cerca de 3.5? (Consejo: primero quita la línea 6, ¡o te van a salir 100 líneas!)\par
\tick Haz un dado de 20 caras para un juego de mesa.\par
\tick Tira \textbf{dos} dados cada vez: \py|tirada = random.randint(1, 6) + random.randint(1, 6)|. ¿Cuál es el total más grande que podrías sacar?\par
\tick Agrega \py|print(tirada == 6)| dentro del bucle. Muestra \texttt{True} cada vez que sale un seis.
\end{challenge}

% =====================================================================
\mission{Impresora de tablas}{multiplicar dentro de un bucle, y los f-strings}

Deja que Python imprima una tabla de multiplicar completa en un instante. Luego cambia un número para tener la tabla que quieras.

\step{Paso 1: Constrúyelo en Scratch}
Crea dos variables, \textbf{tabla} y \textbf{n}, y luego arma esto:

\begin{center}
\begin{scratch}
\blockinit{al presionar \greenflag}
\blockvariable{fijar \selectmenu{tabla} a \ovalnum{7}}
\blockvariable{fijar \selectmenu{n} a \ovalnum{1}}
\blockrepeat{repetir \ovalnum{12}}{
  \blocklook{decir \ovaloperator{unir \ovalvariable{n} \ovaloperator{unir \ovalnum{ x 7 = } \ovaloperator{\ovalvariable{n} * \ovalvariable{tabla}}}} durante \ovalnum{1} segundos}
  \blockvariable{cambiar \selectmenu{n} en \ovalnum{1}}
}
\end{scratch}
\end{center}

¡Son muchos bloques de unir! Python lo hace mucho más fácil.

\step{Paso 2: Escríbelo en Colab}

\begin{typethis}
tabla = 7
n = 1
for i in range(12):
    print(f"{n} x {tabla} = {n * tabla}")
    n = n + 1
\end{typethis}

\step{Paso 3: Léelo línea por línea}

\begin{longtable}{>{\centering\arraybackslash}p{1.1cm} P{5.3cm} P{9.1cm}}
\toprule
\textbf{Línea} & \textbf{Bloque de Scratch} & \textbf{Qué significa} \\
\midrule
1 & \chip[variables]{fijar tabla a 7} & Qué tabla de multiplicar vamos a imprimir. \\
2 & \chip[variables]{fijar n a 1} & \py|n| cuenta 1, 2, 3 \ldots\ hasta 12. \\
3 & \chip[control]{repetir 12} & Haz 12 veces las líneas con sangría. \\
4 & \chip[looks]{decir unir \ldots} & Un \textbf{f-string}. Cada par de llaves es un hueco que Python rellena. ¡\texttt{\{n * tabla\}} hasta hace la cuenta primero! La \texttt{x} es solo una letra x dentro del texto. \\
5 & \chip[variables]{cambiar n en 1} & Súmale 1 a \py|n|, listo para la siguiente línea. \\
\bottomrule
\end{longtable}

\step{Paso 4: Ejecútalo}

\begin{outputbox}
1 x 7 = 7\\
2 x 7 = 14\\
3 x 7 = 21\\
4 x 7 = 28\\
{\sffamily\itshape\color{muted} \ldots\ y así sigue, hasta \ldots}\\
12 x 7 = 84
\end{outputbox}

\begin{think}[La misma línea sin f-string]
También podrías unir con \texttt{+} y \py|str( )|, como en la Misión 1. Funciona, pero queda larguísimo, así que es más fácil en tres líneas:
\begin{codenonum}
a = str(n) + " x "
b = str(tabla) + " = "
print(a + b + str(n * tabla))
\end{codenonum}
Las dos formas están bien. Casi todos los programadores de Python eligen el f-string.
\end{think}

\begin{challenge}
\tick Imprime la tabla del 9. ¿Qué línea cambias?\par
\tick Llega hasta 20 x 7.\par
\tick Imprime los números al cuadrado: cambia la línea 4 para que muestre \texttt{5 x 5 = 25}, y así. (Pista: solo necesitas \py|n|.)\par
\tick Agrega \texttt{\{n * tabla \% 2 == 0\}} al final del f-string. ¿Qué te dice el \texttt{True} o el \texttt{False}?
\end{challenge}

% =====================================================================
\mission{Arte de puntos al azar}{posiciones al azar con \texttt{goto}, \texttt{dot( )}, y elegir un color al azar de una lista}

La tortuga salta a 50 lugares al azar del escenario y deja un punto de color en cada uno. ¡Cada vez que lo ejecutas, sale un dibujo nuevo!

\step{Paso 1: Constrúyelo en Scratch}
Agrega la extensión \textbf{Lápiz} y luego arma esto:

\begin{center}
\begin{scratch}
\blockinit{al presionar \greenflag}
\blockpen{borrar todo}
\blockpen{subir lápiz}
\blocklook{esconder}
\blockrepeat{repetir \ovalnum{50}}{
  \blockmove{ir a x: \ovaloperator{elegir número al azar entre \ovalnum{-240} y \ovalnum{240}} y: \ovaloperator{elegir número al azar entre \ovalnum{-180} y \ovalnum{180}}}
  \blockpen{fijar tamaño de lápiz a \ovaloperator{elegir número al azar entre \ovalnum{10} y \ovalnum{40}}}
  \blockpen{fijar \selectmenu{color} de lápiz a \ovaloperator{elegir número al azar entre \ovalnum{0} y \ovalnum{100}}}
  \blockpen{bajar lápiz}
  \blockpen{subir lápiz}
}
\end{scratch}
\end{center}

\step{Paso 2: Escríbelo en Colab}

\begin{typethis}
from ColabTurtlePlus.Turtle import *
clearscreen()
initializeTurtle((480, 360))
import random

colores = ["red", "gold", "skyblue", "pink", "lime"]
speed(13)
penup()
hideturtle()
for i in range(50):
    x = random.randint(-240, 240)
    y = random.randint(-180, 180)
    goto(x, y)
    medida = random.randint(10, 40)
    tono = colores[random.randint(0, 4)]
    dot(medida, tono)
\end{typethis}

Los colores se escriben en inglés: \py|"red"| es rojo, \py|"gold"| dorado, \py|"skyblue"| azul cielo, \py|"pink"| rosa y \py|"lime"| verde limón.

\step{Paso 3: Léelo línea por línea}

\begin{longtable}{>{\centering\arraybackslash}p{1.1cm} P{5.3cm} P{9.1cm}}
\toprule
\textbf{Línea} & \textbf{Bloque de Scratch} & \textbf{Qué significa} \\
\midrule
1--3 & \chip[events]{al presionar bandera verde}\newline\chip[pen]{borrar todo} & Las tres líneas de inicio. También dejan el escenario limpio. \\
4 & (Agregar Extensión) & Trae las herramientas de azar. \\
6 & (una lista) & Una lista de 5 nombres de colores. En Scratch, los colores del lápiz son números del 0 al 100, pero a Python le gustan los \textbf{nombres} de los colores, así que guardamos nuestros favoritos en una lista. \\
7 & & Dibuja tan rápido como pueda ir la tortuga. \\
8--9 & \chip[pen]{subir lápiz}\newline\chip[looks]{esconder} & No dibujes líneas entre los puntos, y esconde la tortuga. \\
10 & \chip[control]{repetir 50} & Haz 50 puntos. \\
11--13 & \chip[motion]{ir a x: \ldots{} y: \ldots} & Elige una \py|x| al azar de un lado del escenario al otro, y una \py|y| al azar de abajo hacia arriba. Luego ve ahí. \\
14 & \chip[pen]{fijar tamaño de lápiz a}\newline\chip[operators]{elegir número al azar} & Una medida al azar para el punto, del 10 al 40. \\
15 & \chip[pen]{fijar color de lápiz a}\newline\chip[operators]{elegir número al azar} & Elige un lugar al azar de la lista. La lista tiene 5 colores, pero Python cuenta desde 0, así que los lugares son 0, 1, 2, 3 y 4. Por eso es \py|randint(0, 4)| y no \py|randint(1, 5)|. Llamamos \py|tono| al color elegido, porque Python ya usa la palabra \py|color| para la tortuga. \\
16 & \chip[pen]{bajar lápiz}\newline\chip[pen]{subir lápiz} & En Scratch, bajar lápiz y luego subir lápiz hace un punto. \py|dot(medida, tono)| dibuja un punto redondo justo donde está la tortuga. \\
\bottomrule
\end{longtable}

\step{Paso 4: Ejecútalo}
Deberías ver 50 puntos de diferentes tamaños en rojo, dorado, azul cielo, rosa y verde limón, regados por todo el escenario. Ejecútalo otra vez: el dibujo es diferente cada vez.

\begin{challenge}
\tick Agrega \py|bgcolor("black")| después de la línea 3 para tener un cielo de noche. Haz 200 puntos.\par
\tick Agrega más colores a la lista, como \py|"orange"| (naranja) y \py|"purple"| (morado). La línea 15 solo elige los lugares del 0 al 4, así que nunca los va a elegir hasta que la arregles. (Forma astuta: \py|random.randint(0, len(colores) - 1)| siempre le queda a la lista.)\par
\tick Haz que todos los puntos caigan en la mitad de arriba del escenario. ¿Qué números de la línea 12 cambias?\par
\tick \textbf{Extra:} \py|random.choice(colores)| elige un elemento al azar de una lista de una sola vez. Pruébalo en la línea 15.
\end{challenge}

% =====================================================================
\mission{Máquina de nombres}{desarmar palabras con letras y \texttt{len( )}, y armar una palabra al revés}

Python toma un nombre, cuenta sus letras, lo deletrea con una letra por línea y luego lo escribe al revés.

\step{Paso 1: Constrúyelo en Scratch}
Crea tres variables, \textbf{nombre}, \textbf{n} y \textbf{alreves}, y luego arma esto:

\begin{center}
\begin{scratch}
\blockinit{al presionar \greenflag}
\blockvariable{fijar \selectmenu{nombre} a \ovalnum{Matilda}}
\blocklook{decir \ovaloperator{unir \ovalnum{Tu nombre tiene } \ovaloperator{unir \ovaloperator{longitud de \ovalvariable{nombre}} \ovalnum{ letras}}} durante \ovalnum{2} segundos}
\blockvariable{fijar \selectmenu{n} a \ovalnum{1}}
\blockrepeat{repetir \ovaloperator{longitud de \ovalvariable{nombre}}}{
  \blocklook{decir \ovaloperator{letra \ovalvariable{n} de \ovalvariable{nombre}} durante \ovalnum{1} segundos}
  \blockvariable{cambiar \selectmenu{n} en \ovalnum{1}}
}
\blockvariable{fijar \selectmenu{alreves} a \ovalnum{}}
\blockvariable{fijar \selectmenu{n} a \ovalnum{1}}
\blockrepeat{repetir \ovaloperator{longitud de \ovalvariable{nombre}}}{
  \blockvariable{fijar \selectmenu{alreves} a \ovaloperator{unir \ovaloperator{letra \ovalvariable{n} de \ovalvariable{nombre}} \ovalvariable{alreves}}}
  \blockvariable{cambiar \selectmenu{n} en \ovalnum{1}}
}
\blocklook{decir \ovalvariable{alreves}}
\end{scratch}
\end{center}

\step{Paso 2: Escríbelo en Colab}
¡Pon tu propio nombre en la línea 1!

\begin{typethis}
nombre = "Matilda"
print("Tu nombre tiene " + str(len(nombre)) + " letras")
for i in range(len(nombre)):
    print(nombre[i])

alreves = ""
for i in range(len(nombre)):
    alreves = nombre[i] + alreves
print(alreves)
\end{typethis}

\step{Paso 3: Léelo línea por línea}

\begin{longtable}{>{\centering\arraybackslash}p{1.1cm} P{5.3cm} P{9.1cm}}
\toprule
\textbf{Línea} & \textbf{Bloque de Scratch} & \textbf{Qué significa} \\
\midrule
1 & \chip[variables]{fijar nombre a Matilda} & El nombre va entre comillas, porque es un string. \\
2 & \chip[looks]{decir unir Tu nombre tiene \ldots} & \py|len(nombre)| cuenta las letras. \py|str( )| convierte ese número en string, para poder unirlo. \\
3 & \chip[control]{repetir longitud de nombre} & Da una vuelta por cada letra. \\
4 & \chip[looks]{decir letra n de nombre} & Aquí está la parte astuta. La \py|i| del bucle cuenta 0, 1, 2, 3 \ldots\ ¡y esos son justo los números que Python usa para las letras! Así que Python no necesita \texttt{n} ni \chip[variables]{cambiar n en 1}. \\
6 & \chip[variables]{fijar alreves a (vacío)} & Dos comillas sin nada en medio forman un string vacío, como una caja vacía. \\
7 & \chip[control]{repetir longitud de nombre} & Recorre las letras otra vez. \\
8 & \chip[variables]{fijar alreves a unir \ldots} & Pega cada letra nueva \textbf{adelante}. M, luego aM, luego taM \ldots\ hasta que todo el nombre queda al revés. \\
9 & \chip[looks]{decir alreves} & Muestra el nombre al revés, una vez, después del bucle. \\
\bottomrule
\end{longtable}

\step{Paso 4: Ejecútalo}

\begin{outputbox}
Tu nombre tiene 7 letras\\
M\\
a\\
t\\
i\\
l\\
d\\
a\\
adlitaM
\end{outputbox}

\begin{think}[Los espacios también son letras]
Prueba \py|nombre = "Ada Lovelace"|. Python dice que tiene 12 letras, porque también cuenta el espacio. La \chip[operators]{longitud de} de Scratch hace lo mismo.
\end{think}

\begin{challenge}
\tick Imprime el nombre en mayúsculas con \py|print(nombre.upper())|.\par
\tick Imprime la primera y la última letra: \py|nombre[0]| y \py|nombre[-1]|.\par
\tick Numera las letras: cambia la línea 4 por \py|print(str(i + 1) + ": " + nombre[i])|. ¿Por qué dice \py|i + 1|?\par
\tick ¿Tu nombre tiene una \texttt{a}? Agrega \py|print("a" in nombre.lower())|.
\end{challenge}

% =====================================================================
\needspace{22\baselineskip}
\subsection*{Parte 4: Tradúcelo tú}

Aquí hay un programa de Scratch que no has visto. Escríbelo en Python en las líneas, luego escríbelo en Colab y ejecútalo. ¿Hace lo mismo que el programa de Scratch? (La respuesta está al final del libro.)

\begin{tcbraster}[raster columns=2, raster equal height=rows, raster column skip=12pt, raster valign=top]
\begin{tcolorbox}[colback=white, colframe=white, boxrule=0pt]
\begin{scratch}[scale=0.66]
\blockinit{al presionar \greenflag}
\blockrepeat{repetir \ovalnum{3}}{
  \blockvariable{fijar \selectmenu{a} a \ovaloperator{elegir número al azar entre \ovalnum{1} y \ovalnum{10}}}
  \blockvariable{fijar \selectmenu{b} a \ovaloperator{elegir número al azar entre \ovalnum{1} y \ovalnum{10}}}
  \blocklook{decir \ovaloperator{unir \ovalnum{Total: } \ovaloperator{\ovalvariable{a} + \ovalvariable{b}}} durante \ovalnum{1} segundos}
}
\end{scratch}
\end{tcolorbox}
\begin{tcolorbox}[colback=white, colframe=codeframe, arc=5pt, title=\textbf{Mi Python}, coltitle=ink, colbacktitle=codebg]
\vspace{4pt}
\foreach \n in {1,...,7}{\rule{\linewidth}{0.4pt}\par\vspace{9pt}}
\end{tcolorbox}
\end{tcbraster}

\textbf{Dos pistas:} ¿qué necesitas hasta arriba antes de poder elegir números al azar? ¿Y puedes unir \texttt{"Total: "} directamente con un número?

\begin{puzzle}
La contraseña secreta de la tortuga es la \textbf{primera letra} de \texttt{Tortuga} unida al \textbf{número de letras} de \texttt{Tortuga}. Escribe \textbf{dos líneas} de Python que calculen la contraseña y la impriman. No escribas tú la respuesta: ¡haz que Python la calcule!

Cuando funcione, deberías ver:
\begin{outputbox}
T7
\end{outputbox}
\end{puzzle}

\setscratch{scale=1}
\needspace{8\baselineskip}
\subsection*{Lista de la Parte 4}
\tick Ya sé hacer cuentas con \texttt{+ - * /}, y sé que \texttt{/} siempre da un decimal.\par
\tick Puedo elegir un número al azar con \py|random.randint( )|, después de \py|import random|.\par
\tick Puedo comparar con \texttt{> < ==} y con \py|and|, \py|or|, \py|not|, e imprimir \py|True| o \py|False|.\par
\tick Puedo unir strings con \texttt{+}, y usar \py|str( )| o un f-string para los números.\par
\tick Puedo sacar una letra con \py|palabra[0]|, contar letras con \py|len( )| y revisar con \py|in|.\par
\tick Ya sé dónde Python es más exigente que Scratch: \py|round(2.5)|, las mayúsculas y contar desde 0.

% =====================================================================
% PARTE 5: Sensores y decisiones
% =====================================================================
\partopener{sensing}{PARTE 5}{Sensores y decisiones}{Tus programas te escucharán, mirarán el reloj y tomarán decisiones.}

Hasta ahora tus programas hacían lo mismo cada vez que presionabas \playbtn. En esta parte empiezan a \textbf{sentir} cosas y a \textbf{decidir} qué hacer. Los bloques de Sensores son los ojos, los oídos y el reloj del objeto. En Python, \py|input()| son los oídos del programa, \py|distance()| es su cinta para medir y \py|time| es su cronómetro.

Vas a conocer estos bloques de Scratch y sus gemelos en Python:

\begin{itemize}
  \item \chip[sensing]{preguntar y esperar} y \chip[sensing]{respuesta} se convierten en \py|input()|
  \item \chip[control]{si entonces}, \chip[control]{si entonces si no} y \chip[control]{repetir hasta que} se convierten en \py|if|, \py|else| y \py|while|
  \item \chip[sensing]{cronómetro} y \chip[sensing]{año actual} se convierten en \py|time| y \py|datetime|
  \item \chip[sensing]{¿tocando?} y \chip[motion]{si toca un borde, rebotar} se convierten en un poquito de matemáticas con \py|distance()|, \py|xcor()| e \py|if|
\end{itemize}

Ninguno de estos programas necesita la tortuga hasta casi el final, así que la mayoría empieza directamente, sin las líneas de inicio. Recuerda que Python habla inglés: sus palabras (\py|if|, \py|while|, \py|input|\ldots) siempre van en inglés. Lo que tú inventas (los nombres de las cajas y el texto entre comillas) puede ir en español.

% ---------------------------------------------------------------------
\needspace{14\baselineskip}
\subsection*{Hacer una pregunta y esperar}

\begin{tcbraster}[raster columns=2, raster equal height=rows, raster column skip=12pt, raster valign=center, raster force size=false, raster column 1/.style={width=0.36\linewidth}, raster column 2/.style={width=0.64\linewidth-12pt}]
\begin{tcolorbox}[colback=white, colframe=white, boxrule=0pt, valign=center, left=0pt, right=0pt]
\setscratch{scale=0.8}
\begin{scratch}[scale=0.76]
\blockinit{al presionar \greenflag}
\blocksensing{preguntar \ovalnum{¿Cómo te llamas?} y esperar}
\blocklook{decir \ovaloperator{unir \ovalnum{Hola } \ovalsensing{respuesta}}}
\end{scratch}
\end{tcolorbox}
\begin{code}
nombre = input("¿Cómo te llamas? ")
print("Hola " + nombre)
\end{code}
\end{tcbraster}

\textbf{Cómo funciona.} \py|input("...")| hace el trabajo de \chip[sensing]{preguntar y esperar}. Muestra la pregunta y luego aparece una cajita \textbf{debajo de la celda}. Haz clic en la cajita, escribe y presiona \textbf{Enter}. Lo que escribiste sale de \py|input()| y se guarda en la variable que está a la izquierda del \texttt{=}. Scratch guarda la respuesta en un bloque llamado \chip[sensing]{respuesta}; en Python \textbf{tú} eliges el nombre de la caja, así que aquí se llama \py|nombre|.

\begin{outputbox}
¿Cómo te llamas? Ana\\
Hola Ana
\end{outputbox}

\begin{think}[Deja un espacio al final]
Mira el espacio que hay después del signo de interrogación en \texttt{"¿Cómo te llamas? "}. Evita que lo que escribes quede pegado a la pregunta, así: \texttt{¿Cómo te llamas?Ana}.
\end{think}

% ---------------------------------------------------------------------
\needspace{14\baselineskip}
\subsection*{input siempre te da texto}

\begin{center}
\begin{scratch}
\blockinit{al presionar \greenflag}
\blocksensing{preguntar \ovalnum{¿Cuántos años tienes?} y esperar}
\blocklook{decir \ovaloperator{unir \ovalnum{El próximo año tendrás } \ovaloperator{\ovalsensing{respuesta} + \ovalnum{1}}}}
\end{scratch}
\end{center}
\vspace{-4pt}
\begin{code}
edad = input("¿Cuántos años tienes? ")
edad = int(edad)
print("El próximo año tendrás " + str(edad + 1))
\end{code}

\textbf{Cómo funciona.} Todo lo que escribes en \py|input()| es \textbf{texto} para Python, aunque parezca un número. \texttt{"9"} es texto; \texttt{9} es un número. Scratch convierte el texto en número sin decirte nada. Python quiere que tú se lo pidas: \py|int(edad)| convierte el texto \texttt{"9"} en el número \texttt{9}. (\texttt{int} viene de \emph{integer}, que en inglés quiere decir número entero.) Después, \py|str( )| vuelve a convertir el resultado en texto para poder unirlo con \texttt{+}, igual que hiciste en la Parte 4.

Puedes hacer los dos trabajos en una sola línea, y lo vas a ver muchas veces:
\begin{codenonum}
edad = int(input("¿Cuántos años tienes? "))
\end{codenonum}

\needspace{14\baselineskip}
\begin{watchout}[¿Y si escribes una palabra?]
Ejecuta el programa y escribe \texttt{nueve} en lugar de \texttt{9}. Python no puede convertir \texttt{nueve} en un número, así que se detiene y te deja una nota con una pista. Lee primero la \textbf{última línea}:
\begin{errorbox}
ValueError: invalid literal for int() with base 10: 'nueve'
\end{errorbox}
Quiere decir: «a \texttt{int()} le dieron \texttt{'nueve'}, y eso no es un número que yo sepa leer». No se rompió nada. Presiona \playbtn{} otra vez y escribe cifras.

Python siempre escribe sus errores en inglés (no los traduce). La palabra importante es el nombre del error, al principio de la última línea: aquí, \texttt{ValueError}.
\end{watchout}

% ---------------------------------------------------------------------
\needspace{14\baselineskip}
\subsection*{Si pasa esto, haz aquello}

\begin{tcbraster}[raster columns=2, raster equal height=rows, raster column skip=12pt, raster valign=center, raster force size=false, raster column 1/.style={width=0.36\linewidth}, raster column 2/.style={width=0.64\linewidth-12pt}]
\begin{tcolorbox}[colback=white, colframe=white, boxrule=0pt, valign=center, left=0pt, right=0pt]
\setscratch{scale=0.8}
\begin{scratch}[scale=0.63]
\blockinit{al presionar \greenflag}
\blocksensing{preguntar \ovalnum{¿Te gustan los gatos?} y esperar}
\blockif{si \booloperator{\ovalsensing{respuesta} = \ovalnum{si}} entonces}{
  \blocklook{decir \ovalnum{¡A mí también!} durante \ovalnum{2} segundos}
}
\blocklook{decir \ovalnum{¡Adiós!}}
\end{scratch}
\end{tcolorbox}
\begin{code}
respuesta = input("¿Te gustan los gatos? ")
if respuesta == "si":
    print("¡A mí también!")
print("¡Adiós!")
\end{code}
\end{tcbraster}

\textbf{Cómo funciona.} \py|if| funciona igual que el bloque en forma de C \chip[control]{si entonces}. Los \textbf{dos puntos} del final abren la boca de la C, y las líneas con \textbf{cuatro espacios} delante van dentro de ella, igual que en un bucle \py|for|. Python revisa la pregunta una vez. Si es verdad, ejecuta las líneas de adentro; si no, se las salta. \texttt{print("¡Adiós!")} está otra vez pegado al borde izquierdo, así que queda \textbf{debajo} de la C y siempre se ejecuta. Recuerda: \textbf{dos} signos de igual \texttt{==} \textbf{revisan}; uno solo \texttt{=} \textbf{guarda}.

Escribe la respuesta como está en el código: \texttt{si}, sin acento, porque es más fácil de teclear. Si escribes \texttt{sí} con acento, Python no lo reconoce, porque para él son textos distintos.

\begin{watchout}[Las mayúsculas cuentan]
Escribe \texttt{Si} con S mayúscula. Python solo dice \texttt{¡Adiós!}, porque para Python \texttt{"Si"} y \texttt{"si"} son diferentes. ¡Para Scratch son iguales! Aprenderás el truco para arreglarlo, \py|.lower()|, en la Misión 3.
\end{watchout}

% ---------------------------------------------------------------------
\needspace{14\baselineskip}
\subsection*{Si pasa esto, haz aquello; si no, lo otro}

\begin{tcbraster}[raster columns=2, raster equal height=rows, raster column skip=12pt, raster valign=center, raster force size=false, raster column 1/.style={width=0.36\linewidth}, raster column 2/.style={width=0.64\linewidth-12pt}]
\begin{tcolorbox}[colback=white, colframe=white, boxrule=0pt, valign=center, left=0pt, right=0pt]
\setscratch{scale=0.8}
\begin{scratch}[scale=0.69]
\blockinit{al presionar \greenflag}
\blocksensing{preguntar \ovalnum{¿Cuántos años tienes?} y esperar}
\blockifelse{si \booloperator{\ovalsensing{respuesta} > \ovalnum{9}} entonces}{
  \blocklook{decir \ovalnum{¡Ya tienes dos cifras!}}
}{
  \blocklook{decir \ovalnum{Aún no tienes dos cifras}}
}
\end{scratch}
\end{tcolorbox}
\begin{code}
edad = int(input("¿Cuántos años tienes? "))
if edad >= 10:
    print("¡Ya tienes dos cifras!")
else:
    print("Aún no tienes dos cifras")
\end{code}
\end{tcbraster}

\textbf{Cómo funciona.} \py|else:| es la segunda boca del bloque \chip[control]{si entonces si no}. Quiere decir «si no» o «de lo contrario». \py|else| se alinea con \py|if| (sin espacios delante) y necesita sus propios dos puntos. Se ejecuta exactamente \textbf{una} de las dos bocas, nunca las dos. Scratch no tiene un bloque de «mayor o igual», así que Scratch usa \texttt{> 9}; Python puede decir \py|>= 10| («10 o más»). Para números enteros, las dos cosas significan lo mismo.

% ---------------------------------------------------------------------
\needspace{14\baselineskip}
\subsection*{Un atajo de Python: elif}

\begin{center}
\setscratch{scale=0.8}
\begin{scratch}
\blocksensing{preguntar \ovalnum{¿Cuántos puntos sacaste de 10?} y esperar}
\blockifelse{si \booloperator{\ovalsensing{respuesta} = \ovalnum{10}} entonces}{
  \blocklook{decir \ovalnum{¡Perfecto!}}
}{
  \blockifelse{si \booloperator{\ovalsensing{respuesta} > \ovalnum{6}} entonces}{
    \blocklook{decir \ovalnum{¡Muy bien!}}
  }{
    \blockifelse{si \booloperator{\ovalsensing{respuesta} > \ovalnum{3}} entonces}{
      \blocklook{decir \ovalnum{¡Buen intento!}}
    }{
      \blocklook{decir \ovalnum{¡Sigue practicando!}}
    }
  }
}
\end{scratch}
\end{center}
\vspace{-4pt}
\begin{code}
puntos = int(input("¿Cuántos puntos sacaste de 10? "))
if puntos == 10:
    print("¡Perfecto!")
elif puntos >= 7:
    print("¡Muy bien!")
elif puntos >= 4:
    print("¡Buen intento!")
else:
    print("¡Sigue practicando!")
\end{code}

\textbf{Cómo funciona.} En Scratch, para elegir entre más de dos cosas pones un \chip[control]{si entonces si no} dentro de la boca \textbf{si no} de otro, y el programa se hace cada vez más ancho. Python tiene un atajo que Scratch no tiene: \py|elif|, que es la forma corta de «else if» («si no, si\ldots»). Python revisa las preguntas \textbf{desde arriba} y ejecuta solo la \textbf{primera} que sea verdad; luego salta todas las demás. Si ninguna es verdad, se ejecuta la parte del \py|else|. Todas las líneas de \py|if|, \py|elif| y \py|else| están alineadas, así que el código queda ordenado y angosto.

\begin{think}[¿Por qué con 8 no dice también «¡Buen intento!»?]
8 es mayor que 4, así que \py|puntos >= 4| es verdad. Pero Python ya había visto que \py|puntos >= 7| era verdad, escribió \texttt{¡Muy bien!} y se saltó todo lo de abajo. Solo se ejecuta una parte.
\end{think}

% ---------------------------------------------------------------------
\needspace{14\baselineskip}
\subsection*{Repetir hasta que: Python le da la vuelta a la pregunta}

\begin{tcbraster}[raster columns=1, raster row skip=6pt]
\begin{tcolorbox}[colback=white, colframe=white, boxrule=0pt, valign=center, left=0pt, right=0pt]
\setscratch{scale=0.8}
\begin{scratch}
\blockinit{al presionar \greenflag}
\blockrepeat{repetir hasta que \booloperator{\ovalsensing{respuesta} = \ovalnum{Génesis}}}{
  \blocksensing{preguntar \ovalnum{¿Primer libro de la Biblia?} y esperar}
}
\blocklook{decir \ovalnum{¡Muy bien!}}
\end{scratch}
\end{tcolorbox}
\begin{code}
respuesta = ""
while respuesta != "Génesis":
    respuesta = input("¿Primer libro de la Biblia? ")
print("¡Muy bien!")
\end{code}
\end{tcbraster}

\textbf{Cómo funciona.} Python no tiene «repetir hasta que». Tiene \py|while|, que quiere decir «sigue \textbf{mientras} esto sea verdad». Así que le das la vuelta a la pregunta. Scratch dice: repite \textbf{hasta que} la respuesta \textbf{sea} «Génesis». Python dice: sigue \textbf{mientras} la respuesta \textbf{no sea} «Génesis». \texttt{!=} quiere decir «no es igual a». Es como «quédate afuera hasta que llueva» y «quédate afuera mientras no llueva»: el mismo plan, dicho al revés.

También puedes darle la vuelta con la palabra \py|not|, que es igual al bloque verde \chip[operators]{no}:
\begin{codenonum}
while not respuesta == "Génesis":
\end{codenonum}
Las dos líneas hacen lo mismo. Casi todos los programadores de Python usan \texttt{!=} porque es más corto.

\begin{watchout}[Escríbelo exactamente igual]
Tienes que escribir \texttt{Génesis} exactamente como está en el código: con G mayúscula y con acento en la é. \texttt{Genesis} sin acento o \texttt{génesis} con minúscula no cuentan, y el programa vuelve a preguntar. Si te cuesta escribir la é, puedes quitarle el acento en el código, \texttt{"Genesis"}, y entonces escribirlo sin acento al contestar. Lo importante es que el código y lo que escribes sean iguales.
\end{watchout}

\begin{think}[¿Por qué \texttt{respuesta = ""} en la línea 1?]
\py|while| revisa su pregunta \textbf{antes} de la primera vuelta, igual que \chip[control]{repetir hasta que}. Python necesita una caja llamada \py|respuesta| para mirar dentro. \py|""| es texto vacío: una caja que todavía no tiene nada. No es «Génesis», así que el bucle empieza.
\end{think}

% ---------------------------------------------------------------------
\needspace{14\baselineskip}
\subsection*{Tu propio cronómetro}

\begin{center}
\begin{scratch}
\blockinit{al presionar \greenflag}
\blocksensing{reiniciar cronómetro}
\blocksensing{preguntar \ovalnum{Cuenta hasta 10 y presiona Enter} y esperar}
\blocklook{decir \ovaloperator{unir \ovalnum{Tardaste } \ovalsensing{cronómetro}}}
\end{scratch}
\end{center}
\vspace{-4pt}
\begin{code}
import time
inicio = time.time()
input("Cuenta hasta 10 en silencio y presiona Enter ")
tiempo = time.time() - inicio
print("Tardaste " + str(round(tiempo, 1)) + " segundos")
\end{code}

\textbf{Cómo funciona.} Python no tiene un cronómetro que arranque solo. En cambio, \py|time.time()| te dice qué hora es ahora mismo, como un número enorme de segundos. (¡Cuenta desde el 1 de enero de 1970!) Así que tú armas tu propio cronómetro: anotas la hora al principio (\py|inicio = time.time()|, que es \chip[sensing]{reiniciar cronómetro}) y, más tarde, la restas de la hora nueva. \py|time.time() - inicio| es el bloque \chip[sensing]{cronómetro}: los segundos que han pasado desde que empezaste. \py|round(tiempo, 1)| deja una sola cifra después del punto.

\begin{think}[Reiniciar el cronómetro]
Para reiniciar tu cronómetro, solo vuelve a escribir \py|inicio = time.time()|. También puedes tener muchos cronómetros a la vez: \py|inicio1|, \py|inicio2|\ldots{} Scratch tiene solo uno.
\end{think}

% ---------------------------------------------------------------------
\needspace{14\baselineskip}
\subsection*{Esperar hasta que}

\begin{tcbraster}[raster columns=2, raster equal height=rows, raster column skip=12pt, raster valign=center, raster force size=false, raster column 1/.style={width=0.36\linewidth}, raster column 2/.style={width=0.64\linewidth-12pt}]
\begin{tcolorbox}[colback=white, colframe=white, boxrule=0pt, valign=center, left=0pt, right=0pt]
\setscratch{scale=0.8}
\begin{scratch}[scale=0.74]
\blockinit{al presionar \greenflag}
\blocksensing{reiniciar cronómetro}
\blocklook{decir \ovalnum{Espera...}}
\blockcontrol{esperar hasta que \booloperator{\ovalsensing{cronómetro} > \ovalnum{3}}}
\blocklook{decir \ovalnum{¡Pasaron 3 segundos!}}
\end{scratch}
\end{tcolorbox}
\begin{code}
import time
inicio = time.time()
print("Espera...")
while not time.time() - inicio > 3:
    time.sleep(0.1)
print("¡Pasaron 3 segundos!")
\end{code}
\end{tcbraster}

\textbf{Cómo funciona.} \chip[control]{esperar hasta que} es un bucle \py|while not| que no hace nada más que esperar. En cada vuelta, \py|time.sleep(0.1)| toma una siestecita de una décima de segundo, y luego Python vuelve a revisar la pregunta. Cuando el cronómetro pasa de 3, la pregunta \py|not ... > 3| es falsa, el bucle termina y se ejecuta la línea 6. La siesta evita que la computadora pregunte «¿ya llegamos?» millones de veces por segundo.

\begin{watchout}[Algo tiene que cambiar]
En Scratch, otro objeto podría hacer que la pregunta se vuelva verdad. En Python solo se ejecuta un programa, así que un bucle \py|while| espera para siempre si nada de lo que hay adentro puede cambiar la respuesta. Esperar al reloj, o a un \py|input()|, es seguro. Si una celda se queda ejecutando para siempre, presiona \stopbtn.
\end{watchout}

% ---------------------------------------------------------------------
\needspace{14\baselineskip}
\subsection*{¿En qué año estamos? ¿Qué hora es?}

\begin{center}
\setscratch{scale=0.85}
\begin{scratch}
\blockinit{al presionar \greenflag}
\blocklook{decir \ovaloperator{unir \ovalnum{Estamos en el año } \ovalsensing{\selectmenu{año} actual}} durante \ovalnum{2} segundos}
\blockifelse{si \booloperator{\ovalsensing{\selectmenu{hora} actual} < \ovalnum{12}} entonces}{
  \blocklook{decir \ovalnum{¡Buenos días!}}
}{
  \blocklook{decir \ovalnum{¡Buenas tardes!}}
}
\end{scratch}
\end{center}
\vspace{-4pt}
\begin{code}
import datetime
ahora = datetime.datetime.now()
print("Estamos en el año " + str(ahora.year))
if ahora.hour < 12:
    print("¡Buenos días!")
else:
    print("¡Buenas tardes!")
\end{code}

\textbf{Cómo funciona.} \py|import datetime| trae las herramientas de calendario de Python. \py|datetime.datetime.now()| lee el reloj una vez y guarda la respuesta en \py|ahora|. Después, \py|ahora.year| (año), \py|ahora.month| (mes), \py|ahora.day| (día), \py|ahora.hour| (hora), \py|ahora.minute| (minuto) y \py|ahora.second| (segundo) funcionan como el menú del bloque \chip[sensing]{año actual}. (Sí, \texttt{datetime} se escribe dos veces: el primero es la caja de herramientas y el segundo es la herramienta que está dentro.)

\begin{watchout}[El reloj de Colab puede estar en otro país]
Tu código se ejecuta en una computadora de Google, que muchas veces tiene la hora \textbf{UTC} (la hora de Londres en invierno), no la hora de tu casa. Si Colab dice \texttt{¡Buenas tardes!} a la hora del desayuno, tu código está bien: la computadora simplemente está en otro lugar del mundo. El año y la fecha casi siempre salen bien.
\end{watchout}

% ---------------------------------------------------------------------
\needspace{14\baselineskip}
\subsection*{¿Tocando? Usa una cinta para medir}

Ahora regresa la tortuga. Scratch revisa si dos dibujos se enciman. La tortuga no puede ver dibujos, pero sí puede \textbf{medir}: \py|distance(x, y)| te dice a cuántos pasos está un lugar. Si está muy cerca, digamos a menos de 20 pasos, decimos que lo está tocando.

\begin{tcbraster}[raster columns=2, raster equal height=rows, raster column skip=12pt, raster valign=center, raster force size=false, raster column 1/.style={width=0.36\linewidth}, raster column 2/.style={width=0.64\linewidth-12pt}]
\begin{tcolorbox}[colback=white, colframe=white, boxrule=0pt, valign=center, left=0pt, right=0pt]
\setscratch{scale=0.8}
\begin{scratch}[scale=0.68]
\blockinit{al presionar \greenflag}
\blockmove{ir a x: \ovalnum{-200} y: \ovalnum{0}}
\blockrepeat{repetir hasta que \boolsensing{¿tocando \selectmenu{Mancha}?}}{
  \blockmove{mover \ovalnum{10} pasos}
}
\blocklook{decir \ovalnum{¡La toqué!}}
\end{scratch}
\end{tcolorbox}
\begin{code}
from ColabTurtlePlus.Turtle import *
clearscreen()
initializeTurtle((480, 360))
dot(20, "red")
penup()
goto(-200, 0)
pendown()
while distance(0, 0) >= 20:
    forward(10)
print("¡La toqué!")
\end{code}
\end{tcbraster}

\textbf{Cómo funciona.} La línea 4 dibuja una mancha roja en el centro, en $(0, 0)$. Los colores se escriben en inglés: \texttt{"red"} es rojo, \texttt{"blue"} azul, \texttt{"gold"} dorado\ldots{} \py|distance(0, 0)| mide desde la tortuga hasta esa mancha. «Repetir \textbf{hasta que} toque» se voltea y queda «sigue \textbf{mientras} la distancia sea 20 o más». Es como imaginar que cada objeto tiene una burbuja invisible alrededor.

\needspace{20\baselineskip}
Para saber si dos \textbf{objetos} se tocan, crea una segunda tortuga con \py|Turtle()| y mide hasta donde está. \py|perro.pos()| te da la posición del perro, su $x$ y su $y$ juntas.

\begin{tcbraster}[raster columns=2, raster equal height=rows, raster column skip=12pt, raster valign=center, raster force size=false, raster column 1/.style={width=0.36\linewidth}, raster column 2/.style={width=0.64\linewidth-12pt}]
\begin{tcolorbox}[colback=white, colframe=white, boxrule=0pt, valign=center, left=0pt, right=0pt]
\setscratch{scale=0.8}
\begin{scratch}[scale=0.71]
\blockinit{al presionar \greenflag}
\blockrepeat{repetir hasta que \boolsensing{¿tocando \selectmenu{Perro}?}}{
  \blockmove{mover \ovalnum{10} pasos}
}
\blocklook{decir \ovalnum{¡Te atrapé, perro!}}
\end{scratch}
\end{tcolorbox}
\begin{code}
from ColabTurtlePlus.Turtle import *
clearscreen()
initializeTurtle((480, 360))
shape("turtle")
perro = Turtle()
perro.shape("circle")
perro.penup()
perro.goto(150, 0)
while distance(perro.pos()) >= 20:
    forward(10)
print("¡Te atrapé, perro!")
\end{code}
\end{tcbraster}

\textbf{Cómo funciona.} \py|perro = Turtle()| crea un objeto nuevo y le pone el nombre \py|perro|. Los comandos que llevan \py|perro.| delante mueven solo al perro; los comandos solitos mueven a la primera tortuga, como siempre. La línea 9 mide desde la tortuga hasta la posición del perro. Si creas dos objetos con nombre, la revisión se ve así: \py|gato.distance(perro.pos()) < 20|.

\begin{watchout}[No olvides \texttt{.pos()}]
En ColabTurtlePlus tienes que medir hasta la \textbf{posición} del perro, \py|perro.pos()|. Si escribes \py|distance(perro)| sale un \texttt{ValueError}, porque la cinta para medir quiere números, no un objeto entero.
\end{watchout}

% ---------------------------------------------------------------------
\needspace{14\baselineskip}
\subsection*{¿Tecla presionada? ¿Ratón presionado? Mejor por turnos}

\begin{tcbraster}[raster columns=2, raster equal height=rows, raster column skip=12pt, raster valign=center, raster force size=false, raster column 1/.style={width=0.36\linewidth}, raster column 2/.style={width=0.64\linewidth-12pt}]
\begin{tcolorbox}[colback=white, colframe=white, boxrule=0pt, valign=center, left=0pt, right=0pt]
\setscratch{scale=0.7}
\begin{scratch}[scale=0.51]
\blockinit{al presionar \greenflag}
\blockrepeat{repetir \ovalnum{5}}{
  \blockif{si \boolsensing{¿tecla \selectmenu{flecha izquierda} presionada?} entonces}{
    \blockmove{girar \turnleft{} \ovalnum{90} grados}
  }
  \blockif{si \boolsensing{¿tecla \selectmenu{flecha derecha} presionada?} entonces}{
    \blockmove{girar \turnright{} \ovalnum{90} grados}
  }
  \blockmove{mover \ovalnum{50} pasos}
}
\end{scratch}
\end{tcolorbox}
\begin{code}
from ColabTurtlePlus.Turtle import *
clearscreen()
initializeTurtle((480, 360))
shape("turtle")
for i in range(5):
    tecla = input("¿a, i o d? ")
    if tecla == "i":
        left(90)
    elif tecla == "d":
        right(90)
    forward(50)
\end{code}
\end{tcbraster}

\textbf{Cómo funciona.} La tortuga de Colab no puede oír el teclado ni sentir el ratón, así que \chip[sensing]{¿tecla presionada?} y \chip[sensing]{¿ratón presionado?} no funcionan ahí. En cambio, jugamos por \textbf{turnos}: en cada vuelta del bucle, \py|input()| espera a que escribas una letra y presiones Enter, y \py|if| y \py|elif| deciden qué hacer. \texttt{i} es izquierda y \texttt{d} es derecha. Cualquier otra letra, como \texttt{a} (de «adelante»), solo avanza. (Los juegos de verdad con teclas vienen más adelante, en Thonny.)

% ---------------------------------------------------------------------
\needspace{14\baselineskip}
\subsection*{Haz tu propio «si toca un borde, rebotar»}

\begin{tcbraster}[raster columns=2, raster equal height=rows, raster column skip=12pt, raster valign=center, raster force size=false, raster column 1/.style={width=0.28\linewidth}, raster column 2/.style={width=0.72\linewidth-12pt}]
\begin{tcolorbox}[colback=white, colframe=white, boxrule=0pt, valign=center, left=0pt, right=0pt]
\setscratch{scale=0.8}
\begin{scratch}
\blockinit{al presionar \greenflag}
\blockmove{apuntar en dirección \ovalnum{60}}
\blockrepeat{repetir \ovalnum{300}}{
  \blockmove{mover \ovalnum{10} pasos}
  \blockmove{si toca un borde, rebotar}
}
\end{scratch}
\end{tcolorbox}
\begin{code}
from ColabTurtlePlus.Turtle import *
clearscreen()
initializeTurtle((480, 360))
speed(13)
setheading(30)
for i in range(300):
    forward(10)
    if xcor() > 230 or xcor() < -230:
        setheading(180 - heading())
    if ycor() > 170 or ycor() < -170:
        setheading(0 - heading())
\end{code}
\end{tcbraster}

\textbf{Cómo funciona.} La tortuga no tiene bloque de rebotar ni paredes: se saldría del escenario tan tranquila. Así que construimos las paredes con \py|if|. \py|xcor()| es la posición $x$ de la tortuga. Si pasa de 230 (la pared derecha) o de $-230$ (la pared izquierda), la hacemos dar la vuelta. \py|heading()| te dice hacia dónde apunta la tortuga en este momento. La línea 5 usa la regla de 90 menos: la dirección 60 de Scratch es \py|90 - 60| en Python, o sea 30.

\needspace{12\baselineskip}
\begin{think}[Las matemáticas del rebote]
Una pelota que rebota en una pared de un lado sigue subiendo (o bajando), pero cambia derecha por izquierda. Eso es \py|180 - heading()|. Cuando rebota en el piso o en el techo, sigue yendo a la derecha (o a la izquierda), pero cambia arriba por abajo. Eso es \py|0 - heading()|.
\begin{center}
\renewcommand{\arraystretch}{1.25}
\begin{tabular}{lll}
\toprule
\textbf{Apunta hacia} & \textbf{Choca con} & \textbf{Nueva dirección} \\
\midrule
30 (arriba y derecha) & pared derecha & $180 - 30 = 150$ (arriba e izquierda) \\
150 (arriba e izquierda) & techo & $0 - 150 = -150$ (abajo e izquierda) \\
$-150$ (abajo e izquierda) & pared izquierda & $180 - (-150) = 330$ (abajo y derecha) \\
\bottomrule
\end{tabular}
\end{center}
\end{think}

% =====================================================================
\mission{Di hola}{\texttt{input}, a unir palabras y a guardar más de una respuesta}

La computadora te pregunta tu nombre y tu comida favorita, y luego le cuenta a todos lo que te encanta.

\step{Paso 1: Constrúyelo en Scratch}
Crea una variable llamada \textbf{nombre} y luego construye:

\begin{center}
\begin{scratch}
\blockinit{al presionar \greenflag}
\blocksensing{preguntar \ovalnum{¿Cómo te llamas?} y esperar}
\blockvariable{fijar \selectmenu{nombre} a \ovalsensing{respuesta}}
\blocklook{decir \ovaloperator{unir \ovalnum{Hola } \ovalvariable{nombre}} durante \ovalnum{2} segundos}
\blocksensing{preguntar \ovalnum{¿Cuál es tu comida favorita?} y esperar}
\blocklook{decir \ovaloperator{unir \ovalvariable{nombre} \ovaloperator{unir \ovalnum{ adora comer } \ovalsensing{respuesta}}}}
\end{scratch}
\end{center}

¿Para qué la variable \textbf{nombre}? Scratch tiene una sola \chip[sensing]{respuesta}. La segunda pregunta borra la primera respuesta, así que primero tienes que guardar el nombre en un lugar seguro.

\step{Paso 2: Escríbelo en Colab}
Este programa no dibuja, así que no necesitas las líneas de inicio.

\begin{typethis}
nombre = input("¿Cómo te llamas? ")
print("Hola " + nombre)
comida = input("¿Cuál es tu comida favorita? ")
print(nombre + " adora comer " + comida + ".")
\end{typethis}

\step{Paso 3: Léelo línea por línea}

\begin{longtable}{>{\centering\arraybackslash}p{1.1cm} P{5.3cm} P{9.1cm}}
\toprule
\textbf{Línea} & \textbf{Bloque de Scratch} & \textbf{Qué significa} \\
\midrule
--- & \chip[events]{al presionar bandera verde} & No hace falta código. El botón de reproducir \playbtn{} lo pone en marcha. \\
1 & \chip[sensing]{preguntar y esperar}\newline\chip[variables]{fijar nombre a respuesta} & ¡Dos bloques en una sola línea! \py|input()| pregunta y espera, y la respuesta va directo a la caja llamada \py|nombre|. \\
2 & \chip[looks]{decir unir Hola nombre} & \texttt{+} une textos, como \chip[operators]{unir}. \\
3 & \chip[sensing]{preguntar y esperar} & Una segunda pregunta con su propia caja, \py|comida|. Python puede guardar todas las respuestas que quieras. \\
4 & \chip[looks]{decir unir nombre unir \ldots} & Cuatro piezas unidas en fila. Scratch necesita un \chip[operators]{unir} dentro de otro \chip[operators]{unir}; Python solo agrega otro \texttt{+}. \\
\bottomrule
\end{longtable}

\step{Paso 4: Ejecútalo}
Presiona \playbtn. Escribe tu nombre en la cajita y presiona Enter; luego escribe una comida y presiona Enter.

\begin{outputbox}
¿Cómo te llamas? Ana\\
Hola Ana\\
¿Cuál es tu comida favorita? pizza\\
Ana adora comer pizza.
\end{outputbox}

\begin{think}
Mira los espacios dentro de \texttt{" adora comer "}: uno antes y uno después. Sin ellos te saldría \texttt{Anaadora comerpizza.}
\end{think}

\begin{challenge}
\tick Haz una tercera pregunta, por ejemplo tu animal favorito, y agrégala a la última oración.\par
\tick Pregunta \texttt{edad = int(input("¿Cuántos años tienes? "))} y escribe cuántos años tendrás dentro de 10 años. ¡No olvides \py|str( )| cuando unas un número!\par
\tick Usa un \py|if| para dar una respuesta especial cuando la comida sea \texttt{"brócoli"}.\par
\tick \textbf{Súper reto:} quita el \py|int( )| de la pregunta de la edad y ejecútalo. Lee la última línea del texto rojo. ¿Qué tipo de error es?
\end{challenge}

% =====================================================================
\mission{Adivina mi número}{\texttt{random}, \texttt{while}, \texttt{if} y \texttt{elif}, \texttt{int}, y a contar intentos}

La computadora elige un número secreto del 1 al 10. Tú sigues adivinando hasta acertar. Te dice «muy bajo» o «muy alto» y, al final, te dice cuántos intentos necesitaste.

\step{Paso 1: Constrúyelo en Scratch}
Crea tres variables llamadas \textbf{secreto}, \textbf{intento} y \textbf{veces}, y luego construye:

\begin{center}
\setscratch{scale=0.8}
\begin{scratch}
\blockinit{al presionar \greenflag}
\blockvariable{fijar \selectmenu{secreto} a \ovaloperator{elegir número al azar entre \ovalnum{1} y \ovalnum{10}}}
\blockvariable{fijar \selectmenu{intento} a \ovalnum{0}}
\blockvariable{fijar \selectmenu{veces} a \ovalnum{0}}
\blockrepeat{repetir hasta que \booloperator{\ovalvariable{intento} = \ovalvariable{secreto}}}{
  \blocksensing{preguntar \ovalnum{Adivina del 1 al 10:} y esperar}
  \blockvariable{fijar \selectmenu{intento} a \ovalsensing{respuesta}}
  \blockvariable{cambiar \selectmenu{veces} en \ovalnum{1}}
  \blockifelse{si \booloperator{\ovalvariable{intento} < \ovalvariable{secreto}} entonces}{
    \blocklook{decir \ovalnum{Muy bajo} durante \ovalnum{1} segundos}
  }{
    \blockif{si \booloperator{\ovalvariable{intento} > \ovalvariable{secreto}} entonces}{
      \blocklook{decir \ovalnum{Muy alto} durante \ovalnum{1} segundos}
    }
  }
}
\blocklook{decir \ovaloperator{unir \ovalnum{¡Lo adivinaste en } \ovaloperator{unir \ovalvariable{veces} \ovalnum{ intentos!}}}}
\end{scratch}
\end{center}

\step{Paso 2: Escríbelo en Colab}

\begin{typethis}
import random
secreto = random.randint(1, 10)
intento = 0
veces = 0
while intento != secreto:
    intento = int(input("Adivina del 1 al 10: "))
    veces += 1
    if intento < secreto:
        print("Muy bajo")
    elif intento > secreto:
        print("Muy alto")
print("¡Lo adivinaste en " + str(veces) + " intentos!")
\end{typethis}

\step{Paso 3: Léelo línea por línea}

\begin{longtable}{>{\centering\arraybackslash}p{1.1cm} P{5.3cm} P{9.1cm}}
\toprule
\textbf{Línea} & \textbf{Bloque de Scratch} & \textbf{Qué significa} \\
\midrule
1 & (Agregar Extensión) & Trae las herramientas de números al azar de Python. \\
2 & \chip[variables]{fijar secreto a}\newline\chip[operators]{elegir número al azar} & \py|random.randint(1, 10)| elige un número entero del 1 al 10. Puede elegir el 1 o el 10. \\
3--4 & \chip[variables]{fijar intento a 0}\newline\chip[variables]{fijar veces a 0} & Crea las dos cajas. \py|intento| empieza en 0, que nunca puede ser el secreto, así que el bucle seguro empieza. \\
5 & \chip[control]{repetir hasta que}\newline\chip[operators]{intento = secreto} & ¡Al revés! Sigue \textbf{mientras} el intento \textbf{no sea} el secreto. \\
6 & \chip[sensing]{preguntar}\newline\chip[variables]{fijar intento a respuesta} & Pregunta, convierte el texto en número con \py|int( )| y lo guarda en \py|intento|. \\
7 & \chip[variables]{cambiar veces en 1} & \py|veces += 1| es la forma corta de escribir \py|veces = veces + 1|. \\
8--9 & \chip[control]{si intento < secreto entonces} & ¿Muy chico? Dilo. \\
10--11 & \chip[control]{si no} \chip[control]{si intento > secreto entonces} & \py|elif| quiere decir «si no, si\ldots». Si el intento no es ni muy bajo ni muy alto, no se escribe nada, ¡porque tiene que ser el correcto! \\
12 & \chip[looks]{decir unir ¡Lo adivinaste en \ldots} & No tiene espacios delante, así que se ejecuta una sola vez, después del bucle. \py|str(veces)| convierte el número en texto para poder unirlo. \\
\bottomrule
\end{longtable}

\step{Paso 4: Ejecútalo}
Presiona \playbtn{} y escribe un número en la cajita. Presiona Enter después de cada intento.

\begin{outputbox}[Podrías ver]
Adivina del 1 al 10: 5\\
Muy bajo\\
Adivina del 1 al 10: 8\\
Muy alto\\
Adivina del 1 al 10: 7\\
¡Lo adivinaste en 3 intentos!
\end{outputbox}

\begin{think}[¿Puedes ganar siempre en 4 intentos?]
Siempre adivina el número de en medio de los que quedan. Del 1 al 10: di 5. ¿Muy bajo? Ahora quedan del 6 al 10: di 8. Este truco se llama \textbf{búsqueda binaria}, y los programadores de verdad lo usan todo el tiempo.
\end{think}

\begin{challenge}
\tick Hazlo más difícil: elige un número del 1 al \textbf{100}. ¡Cambia también la pregunta!\par
\tick Cambia la última línea para que diga \texttt{intento} (no \texttt{intentos}) cuando aciertas a la primera. Usa \py|if veces == 1:| y \py|else:|.\par
\tick Dale al jugador solo 5 intentos. Cambia la línea 5 por \texttt{while intento != secreto and veces < 5:}. Después del bucle, usa \py|if|/\py|else| para decir si ganó o si se le acabaron los intentos, y dile cuál era el secreto.\par
\tick Escribe \texttt{gato} como intento. Lee la última línea del texto rojo. ¿Qué salió mal?
\end{challenge}

% =====================================================================
\mission{Concurso con puntos}{\texttt{if}/\texttt{else}, \texttt{.lower()} y a llevar la cuenta con \texttt{+=}}

Haz un concurso de tres preguntas. Cada respuesta correcta suma un punto y, al final, la computadora te dice cómo te fue.

\step{Paso 1: Constrúyelo en Scratch}
Crea una variable llamada \textbf{puntos}. Aquí está la primera pregunta. Las preguntas 2 y 3 se construyen igual, una debajo de la otra.

\begin{center}
\setscratch{scale=0.8}
\begin{scratch}
\blockinit{al presionar \greenflag}
\blockvariable{fijar \selectmenu{puntos} a \ovalnum{0}}
\blocksensing{preguntar \ovalnum{¿Cuánto es 6 x 7?} y esperar}
\blockifelse{si \booloperator{\ovalsensing{respuesta} = \ovalnum{42}} entonces}{
  \blocklook{decir \ovalnum{¡Correcto!} durante \ovalnum{1} segundos}
  \blockvariable{cambiar \selectmenu{puntos} en \ovalnum{1}}
}{
  \blocklook{decir \ovalnum{¡Uy! Era 42.} durante \ovalnum{1} segundos}
}
\blockspace[0.6]
\blocklook{decir \ovaloperator{unir \ovalnum{Sacaste } \ovaloperator{unir \ovalvariable{puntos} \ovalnum{ de 3}}}}
\end{scratch}
\end{center}

\step{Paso 2: Escríbelo en Colab}

\begin{typethis}
puntos = 0

respuesta = int(input("¿Cuánto es 6 x 7? "))
if respuesta == 42:
    print("¡Correcto!")
    puntos += 1
else:
    print("¡Uy! Era 42.")

respuesta = input("¿De qué color es un plátano? ")
if respuesta.lower() == "amarillo":
    print("¡Correcto!")
    puntos += 1
else:
    print("¡Uy! Era amarillo.")

respuesta = input("¿Qué animal hace muu? ")
if respuesta.lower() == "vaca":
    print("¡Correcto!")
    puntos += 1
else:
    print("¡Uy! Era una vaca.")

print("Sacaste " + str(puntos) + " de 3")
\end{typethis}

\step{Paso 3: Léelo línea por línea}

\begin{longtable}{>{\centering\arraybackslash}p{1.1cm} P{5.3cm} P{9.1cm}}
\toprule
\textbf{Línea} & \textbf{Bloque de Scratch} & \textbf{Qué significa} \\
\midrule
1 & \chip[variables]{fijar puntos a 0} & Los puntos empiezan en 0. Las líneas vacías (2, 9, 16, 23) no hacen nada; solo hacen que el código sea más fácil de leer. \\
3 & \chip[sensing]{preguntar y esperar} & La respuesta debe ser un número, así que \py|int( )| la convierte en uno. \\
4 & \chip[control]{si respuesta = 42 entonces} & ¿El número es 42? Dos signos de igual para revisar. \\
5--6 & \chip[looks]{decir ¡Correcto!}\newline\chip[variables]{cambiar puntos en 1} & Las dos líneas están dentro del \py|if|. \py|puntos += 1| suma un punto. \\
7--8 & \chip[control]{si no} & Si no, le dice cuál era la respuesta correcta. \\
10 & \chip[sensing]{preguntar} & Esta vez la respuesta es una palabra, así que no hay \py|int( )|. Volvemos a usar la caja \py|respuesta|; la respuesta nueva reemplaza a la vieja, igual que en Scratch. \\
11 & \chip[control]{si ... entonces}\newline\chip[operators]{respuesta = amarillo} & \py|respuesta.lower()| primero pasa la respuesta a minúsculas, así que \texttt{Amarillo}, \texttt{AMARILLO} y \texttt{amarillo} cuentan todas. \\
12--15 & & El mismo patrón que las líneas 5 a 8. \\
17--22 & & La pregunta 3, construida igual. \\
24 & \chip[looks]{decir unir Sacaste \ldots} & Está pegada al borde izquierdo, así que siempre se ejecuta al final. \\
\bottomrule
\end{longtable}

\begin{watchout}[¿Para qué sirve \texttt{.lower()}?]
Scratch piensa que \texttt{Amarillo} y \texttt{amarillo} son iguales. Python no: ¡las mayúsculas cuentan! \py|.lower()| hace una copia del texto en minúsculas, así que no importa cómo escriba el jugador. Asegúrate de que la palabra con la que comparas, como \py|"amarillo"|, también esté en minúsculas, o nunca va a coincidir.
\end{watchout}

\step{Paso 4: Ejecútalo}

\begin{outputbox}[Podrías ver]
¿Cuánto es 6 x 7? 42\\
¡Correcto!\\
¿De qué color es un plátano? Amarillo\\
¡Correcto!\\
¿Qué animal hace muu? oveja\\
¡Uy! Era una vaca.\\
Sacaste 2 de 3
\end{outputbox}

\begin{challenge}
\tick Cambia las preguntas por algo que te encante: futbol, dinosaurios, el espacio\ldots\par
\tick Agrega una cuarta pregunta. Acuérdate de cambiar el \texttt{3} por \texttt{4} en la última línea.\par
\tick Al final, usa \py|if|, \py|elif| y \py|else| para dar un mensaje: \texttt{¡Campeón del concurso!} si sacas 3, \texttt{¡Bien hecho!} si sacas 1 o 2, y \texttt{¡Más suerte la próxima vez!} si sacas 0.\par
\tick \textbf{Súper reto:} mide con un cronómetro cuánto tardas en todo el concurso, y escribe cuántos segundos fueron.
\end{challenge}

% =====================================================================
\mission{Prueba de reflejos}{el cronómetro, \texttt{time.sleep}, \texttt{random} e \texttt{if}/\texttt{else}}

¿Qué tan rápido eres? La computadora espera un tiempo sorpresa y luego grita ¡YA! Presiona Enter lo más rápido que puedas, y ella mide cuánto tardaste.

\step{Paso 1: Constrúyelo en Scratch}
Crea una variable llamada \textbf{tiempo} y luego construye:

\begin{center}
\setscratch{scale=0.8}
\begin{scratch}
\blockinit{al presionar \greenflag}
\blocksensing{preguntar \ovalnum{Presiona Enter para empezar} y esperar}
\blocklook{decir \ovalnum{Prepárate...}}
\blockcontrol{esperar \ovaloperator{elegir número al azar entre \ovalnum{2} y \ovalnum{5}} segundos}
\blocksensing{reiniciar cronómetro}
\blocksensing{preguntar \ovalnum{¡YA! ¡Presiona Enter!} y esperar}
\blockvariable{fijar \selectmenu{tiempo} a \ovalsensing{cronómetro}}
\blocklook{decir \ovaloperator{unir \ovalnum{Tu tiempo: } \ovalvariable{tiempo}} durante \ovalnum{2} segundos}
\blockifelse{si \booloperator{\ovalvariable{tiempo} < \ovalnum{0.5}} entonces}{
  \blocklook{decir \ovalnum{¡Rápido como un rayo!}}
}{
  \blocklook{decir \ovalnum{¡Bien! Juega otra vez para mejorar.}}
}
\end{scratch}
\end{center}

\step{Paso 2: Escríbelo en Colab}

\begin{typethis}
import time
import random
input("Presiona Enter para empezar ")
print("Prepárate...")
time.sleep(random.randint(2, 5))
inicio = time.time()
input("¡YA! ¡Presiona Enter! ")
tiempo = time.time() - inicio
tiempo = round(tiempo, 2)
print("Tu tiempo: " + str(tiempo) + " segundos")
if tiempo < 0.5:
    print("¡Rápido como un rayo!")
else:
    print("¡Bien! Juega otra vez para mejorar.")
\end{typethis}

\step{Paso 3: Léelo línea por línea}

\begin{longtable}{>{\centering\arraybackslash}p{1.1cm} P{5.3cm} P{9.1cm}}
\toprule
\textbf{Línea} & \textbf{Bloque de Scratch} & \textbf{Qué significa} \\
\midrule
1--2 & (Agregar Extensión) & Traen las herramientas del reloj y las de números al azar. \\
3 & \chip[sensing]{preguntar y esperar} & No guardamos la respuesta, así que no hay \py|respuesta =|. Usamos \py|input()| solo para que espere. \\
4 & \chip[looks]{decir Prepárate...} & \\
5 & \chip[control]{esperar ... segundos}\newline\chip[operators]{elegir número al azar} & \py|time.sleep( )| es \chip[control]{esperar}. Con una espera al azar no puedes adivinar cuándo viene el ¡YA! \\
6 & \chip[sensing]{reiniciar cronómetro} & Arranca el cronómetro. \\
7 & \chip[sensing]{preguntar ¡YA! y esperar} & Espera hasta que presiones Enter. \\
8 & \chip[variables]{fijar tiempo a cronómetro} & La hora de ahora menos la del inicio: los segundos que tardaste. \\
9 & & \py|round(tiempo, 2)| deja dos cifras después del punto, así que \texttt{0.4371825} se vuelve \texttt{0.44}. \\
10 & \chip[looks]{decir unir Tu tiempo: tiempo} & Convierte el número en texto y lo une. \\
11--14 & \chip[control]{si entonces si no} & ¿Menos de medio segundo? ¡Como un rayo! Si no, te anima a jugar otra vez. \\
\bottomrule
\end{longtable}

\step{Paso 4: Ejecútalo}
Presiona \playbtn{} y luego presiona Enter en la cajita. Prepárate y, cuando veas ¡YA!, presiona Enter lo más rápido que puedas.

\begin{outputbox}[Podrías ver]
Presiona Enter para empezar\\
Prepárate...\\
¡YA! ¡Presiona Enter!\\
Tu tiempo: 0.62 segundos\\
¡Bien! Juega otra vez para mejorar.
\end{outputbox}

\begin{think}[¿Por qué Colab es un poco lento?]
Colab se ejecuta en una computadora que está lejos. Cuando presionas Enter, el aviso tiene que viajar hasta allá por internet, así que tu tiempo incluye el viaje. Compara tus tiempos entre sí, no con los de un amigo en otra computadora. Si \texttt{0.5} es demasiado difícil, cambia la línea 11 a \texttt{0.8}.
\end{think}

\begin{challenge}
\tick Primero pregunta el nombre del jugador, y escríbelo junto con su tiempo.\par
\tick Agrega un \py|elif| para que haya tres mensajes: menos de 0.5, menos de 1, y todo lo demás.\par
\tick \textbf{El mejor de tres:} pon las líneas 3 a 8 dentro de \py|for i in range(3):| (con cuatro espacios más). Antes del bucle, escribe \py|mejor = 100|. Adentro, después de la línea 8, agrega \py|if tiempo < mejor:| y \py|mejor = tiempo|. Al final, escribe el mejor tiempo.\par
\tick Hazlo más tramposo: cambia la espera a \py|random.randint(1, 8)|.
\end{challenge}

% =====================================================================
\needspace{30\baselineskip}
\subsection*{Tradúcelo tú}
Aquí hay un programa de Scratch que no has visto antes. Escribe el Python en las líneas, luego escríbelo en Colab y ejecútalo. ¿Hace lo mismo?

\begin{tcbraster}[raster columns=2, raster equal height=rows, raster column skip=12pt, raster valign=top]
\begin{tcolorbox}[colback=white, colframe=white, boxrule=0pt]
\setscratch{scale=0.85}
\begin{scratch}[scale=0.67]
\blockinit{al presionar \greenflag}
\blocksensing{preguntar \ovalnum{¿Cuál es la contraseña?} y esperar}
\blockrepeat{repetir hasta que \booloperator{\ovalsensing{respuesta} = \ovalnum{banana}}}{
  \blocklook{decir \ovalnum{¡Mal! Intenta otra vez.} durante \ovalnum{1} segundos}
  \blocksensing{preguntar \ovalnum{¿Cuál es la contraseña?} y esperar}
}
\blocklook{decir \ovalnum{¡Adelante!}}
\end{scratch}
\end{tcolorbox}
\begin{tcolorbox}[colback=white, colframe=codeframe, arc=5pt, title=\textbf{Mi Python}, coltitle=ink, colbacktitle=codebg]
\vspace{4pt}
\foreach \n in {1,...,7}{\rule{\linewidth}{0.4pt}\par\vspace{9pt}}
\end{tcolorbox}
\end{tcbraster}

\begin{think}[Dos pistas]
\chip[control]{repetir hasta que} se voltea y queda \py|while| \ldots\ \texttt{!=}. Y para \chip[looks]{decir ¡Mal! Intenta otra vez.} basta con un \py|print|.
\end{think}

% ---------------------------------------------------------------------
\needspace{12\baselineskip}
\begin{puzzle}
Escribe \textbf{dos líneas} de Python que pregunten «¿Cuántos años tienes?» y luego escriban cuántos años tendrás el próximo año, por ejemplo \texttt{El próximo año tendrás 10}.

\vspace{4pt}
\rule{\linewidth}{0.4pt}\par\vspace{12pt}
\rule{\linewidth}{0.4pt}
\end{puzzle}

% ---------------------------------------------------------------------
\needspace{10\baselineskip}
\subsection*{Lista de la parte 5}
\begingroup\setlength{\parskip}{5pt}
\tick Puedo hacer una pregunta con \py|input()| y guardar la respuesta en una variable.\par
\tick Puedo convertir texto en número con \py|int( )|, y sé que un \texttt{ValueError} quiere decir que escribí una palabra.\par
\tick Puedo tomar decisiones con \py|if|, \py|else| y \py|elif|, con dos puntos y cuatro espacios.\par
\tick Puedo voltear un «repetir hasta que» y convertirlo en un bucle \py|while| con \texttt{!=} o con \py|not|.\par
\tick Puedo hacer un cronómetro con \py|time.time()| y leer el año con \py|datetime|.\par
\tick Ya sé revisar si algo se «toca» con \py|distance()| y hacer que una tortuga rebote en las paredes.\par
\endgroup

% =====================================================================
% PARTE 6: Mis Bloques (funciones, return, objetos, clones)
% =====================================================================
\partopener{myblocks}{PARTE 6}{Mis Bloques}{Inventa tus propios bloques de Python, haz que te devuelvan respuestas y llena el escenario de objetos.}

En Scratch, el botón rosa \textbf{Mis Bloques} te deja inventar un bloque totalmente nuevo, como \chip[myblocks]{cuadrado} o \chip[myblocks]{saltar}. Lo construyes una vez y luego puedes usarlo una y otra vez. En Python, un bloque hecho en casa se llama \textbf{función}. Piensa en él como una \textbf{tarjeta de receta}: escribes los pasos una sola vez, le pones un nombre a la tarjeta y después solo dices el nombre cada vez que quieras esos pasos.

En esta parte vas a conocer:
\begin{itemize}
  \item \chip[myblocks]{definir cuadrado (lado)} \quad se convierte en \quad \py|def cuadrado(lado):|
  \item \chip[myblocks]{cuadrado 100} \quad se convierte en \quad \py|cuadrado(100)|
  \item \chip[control]{crear clon de mí mismo} \quad se convierte en \quad \py|s = Turtle()| dentro de un bucle
\end{itemize}

Y vas a aprender un superpoder de Python que los bloques rosas de Scratch no tienen: un bloque que \textbf{te devuelve una respuesta}.

\setscratch{scale=0.85}

% ---------------------------------------------------------------------
\needspace{14\baselineskip}
\subsection*{Crea un bloque: \texttt{def}}

\begin{tcbraster}[raster columns=2, raster equal height=rows, raster column skip=12pt, raster valign=center, raster force size=false, raster column 1/.style={add to width=-2cm}, raster column 2/.style={add to width=2cm}]
\begin{tcolorbox}[colback=white, colframe=white, boxrule=0pt, valign=center]
\begin{scratch}[scale=0.71]
\initmoreblocks{definir \namemoreblocks{hola}}
\blocklook{decir \ovalnum{¡Hola!} durante \ovalnum{1} segundos}
\blocklook{decir \ovalnum{¡Mi bloque nuevo!}}
\end{scratch}
\end{tcolorbox}
\begin{code}
def hola():
    print("¡Hola!")
    print("¡Mi bloque nuevo!")
\end{code}
\end{tcbraster}

\textbf{Cómo funciona.} \py|def| es la forma corta de \textbf{define} («definir» en inglés), así que \py|def hola():| es el sombrero rosa \chip[myblocks]{definir hola}. Python habla inglés: sus palabras (\py|def|, \py|for|, \py|print|, \py|forward|...) siempre van en inglés, pero el nombre de tu bloque lo inventas tú, ¡y puede ser en español! Después viene el nombre y luego un par de paréntesis (vacíos por ahora, porque este bloque no tiene entradas). Los \textbf{dos puntos} abren la boca del bloque, igual que los dos puntos después de \py|for| o \py|if|. Cada línea con cuatro espacios delante está \textbf{dentro} del bloque. La primera línea que vuelve al borde izquierdo ya no forma parte de él.

\begin{think}[¿No pasó nada?]
Si ejecutas solo estas tres líneas, \textbf{no aparece nada}. ¡Y eso está bien! En Scratch, un sombrero \chip[myblocks]{definir} solo tampoco hace nada. Ya escribiste la tarjeta de receta, pero nadie la ha cocinado todavía. Sigue leyendo para descubrir cómo.
\end{think}

% ---------------------------------------------------------------------
\needspace{14\baselineskip}
\subsection*{Usa tu bloque: los paréntesis lo ponen en marcha}

\begin{tcbraster}[raster columns=2, raster equal height=rows, raster column skip=12pt, raster valign=center, raster force size=false, raster column 1/.style={add to width=-2cm}, raster column 2/.style={add to width=2cm}]
\begin{tcolorbox}[colback=white, colframe=white, boxrule=0pt, valign=center]
\begin{scratch}
\blockinit{al presionar \greenflag}
\blockmoreblocks{hola}
\blockmoreblocks{hola}
\end{scratch}
\end{tcolorbox}
\begin{code}
def hola():
    print("¡Hola!")
    print("¡Mi bloque nuevo!")

hola()
hola()
\end{code}
\end{tcbraster}

\textbf{Cómo funciona.} Escribir el nombre con paréntesis, \py|hola()|, \textbf{llama} al bloque. Es lo mismo que encajar el bloque rosa \chip[myblocks]{hola} en un programa. Python salta hasta el \py|def|, hace cada línea de adentro y luego regresa para seguir. Aquí se ejecuta dos veces, así que ves las dos líneas dos veces.

\begin{watchout}[¡Cuidado: sin paréntesis, no hay acción!]
Si escribes \py|hola| sin los paréntesis, Python \textbf{no hace nada}, ¡y tampoco muestra ningún error! Los paréntesis son el botón de «¡ya!». Sin ellos, solo dijiste el nombre del bloque, como si señalaras la tarjeta de receta sin cocinar. Cuando parezca que Python ignora un bloque, revisa los \texttt{()}.
\end{watchout}

% ---------------------------------------------------------------------
\needspace{14\baselineskip}
\subsection*{Una entrada numérica}

\begin{tcbraster}[raster columns=2, raster equal height=rows, raster column skip=12pt, raster valign=center, raster force size=false, raster column 1/.style={add to width=-2cm}, raster column 2/.style={add to width=2cm}]
\begin{tcolorbox}[colback=white, colframe=white, boxrule=0pt, valign=center]
\begin{scratch}
\initmoreblocks{definir \namemoreblocks{cuadrado \ovalmoreblocks{lado}}}
\blockrepeat{repetir \ovalnum{4}}{
  \blockmove{mover \ovalmoreblocks{lado} pasos}
  \blockmove{girar \turnright{} \ovalnum{90} grados}
}
\end{scratch}
\par\medskip
\begin{scratch}
\blockinit{al presionar \greenflag}
\blockmoreblocks{cuadrado \ovalnum{50}}
\blockmoreblocks{cuadrado \ovalnum{100}}
\end{scratch}
\end{tcolorbox}
\begin{code}
def cuadrado(lado):
    for i in range(4):
        forward(lado)
        right(90)

cuadrado(50)
cuadrado(100)
\end{code}
\end{tcbraster}

\textbf{Cómo funciona.} La palabra dentro de los paréntesis de la línea \py|def| es una \textbf{entrada}. Python la llama \textbf{parámetro}. Es como el óvalo rosa \chip[myblocks]{lado} de Scratch: una caja con etiqueta que se llena cada vez que se usa el bloque. \py|cuadrado(50)| pone 50 en \py|lado|, y \py|cuadrado(100)| pone 100. Dentro del bloque, usas \py|lado| en cada lugar donde arrastrarías el óvalo rosa. (Pon las tres líneas de inicio arriba de este código cuando lo ejecutes).

% ---------------------------------------------------------------------
\needspace{14\baselineskip}
\subsection*{Dos entradas o más}

\begin{tcbraster}[raster columns=2, raster equal height=rows, raster column skip=12pt, raster valign=center, raster force size=false, raster column 1/.style={add to width=-0.6cm}, raster column 2/.style={add to width=0.6cm}]
\begin{tcolorbox}[colback=white, colframe=white, boxrule=0pt, valign=center]
\begin{scratch}
\initmoreblocks{definir \namemoreblocks{poligono \ovalmoreblocks{lados} \ovalmoreblocks{medida}}}
\blockrepeat{repetir \ovalmoreblocks{lados}}{
  \blockmove{mover \ovalmoreblocks{medida} pasos}
  \blockmove{girar \turnright{} \ovaloperator{\ovalnum{360} / \ovalmoreblocks{lados}} grados}
}
\end{scratch}
\par\medskip
\begin{scratch}
\blockmoreblocks{poligono \ovalnum{3} \ovalnum{100}}
\end{scratch}
\end{tcolorbox}
\begin{code}
def poligono(lados, medida):
    for i in range(lados):
        forward(medida)
        right(360 / lados)

poligono(3, 100)
poligono(6, 50)
\end{code}
\end{tcbraster}

\textbf{Cómo funciona.} Pon una coma entre las entradas. Cuando llamas al bloque, los números entran en las entradas \textbf{en orden}: el primer número va a \py|lados| y el segundo a \py|medida|. Así, \py|poligono(3, 100)| es un triángulo grande y \py|poligono(6, 50)| es un hexágono pequeño. (Escribimos \texttt{poligono} sin acento: los nombres que inventas van sin acentos ni ñ, porque son difíciles de teclear).

\begin{watchout}[¡Cuidado: el orden importa, y no puede faltar ninguna!]
\py|poligono(100, 3)| las intercambia: 100 lados de 3 pasos cada uno. ¡Te sale un circulito! Y si olvidas una entrada, Scratch simplemente usa un espacio vacío, pero Python se detiene y te dice:
\begin{errorbox}
TypeError: poligono() missing 1 required positional argument: 'medida'
\end{errorbox}
Python siempre escribe sus errores en inglés; lo más importante es el nombre del error (aquí, \texttt{TypeError}). Léelo como la pista de un detective: \textit{missing} quiere decir que al bloque \texttt{poligono} le \textbf{falta} la entrada llamada \texttt{medida}. Dale dos números, como \py|poligono(6, 50)|.
\end{watchout}

\begin{think}[Las entradas viven dentro del bloque]
En Scratch, si sacas el óvalo rosa \chip[myblocks]{lado} del sombrero definir y lo usas en otro lugar, solo reporta 0. En Python, una entrada solo existe \textbf{dentro} de su bloque. Prueba \py|print(lados)| en una línea pegada al borde izquierdo, después del bloque: Python dice \texttt{NameError}, porque \py|lados| le pertenece solo a \py|poligono|.
\end{think}

% ---------------------------------------------------------------------
\needspace{10\baselineskip}
\subsection*{Defínelo antes de usarlo}

En Scratch, el sombrero \chip[myblocks]{definir} puede estar en cualquier lugar del área de programación. Python es diferente: lee de \textbf{arriba hacia abajo}, como cuando tú lees un libro. Si una línea usa un bloque que Python todavía no ha leído, se detiene:

\begin{code}
flor()

def flor():
    print("¡Una flor!")
\end{code}
\begin{errorbox}
NameError: name 'flor' is not defined
\end{errorbox}

La pista dice que Python nunca ha oído hablar de \texttt{flor} (\textit{is not defined} significa «no está definido»). La solución: mueve la parte del \py|def| \textbf{arriba} de la línea que lo usa.

\begin{watchout}[¡Cuidado en Colab: ejecuta primero la celda del def!]
Si tu \py|def| está en una celda y usas el bloque en otra, ejecuta \textbf{primero} la celda del \py|def|. Y si cambias un \py|def|, vuelve a ejecutar su celda, o Colab seguirá usando la versión vieja.
\end{watchout}

% ---------------------------------------------------------------------
\needspace{14\baselineskip}
\subsection*{Entradas de texto y entradas de verdadero o falso}

\begin{tcbraster}[raster columns=2, raster equal height=rows, raster column skip=12pt, raster valign=center, raster force size=false, raster column 1/.style={add to width=-2cm}, raster column 2/.style={add to width=2cm}]
\begin{tcolorbox}[colback=white, colframe=white, boxrule=0pt, valign=center]
\setscratch{scale=0.85}
\begin{scratch}[scale=0.77]
\initmoreblocks{definir \namemoreblocks{saludar \ovalmoreblocks{nombre} \boolmoreblocks{gritar}}}
\blockifelse{si \boolmoreblocks{gritar} entonces}{
  \blocklook{decir \ovaloperator{unir \ovalnum{HOLA } \ovalmoreblocks{nombre}}}
}{
  \blocklook{decir \ovaloperator{unir \ovalnum{Hola } \ovalmoreblocks{nombre}}}
}
\end{scratch}
\par\medskip
\begin{scratch}
\blockmoreblocks{saludar \ovalnum{Ada} \booloperator{\ovalnum{1} = \ovalnum{2}}}
\blockmoreblocks{saludar \ovalnum{Sam} \booloperator{\ovalnum{1} = \ovalnum{1}}}
\end{scratch}
\end{tcolorbox}
\begin{code}
def saludar(nombre, gritar):
    if gritar:
        print("HOLA " + nombre)
    else:
        print("Hola " + nombre)

saludar("Ada", False)
saludar("Sam", True)
\end{code}
\end{tcbraster}

\begin{outputbox}
Hola Ada\\
HOLA Sam
\end{outputbox}

\textbf{Cómo funciona.} Una entrada también puede guardar palabras. Las palabras necesitan comillas cuando llamas al bloque: \py|saludar("Ada", False)|. La entrada en forma de hexágono puntiagudo de Scratch es una entrada de \textbf{verdadero o falso} (Scratch la llama \textit{booleano}). Scratch no tiene un bloque «verdadero», así que los scratchers meten una prueba como \chip[operators]{1 = 1}. Python te deja escribir simplemente \py|True| (verdadero) o \py|False| (falso), en inglés, con mayúscula y sin comillas. Luego \py|if gritar:| pregunta «¿\texttt{gritar} es True?».

% ---------------------------------------------------------------------
\needspace{14\baselineskip}
\subsection*{Ejecutar sin actualizar la pantalla}

\begin{tcbraster}[raster columns=2, raster equal height=rows, raster column skip=12pt, raster valign=center, raster force size=false, raster column 1/.style={add to width=-2cm}, raster column 2/.style={add to width=2cm}]
\begin{tcolorbox}[colback=white, colframe=white, boxrule=0pt, valign=center]
\begin{scratch}
\initmoreblocks{definir \namemoreblocks{flor}}
\blockrepeat{repetir \ovalnum{36}}{
  \blockmoreblocks{cuadrado \ovalnum{80}}
  \blockmove{girar \turnright{} \ovalnum{10} grados}
}
\end{scratch}
\par\smallskip
{\small $\boxtimes$ Ejecutar sin actualizar la pantalla}
\end{tcolorbox}
\begin{code}
def flor():
    animationOff()
    for i in range(36):
        cuadrado(80)
        right(10)
    animationOn()

flor()
\end{code}
\end{tcbraster}

\textbf{Cómo funciona.} Cuando creas un bloque en Scratch, puedes marcar \textbf{Ejecutar sin actualizar la pantalla} para que dibuje en un instante. Python tiene dos maneras de apurar a la tortuga:
\begin{itemize}
  \item \py|speed(13)| es la velocidad más rápida en la que todavía ves cómo se dibujan las líneas.
  \item \py|animationOff()| hace que la tortuga \textbf{salte} en vez de caminar, así que un dibujo grande aparece muy rápido. \py|animationOn()| vuelve a activar la caminata.
\end{itemize}
(Este programa también necesita las tres líneas de inicio y el \py|def cuadrado| de la página de la entrada numérica, arriba de él).

\begin{watchout}[¡Cuidado: no uses speed(0)!]
En Colab, \py|speed(0)| no muestra \textbf{absolutamente nada} hasta el final, y solo si agregas \py|done()|. Parece que tu programa está descompuesto. Quédate con \py|speed(13)| o \py|animationOff()|.
\end{watchout}

% ---------------------------------------------------------------------
\needspace{14\baselineskip}
\subsection*{El superpoder de Python: \texttt{return}}

Los bloques rosas de Scratch pueden \textbf{hacer} cosas, pero no pueden \textbf{devolver una respuesta}. No existe un bloque reportero rosa y redondo. Por eso los scratchers guardan la respuesta en una variable y la leen después. En Python, un bloque puede entregar la respuesta directamente con \py|return| («devolver» en inglés).

\begin{tcbraster}[raster columns=2, raster equal height=rows, raster column skip=12pt, raster valign=center, raster force size=false, raster column 1/.style={add to width=-2cm}, raster column 2/.style={add to width=2cm}]
\begin{tcolorbox}[colback=white, colframe=white, boxrule=0pt, valign=center]
\begin{scratch}
\initmoreblocks{definir \namemoreblocks{doble \ovalmoreblocks{n}}}
\blockvariable{fijar \selectmenu{resultado} a \ovaloperator{\ovalmoreblocks{n} * \ovalnum{2}}}
\end{scratch}
\par\medskip
\begin{scratch}
\blockinit{al presionar \greenflag}
\blockmoreblocks{doble \ovalnum{5}}
\blocklook{decir \ovalvariable{resultado}}
\end{scratch}
\end{tcolorbox}
\begin{code}
def doble(n):
    return n * 2

resultado = doble(5)
print(resultado)
print(doble(doble(3)))
\end{code}
\end{tcbraster}

\begin{outputbox}
10\\
12
\end{outputbox}

\textbf{Cómo funciona.} \py|return n * 2| significa «calcula \py|n * 2| y entrégaselo a quien lo pidió». Así, \py|doble(5)| se convierte en \texttt{10}, y puedes guardarlo en una variable, imprimirlo o hasta meterlo en otro \py|doble|. \py|doble(doble(3))| funciona de adentro hacia afuera: \py|doble(3)| es 6, y \py|doble(6)| es 12. ¡Tu bloque ahora funciona como un reportero verde y redondo, por ejemplo \chip[operators]{elegir número al azar entre 1 y 10}!

\needspace{14\baselineskip}
\begin{watchout}[¡Cuidado: print no es return!]
\py|print| solo \textbf{muestra} un número en la pantalla. No lo devuelve. Mira:
\begin{code}
def mostrar_doble(n):
    print(n * 2)

x = mostrar_doble(5)
print(x)
\end{code}
Esto muestra \texttt{10} y luego \texttt{None}. \texttt{None} es la palabra de Python para «nada». Un bloque sin \py|return| no devuelve nada, así que \py|x| queda vacía. Si quieres \textbf{usar} la respuesta, usa \py|return|.
\end{watchout}

% ---------------------------------------------------------------------
\needspace{14\baselineskip}
\subsection*{Más objetos: \texttt{Turtle()}}

\begin{tcbraster}[raster columns=2, raster equal height=rows, raster column skip=12pt, raster valign=center, raster force size=false, raster column 1/.style={add to width=-1.2cm}, raster column 2/.style={add to width=1.2cm}]
\begin{tcolorbox}[colback=white, colframe=white, boxrule=0pt, valign=center]
\setscratch{scale=0.7}{\small\bfseries Objeto Gato}\par
\begin{scratch}
\blockinit{al presionar \greenflag}
\blockpen{fijar color de lápiz a \ovalnum{naranja}}
\blockmove{mover \ovalnum{100} pasos}
\end{scratch}
\par\smallskip
{\small\bfseries Objeto Perro}\par
\begin{scratch}
\blockinit{al presionar \greenflag}
\blockpen{fijar color de lápiz a \ovalnum{azul}}
\blockmove{girar \turnleft{} \ovalnum{90} grados}
\blockmove{mover \ovalnum{100} pasos}
\end{scratch}
\end{tcolorbox}
\begin{code}
gato = Turtle()
perro = Turtle()
gato.pencolor("orange")
perro.pencolor("blue")
gato.forward(100)
perro.left(90)
perro.forward(100)
\end{code}
\end{tcbraster}

\textbf{Cómo funciona.} Hasta ahora, tus comandos simples han movido a una sola tortuga. \py|gato = Turtle()| crea un \textbf{objeto nuevo} y le pone el nombre \py|gato|. Cada objeto nuevo empieza en el centro, mirando a la derecha, con el lápiz abajo. Para decirle a un objeto qué hacer, escribe su nombre, un punto y luego el comando: \py|gato.forward(100)|. El nombre antes del punto es como hacer clic en ese objeto en la lista de objetos de Scratch. Los colores se escriben en inglés: \py|"orange"| es naranja, \py|"blue"| azul, \py|"gold"| dorado, \py|"purple"| morado.


% ---------------------------------------------------------------------
\needspace{8\baselineskip}
\subsection*{Dale un objeto a tu bloque}

En Scratch, un bloque rosa le pertenece a \textbf{un} objeto. Si el perro quiere el bloque \chip[myblocks]{cuadrado} del gato, tienes que arrastrar una copia hasta él. En Python, un mismo bloque puede servirle a todos los objetos: ¡le pasas el objeto como entrada!

\begin{tcbraster}[raster columns=2, raster equal height=rows, raster column skip=12pt, raster valign=center, raster force size=false, raster column 1/.style={add to width=-2cm}, raster column 2/.style={add to width=2cm}]
\begin{tcolorbox}[colback=white, colframe=white, boxrule=0pt, valign=center]
\begin{scratch}
\initmoreblocks{definir \namemoreblocks{cuadrado \ovalmoreblocks{lado}}}
\blockrepeat{repetir \ovalnum{4}}{
  \blockmove{mover \ovalmoreblocks{lado} pasos}
  \blockmove{girar \turnright{} \ovalnum{90} grados}
}
\end{scratch}
\par\smallskip
{\small (una copia en el Gato y otra copia en el Perro)}
\end{tcolorbox}
\begin{code}
def cuadrado(t, lado):
    for i in range(4):
        t.forward(lado)
        t.right(90)

cuadrado(gato, 60)
cuadrado(perro, 100)
\end{code}
\end{tcbraster}

\textbf{Cómo funciona.} La entrada \py|t| guarda un objeto (\texttt{t} de tortuga). Cuando llamas \py|cuadrado(gato, 60)|, el gato entra en \py|t|, así que \py|t.forward(lado)| mueve al gato. Llama \py|cuadrado(perro, 100)| y las mismas líneas mueven al perro. (Primero crea a \py|gato| y a \py|perro|, como arriba).

\begin{think}[La velocidad es de cada objeto]
\py|speed(13)| solo acelera a la tortuga de los comandos simples. Cada objeto nuevo tiene su propia velocidad, así que escribe \py|gato.speed(13)| para el gato.
\end{think}

% ---------------------------------------------------------------------
\needspace{14\baselineskip}
\subsection*{Clones: una lista de objetos}

\begin{tcbraster}[raster columns=2, raster equal height=rows, raster column skip=12pt, raster valign=center, raster force size=false, raster column 1/.style={add to width=-2cm}, raster column 2/.style={add to width=2cm}]
\begin{tcolorbox}[colback=white, colframe=white, boxrule=0pt, valign=center]
\setscratch{scale=0.85}
\begin{scratch}[scale=0.83]
\blockinit{al presionar \greenflag}
\blocklook{esconder}
\blockpen{subir lápiz}
\blockmove{ir a x: \ovalnum{-200} y: \ovalnum{120}}
\blockrepeat{repetir \ovalnum{5}}{
  \blockcontrol{crear clon de \selectmenu{mí mismo}}
  \blockmove{cambiar y en \ovalnum{-60}}
}
\end{scratch}
\par\smallskip
\begin{scratch}
\blockinitclone{al comenzar como clon}
\blocklook{mostrar}
\blockmove{mover \ovalnum{100} pasos}
\end{scratch}
\end{tcolorbox}
\begin{code}
ejercito = []
for i in range(5):
    s = Turtle()
    s.shape("turtle")
    s.penup()
    s.goto(-200, 120 - i * 60)
    ejercito.append(s)

for s in ejercito:
    s.forward(100)
\end{code}
\end{tcbraster}

\textbf{Cómo funciona.} Python no tiene un bloque «crear clon». En su lugar, creas objetos nuevos dentro de un bucle y los guardas en una \textbf{lista}, como una mochila llena de objetos. En cada vuelta, \py|Turtle()| crea un objeto nuevo llamado \py|s|, lo ponemos en su lugar y \py|ejercito.append(s)| lo mete en la mochila. \py|i| vale 0, 1, 2, 3, 4, así que \py|120 - i * 60| da las alturas 120, 60, 0, $-60$ y $-120$. Luego, \py|for s in ejercito:| saca los objetos uno por uno y mueve a cada uno. Ese segundo bucle es tu programa \chip[control]{al comenzar como clon}.

\begin{think}[Por turnos]
Los clones de Scratch parecen moverse todos al mismo tiempo. Python los mueve \textbf{uno después de otro}, muy rápido. Mantén pequeño tu ejército (hasta unos 20 objetos), o Colab se pone lento.
\end{think}

% ---------------------------------------------------------------------
\needspace{14\baselineskip}
\subsection*{Extra: el cortador de galletas}

\textit{Esta página es extra. Sáltatela si quieres y regresa cuando estés listo.}

Una \textbf{clase} (\textit{class}) es un cortador de galletas. Diseñas el cortador una vez y luego sacas todas las galletas que quieras. Cada galleta es su propio objeto, con su propia posición.

\begin{code}
class Hormiga:
    def __init__(self, y):
        self.t = Turtle()
        self.t.shape("turtle")
        self.t.penup()
        self.t.goto(-200, y)

    def mover(self):
        self.t.forward(30)

hormigas = []
for i in range(3):
    hormigas.append(Hormiga(i * 60))

for paso in range(5):
    for h in hormigas:
        h.mover()
\end{code}

\textbf{Cómo funciona.} \py|class Hormiga:| empieza el cortador de galletas. \py|__init__| (dos guiones bajos de cada lado) se ejecuta una vez, en el momento en que se crea una hormiga nueva, igual que \chip[control]{al comenzar como clon}. \py|self| significa «esta hormiga, la que se está creando ahora mismo», así que \py|self.t| es la tortuga propia de esta hormiga. \py|mover| es un bloque que todas las hormigas conocen. \py|Hormiga(60)| saca una hormiga nueva a la altura 60, y \py|h.mover()| le dice a una hormiga que se mueva.

% ---------------------------------------------------------------------
\needspace{14\baselineskip}
\subsection*{Reto extra: un bloque que se usa a sí mismo}

\textit{Extra, para programadores valientes.} ¡Un bloque puede llamarse a \textbf{sí mismo}! Esto se llama \textbf{recursión}. Cada rama de este árbol hace crecer dos árboles más pequeños.

\begin{code}
speed(13)

def arbol(medida):
    if medida < 10:
        return
    forward(medida)
    left(30)
    arbol(medida * 0.7)
    right(60)
    arbol(medida * 0.7)
    left(30)
    back(medida)

penup()
goto(0, -170)
setheading(90)
pendown()
arbol(80)
\end{code}

\textbf{Cómo funciona.} Las líneas 4 y 5 son la \textbf{regla para parar}: cuando la rama es diminuta, \py|return| termina el bloque de inmediato (como \chip[control]{detener este programa}). Sin ella, el árbol crecería para siempre, y Python se detendría con un \texttt{RecursionError}. Cada árbol nuevo mide 0.7 veces lo que medía el anterior. La última línea, \py|back(medida)|, camina de regreso por la rama, así que la tortuga siempre termina donde empezó. (Agrega las tres líneas de inicio hasta arriba).

% =====================================================================
\setscratch{scale=1}
\mission{Haz tu propio bloque}{a usar \texttt{def}, una entrada numérica y un mismo bloque muchas veces}

En esta misión creas un bloque \chip[myblocks]{cuadrado} con una entrada \texttt{lado} y luego lo usas tres veces.

\step{Paso 1: Constrúyelo en Scratch}
Haz clic en \textbf{Mis Bloques} y luego en \textbf{Crear un bloque}. Llámalo \texttt{cuadrado} y agrega una entrada de número llamada \texttt{lado} (\textbf{Agregar una entrada}, \textit{número o texto}). Agrega también la extensión \textbf{Lápiz} con el botón \textbf{Agregar Extensión}.

\begin{center}
\begin{tabular}{c@{\hspace{1.2cm}}c}
\begin{scratch}
\initmoreblocks{definir \namemoreblocks{cuadrado \ovalmoreblocks{lado}}}
\blockrepeat{repetir \ovalnum{4}}{
  \blockmove{mover \ovalmoreblocks{lado} pasos}
  \blockmove{girar \turnright{} \ovalnum{90} grados}
}
\end{scratch}
&
\begin{scratch}
\blockinit{al presionar \greenflag}
\blockpen{bajar lápiz}
\blockmoreblocks{cuadrado \ovalnum{40}}
\blockmoreblocks{cuadrado \ovalnum{80}}
\blockmoreblocks{cuadrado \ovalnum{120}}
\end{scratch}
\end{tabular}
\end{center}

\step{Paso 2: Escríbelo en Colab}

\begin{typethis}
from ColabTurtlePlus.Turtle import *
clearscreen()
initializeTurtle((480, 360))

def cuadrado(lado):
    for i in range(4):
        forward(lado)
        right(90)

cuadrado(40)
cuadrado(80)
cuadrado(120)
\end{typethis}

\step{Paso 3: Léelo línea por línea}

\begin{longtable}{>{\centering\arraybackslash}p{1.1cm} P{5.3cm} P{9.1cm}}
\toprule
\textbf{Línea} & \textbf{Bloque de Scratch} & \textbf{Qué significa} \\
\midrule
1--3 & \chip[events]{al presionar bandera verde} & Las tres líneas de inicio. El lápiz de la tortuga empieza abajo, así que no necesitas \chip[pen]{bajar lápiz}. \\
5 & \chip[myblocks]{definir cuadrado (lado)} & \py|def cuadrado(lado):| crea un bloque nuevo llamado \py|cuadrado| con una entrada llamada \py|lado|. Los dos puntos lo abren. \\
6 & \chip[control]{repetir 4} & Esta línea tiene cuatro espacios, así que está dentro del bloque. \\
7--8 & \chip[motion]{mover lado pasos}\newline\chip[motion]{girar $\circlearrowright$ 90 grados} & ¡\textbf{Ocho} espacios, porque están dentro del bucle, que está dentro del bloque! \py|lado| se usa donde va el óvalo rosa. \\
10 & \chip[myblocks]{cuadrado 40} & Esto \textbf{usa} tu bloque. El 40 entra en \py|lado|. \\
11--12 & \chip[myblocks]{cuadrado 80}\newline\chip[myblocks]{cuadrado 120} & Úsalo otra vez con lados más grandes. ¡Escribiste el código del cuadrado una sola vez, pero lo usaste tres veces! \\
\bottomrule
\end{longtable}

\step{Paso 4: Ejecútalo}
Presiona \playbtn. Deberías ver tres cuadrados de distintos tamaños, que empiezan todos en la misma esquina, en el centro del escenario, y crecen hacia abajo y a la derecha.

\begin{think}[¿Por qué la línea 9 está vacía?]
Python no la necesita. Te ayuda a \textbf{ti} a ver dónde termina el bloque.
\end{think}

\enlargethispage{\baselineskip}
\begin{challenge}\small
\tick Crea un bloque \py|triangulo(lado)|. Funciona como \texttt{cuadrado}, pero repite 3 veces y gira 120.\par
\tick \textbf{Dibuja una flor.} Agrega esto al final y ejecútalo. (Pon \py|speed(13)| después de la línea 3, para que no tarde una eternidad).
\begin{codenonum}
for i in range(12):
    cuadrado(80)
    right(30)
\end{codenonum}
\tick Dale a tu bloque una segunda entrada para el color: \py|def cuadrado(lado, tinta):|, y luego agrega \py|pencolor(tinta)| como primera línea de adentro. Úsalo así: \py|cuadrado(80, "red")|.\par
\tick Escribe \py|cuadrado| solo en una línea, sin paréntesis, y ejecútalo. ¿Qué pasa? ¿Por qué?
\end{challenge}

% =====================================================================
\mission{Máquina de figuras}{bloques con dos entradas, y un bloque que usa otro bloque}

Un solo bloque puede dibujar \textbf{cualquier} figura. Luego, un segundo bloque usa el primero para dibujar una flor.

\step{Paso 1: Constrúyelo en Scratch}

\begin{center}
\begin{tabular}{l@{\hspace{0.6cm}}l}
\begin{tabular}[t]{@{}l@{}}
\begin{scratch}
\initmoreblocks{definir \namemoreblocks{poligono \ovalmoreblocks{lados} \ovalmoreblocks{medida}}}
\blockrepeat{repetir \ovalmoreblocks{lados}}{
  \blockmove{mover \ovalmoreblocks{medida} pasos}
  \blockmove{girar \turnright{} \ovaloperator{\ovalnum{360} / \ovalmoreblocks{lados}} grados}
}
\end{scratch}
\\ \rule{0pt}{0.7cm}\\
\begin{scratch}
\initmoreblocks{definir \namemoreblocks{flor \ovalmoreblocks{petalos} \ovalmoreblocks{lados} \ovalmoreblocks{medida}}}
\blockrepeat{repetir \ovalmoreblocks{petalos}}{
  \blockmoreblocks{poligono \ovalmoreblocks{lados} \ovalmoreblocks{medida}}
  \blockmove{girar \turnright{} \ovaloperator{\ovalnum{360} / \ovalmoreblocks{petalos}} grados}
}
\end{scratch}
\end{tabular}
&
\begin{tabular}[t]{@{}l@{}}
\begin{scratch}
\blockinit{al presionar \greenflag}
\blockpen{subir lápiz}
\blockmove{ir a x: \ovalnum{-170} y: \ovalnum{60}}
\blockpen{bajar lápiz}
\blockpen{fijar color de lápiz a \ovalnum{azul}}
\blockmoreblocks{poligono \ovalnum{5} \ovalnum{70}}
\blockpen{subir lápiz}
\blockmove{ir a x: \ovalnum{90} y: \ovalnum{0}}
\blockpen{bajar lápiz}
\blockpen{fijar color de lápiz a \ovalnum{morado}}
\blockmoreblocks{flor \ovalnum{12} \ovalnum{6} \ovalnum{40}}
\end{scratch}
\end{tabular}
\end{tabular}
\end{center}

\step{Paso 2: Escríbelo en Colab}

\begin{typethis}
from ColabTurtlePlus.Turtle import *
clearscreen()
initializeTurtle((480, 360))
speed(13)

def poligono(lados, medida):
    for i in range(lados):
        forward(medida)
        right(360 / lados)

def flor(petalos, lados, medida):
    for i in range(petalos):
        poligono(lados, medida)
        right(360 / petalos)

penup()
goto(-170, 60)
pendown()
pencolor("blue")
poligono(5, 70)

penup()
goto(90, 0)
pendown()
pencolor("purple")
flor(12, 6, 40)
\end{typethis}

\step{Paso 3: Léelo línea por línea}

\begin{longtable}{>{\centering\arraybackslash}p{1.1cm} P{5.3cm} P{9.1cm}}
\toprule
\textbf{Línea} & \textbf{Bloque de Scratch} & \textbf{Qué significa} \\
\midrule
1--3 & \chip[events]{al presionar bandera verde} & Las tres líneas de inicio. \\
4 & $\boxtimes$ Ejecutar sin actualizar la pantalla & \py|speed(13)| dibuja lo más rápido posible mientras todavía lo ves. \\
6 & \chip[myblocks]{definir poligono (lados)}\newline\chip[myblocks]{(medida)} & Un bloque con \textbf{dos} entradas, con una coma entre ellas. \\
7 & \chip[control]{repetir lados} & \py|range(lados)| repite tantas veces como el número que hay en \py|lados|. \\
8--9 & \chip[motion]{mover medida pasos}\newline\chip[motion]{girar $\circlearrowright$ 360 / lados} & Los giros de cualquier figura suman 360, así que cada giro es \py|360 / lados|. \\
11 & \chip[myblocks]{definir flor (petalos)}\newline\chip[myblocks]{(lados) (medida)} & Un segundo bloque, con tres entradas. \\
12--14 & \chip[myblocks]{poligono lados medida} & El bloque \texttt{flor} \textbf{usa} el bloque \texttt{poligono} y luego gira un poquito. ¡Los bloques pueden usar otros bloques! \\
16--18 & \chip[pen]{subir lápiz}\newline\chip[motion]{ir a x: -170 y: 60}\newline\chip[pen]{bajar lápiz} & Va al lado izquierdo sin dibujar una línea. \\
19--20 & \chip[myblocks]{poligono 5 70} & Un pentágono azul: 5 lados, cada uno de 70 pasos. \\
22--24 & \chip[motion]{ir a x: 90 y: 0} & Va al lado derecho sin dibujar. \\
25--26 & \chip[myblocks]{flor 12 6 40} & Una flor morada: 12 pétalos, y cada pétalo es un hexágono con lados de 40. \\
\bottomrule
\end{longtable}

\step{Paso 4: Ejecútalo}
Deberías ver un pentágono azul a la izquierda y, a la derecha, una flor morada hecha de 12 hexágonos encimados.

\begin{think}[El orden dentro de los paréntesis]
¡\py|poligono(5, 70)| y \py|poligono(70, 5)| son muy diferentes! El primer número siempre va a la primera entrada. Prueba \py|poligono(70, 5)| y mira qué te sale.
\end{think}

\enlargethispage{\baselineskip}
\begin{challenge}\small
\tick Cambia la línea 26 por \py|flor(8, 4, 60)|. ¿Cuántos pétalos puedes contar? ¿Qué figura es cada pétalo?\par
\tick Agrega a \texttt{flor} una cuarta entrada llamada \texttt{tinta}, y usa \py|pencolor(tinta)| adentro.\par
\tick Dibuja un triángulo, un cuadrado y un hexágono en fila, usando solo \py|poligono|, \py|penup|, \py|goto| y \py|pendown|.\par
\tick Cambia la línea 4 por \py|animationOff()|. ¿Es más rápido?
\end{challenge}

% =====================================================================
\mission{Máquina de respuestas}{\texttt{return}: bloques que devuelven una respuesta}

Ahora vas a construir bloques que calculan respuestas y te las entregan, como tus propios bloques reporteros. Este programa no dibuja, así que no necesitas las tres líneas de inicio.

\step{Paso 1: Constrúyelo en Scratch}
Scratch no puede crear un bloque reportero, así que estos bloques rosas guardan su respuesta en una variable llamada \textbf{resultado}.

\begin{center}
\setscratch{scale=0.75}
\begin{tabular}{l@{\hspace{0.8cm}}l}
\begin{tabular}[t]{@{}l@{}}
\begin{scratch}
\initmoreblocks{definir \namemoreblocks{doble \ovalmoreblocks{n}}}
\blockvariable{fijar \selectmenu{resultado} a \ovaloperator{\ovalmoreblocks{n} * \ovalnum{2}}}
\end{scratch}
\\ \rule{0pt}{0.6cm}\\
\begin{scratch}
\initmoreblocks{definir \namemoreblocks{area \ovalmoreblocks{ancho} \ovalmoreblocks{alto}}}
\blockvariable{fijar \selectmenu{resultado} a \ovaloperator{\ovalmoreblocks{ancho} * \ovalmoreblocks{alto}}}
\end{scratch}
\end{tabular}
&
\begin{tabular}[t]{@{}l@{}}
\begin{scratch}
\blockinit{al presionar \greenflag}
\blockmoreblocks{doble \ovalnum{4}}
\blocklook{decir \ovalvariable{resultado} durante \ovalnum{2} segundos}
\blockmoreblocks{area \ovalnum{5} \ovalnum{3}}
\blocklook{decir \ovaloperator{unir \ovalnum{El tapete cubre } \ovalvariable{resultado}}}
\end{scratch}
\end{tabular}
\end{tabular}
\end{center}

\step{Paso 2: Escríbelo en Colab}

\begin{typethis}
def doble(n):
    return n * 2

def area(ancho, alto):
    return ancho * alto

print(doble(4))
print(doble(21))
tapete = area(5, 3)
print("El tapete cubre " + str(tapete) + " cuadros")
print(area(doble(2), 10))

w = int(input("¿Qué tan ancho? "))
h = int(input("¿Qué tan alto? "))
print("Área: " + str(area(w, h)))
\end{typethis}

\step{Paso 3: Léelo línea por línea}

\begin{longtable}{>{\centering\arraybackslash}p{1.1cm} P{5.3cm} P{9.1cm}}
\toprule
\textbf{Línea} & \textbf{Bloque de Scratch} & \textbf{Qué significa} \\
\midrule
1--2 & \chip[myblocks]{definir doble (n)} & \py|return n * 2| devuelve el doble del número. ¡No hace falta ninguna variable! \\
4--5 & \chip[myblocks]{definir area (ancho) (alto)} & El área de un rectángulo es el ancho por el alto. \py|return| la devuelve. \\
7 & \chip[myblocks]{doble 4} \chip[looks]{decir resultado} & \py|doble(4)| se convierte en 8, y \py|print| lo muestra. Una línea en vez de dos bloques. \\
8 & & El mismo bloque con otro número da 42. \\
9 & \chip[myblocks]{area 5 3} & La respuesta, 15, se guarda en una variable llamada \py|tapete|. \\
10 & \chip[looks]{decir unir El tapete\ldots} & \py|str(tapete)| convierte el número en texto, para poder unirlo con \texttt{+}. \\
11 & & Bloques dentro de bloques: \py|doble(2)| es 4, así que esto es \py|area(4, 10)|, que da 40. \\
13--14 & \chip[sensing]{preguntar ¿Qué tan ancho?}\newline\chip[sensing]{y esperar} & Pide dos números. \py|int( )| convierte cada respuesta en número. \\
15 & \chip[looks]{decir} & Calcula el área de \textbf{tu} rectángulo y la muestra. \\
\bottomrule
\end{longtable}

\step{Paso 4: Ejecútalo}
Presiona \playbtn. Cuando te pregunte, escribe un ancho y un alto, y presiona Enter después de cada uno.

\begin{outputbox}[Podrías ver]
8\\
42\\
El tapete cubre 15 cuadros\\
40\\
¿Qué tan ancho? 7\\
¿Qué tan alto? 4\\
Área: 28
\end{outputbox}

\begin{watchout}[¡Cuidado: return termina el bloque!]
En cuanto Python llega a \py|return|, sale del bloque de inmediato. Las líneas que estén después de \py|return| dentro del bloque nunca se ejecutan.
\end{watchout}

\begin{challenge}
\tick Escribe \py|def mitad(n):| que devuelva la mitad de un número. Prueba \py|print(mitad(9))|.\par
\tick Escribe \py|def perimetro(ancho, alto):| que devuelva la distancia de toda la vuelta alrededor de un rectángulo (o sea, \py|2 * ancho + 2 * alto|).\par
\tick Cambia la línea 2 por \py|print(n * 2)|. Ejecútalo. ¿Por qué ves \texttt{None}, y por qué ahora la línea 11 sale mal?\par
\tick \textbf{Súper reto:} escribe \py|def es_par(n):| que devuelva \py|n % 2 == 0|. Luego prueba \py|if es_par(7):|. ¡Scratch no puede crear un bloque rosa de verdadero o falso para nada!
\end{challenge}

% =====================================================================
\mission{Ejército de objetos}{a crear muchos objetos con un bucle, guardarlos en una lista y moverlos juntos}

Vas a crear cinco objetos tortuga, cada uno de un color diferente, y hacerlos marchar juntos por el escenario. Son la versión de Python de los clones.

\step{Paso 1: Constrúyelo en Scratch}

\begin{center}
\setscratch{scale=0.85}
\begin{tabular}{c@{\hspace{0.8cm}}c}
\begin{scratch}
\blockinit{al presionar \greenflag}
\blocklook{esconder}
\blockpen{subir lápiz}
\blockmove{ir a x: \ovalnum{-200} y: \ovalnum{120}}
\blockrepeat{repetir \ovalnum{5}}{
  \blockcontrol{crear clon de \selectmenu{mí mismo}}
  \blockmove{cambiar y en \ovalnum{-60}}
  \blocklook{cambiar el efecto \selectmenu{color} en \ovalnum{25}}
}
\end{scratch}
&
\begin{scratch}
\blockinitclone{al comenzar como clon}
\blocklook{mostrar}
\blockpen{bajar lápiz}
\blockrepeat{repetir \ovalnum{4}}{
  \blockmove{mover \ovalnum{80} pasos}
}
\end{scratch}
\end{tabular}
\end{center}

\step{Paso 2: Escríbelo en Colab}

\begin{typethis}
from ColabTurtlePlus.Turtle import *
clearscreen()
initializeTurtle((480, 360))

colores = ["red", "orange", "gold", "green", "blue"]
ejercito = []
for i in range(5):
    s = Turtle()
    s.shape("turtle")
    s.color(colores[i])
    s.speed(13)
    s.penup()
    s.goto(-200, 120 - i * 60)
    s.pendown()
    ejercito.append(s)

def marchar(pasos):
    for s in ejercito:
        s.forward(pasos)

for paso in range(4):
    marchar(80)
\end{typethis}

\par\medskip{\large\bfseries\color{ink}Paso 3: Léelo línea por línea}\par

\begin{longtable}{>{\centering\arraybackslash}p{1.1cm} P{5.3cm} P{9.1cm}}
\toprule
\textbf{Línea} & \textbf{Bloque de Scratch} & \textbf{Qué significa} \\
\midrule
1--3 & \chip[events]{al presionar bandera verde} & Las tres líneas de inicio. \\
5 & \chip[looks]{cambiar el efecto color en 25} & Una lista de cinco colores, uno para cada objeto (en inglés: rojo, naranja, dorado, verde y azul). \\
6 & & Una lista vacía: la mochila que va a guardar al ejército. \\
7 & \chip[control]{repetir 5} & Da 5 vueltas. \py|i| cuenta 0, 1, 2, 3, 4. \\
8 & \chip[control]{crear clon de mí mismo} & \py|Turtle()| crea un objeto totalmente nuevo, que por ahora se llama \py|s|. \\
9--11 & \chip[looks]{mostrar} & Le da el disfraz de tortuga, su propio color (\py|colores[i]|) y la velocidad más rápida. \py|color| fija al mismo tiempo el color del lápiz y el del disfraz. \\
12--14 & \chip[pen]{subir lápiz}\newline\chip[motion]{ir a x: -200 y: 120}\newline\chip[pen]{bajar lápiz} & Salta a la línea de salida sin dibujar. Cada objeto queda 60 pasos más abajo que el anterior. \\
15 & & Mete el objeto nuevo en la mochila. \\
17--19 & \chip[myblocks]{definir marchar (pasos)} & Un bloque que mueve hacia adelante a \textbf{todos} los objetos del ejército. \\
21--22 & \chip[control]{repetir 4}\newline\chip[motion]{mover 80 pasos} & Marcha 4 veces. Cada vez, todos los objetos dan 80 pasos. \\
\bottomrule
\end{longtable}

\step{Paso 4: Ejecútalo}
Cinco tortugas de colores aparecen a la izquierda, una arriba de la otra. Luego marchan juntas hacia la derecha, paso a paso, y cada una deja una línea de su propio color. Se detienen en $x = 120$.

\begin{think}[¿Por qué no cinco nombres?]
Podrías escribir \py|rojo = Turtle()|, \py|naranja = Turtle()| y así. Pero con una lista, 5 objetos ocupan el mismo número de líneas que 50. ¡Solo cambia el \texttt{5} de la línea 7 (y agrega más colores)!
\end{think}

\begin{challenge}
\tick Dale 8 objetos al ejército. Agrega tres colores más a la lista y cambia \py|i * 60| por \py|i * 35| para que quepan en el escenario.\par
\tick Crea un bloque \py|girar_todos(angulo)| que gire a todos los objetos, y úsalo para que el ejército marche hacia arriba del escenario.\par
\tick Importa \py|random| y haz que cada objeto dé \py|random.randint(10, 80)| pasos en vez de todos el mismo número. ¡Ahora es una carrera!\par
\tick Usa \py|ejercito[0]| para mover solo al primer objeto. ¿Cuál es?
\end{challenge}

% =====================================================================
\banner{pen}{TU TURNO}{Tradúcelo tú}

Aquí tienes un proyecto de Scratch con un bloque hecho en casa. Escribe el Python en las líneas (con las tres líneas de inicio arriba). Luego escríbelo en Colab y comprueba que la tortuga dibuje lo mismo. Pista: no olvides \textbf{llamar} a tu bloque, con paréntesis, debajo del \py|def|.

\needspace{20\baselineskip}
\begin{tcbraster}[raster columns=2, raster equal height=rows, raster column skip=12pt, raster valign=top]
\begin{tcolorbox}[colback=white, colframe=white, boxrule=0pt]
\begin{scratch}
\initmoreblocks{definir \namemoreblocks{estrella \ovalmoreblocks{medida}}}
\blockrepeat{repetir \ovalnum{5}}{
  \blockmove{mover \ovalmoreblocks{medida} pasos}
  \blockmove{girar \turnright{} \ovalnum{144} grados}
}
\end{scratch}
\par\bigskip
\begin{scratch}
\blockinit{al presionar \greenflag}
\blockpen{fijar color de lápiz a \ovalnum{dorado}}
\blockmoreblocks{estrella \ovalnum{100}}
\blockmoreblocks{estrella \ovalnum{50}}
\end{scratch}
\end{tcolorbox}
\begin{tcolorbox}[colback=white, colframe=codeframe, arc=5pt, title=\textbf{Mi Python}, coltitle=ink, colbacktitle=codebg]
\vspace{4pt}
\foreach \n in {1,...,11}{\rule{\linewidth}{0.4pt}\par\vspace{9pt}}
\end{tcolorbox}
\end{tcbraster}


\needspace{10\baselineskip}
\begin{puzzle}
Escribe un bloque llamado \py|triple| que \textbf{devuelva} el triple del número que le das. Ocupa exactamente \textbf{dos líneas}. Cuando funcione, esta línea debería mostrar \texttt{18}:
\begin{codenonum}
print(triple(triple(2)))
\end{codenonum}
\vspace{4pt}
\foreach \n in {1,...,2}{\rule{\linewidth}{0.4pt}\par\vspace{9pt}}
\end{puzzle}

\needspace{14\baselineskip}
\subsection*{Lista de la parte 6}

Ya sé\ldots
\begin{itemize}[label={}, leftmargin=0pt]
  \item \tick crear mi propio bloque con \py|def| y dos puntos, y llamarlo con paréntesis.
  \item \tick darle a un bloque una o más entradas, en el orden correcto.
  \item \tick explicar por qué el \py|def| debe ir antes de la llamada (y arreglar un \texttt{NameError}).
  \item \tick usar \py|return| para devolver una respuesta, y distinguir \py|print| de \py|return|.
  \item \tick crear un objeto nuevo con \py|Turtle()| y pasárselo a un bloque.
  \item \tick crear un ejército de objetos con un bucle y moverlos todos con otro bucle.
\end{itemize}

% =====================================================================
% PARTE 7: Apariencia, sonido y eventos
% =====================================================================
\partopener{events}{PARTE 7}{Apariencia, sonido y eventos}{Haz que tus objetos se disfracen, canten una melodía y se pasen mensajes, y luego hazlos competir a todos en un solo bucle.}

Esta es la última parte, ¡y es muy divertida! Vas a disfrazar a tu tortuga, hacer que Colab toque música y aprender cómo hace Python el trabajo de los bloques amarillos con sombrero de Scratch. Estos son algunos de los bloques que vas a conocer:

\begin{center}
\chip[looks]{decir ¡Hola! durante 2 segundos}\quad
\chip[looks]{siguiente disfraz}\par\vspace{6pt}
\chip[looks]{fijar tamaño a 200 \%}\quad
\chip[looks]{cambiar fondo a cielo azul}\par\vspace{6pt}
\chip[sound]{iniciar sonido Miau}\quad
\chip[sound]{tocar sonido Miau hasta que termine}\par\vspace{6pt}
\chip[events]{al presionar bandera verde}\quad
\chip[events]{enviar mensaje1}\quad
\chip[events]{al recibir mensaje1}
\end{center}

Algunos funcionan casi igual en Python. Otros necesitan un truco ingenioso. Y hay una gran idea que es distinta: Scratch parece ejecutar muchos programas a la vez, pero Python hace \textbf{una cosa a la vez}. Al final de esta parte sabrás cómo hacer que un solo bucle mueva a todos los objetos.

% ---------------------------------------------------------------------
\needspace{18\baselineskip}
\subsection*{Decir cosas: \texttt{print}}

Conociste \py|print| al principio del libro. Es el bloque \chip[looks]{decir} de Python. Las palabras aparecen debajo de la celda de código en lugar de en un globo de diálogo.

\begin{tcbraster}[raster columns=2, raster equal height=rows, raster column skip=12pt, raster valign=center]
\begin{tcolorbox}[colback=white, colframe=white, boxrule=0pt, valign=center]
\begin{scratch}
\blockinit{al presionar \greenflag}
\blocklook{decir \ovalnum{¡Hola!} durante \ovalnum{2} segundos}
\blocklook{decir \ovalnum{¡Adiós!}}
\end{scratch}
\end{tcolorbox}
\begin{code}
import time
print("¡Hola!")
time.sleep(2)
print("¡Adiós!")
\end{code}
\end{tcbraster}

\textbf{Cómo funciona.} \py|print("¡Hola!")| es la parte de \chip[looks]{decir}. La parte de «durante 2 segundos» es una espera, así que se convierte en \py|time.sleep(2)|, que ya conociste con el bloque \chip[control]{esperar}. \py|import time| arriba trae las herramientas del reloj. En Scratch, el globo desaparece después de 2 segundos. Las palabras impresas se quedan en la página, así que verás \texttt{¡Hola!} y luego, 2 segundos después, \texttt{¡Adiós!} debajo.

\begin{think}[¿Y pensar?]
Python no tiene globo de pensamiento, así que \chip[looks]{pensar Mmm...} también es solo \py|print("Mmm...")|. Si quieres que se sepa que es un pensamiento, puedes escribir \py|print("(piensa) Mmm...")|.
\end{think}

% ---------------------------------------------------------------------
\needspace{14\baselineskip}
\subsection*{Palabras en el escenario: una tortuga globo}

¿Quieres que las palabras aparezcan \textbf{en el escenario}, junto a tu objeto? La tortuga tiene un comando \py|write()| que escribe palabras donde está parada. El truco es crear una segunda tortuga, \textbf{invisible}, cuyo único trabajo es escribir. Vamos a llamarla \py|globo|, como el globo de diálogo de Scratch.

\begin{center}
\begin{scratch}
\blockinit{al presionar \greenflag}
\blocklook{cambiar disfraz a \selectmenu{tortuga}}
\blocklook{decir \ovalnum{¡Hola!} durante \ovalnum{2} segundos}
\end{scratch}
\end{center}

\begin{typethis}
from ColabTurtlePlus.Turtle import *
clearscreen()
initializeTurtle((480, 360))
import time

shape("turtle")
globo = Turtle()
globo.hideturtle()
globo.penup()
globo.goto(0, 30)
globo.write("¡Hola!", align="center",
            font=("Arial", 20, "normal"))
time.sleep(2)
globo.clear()
\end{typethis}

\textbf{Cómo funciona.} Las líneas 7 a 10 crean un objeto nuevo llamado \py|globo|, lo esconden, levantan su lápiz y lo mueven justo encima de la tortuga. Las líneas 11 y 12 son la parte de \chip[looks]{decir}: \py|write| escribe las palabras donde está parado \py|globo|. \py|align="center"| pone ahí el centro de las palabras, y \py|font=| da el tipo de letra, el tamaño (20) y \py|"normal"|. Es un solo comando partido en dos líneas: Python sabe que no ha terminado hasta que se cierra el paréntesis. La línea 13 espera, y la línea 14 hace desaparecer el «globo»: \py|clear()| borra todo lo que \py|globo| ha escrito o dibujado.

\begin{watchout}[¿Por qué una tortuga globo aparte?]
\py|clear()| borra \textbf{todo} lo que dibujó esa tortuga, no solo las palabras. Si tu tortuga principal escribiera las palabras y luego las borrara, ¡sus líneas también desaparecerían! Una tortuga globo es como un amigo que sostiene un pizarrón: para cambiar el letrero, el amigo borra el pizarrón y vuelve a escribir, y tu dibujo queda a salvo. Para un pensamiento, cambia \py|"normal"| por \py|"italic"| y las palabras salen inclinadas, como cuando piensas.
\end{watchout}

% ---------------------------------------------------------------------
\needspace{20\baselineskip}
\subsection*{Disfraces: \texttt{shape}}

En Scratch, un objeto puede cambiar de disfraz. La tortuga también, con \py|shape()|. Le das el nombre de un disfraz entre comillas.

\begin{tcbraster}[raster columns=2, raster equal height=rows, raster column skip=12pt, raster valign=center]
\begin{tcolorbox}[colback=white, colframe=white, boxrule=0pt, valign=center]
\begin{scratch}
\blocklook{cambiar disfraz a \selectmenu{tortuga}}
\end{scratch}
\end{tcolorbox}
\begin{code}
shape("turtle")
\end{code}
\end{tcbraster}

La tortuga tiene nueve disfraces. ¡Pruébalos todos! Python habla inglés: sus palabras (\py|shape|, \py|print|, \py|forward|...) siempre van en inglés, y los nombres de los disfraces también.

\begin{center}
\renewcommand{\arraystretch}{1.3}
\begin{tabular}{lll}
\toprule
\py|"classic"| (el del principio) & \py|"arrow"| (flecha) & \py|"triangle"| (triángulo) \\
\py|"square"| (cuadrado) & \py|"circle"| (círculo) & \py|"turtle"| (tortuga) \\
\py|"turtle2"| (tortuga puntiaguda) & \py|"ring"| (anillo) & \py|"blank"| (invisible) \\
\bottomrule
\end{tabular}
\end{center}

\textbf{Cómo funciona.} El nombre del disfraz es texto, así que necesita comillas, y debe estar escrito exactamente igual. Si escribes un nombre que no está en la lista, Python se detiene con un \texttt{ValueError} y te muestra la lista de nombres que conoce. ¡Esa es una nota de pista muy útil! Python siempre escribe sus errores en inglés; lo importante es el nombre del error (\texttt{ValueError}, \texttt{NameError}...).

\begin{watchout}[En Colab no hay disfraces con dibujos]
En Scratch puedes pintar un disfraz o elegir un gato. La tortuga de Colab solo tiene las formas de la tabla: no puede cargar imágenes. En su lugar, haz que tus disfraces sean diferentes cambiando su \textbf{forma}, su \textbf{color} y su \textbf{tamaño}. De eso trata el resto de esta página.
\end{watchout}

% ---------------------------------------------------------------------
\needspace{12\baselineskip}
\subsection*{Siguiente disfraz: una lista y un gancho}

Scratch te guarda un ropero de disfraces. En Python, tú construyes el ropero, como una \textbf{lista}. Un número llamado \py|n| es el gancho al que estás apuntando.

\begin{tcbraster}[raster columns=2, raster equal height=rows, raster column skip=12pt, raster valign=center, raster force size=false, raster column 1/.style={add to width=-2.2cm}, raster column 2/.style={add to width=2.2cm}]
\begin{tcolorbox}[colback=white, colframe=white, boxrule=0pt, valign=center]
\begin{scratch}
\blockinit{al presionar \greenflag}
\blockrepeat{repetir \ovalnum{6}}{
  \blockcontrol{esperar \ovalnum{1} segundos}
  \blocklook{siguiente disfraz}
}
\end{scratch}
\end{tcolorbox}
\begin{code}
disfraces = ["turtle", "circle",
             "square"]
n = 0
shape(disfraces[n])
for i in range(6):
    time.sleep(1)
    n = (n + 1) % len(disfraces)
    shape(disfraces[n])
\end{code}
\end{tcbraster}

\textbf{Cómo funciona.} \py|disfraces[n]| toma el disfraz del gancho \py|n|. Recuerda que las listas empiezan a contar en 0, así que el gancho 0 es \py|"turtle"|. La línea 7 es la parte ingeniosa. \py|n + 1| pasa al siguiente gancho. \py|len(disfraces)| es 3, el número de disfraces. El signo \texttt{\%} es el bloque \chip[operators]{módulo} de Scratch: da el residuo de una división. Así que cuando \py|n + 1| llega a 3, \py|3 % 3| es 0 y vuelves al primer gancho, igual que Scratch cuando pasa del último disfraz al primero.

\begin{center}
\renewcommand{\arraystretch}{1.25}
\begin{tabular}{cccc}
\toprule
\py|n| antes & \py|n + 1| & \py|(n + 1) % 3| & disfraz \\
\midrule
0 & 1 & 1 & \py|"circle"| \\
1 & 2 & 2 & \py|"square"| \\
2 & 3 & 0 & \py|"turtle"| \\
\bottomrule
\end{tabular}
\end{center}

El bloque \chip[looks]{disfraz número} de Scratch cuenta desde 1, así que en Python es \py|n + 1|.

% ---------------------------------------------------------------------
\needspace{16\baselineskip}
\subsection*{Tamaño: \texttt{shapesize}}

\begin{tcbraster}[raster columns=2, raster equal height=rows, raster column skip=12pt, raster valign=center]
\begin{tcolorbox}[colback=white, colframe=white, boxrule=0pt, valign=center]
\begin{scratch}
\blocklook{fijar tamaño a \ovalnum{200} \%}
\blockrepeat{repetir \ovalnum{5}}{
  \blocklook{cambiar tamaño en \ovalnum{50}}
}
\end{scratch}
\end{tcolorbox}
\begin{code}
shapesize(2)
medida = 200
for i in range(5):
    medida = medida + 50
    shapesize(medida / 100)
\end{code}
\end{tcbraster}

\textbf{Cómo funciona.} Scratch mide el tamaño en porcentaje, y 100\,\% es lo normal. Python usa «cuántas veces más grande», así que \py|shapesize(1)| es normal (100\,\%), \py|shapesize(2)| es el doble (200\,\%) y \py|shapesize(0.5)| es la mitad (50\,\%). Para cambiar el tamaño poco a poco, guarda tu propia variable \py|medida| en porcentaje, igual que en Scratch, y divídela entre 100 cuando se la des a \py|shapesize|.

\begin{watchout}[Nunca tamaño 0]
\py|shapesize(0)| hace que Python se detenga con un \texttt{ValueError}. Scratch evita en silencio que tu objeto se haga demasiado pequeño, pero Python no. Si tu tamaño baja dentro de un bucle, asegúrate de que nunca llegue a 0. Usa 0.1 para una tortuga diminuta.
\end{watchout}

% ---------------------------------------------------------------------
\needspace{16\baselineskip}
\subsection*{Mostrar y esconder}

\begin{tcbraster}[raster columns=2, raster equal height=rows, raster column skip=12pt, raster valign=center]
\begin{tcolorbox}[colback=white, colframe=white, boxrule=0pt, valign=center]
\begin{scratch}
\blockrepeat{repetir \ovalnum{3}}{
  \blocklook{esconder}
  \blockcontrol{esperar \ovalnum{0.5} segundos}
  \blocklook{mostrar}
  \blockcontrol{esperar \ovalnum{0.5} segundos}
}
\end{scratch}
\end{tcolorbox}
\begin{code}
for i in range(3):
    hideturtle()
    time.sleep(0.5)
    showturtle()
    time.sleep(0.5)
\end{code}
\end{tcbraster}

\textbf{Cómo funciona.} \py|hideturtle()| hace invisible a la tortuga y \py|showturtle()| la trae de vuelta. Los paréntesis están vacíos porque no hay nada que decirle, pero igual los necesitas: son los que hacen que Python \textbf{haga} el comando. Una tortuga escondida puede seguir moviéndose y dibujando. ¡Así funciona la tortuga globo!

% ---------------------------------------------------------------------
\needspace{18\baselineskip}
\subsection*{Fondos y colores}

\begin{tcbraster}[raster columns=2, raster equal height=rows, raster column skip=12pt, raster valign=center]
\begin{tcolorbox}[colback=white, colframe=white, boxrule=0pt, valign=center]
\begin{scratch}[scale=0.94]
\blocklook{cambiar fondo a \selectmenu{cielo azul}}
\blockrepeat{repetir \ovalnum{4}}{
  \blocklook{cambiar el efecto \selectmenu{color} en \ovalnum{25}}
  \blockcontrol{esperar \ovalnum{0.5} segundos}
}
\end{scratch}
\end{tcolorbox}
\begin{code}
bgcolor("skyblue")
arcoiris = ["red", "orange",
            "gold", "green"]
for c in arcoiris:
    color(c)
    time.sleep(0.5)
\end{code}
\end{tcbraster}

\textbf{Cómo funciona.} \py|bgcolor()| («color de fondo», en inglés \emph{background color}) pinta todo el escenario. En Colab los fondos son solo colores lisos, no imágenes, así que elige un color que vaya bien: \py|"skyblue"| para el cielo, \py|"black"| para el espacio, \py|"lightgreen"| para un prado. \py|color(c)| colorea a la tortuga misma. El bucle recorre la lista \py|arcoiris| y le da a la tortuga cada color por turno, que es una forma sencilla de hacer el efecto de color de Scratch. Los colores se escriben en inglés: \py|"red"| es rojo, \py|"orange"| naranja, \py|"gold"| dorado, \py|"green"| verde y \py|"skyblue"| azul cielo. Y por suerte, en Python \texttt{color} se escribe igual que en español. Recuerda: los nombres de los colores van todos en minúsculas.

% ---------------------------------------------------------------------
\needspace{16\baselineskip}
\subsection*{Sonido en Colab: la celda de sonido}

Colab funciona en las computadoras de Google, muy lejos de ti, así que Python no puede llegar a tus bocinas. Pero \textbf{sí} puede mandar un sonido a la página web, y tu navegador lo toca. Para eso, ejecuta esta \textbf{celda de sonido} una vez, después de la celda «En el principio». Crea un bloque nuevo llamado \py|pitido| (porque hace un pitido).

\begin{typethis}[Celda de sonido: ejecútala una vez cada vez que abras Colab]
import numpy as np
from IPython.display import Audio, display

def pitido(frec=440, segs=0.5):
    instantes = np.linspace(0, segs, int(22050 * segs))
    sonido = np.sin(2 * np.pi * frec * instantes)
    display(Audio(sonido, rate=22050, autoplay=True))
\end{typethis}

No necesitas entender cada línea. Es una tarjeta de receta que escribió otra persona, y tú solo la usas, como un bloque de una extensión. Esta es la idea. Las líneas 1 y 2 traen dos extensiones: \py|numpy| (matemáticas superrápidas) y las herramientas que ponen un reproductor de sonido en la página. La línea 4 crea tu propio bloque, \py|pitido|, con dos entradas: qué tan aguda es la nota (\py|frec|, abreviatura de \emph{frecuencia}) y cuánto dura (\py|segs|, de \emph{segundos}). \py|frec=440| significa «si nadie dice otra cosa, usa 440». Las líneas 5 y 6 calculan la forma de la onda de sonido, y la línea 7 pone un pequeño reproductor de sonido en la página y lo toca enseguida (\py|autoplay=True|).

\begin{watchout}[¿No se oye nada?]
\begin{itemize}
  \item Revisa que el volumen de tu computadora esté alto.
  \item Algunos navegadores bloquean los sonidos que suenan solos. Si ves un pequeño reproductor con un botón \playbtn{} debajo de la celda pero no oyes nada, presiona su botón de reproducir.
  \item ¿Sale \texttt{NameError: name 'pitido' is not defined}? Ejecuta primero la celda de sonido.
\end{itemize}
\end{watchout}

% ---------------------------------------------------------------------
\needspace{18\baselineskip}
\subsection*{Iniciar sonido, y tocar hasta que termine}

\begin{tcbraster}[raster columns=2, raster equal height=rows, raster column skip=12pt, raster valign=center]
\begin{tcolorbox}[colback=white, colframe=white, boxrule=0pt, valign=center]
\begin{scratch}[scale=0.8]
\blocksound{tocar sonido \selectmenu{Do grave} hasta que termine}
\blocksound{tocar sonido \selectmenu{Mi} hasta que termine}
\blocksound{iniciar sonido \selectmenu{Sol largo}}
\end{scratch}
\end{tcolorbox}
\begin{code}
import time
pitido(262)
time.sleep(0.5)
pitido(330)
time.sleep(0.5)
pitido(392, 1)
\end{code}
\end{tcbraster}

\textbf{Cómo funciona.} \py|pitido(262)| es \chip[sound]{iniciar sonido}: empieza la nota y Python sigue de inmediato con la siguiente línea. \chip[sound]{tocar sonido hasta que termine} espera a que el sonido acabe. Colab no sabe cuándo terminó un sonido, así que esperas tú: la nota dura 0.5 segundos, así que \py|time.sleep(0.5)| espera 0.5 segundos. \py|pitido(392, 1)| también da la segunda entrada, así que esa nota dura 1 segundo.

\begin{think}[El timbre y el radio]
\chip[sound]{iniciar sonido} es como prender el radio e irte a hacer otra cosa. \chip[sound]{tocar sonido hasta que termine} es como quedarte junto al radio hasta que acabe la canción. Quita las dos líneas de \py|time.sleep| y ejecútalo otra vez. ¡Todas las notas empiezan a la vez y suenan unas encima de otras!
\end{think}

% ---------------------------------------------------------------------
\needspace{14\baselineskip}
\subsection*{Las notas son números}

El número que le das a \py|pitido| decide qué tan aguda es la nota. Los números más grandes son más agudos. Estas son las notas de una escala, empezando en el Do central. En inglés las notas se nombran con letras: Do es C, Re es D, y así.

\begin{center}
\renewcommand{\arraystretch}{1.3}
\begin{tabular}{l*{8}{c}}
\toprule
\textbf{Nota} & Do & Re & Mi & Fa & Sol & La & Si & Do agudo \\
\textbf{Letra en inglés} & C & D & E & F & G & A & B & C \\
\textbf{Número} & 262 & 294 & 330 & 349 & 392 & 440 & 494 & 523 \\
\textbf{Nota en Música} & 60 & 62 & 64 & 65 & 67 & 69 & 71 & 72 \\
\bottomrule
\end{tabular}
\end{center}

Una melodía es una lista de notas, así que puedes tocarla con un bucle \py|for|:

\begin{code}
notas = [262, 294, 330, 349, 392, 440, 494, 523]
for nota in notas:
    pitido(nota, 0.3)
    time.sleep(0.4)
\end{code}

Esto toca toda la escala, desde Do hasta Do agudo. Cada nota dura 0.3 segundos, y la espera de 0.4 segundos deja un huequito antes de la siguiente.

\begin{think}[Para los adultos: si todas las notas suenan al final]
En la mayoría de las computadoras, cada nota suena en cuanto se ejecuta \py|pitido|. Si la tuya las guarda todas y las toca juntas cuando termina la celda, une las notas en un solo sonido y toca ese. Pega esto debajo de la celda de sonido y usa \py|tocar_melodia([262, 294, 330])|:
\begin{codenonum}
def tocar_melodia(notas, segs=0.4):
    partes = []
    for nota in notas:
        instantes = np.linspace(0, segs, int(22050 * segs))
        partes.append(np.sin(2 * np.pi * nota * instantes))
    cancion = np.concatenate(partes)
    display(Audio(cancion, rate=22050, autoplay=True))
\end{codenonum}
\end{think}

% ---------------------------------------------------------------------
\needspace{14\baselineskip}
\subsection*{La bandera verde: ejecutar la celda}

Todo programa de Scratch empieza con un bloque de sombrero. El que más usas es \chip[events]{al presionar bandera verde}. En Colab, la bandera verde es el botón de reproducir \playbtn: al ejecutar una celda empieza todo lo que tiene, desde arriba. Por eso tus programas de dibujo empiezan con las tres líneas de inicio. Hacen lo que hace un programa de bandera verde al empezar: limpiar el escenario y preparar al objeto.

Los otros bloques de sombrero de Scratch necesitan trucos diferentes en Python. Aquí están:

\begin{center}
\renewcommand{\arraystretch}{1.4}
\begin{tabular}{P{5.6cm} P{9.6cm}}
\toprule
\textbf{Scratch} & \textbf{Python en Colab} \\
\midrule
\chip[events]{al presionar bandera verde} & Presiona \playbtn{} para ejecutar la celda. \\
\chip[events]{al recibir mensaje1} & Crea una función: \py|def mensaje1():| \\
\chip[events]{enviar mensaje1} & Llama a la función: \py|mensaje1()| \\
\chip[events]{al presionar tecla espacio} & No existe en Colab. Pregunta con \py|input()| o usa Thonny (página al final de esta parte). \\
\chip[events]{al hacer clic en este objeto} & No existe en Colab. Pregunta con \py|input()|. \\
\bottomrule
\end{tabular}
\end{center}

% ---------------------------------------------------------------------
\needspace{12\baselineskip}
\subsection*{Enviar mensajes es llamar funciones}

En Scratch, un programa puede \textbf{enviar} un mensaje, y otro programa que empieza con \chip[events]{al recibir} hace su trabajo. En Python, el programa que recibe es una \textbf{función}, y enviar el mensaje es \textbf{llamarla}.

\begin{center}
\begin{tabular}{c@{\hspace{1cm}}c}
\begin{scratch}
\blockinit{al presionar \greenflag}
\blocklook{decir \ovalnum{Listos...}}
\blockevent{enviar \selectmenu{vamos} y esperar}
\blocklook{decir \ovalnum{¡Listo!}}
\end{scratch}
&
\begin{scratch}
\blockinit{al recibir \selectmenu{vamos}}
\blocklook{decir \ovalnum{¡Vamos, vamos, vamos!}}
\end{scratch}
\end{tabular}
\end{center}

\begin{code}
def vamos():
    print("¡Vamos, vamos, vamos!")

print("Listos...")
vamos()
print("¡Listo!")
\end{code}

\begin{outputbox}
Listos...\\
¡Vamos, vamos, vamos!\\
¡Listo!
\end{outputbox}

\textbf{Cómo funciona.} Las líneas 1--2 son el programa \chip[events]{al recibir vamos}. Python aprende la receta, pero todavía no la hace. La línea 4 se ejecuta primero. La línea 5 es el envío del mensaje: Python salta dentro de \py|vamos|, hace todo lo que tiene adentro, luego regresa y sigue con la línea 6.

\begin{think}[Una llamada siempre espera]
Scratch tiene dos bloques para enviar mensajes. \chip[events]{enviar} manda el mensaje y sigue de inmediato. \chip[events]{enviar y esperar} espera hasta que el otro programa haya terminado. Una llamada a una función en Python \textbf{siempre espera}: Python termina \py|vamos()| antes de pasar a la siguiente línea. Así que una llamada normal es como \chip[events]{enviar y esperar}. Buena noticia: así es fácil saber qué pasa y en qué orden.
\end{think}

% ---------------------------------------------------------------------
\needspace{12\baselineskip}
\subsection*{Las teclas se vuelven preguntas}

La tortuga de Colab no puede oír el teclado ni el ratón mientras corre un programa. Así que en lugar de \chip[events]{al presionar tecla a}, \textbf{preguntas} qué tecla, con \py|input()|. El jugador escribe una letra y presiona Enter. Es como un juego de mesa: una jugada por turno. Usaremos \texttt{a} para avanzar, \texttt{i} para izquierda y \texttt{d} para derecha.

\begin{center}
\begin{tabular}{c@{\hspace{0.6cm}}c@{\hspace{0.6cm}}c}
\begin{scratch}
\blockinit{al presionar tecla \selectmenu{a}}
\blockmove{mover \ovalnum{50} pasos}
\end{scratch}
&
\begin{scratch}
\blockinit{al presionar tecla \selectmenu{i}}
\blockmove{girar \turnleft{} \ovalnum{90} grados}
\end{scratch}
&
\begin{scratch}
\blockinit{al presionar tecla \selectmenu{d}}
\blockmove{girar \turnright{} \ovalnum{90} grados}
\end{scratch}
\end{tabular}
\end{center}

\begin{typethis}
from ColabTurtlePlus.Turtle import *
clearscreen()
initializeTurtle((480, 360))

shape("turtle")
for turno in range(5):
    tecla = input("¿a, i o d? ")
    if tecla == "a":
        forward(50)
    elif tecla == "i":
        left(90)
    elif tecla == "d":
        right(90)
\end{typethis}

\textbf{Cómo funciona.} Tres programas de Scratch se convierten en un solo bucle con cinco turnos. En cada turno, \py|input| espera una letra. Luego el \py|if| revisa qué letra fue y hace el trabajo del bloque de sombrero correspondiente. \py|elif| significa «si no, si...» (en inglés \emph{else if}): solo se revisa cuando las líneas de arriba no fueron verdaderas. Prueba a escribir \texttt{a}, \texttt{i}, \texttt{a}, \texttt{d}, \texttt{a}. ¡La tortuga dibuja un escalón en zigzag!

\begin{watchout}[¡Cuidado! Dale siempre una salida a tu bucle]
En Scratch, \chip[control]{por siempre} es seguro porque puedes hacer clic en la señal roja de alto. En Colab, \py|while True:| sin salida hace que la celda corra para siempre, y no puede pasar nada más. Usa un bucle \py|for| como \py|for turno in range(5):|, o un \py|while| cuya condición llegue a ser falsa, como el de la Carrera de tortugas más adelante en esta parte. Si una celda se queda atorada, presiona el botón de detener \stopbtn.
\end{watchout}

% ---------------------------------------------------------------------
\needspace{12\baselineskip}
\subsection*{Muchos objetos, un solo bucle}

En Scratch, cada objeto puede tener su propio programa \chip[events]{al presionar bandera verde}, y todos parecen correr al mismo tiempo. Aquí un gato camina a la derecha mientras un perro camina hacia arriba:

\begin{center}
\begin{tabular}{c@{\hspace{1cm}}c}
\textbf{Objeto Gato} & \textbf{Objeto Perro} \\[4pt]
\begin{scratch}
\blockinit{al presionar \greenflag}
\blockrepeat{repetir \ovalnum{20}}{
  \blockmove{mover \ovalnum{10} pasos}
}
\end{scratch}
&
\begin{scratch}
\blockinit{al presionar \greenflag}
\blockmove{apuntar en dirección \ovalnum{0}}
\blockrepeat{repetir \ovalnum{20}}{
  \blockmove{mover \ovalnum{5} pasos}
}
\end{scratch}
\end{tabular}
\end{center}

Python hace \textbf{una cosa a la vez}. Si escribieras el bucle del gato y después el bucle del perro, el gato terminaría toda su caminata antes de que el perro siquiera empezara. En lugar de eso, escribe \textbf{un solo bucle que mueva un poquito a cada objeto}:

\begin{typethis}
from ColabTurtlePlus.Turtle import *
clearscreen()
initializeTurtle((480, 360))

hideturtle()
gato = Turtle()
gato.shape("turtle")
perro = Turtle()
perro.shape("circle")
perro.left(90)
for paso in range(20):
    gato.forward(10)
    perro.forward(5)
\end{typethis}

\textbf{Cómo funciona.} Cada vuelta del bucle, el gato da un pasito y luego el perro da un pasito. Pasa tan rápido que parecen moverse juntos. Esto se llama un \textbf{bucle de juego}, y todos los videojuegos de verdad tienen uno. Scratch también tiene uno, escondido por dentro, como un malabarista que le da un empujoncito a cada pelota.


% =====================================================================
\mission{Desfile de disfraces}{listas de disfraces y colores, \texttt{shapesize}, \texttt{\%} y \texttt{time.sleep}}

¡Tu tortuga va a dar un desfile de moda! Cambia de disfraz, de color y de tamaño cada medio segundo, luego hace una reverencia y desaparece.

\step{Paso 1: Constrúyelo en Scratch}
Dale cuatro disfraces a un objeto y luego construye:

\begin{center}
\begin{scratch}
\blockinit{al presionar \greenflag}
\blocklook{cambiar fondo a \selectmenu{lavanda}}
\blocklook{cambiar disfraz a \selectmenu{tortuga}}
\blocklook{fijar tamaño a \ovalnum{100} \%}
\blockrepeat{repetir \ovalnum{8}}{
  \blocklook{cambiar el efecto \selectmenu{color} en \ovalnum{25}}
  \blockcontrol{esperar \ovalnum{0.5} segundos}
  \blocklook{siguiente disfraz}
  \blocklook{cambiar tamaño en \ovalnum{50}}
}
\blocklook{esconder}
\blocklook{decir \ovalnum{¡Tarán!}}
\end{scratch}
\end{center}

\step{Paso 2: Escríbelo en Colab}

\begin{typethis}
from ColabTurtlePlus.Turtle import *
clearscreen()
initializeTurtle((480, 360))
import time

disfraces = ["turtle", "circle", "square", "triangle"]
colores = ["green", "red", "blue", "orange"]
n = 0
medida = 100
bgcolor("lavender")
for i in range(8):
    shape(disfraces[n])
    color(colores[n])
    shapesize(medida / 100)
    print("Disfraz", n + 1, disfraces[n])
    time.sleep(0.5)
    n = (n + 1) % len(disfraces)
    medida = medida + 50
hideturtle()
print("¡Tarán!")
\end{typethis}

\step{Paso 3: Léelo línea por línea}

\begin{longtable}{>{\centering\arraybackslash}p{1.1cm} P{5.3cm} P{9.1cm}}
\toprule
\textbf{Línea} & \textbf{Bloque de Scratch} & \textbf{Qué significa} \\
\midrule
1--4 & \chip[events]{al presionar bandera verde} & Las tres líneas de inicio, más las herramientas del reloj. \\
6 & (la pestaña Disfraces) & El ropero: una lista de cuatro disfraces. \\
7 & \chip[looks]{cambiar el efecto color} & Una lista de cuatro colores, uno para cada disfraz. \\
8 & \chip[looks]{cambiar disfraz a tortuga} & El gancho empieza en 0, el primer disfraz. \\
9 & \chip[looks]{fijar tamaño a 100 \%} & Tu propia variable de tamaño, en porcentaje como en Scratch. \\
10 & \chip[looks]{cambiar fondo a lavanda} & Pinta el escenario de morado clarito (\texttt{lavender} es lavanda). \\
11 & \chip[control]{repetir 8} & Haz 8 veces las líneas de adentro. \\
12--14 & & Ponte el disfraz número \py|n|, con el color número \py|n|, en el tamaño correcto. \\
15 & (el número de disfraz) & Di qué disfraz es. \py|print| puede mostrar varias cosas a la vez si pones comas entre ellas, y agrega los espacios por ti. \py|n + 1| porque Scratch cuenta los disfraces desde 1. \\
16 & \chip[control]{esperar 0.5 segundos} & Pausa para que puedas ver el disfraz. \\
17 & \chip[looks]{siguiente disfraz} & Pasa al siguiente gancho, y vuelve al 0 después del último. \\
18 & \chip[looks]{cambiar tamaño en 50} & Crece un 50\,\% para la próxima vez. \\
19 & \chip[looks]{esconder} & Desaparece\ldots \\
20 & \chip[looks]{decir ¡Tarán!} & \ldots y hace una reverencia. \\
\bottomrule
\end{longtable}

\step{Paso 4: Ejecútalo}
El escenario se vuelve morado clarito. La tortuga cambia de disfraz y de color cada medio segundo, y se hace más grande cada vez. Luego desaparece. Debajo de la celda verás:

\begin{outputbox}
Disfraz 1 turtle\\
Disfraz 2 circle\\
Disfraz 3 square\\
Disfraz 4 triangle\\
Disfraz 1 turtle\\
Disfraz 2 circle\\
Disfraz 3 square\\
Disfraz 4 triangle\\
¡Tarán!
\end{outputbox}

\begin{think}
El disfraz y el color siempre van juntos: el disfraz 2 siempre es rojo. Eso es porque las dos listas usan el mismo número de gancho \py|n|. Las dos listas son como dos percheros de ropa, uno al lado del otro.
\end{think}

\enlargethispage{\baselineskip}
\begin{challenge}\small
\tick Agrega un quinto disfraz, \py|"arrow"|, al final de la lista de disfraces. ¿Qué más tienes que agregar para que la línea 13 no truene? (Pista: cuenta los colores.)\par
\tick Haz que la tortuga se \textbf{encoja}: empieza con \py|medida = 400| y réstale 40 cada vez. ¡Revisa que nunca llegue a 0!\par
\tick Agrega una lista de colores de fondo, como \py|["lavender", "pink", "skyblue", "lightyellow"]|, y cambia también el fondo cada vez.\par
\tick Usa una tortuga globo para escribir el nombre del disfraz en el escenario, encima de la tortuga, en lugar de imprimirlo. Acuérdate de borrar primero el nombre anterior con \py|clear()|.
\end{challenge}

% =====================================================================
\mission{Hacer música}{a tocar una melodía con una lista de notas, y \texttt{range(len(...))}}

Ahora vamos a hacer que Colab cante. Vas a tocar el principio de «Estrellita, ¿dónde estás?». Antes de empezar, ejecuta la celda «En el principio» y la \textbf{celda de sonido}.

\step{Paso 1: Constrúyelo en Scratch}
Haz clic en el botón \textbf{Agregar Extensión} y agrega \textbf{Música}. Sus bloques son verdes en Scratch. Este libro los dibuja en rosa, junto con los bloques de sonido.

\begin{center}
\begin{scratch}
\blockinit{al presionar \greenflag}
\blocksound{fijar tempo a \ovalnum{150}}
\blocksound{tocar nota \ovalnum{60} durante \ovalnum{1} tiempos}
\blocksound{tocar nota \ovalnum{60} durante \ovalnum{1} tiempos}
\blocksound{tocar nota \ovalnum{67} durante \ovalnum{1} tiempos}
\blocksound{tocar nota \ovalnum{67} durante \ovalnum{1} tiempos}
\blocksound{tocar nota \ovalnum{69} durante \ovalnum{1} tiempos}
\blocksound{tocar nota \ovalnum{69} durante \ovalnum{1} tiempos}
\blocksound{tocar nota \ovalnum{67} durante \ovalnum{2} tiempos}
\blocklook{decir \ovalnum{Estrellita, ¿dónde estás?}}
\end{scratch}
\end{center}

\step{Paso 2: Escríbelo en Colab}

\begin{typethis}
import time
Do = 262
Re = 294
Mi = 330
Fa = 349
Sol = 392
La = 440

melodia = [Do, Do, Sol, Sol, La, La, Sol]
tiempos = [1, 1, 1, 1, 1, 1, 2]
for i in range(len(melodia)):
    segs = tiempos[i] * 0.4
    pitido(melodia[i], segs)
    time.sleep(segs + 0.1)
print("Estrellita, ¿dónde estás?")
\end{typethis}

\step{Paso 3: Léelo línea por línea}

\begin{longtable}{>{\centering\arraybackslash}p{1.1cm} P{5.3cm} P{9.1cm}}
\toprule
\textbf{Línea} & \textbf{Bloque de Scratch} & \textbf{Qué significa} \\
\midrule
--- & \chip[events]{al presionar bandera verde} & No hace falta la tortuga, así que no hay líneas de inicio. La celda de sonido ya debe estar ejecutada. \\
1 & & Las herramientas del reloj, para esperar. \\
2--7 & (los números de las notas) & Seis variables, una para cada nota. Ahora puedes escribir \py|Sol| en lugar de acordarte de 392. Las mayúsculas sirven en los nombres de variables, y así parecen notas de verdad. \\
9 & los siete bloques \chip[sound]{tocar nota} & La melodía es una lista de notas, en orden. \\
10 & los tiempos & Cuántos tiempos dura cada nota. La última nota dura el doble. \\
11 & \chip[control]{repetir 7} & \py|len(melodia)| es 7, así que \py|range(7)| da \py|i| = 0, 1, 2 \ldots{} 6: un número para cada nota. \\
12 & \chip[sound]{fijar tempo a 150} & Un tiempo dura 0.4 segundos. Multiplica para saber cuánto dura esta nota. \\
13 & \chip[sound]{tocar nota durante tiempos} & Toca la nota número \py|i| durante el tiempo correcto. \\
14 & («hasta que termine») & Espera a que la nota termine, más un huequito para que oigas dos Do, y no uno largo. \\
15 & \chip[looks]{decir} & Se imprime cuando termina la melodía. \\
\bottomrule
\end{longtable}

\begin{think}[¿Por qué \texttt{range(len(melodia))}?]
La línea 11 cuenta las posiciones, de 0 a 6, en lugar de las notas mismas. Eso es porque necesitas la \textbf{misma posición} de dos listas: \py|melodia[i]| es la nota y \py|tiempos[i]| es cuánto dura. Es el mismo truco de los dos percheros del Desfile de disfraces.
\end{think}

\step{Paso 4: Ejecútalo}
Deberías oír «Estrellita, ¿dónde estás?» y ver aparecer un pequeño reproductor de sonido por cada nota. Luego aparecen las palabras debajo de la celda:

\begin{outputbox}
Estrellita, ¿dónde estás?
\end{outputbox}

Si no oyes nada, vuelve a mirar el recuadro \textbf{¿No se oye nada?} en la página de la celda de sonido.

\begin{challenge}
\tick Hazla más rápida: cambia \py|0.4| en la línea 12 por \py|0.25|. Es como \chip[sound]{fijar tempo a 240}.\par
\tick Agrega la siguiente línea de la canción, «me pregunto qué serás»: las notas son Fa Fa Mi Mi Re Re Do, y la última es larga. Agrégalas al final de las dos listas.\par
\tick \textbf{¡Feliz cumpleaños!} Agrega \py|Si = 494| y \py|Do2 = 523|, y luego toca \py|[Sol, Sol, La, Sol, Do2, Si]| con los tiempos \py|[0.75, 0.25, 1, 1, 1, 2]|.\par
\tick Haz una tortuga bailarina: agrega las tres líneas de inicio arriba y pon \py|right(30)| dentro del bucle, para que la tortuga gire con cada nota.
\end{challenge}

% =====================================================================
\mission{Carrera de tortugas}{dos objetos en un bucle, una función que crea objetos, e \texttt{if}/\texttt{else}}

Dos tortugas corren a través del escenario. En cada paso del bucle, las dos avanzan una cantidad al azar. ¿Quién ganará? ¡Ni la computadora lo sabe!

\step{Paso 1: Constrúyelo en Scratch}
Crea dos objetos llamados \textbf{Rojo} y \textbf{Azul}. El rojo lleva este programa, y el azul lleva el mismo programa con $y = -50$:

\begin{center}
\begin{scratch}
\blockinit{al presionar \greenflag}
\blocklook{cambiar disfraz a \selectmenu{tortuga}}
\blockpen{subir lápiz}
\blockmove{ir a x: \ovalnum{-200} y: \ovalnum{50}}
\blockpen{bajar lápiz}
\blockrepeat{repetir hasta que \booloperator{\ovalmove{posición en x} > \ovalnum{200}}}{
  \blockmove{cambiar x en \ovaloperator{elegir número al azar entre \ovalnum{1} y \ovalnum{20}}}
}
\blocklook{decir \ovalnum{¡Gané!}}
\end{scratch}
\end{center}

En Scratch, cada objeto corre su propio programa \textbf{al mismo tiempo}. En Python, \textbf{un solo bucle mueve a las dos tortugas}, un paso cada una, una y otra vez. Y como los dos objetos se preparan de la misma forma, una función puede prepararlos por ti.

\step{Paso 2: Escríbelo en Colab}

\begin{typethis}
from ColabTurtlePlus.Turtle import *
clearscreen()
initializeTurtle((480, 360))
import random

def crear_corredor(su_color, y):
    corredor = Turtle()
    corredor.shape("turtle")
    corredor.color(su_color)
    corredor.penup()
    corredor.goto(-200, y)
    corredor.pendown()
    return corredor

hideturtle()
rojo = crear_corredor("red", 50)
azul = crear_corredor("blue", -50)

while rojo.xcor() < 200 and azul.xcor() < 200:
    rojo.forward(random.randint(1, 20))
    azul.forward(random.randint(1, 20))

if rojo.xcor() > azul.xcor():
    print("¡Gana Rojo!")
    rojo.write("¡Gané!", font=("Arial", 16, "bold"))
else:
    print("¡Gana Azul!")
    azul.write("¡Gané!", font=("Arial", 16, "bold"))
\end{typethis}

\step{Paso 3: Léelo línea por línea}

\begin{longtable}{>{\centering\arraybackslash}p{1.1cm} P{5.3cm} P{9.1cm}}
\toprule
\textbf{Línea} & \textbf{Bloque de Scratch} & \textbf{Qué significa} \\
\midrule
1--4 & \chip[events]{al presionar bandera verde} & Las tres líneas de inicio, más las herramientas de números al azar. \\
6 & (una fábrica de objetos) & Una función con dos entradas: el color del corredor y a qué altura empieza. \\
7 & (un objeto nuevo) & Crea un objeto nuevecito y llámalo \py|corredor| por ahora. \\
8--9 & \chip[looks]{cambiar disfraz a tortuga} & El disfraz de tortuga, en el color que pediste. \\
10--12 & \chip[pen]{subir lápiz}\newline\chip[motion]{ir a x: -200 y: 50}\newline\chip[pen]{bajar lápiz} & Ve a la línea de salida sin dibujar, y luego baja el lápiz para que deje un rastro. \\
13 & & \py|return| le entrega el objeto terminado a quien lo pidió. \\
15 & \chip[looks]{esconder} & Esconde la tortuga normal del centro: ella no corre. \\
16--17 & (dos objetos) & Usa la fábrica dos veces. Cada llamada crea un objeto nuevo, y tú le pones nombre. \\
19 & \chip[control]{repetir hasta que}\newline\chip[operators]{posición en x > 200} & Sigue \textbf{mientras} (\emph{while}) rojo no haya llegado a 200 \textbf{y} (\emph{and}) azul no haya llegado a 200. Este es el bucle de juego, y tiene una salida: en cuanto una tortuga cruza la meta, se detiene. \\
20--21 & \chip[motion]{cambiar x en}\newline\chip[operators]{elegir número al azar} & Un pasito para cada tortuga, en cada vuelta. \\
23 & \chip[control]{si entonces si no} & ¿Está rojo más a la derecha que azul? \\
24--25 & \chip[looks]{decir ¡Gané!} & Gana rojo: imprímelo y escríbelo en el escenario junto a rojo, en negritas. \\
26--28 & (la parte del si no) & Si no, gana azul. \\
\bottomrule
\end{longtable}

\begin{think}[¿Por qué una función?]
Sin \py|crear_corredor|, tendrías que escribir dos veces las seis líneas de preparación, una para rojo y otra para azul. Con ella, agregar un corredor verde es solo una línea más: \py|verde = crear_corredor("green", 0)|. Una función es una tarjeta de receta que puedes usar una y otra vez.
\end{think}

\step{Paso 4: Ejecútalo}
¡Mira cómo corren las tortugas! La ganadora escribe \textbf{¡Gané!} en la meta, y debajo del dibujo verás uno de estos:

\begin{outputbox}
¡Gana Rojo!
\end{outputbox}

Ejecútalo otra vez. ¿Gana la misma cada vez?

\begin{challenge}
\tick Agrega un tercer corredor: \py|verde = crear_corredor("green", 0)|. Luego cambia la línea 19 para que la carrera se detenga cuando \textbf{cualquiera} de los tres cruce la meta, y agrega una línea de \py|forward| para verde. (Pista: tres tortugas significan \texttt{and} dos veces.)\par
\tick Haz que una tortuga haga trampa: que avance \py|random.randint(1, 25)|. Corre la carrera 5 veces. ¿Quién gana más?\par
\tick Dale a \py|crear_corredor| una tercera entrada, \py|disfraz|, para que un corredor pueda ser un \py|"circle"| y el otro un \py|"triangle"|.\par
\tick Toca un sonido cuando termine la carrera: ejecuta primero la celda de sonido y luego agrega \py|pitido(523, 1)| justo al final.
\end{challenge}

% =====================================================================
\mission{Cuento con mensajes}{funciones como mensajes, y decisiones con \texttt{input()}}

En Scratch, los cuentos grandes muchas veces se construyen con mensajes. El escenario envía \texttt{escena 1}, luego hace una pregunta y luego envía un mensaje distinto según la respuesta. En Python, cada mensaje es una función, y cada envío es una llamada a la función.

\step{Paso 1: Constrúyelo en Scratch}
Estos programas pueden ir en el escenario o en un objeto. El de la izquierda dirige el cuento. Los otros esperan sus mensajes. (Los programas \texttt{camino izquierda} y \texttt{el fin} se parecen al de \texttt{camino derecha}.)

\begin{center}
\begin{tabular}{c@{\hspace{0.5cm}}c}
\begin{scratch}[scale=0.75]
\blockinit{al presionar \greenflag}
\blockevent{enviar \selectmenu{escena 1} y esperar}
\blocksensing{preguntar \ovalnum{¿Izquierda o derecha?} y esperar}
\blockifelse{si \booloperator{\ovalsensing{respuesta} = \ovalnum{izquierda}} entonces}{
  \blockevent{enviar \selectmenu{camino izquierda} y esperar}
}{
  \blockevent{enviar \selectmenu{camino derecha} y esperar}
}
\blockevent{enviar \selectmenu{el fin} y esperar}
\end{scratch}
&
\begin{tabular}{c}
\begin{scratch}[scale=0.75]
\blockinit{al recibir \selectmenu{escena 1}}
\blocklook{decir \ovalnum{Una tortuguita despierta...} durante \ovalnum{2} segundos}
\blocklook{decir \ovalnum{Dos caminos salen de aquí...}}
\end{scratch}
\\[8pt]
\begin{scratch}[scale=0.75]
\blockinit{al recibir \selectmenu{camino derecha}}
\blocklook{decir \ovalnum{¡Conoces a un dragón amistoso!}}
\end{scratch}
\end{tabular}
\end{tabular}
\end{center}

\step{Paso 2: Escríbelo en Colab}

\begin{typethis}
import time

def escena1():
    print("Una tortuguita despierta en un bosque.")
    time.sleep(2)
    print("Dos caminos salen de aquí.")

def camino_izquierda():
    print("Encuentras un río lleno de peces. ¡Ñam!")

def camino_derecha():
    print("¡Conoces a un dragón amistoso!")

def el_fin():
    time.sleep(1)
    print("FIN")

escena1()
eleccion = input("¿Izquierda o derecha? ")
if eleccion == "izquierda":
    camino_izquierda()
else:
    camino_derecha()
el_fin()
\end{typethis}

\step{Paso 3: Léelo línea por línea}

\begin{longtable}{>{\centering\arraybackslash}p{1.1cm} P{5.3cm} P{9.1cm}}
\toprule
\textbf{Línea} & \textbf{Bloque de Scratch} & \textbf{Qué significa} \\
\midrule
1 & & Las herramientas del reloj. No hace falta la tortuga: este cuento es solo de palabras. \\
3--6 & \chip[events]{al recibir escena 1} & El programa \texttt{escena 1}. Python lo aprende ahora, y lo hace después, cuando lo llamen. \\
8--9 & \chip[events]{al recibir camino izquierda} & El programa \texttt{camino izquierda}. Los nombres de funciones no pueden tener espacios, así que usamos un guion bajo: \py|camino_izquierda|. \\
11--12 & \chip[events]{al recibir camino derecha} & El programa \texttt{camino derecha}. \\
14--16 & \chip[events]{al recibir el fin} & El programa \texttt{el fin}. \\
18 & \chip[events]{enviar escena 1 y esperar} & ¡Aquí empieza el cuento! Llama a \py|escena1| y espera a que termine. \\
19 & \chip[sensing]{preguntar y esperar} & Haz la pregunta y guarda la respuesta en \py|eleccion|. \\
20--23 & \chip[control]{si entonces si no} con dos envíos & Llama a la función que corresponde a la respuesta. \\
24 & \chip[events]{enviar el fin y esperar} & Esta línea está otra vez pegada a la orilla, así que pasa sin importar qué camino elegiste. \\
\bottomrule
\end{longtable}

\begin{watchout}[El cuento empieza en la línea 18, no en la línea 1]
Las líneas 3 a 16 solo le \textbf{enseñan} las escenas a Python. No se imprime nada hasta que la línea 18 llama a la primera. Si olvidas las llamadas de abajo, ejecutarás la celda y verás\ldots{} ¡nada de nada! Es como construir programas \chip[events]{al recibir} y nunca enviar los mensajes.
\end{watchout}

\step{Paso 4: Ejecútalo}
Escribe \texttt{izquierda} o \texttt{derecha} cuando aparezca la pregunta (sin acento ni mayúscula), y luego presiona Enter.

\begin{outputbox}[Podrías ver]
Una tortuguita despierta en un bosque.\\
Dos caminos salen de aquí.\\
¿Izquierda o derecha? izquierda\\
Encuentras un río lleno de peces. ¡Ñam!\\
FIN
\end{outputbox}

\begin{think}
¿Qué pasa si escribes \texttt{Izquierda} con I mayúscula? Python piensa que \texttt{"Izquierda"} y \texttt{"izquierda"} son distintas, ¡así que te vas a la derecha! Puedes arreglarlo cambiando la línea 19 por \py|eleccion = input("¿Izquierda o derecha? ").lower()|, que convierte la respuesta a minúsculas.
\end{think}

\enlargethispage{\baselineskip}
\begin{challenge}\small
\tick Agrega una tercera opción, \texttt{derecho}, con su propia función. Vas a necesitar un \py|elif|.\par
\tick Dentro de \py|camino_derecha|, haz otra pregunta («¿Quieres ser su amigo? si o no») y llama a una función diferente para cada respuesta. Usa \texttt{si} sin acento, para que sea fácil de escribir.\par
\tick Haz que el cuento sea musical: ejecuta la celda de sonido y luego pon un \py|pitido| dentro de cada escena. Una nota grave (\py|pitido(262)|) para el miedo y una aguda (\py|pitido(523)|) para la alegría.\par
\tick Agrega un fondo: pon las tres líneas de inicio arriba, y usa \py|bgcolor("darkgreen")| en \py|escena1| (el bosque) y \py|bgcolor("skyblue")| en \py|camino_izquierda| (el río).
\end{challenge}

% =====================================================================
\banner{events}{EXTRA}{Teclas en tu computadora: Thonny}

La tortuga de Colab no puede oír las teclas, pero Python en tu propia computadora sí. \textbf{Thonny} (\texttt{thonny.org}) es un editor de Python gratuito hecho para principiantes. Funciona sin internet y no necesita cuenta. Pídele a un adulto que lo instale. Thonny trae la \texttt{turtle} propia de Python, que abre su propia ventana y \textbf{sí} puede escuchar el teclado.

Este es \chip[events]{al presionar tecla espacio} en Thonny:

\begin{center}
\begin{tabular}{c}
\begin{scratch}
\blockinit{al presionar tecla \selectmenu{espacio}}
\blockmove{girar \turnleft{} \ovalnum{90} grados}
\blockmove{mover \ovalnum{50} pasos}
\blockmove{mover \ovalnum{-50} pasos}
\blockmove{girar \turnright{} \ovalnum{90} grados}
\end{scratch}
\end{tabular}
\end{center}

\begin{typethis}[Escribe esto en Thonny y presiona el botón verde de ejecutar (Run)]
import turtle

t = turtle.Turtle()
t.shape("turtle")
pantalla = turtle.Screen()

def saltar():
    t.left(90)
    t.forward(50)
    t.back(50)
    t.right(90)

pantalla.onkeypress(saltar, "space")
pantalla.listen()
turtle.mainloop()
\end{typethis}

\begin{longtable}{>{\centering\arraybackslash}p{1.1cm} P{14.4cm}}
\toprule
\textbf{Línea} & \textbf{Qué significa} \\
\midrule
1--5 & Las líneas de inicio de Thonny. No hay \texttt{pip install}: la tortuga ya viene incluida. La tortuga se llama \py|t|, así que cada comando empieza con \py|t.|, y \py|pantalla| es la ventana. \\
7--11 & Una función: el programa que va debajo del bloque de sombrero. \\
13 & «Cuando se presione la tecla espacio, ejecuta \py|saltar|». ¡Este es el bloque de sombrero! Los nombres de las teclas son quisquillosos y van en inglés: \py|"space"| (espacio), \py|"Right"|, \py|"Left"|, \py|"Up"|, \py|"Down"| (las flechas), o una letra como \py|"a"|. \\
14 & Dile a la ventana que empiece a escuchar el teclado. \\
15 & Mantén la ventana abierta y sigue esperando teclas. Esta debe ser la \textbf{última} línea. \\
\bottomrule
\end{longtable}

Ejecútalo, luego \textbf{haz clic en la ventana de la tortuga} para que sepa que le estás hablando a ella, y presiona la barra espaciadora. ¡La tortuga salta! ¿Puedes agregar una función \py|ir_derecha| para la flecha \py|"Right"|?

\begin{watchout}[Tres cosas que le pasan a todo el mundo]
\begin{itemize}
  \item \textbf{Sin paréntesis después del nombre de la función.} Escribe \py|pantalla.onkeypress(saltar, "space")|, no \py|saltar()|. Le estás diciendo a Python \emph{qué} bloque ejecutar después, no lo estás ejecutando ahora.
  \item \textbf{Usa \texttt{onkeypress}, no \texttt{onkey}.} \py|onkey| espera hasta que sueltas la tecla, y se siente lento. \py|onkeypress| funciona en el momento en que la presionas, como en Scratch.
  \item \textbf{Haz clic primero en la ventana de la tortuga}, o tus teclas irán al código.
\end{itemize}
\end{watchout}

En Thonny, \py|t.setheading()| usa la misma brújula que Colab, así que la regla de \textbf{90 menos} sigue funcionando.

% =====================================================================
\needspace{30\baselineskip}
\subsection*{Tradúcelo tú}

Este es un programa de Scratch que no has visto antes. Escribe el Python en las líneas. Empieza con las tres líneas de inicio e \py|import time|. Luego escríbelo en Colab y revisa que la tortuga haga lo mismo que el objeto de Scratch.

\begin{tcbraster}[raster columns=2, raster equal height=rows, raster column skip=12pt, raster valign=top]
\begin{tcolorbox}[colback=white, colframe=white, boxrule=0pt]
\begin{scratch}
\blockinit{al presionar \greenflag}
\blocklook{cambiar disfraz a \selectmenu{tortuga}}
\blocklook{fijar tamaño a \ovalnum{300} \%}
\blockrepeat{repetir \ovalnum{3}}{
  \blocklook{esconder}
  \blockcontrol{esperar \ovalnum{0.5} segundos}
  \blocklook{mostrar}
  \blockcontrol{esperar \ovalnum{0.5} segundos}
}
\blocklook{decir \ovalnum{¡Bu!} durante \ovalnum{2} segundos}
\end{scratch}
\end{tcolorbox}
\begin{tcolorbox}[colback=white, colframe=codeframe, arc=5pt, title=\textbf{Mi Python}, coltitle=ink, colbacktitle=codebg]
\vspace{4pt}
\foreach \n in {1,...,13}{\rule{\linewidth}{0.4pt}\par\vspace{9pt}}
\end{tcolorbox}
\end{tcbraster}

% ---------------------------------------------------------------------
\needspace{12\baselineskip}
\subsection*{Acertijo de dos líneas}

\begin{puzzle}
Escribe el Python para este programa de Scratch. ¡Solo necesita dos líneas!
\begin{center}
\begin{scratch}
\blockinit{al recibir \selectmenu{festejo}}
\blocklook{decir \ovalnum{¡Hurra!}}
\end{scratch}
\end{center}
Luego agrega una tercera línea que \textbf{envíe} el mensaje \texttt{festejo}, y ejecútalo.
\end{puzzle}

% ---------------------------------------------------------------------
\needspace{14\baselineskip}
\subsection*{Lista de la parte 7}

\tick Puedo hacer que la tortuga diga cosas con \py|print|, y en el escenario con una tortuga globo.\par
\tick Puedo cambiar el disfraz, el tamaño y el color de la tortuga y el fondo, y hacer que se esconda y se muestre.\par
\tick Puedo usar una lista y \texttt{\%} para hacer \chip[looks]{siguiente disfraz}.\par
\tick Puedo tocar notas y una melodía en Colab con \py|pitido| y \py|time.sleep|.\par
\tick Puedo convertir \chip[events]{al recibir} en una función y \chip[events]{enviar} en una llamada.\par
\tick Puedo mover dos objetos con un solo bucle, y siempre le doy una salida a mis bucles.

% =====================================================================
% 08-projects.tex : GRANDES PROYECTOS (después de las Partes 1 a 7)
% Ocho proyectos en orden de dificultad y luego «Tu propio proyecto».
% Todos los programas probados: juegos de texto con entrada por tubería
% (python3 -I), juegos de dibujo con la biblioteca real. Los programas
% de dibujo también están en 08-colab-tests.py.
% =====================================================================

\partopener{pen}{GRANDES PROYECTOS}{Grandes proyectos}{Ocho proyectos de verdad que usan todo lo que sabes, y luego uno tuyo.}

Terminaste las siete Partes. Sabes dibujar, repetir con bucles, guardar cosas en variables y listas, hacer cuentas, elegir números al azar, hacer preguntas, tomar decisiones con \py|if|, repetir con \py|while|, crear tus propias funciones y hacer objetos extra. ¡Eso es una caja de herramientas completa! Ahora toca construir \textbf{cosas grandes} con ella.

Cada proyecto de aquí es un proyecto clásico de Scratch, reconstruido en Python. Se van poniendo un poco más difíciles, así que hazlos en orden:

\begin{center}
\begin{tabular}{rll}
\toprule
& \textbf{Proyecto} & \textbf{La gran idea} \\
\midrule
1 & Anima tu nombre & bucles que recorren letras, \py|write()| y colores \\
2 & Flor de lápiz & una función dentro de un bucle, cuentas en los giros \\
3 & Concurso de tablas & una función que devuelve puntos, un cronómetro \\
4 & Adivina el número de lujo & niveles, intentos y un bucle de «otra vez» \\
5 & Elige tu propia historia & funciones como escenas \\
6 & Mascota virtual & variables que cambian en cada turno \\
7 & Juego de persecución & tres objetos, \py|towards()| y \py|distance()| \\
8 & Juego de atrapar & un objeto que cae y un bucle de juego de verdad \\
\bottomrule
\end{tabular}
\end{center}

\subsection*{Cómo hacer un gran proyecto}
Cada proyecto tiene cuatro partes. Trabájalas así:
\begin{enumerate}[label=\textbf{\arabic*.}]
  \item \textbf{La versión de Scratch.} Mira el programa de Scratch. Seguro que ya conoces un proyecto parecido. Si quieres, constrúyelo primero en Scratch.
  \item \textbf{La versión de Python.} \textbf{Predice} qué va a hacer. Luego escríbelo y \textbf{ejecútalo}. Es más largo que las misiones, así que tómate tu tiempo y revisa cada dos puntos y cada paréntesis.
  \item \textbf{Cómo funciona.} Lee la tabla. Señala cada línea de tu código mientras lees sobre ella.
  \item \textbf{Sube de nivel.} ¡\textbf{Cámbialo}! Elige al menos un reto. Luego agrega una idea tuya. Esta es la parte más importante: el proyecto se vuelve \emph{tuyo}.
\end{enumerate}

\begin{think}[En Colab los juegos van por turnos]
Colab no puede sentir tu teclado mientras un programa se ejecuta, así que no hay \chip[events]{al presionar tecla espacio}. En su lugar, los juegos de aquí te preguntan qué hacer con \py|input()|, luego mueven todo y luego vuelven a preguntar. Es como un juego de mesa: tú juegas tu turno y luego la computadora juega el suyo. Las ideas del juego son exactamente las mismas que en Scratch: moverse, revisar si algo se toca, cambiar los puntos.
\end{think}

\begin{watchout}[Un programa largo se escribe poco a poco]
No escribas 40 líneas para luego presionar \playbtn{} por primera vez. Escribe la primera parte (hasta una línea en blanco), ejecútala y revisa que funcione. Luego agrega la siguiente parte. Los programadores hacen esto todo el tiempo. Si aparece un error, está en la parte que acabas de agregar, así que es fácil de encontrar.
\end{watchout}

\needspace{12\baselineskip}
\subsection*{Mis grandes proyectos}
Marca cada proyecto cuando funcione, y otra vez cuando lo hayas subido de nivel.

\begin{center}
\begin{tabular}{lcc}
\toprule
\textbf{Proyecto} & \textbf{¡Funciona!} & \textbf{Lo subí de nivel} \\
\midrule
1 Anima tu nombre & \tick & \tick \\
2 Flor de lápiz & \tick & \tick \\
3 Concurso de tablas & \tick & \tick \\
4 Adivina el número de lujo & \tick & \tick \\
5 Elige tu propia historia & \tick & \tick \\
6 Mascota virtual & \tick & \tick \\
7 Juego de persecución & \tick & \tick \\
8 Juego de atrapar & \tick & \tick \\
Mi propio proyecto & \tick & \tick \\
\bottomrule
\end{tabular}
\end{center}

% =====================================================================
\mission{Anima tu nombre}{a recorrer las letras de una palabra con un bucle, \texttt{write()}, colores al azar y \texttt{clear()}}

Tu nombre baila por el escenario en colores que destellan. ¡Cada letra salta arriba y abajo por su cuenta!

\step{La versión de Scratch}
En Scratch harías un objeto para cada letra. Cada objeto-letra recibe un programa como este:

\begin{center}
\begin{scratch}
\blockinit{al presionar \greenflag}
\blockrepeat{repetir \ovalnum{8}}{
  \blocklook{fijar efecto \selectmenu{color} a \ovaloperator{elegir número al azar entre \ovalnum{0} y \ovalnum{200}}}
  \blockmove{ir a x: \ovalmove{posición en x} y: \ovaloperator{elegir número al azar entre \ovalnum{-20} y \ovalnum{20}}}
  \blockcontrol{esperar \ovalnum{0.3} segundos}
}
\end{scratch}
\end{center}

En Python no necesitas un objeto para cada letra. Un bucle \py|for| visita las letras una por una, y la tortuga \textbf{escribe} cada letra en su lugar.

\begin{think}[Primero predice]
Antes de escribir, mira el Python de la página siguiente. Encuentra la línea que pregunta tu nombre. Encuentra los dos bucles \py|for|: uno son los 8 cuadros de la animación y el otro son las letras. ¿Cuál es cuál? ¿Qué crees que pasa si tu nombre es muy largo?
\end{think}

\needspace{40\baselineskip}
\step{La versión de Python}
\begin{typethis}
from ColabTurtlePlus.Turtle import *
clearscreen()
initializeTurtle((480, 360))
import random
import time

bgcolor("black")
hideturtle()
penup()
speed(13)
colores = ["red", "orange", "gold", "lime", "skyblue",
           "violet", "hotpink"]

nombre = input("¿Cómo te llamas? ")
espacio = 40
if len(nombre) > 10:
    espacio = 400 / len(nombre)
inicio = 0 - (len(nombre) - 1) * espacio / 2

for cuadro in range(8):
    clear()
    x = inicio
    for letra in nombre:
        color(random.choice(colores))
        goto(x, random.randint(-20, 20))
        write(letra, align="center",
              font=("Arial", 40, "bold"))
        x = x + espacio
    time.sleep(0.3)

goto(0, -120)
color("white")
write("¡Hola, " + nombre + "!", align="center",
      font=("Arial", 24, "italic"))
\end{typethis}

\begin{think}[Python habla inglés]
Python habla inglés: sus palabras (\py|for|, \py|print|, \py|write|, \py|goto|...) siempre van en inglés. Lo que \textbf{tú} inventas sí puede ir en español: \py|nombre|, \py|letra|, \py|colores|. Los colores se escriben en inglés: \py|"red"| es rojo, \py|"orange"| naranja, \py|"gold"| dorado, \py|"lime"| verde lima, \py|"skyblue"| azul cielo, \py|"violet"| violeta, \py|"hotpink"| rosa fuerte y \py|"black"| negro.
\end{think}

\needspace{14\baselineskip}
\step{Cómo funciona}
\begin{longtable}{>{\centering\arraybackslash}p{1.1cm} P{5.3cm} P{9.1cm}}
\toprule
\textbf{Línea} & \textbf{Bloque de Scratch} & \textbf{Qué significa} \\
\midrule
1--5 & \chip[events]{al presionar bandera verde} & Las tres líneas de inicio, y luego \py|random| y \py|time|. \\
7--10 & \chip[looks]{esconder} & Un escenario negro, y la tortuga se esconde y sube su lápiz. Solo queremos las letras, no la tortuga ni líneas. \\
11--12 & (una lista) & Los colores para elegir. Una lista larga puede seguir en la línea siguiente, porque Python sabe que el \texttt{[} todavía no se ha cerrado. \\
15--18 & \chip[operators]{longitud de} & \py|len(nombre)| cuenta las letras. Los nombres largos reciben un espacio más pequeño para que quepan en el escenario. \py|inicio| es donde va la primera letra, para que el nombre quede en el centro. \\
20--21 & \chip[control]{repetir 8}\newline\chip[pen]{borrar todo} & Ocho cuadros de animación. \py|clear()| borra las letras viejas de la tortuga, lista para escribirlas otra vez. \\
23 & & \py|for letra in nombre:| visita las letras una a la vez, como cuando lees una palabra con el dedo. \\
24--25 & \chip[looks]{fijar efecto color a}\newline\chip[motion]{ir a x: y:} & Un color al azar, y luego salta a una altura al azar. \\
26--27 & \chip[looks]{decir} & \py|write()| escribe la letra en el escenario. \py|align="center"| pone el centro de la letra donde está la tortuga; \py|font| da el nombre del tipo de letra, el tamaño y el estilo. \\
28 & \chip[motion]{cambiar x en 40} & Avanza para la siguiente letra. \\
29 & \chip[control]{esperar 0.3 segundos} & Una pausa corta, para que puedas ver cada cuadro. \\
31--34 & \chip[looks]{decir ¡Hola!} & Un saludo blanco abajo. \\
\bottomrule
\end{longtable}

\begin{think}[Otra vez, la ortografía del código]
La lista se llama \py|colores| (en español, porque el nombre lo inventamos nosotros), pero la orden es \py|color()|. Esa es la palabra de Python en inglés, y Python solo la entiende así.
\end{think}

\begin{challenge}[Sube de nivel]
\tick Haz 20 cuadros en vez de 8, y prueba \py|time.sleep(0.1)| para un baile más rápido.\par
\tick Agrega tus colores favoritos a la lista (solo en minúsculas: \py|"deeppink"|, \py|"aqua"|).\par
\tick Haz que las letras crezcan: escribe cada cuadro con tamaño \py|20 + cuadro * 5| en vez de 40.\par
\tick Haz un arcoíris: en vez de un color al azar, usa \py|colores[cuadro % 7]|, así todas las letras de un cuadro tienen el mismo color, y el color cambia en cada cuadro.
\end{challenge}

% =====================================================================
\mission{Flor de lápiz}{una función dentro de un bucle, cuentas en un giro, \texttt{\%} para repetir una lista y \texttt{circle()}}

Una flor con tantos pétalos como quieras, en colores de arcoíris, con un dibujo de espirógrafo que gira en el centro.

\step{La versión de Scratch}
Usa la extensión \textbf{Lápiz} y un bloque de \textbf{Mis Bloques} llamado \textbf{pétalo}:

\begin{center}
\begin{scratch}
\blockinit{al presionar \greenflag}
\blockpen{borrar todo}
\blockpen{bajar lápiz}
\blocksensing{preguntar \ovalnum{¿Cuántos pétalos?} y esperar}
\blockrepeat{repetir \ovalsensing{respuesta}}{
  \blockpen{cambiar \selectmenu{color} de lápiz en \ovalnum{10}}
  \blockmoreblocks{pétalo \ovalnum{100}}
  \blockmove{girar \turnleft{} \ovaloperator{\ovalnum{360} / \ovalsensing{respuesta}} grados}
}
\end{scratch}
\end{center}

\begin{think}[Primero predice]
Si escribes 6 pétalos, ¿cuánto girará la tortuga después de cada pétalo? (Calcula $360 \div 6$.) ¿Y con 8 pétalos? Busca la línea del Python que hace esta cuenta por ti.
\end{think}

\needspace{40\baselineskip}
\step{La versión de Python}
\begin{typethis}
from ColabTurtlePlus.Turtle import *
clearscreen()
initializeTurtle((480, 360))

bgcolor("black")
speed(13)
pensize(2)
colores = ["red", "orange", "gold", "lime",
           "skyblue", "violet"]

def petalo(medida):
    circle(medida, 60)
    left(120)
    circle(medida, 60)
    left(120)

petalos = int(input("¿Cuántos pétalos? (3 a 20) "))
for p in range(petalos):
    color(colores[p % 6])
    petalo(100)
    left(360 / petalos)

color("white")
for anillo in range(36):
    circle(40)
    left(10)

penup()
goto(0, 0)
dot(25, "gold")
hideturtle()
\end{typethis}

\needspace{14\baselineskip}
\step{Cómo funciona}
\begin{longtable}{>{\centering\arraybackslash}p{1.1cm} P{5.3cm} P{9.1cm}}
\toprule
\textbf{Línea} & \textbf{Bloque de Scratch} & \textbf{Qué significa} \\
\midrule
1--3 & \chip[events]{al presionar bandera verde} & Las tres líneas de inicio. \\
5--7 & \chip[pen]{fijar tamaño de lápiz a 2} & Escenario negro, la velocidad más rápida, un lápiz más grueso. \\
11--15 & \chip[myblocks]{definir pétalo (medida)} & Nuestro propio bloque. \py|circle(medida, 60)| dibuja un sexto de círculo (60 grados del círculo). Dos de esas curvas con giros forman un pétalo puntiagudo. Escribimos \py|petalo| sin acento: en los nombres de Python es más fácil sin acentos. \\
17 & \chip[sensing]{preguntar y esperar} & Pregunta un número. \py|int()| convierte el texto escrito en un número con el que podemos hacer cuentas. \\
18 & \chip[control]{repetir (respuesta)} & Una vuelta del bucle por cada pétalo. \py|p| cuenta 0, 1, 2, \ldots \\
19 & \chip[pen]{cambiar color de lápiz en 10} & \py|p % 6| es el residuo de dividir entre 6. Va 0, 1, 2, 3, 4, 5, 0, 1, \ldots así los colores dan vueltas y vueltas por la lista y nunca se acaban. \\
20 & \chip[myblocks]{pétalo 100} & Usa nuestro bloque. \\
21 & \chip[motion]{girar (360 / respuesta)} & Reparte una vuelta completa (360 grados) en partes iguales entre los pétalos. \\
23--26 & \chip[control]{repetir 36} & 36 círculos, girando 10 grados cada vez: $36 \times 10 = 360$, la vuelta completa. Eso forma el espirógrafo. \\
28--31 & \chip[pen]{subir lápiz} & Vuelve al centro, pon ahí un punto dorado y esconde la tortuga. \\
\bottomrule
\end{longtable}

\begin{challenge}[Sube de nivel]
\tick Prueba con 3, 7 y 20 pétalos. ¿Cuál se ve mejor?\par
\tick Cambia \py|petalo(100)| por \py|petalo(150)|. ¿Qué pasa? ¿Puedes hacer una flor pequeña dentro de la grande?\par
\tick Dale al espirógrafo sus propios colores: dentro del bucle de \py|anillo|, agrega \py|color(colores[anillo % 6])|.\par
\tick Haz un jardín entero: pon las líneas 18--21 dentro de una función \py|flor(x, y)| que primero vaya a \py|(x, y)|, y dibuja tres flores con pétalos más pequeños.
\end{challenge}

% =====================================================================
\mission{Concurso de tablas}{una función que devuelve una respuesta con \texttt{return}, llevar los puntos y un cronómetro}

Cinco preguntas de las tablas de multiplicar. La computadora lleva los puntos y te toma el tiempo. ¿Puedes sacar 5 de 5 en menos de 30 segundos?

\step{La versión de Scratch}
\begin{center}
\begin{scratch}
\blockinit{al presionar \greenflag}
\blockvariable{fijar \selectmenu{puntos} a \ovalnum{0}}
\blocksensing{reiniciar cronómetro}
\blockrepeat{repetir \ovalnum{5}}{
  \blockvariable{fijar \selectmenu{a} a \ovaloperator{elegir número al azar entre \ovalnum{2} y \ovalnum{12}}}
  \blockvariable{fijar \selectmenu{b} a \ovaloperator{elegir número al azar entre \ovalnum{2} y \ovalnum{12}}}
  \blocksensing{preguntar \ovaloperator{unir \ovalvariable{a} \ovaloperator{unir \ovalnum{ x } \ovalvariable{b}}} y esperar}
  \blockifelse{si \booloperator{\ovalsensing{respuesta} = \ovaloperator{\ovalvariable{a} * \ovalvariable{b}}} entonces}{
    \blockvariable{cambiar \selectmenu{puntos} en \ovalnum{1}}
    \blocklook{decir \ovalnum{¡Correcto!} durante \ovalnum{1} segundos}
  }{
    \blocklook{decir \ovalnum{¡Uy!} durante \ovalnum{1} segundos}
  }
}
\blocklook{decir \ovaloperator{unir \ovalvariable{puntos} \ovalsensing{cronómetro}}}
\end{scratch}
\end{center}

\needspace{44\baselineskip}
\begin{think}[Primero predice]
¿Qué crees que pasa si escribes \texttt{gato} en vez de un número? ¿El programa se va a romper? Lee la línea 6 del Python y la tabla para averiguarlo.
\end{think}

\step{La versión de Python}
\begin{typethis}
import random
import time

def pregunta(a, b):
    respuesta = input(str(a) + " x " + str(b) + " = ")
    if respuesta == str(a * b):
        print("¡Correcto!")
        return 1
    else:
        print("¡Uy! Era", a * b)
        return 0

print("Tablas de multiplicar: 5 preguntas. ¡Ya!")
puntos = 0
inicio = time.time()
for i in range(5):
    a = random.randint(2, 12)
    b = random.randint(2, 12)
    puntos = puntos + pregunta(a, b)

segundos = round(time.time() - inicio)
print("Sacaste", puntos, "de 5")
print("Tardaste", segundos, "segundos")
if puntos == 5 and segundos < 30:
    print("¡Eres un rayo con las tablas!")
elif puntos >= 3:
    print("¡Buen trabajo!")
else:
    print("Sigue practicando. ¡Lo vas a lograr!")
\end{typethis}

\needspace{14\baselineskip}
\step{Cómo funciona}
\begin{longtable}{>{\centering\arraybackslash}p{1.1cm} P{5.3cm} P{9.1cm}}
\toprule
\textbf{Línea} & \textbf{Bloque de Scratch} & \textbf{Qué significa} \\
\midrule
4 & \chip[myblocks]{definir pregunta (a) (b)} & Nuestro propio bloque, con dos entradas. \\
5 & \chip[sensing]{preguntar (unir a " x " b)}\newline\chip[sensing]{y esperar} & \py|str()| convierte un número en texto para poder unirlo con \py|+|. \\
6 & \chip[operators]{respuesta = a * b} & \py|input()| siempre da texto, así que convertimos también la respuesta correcta en texto y comparamos texto con texto. Ahora escribir \texttt{gato} no puede romper el concurso: ¡simplemente está mal! \\
8, 11 & (¡no hay bloque de Scratch!) & \py|return 1| le devuelve el número 1 a quien usó el bloque. Los bloques de Mis Bloques de Scratch no pueden hacer esto. Es un superpoder de Python. \\
14--15 & \chip[variables]{fijar puntos a 0}\newline\chip[sensing]{reiniciar cronómetro} & \py|time.time()| lee el reloj (en segundos). Lo guardamos en \py|inicio|, como cuando presionas un cronómetro. \\
19 & \chip[variables]{cambiar puntos en 1} & Haz una pregunta y suma lo que regresa: 1 si está bien, 0 si está mal. \\
21 & \chip[sensing]{cronómetro} & El reloj ahora menos el reloj al inicio = cuánto tardaste. \py|round()| lo vuelve un número entero. \\
24--29 & \chip[control]{si entonces si no} & Tres mensajes diferentes. \py|elif| quiere decir «si no, si». \\
\bottomrule
\end{longtable}

\needspace{9\baselineskip}
\begin{outputbox}[Podrías ver (tus preguntas serán diferentes)]
Tablas de multiplicar: 5 preguntas. ¡Ya!\\
12 x 4 = 48\\
¡Correcto!\\
9 x 11 = 98\\
¡Uy! Era 99\\
\ldots
\end{outputbox}

\begin{challenge}[Sube de nivel]
\tick Haz que sean 10 preguntas. ¿Qué dos líneas tienen que cambiar?\par
\tick Practica solo tu tabla más difícil: haz que \py|a = 7| siempre.\par
\tick Haz el cronómetro más estricto: rayo de las tablas solo si \py|segundos < 20|.\par
\tick Si hiciste \py|pitido()| en la Parte 7, ejecuta esa celda primero, y luego toca una nota alta y feliz en la parte del \py|if| de \py|pregunta| y una nota baja en la parte del \py|else|.
\end{challenge}

% =====================================================================
\mission{Adivina el número de lujo}{bucles \texttt{while} uno dentro de otro, \texttt{and}/\texttt{not}, niveles de dificultad y «otra vez»}

El juego de adivinar, con tres niveles, un número limitado de intentos y la pregunta «¿jugar otra vez?» al final.

\step{La versión de Scratch}
Esta es la parte importante: un «repetir hasta que» dentro de otro.

\begin{center}
\begin{scratch}
\blockinit{al presionar \greenflag}
\blockrepeat{repetir hasta que \booloperator{\ovalsensing{respuesta} = \ovalnum{n}}}{
  \blockvariable{fijar \selectmenu{secreto} a \ovaloperator{elegir número al azar entre \ovalnum{1} y \ovalvariable{maximo}}}
  \blockrepeat{repetir hasta que \booloperator{\booloperator{\ovalvariable{intentos} = \ovalnum{0}} o \booloperator{\ovalvariable{ganaste} = \ovalnum{sí}}}}{
    \blocksensing{preguntar \ovalnum{Tu número:} y esperar}
    \blockvariable{cambiar \selectmenu{intentos} en \ovalnum{-1}}
  }
  \blocksensing{preguntar \ovalnum{¿Jugar otra vez?} y esperar}
}
\end{scratch}
\end{center}

\begin{think}[Primero predice]
El Python tiene dos bucles \py|while|, uno dentro del otro, igual que el programa de Scratch. Encuéntralos. ¿Cuál es un juego completo y cuál es un solo intento? ¿Qué crees que hace \py|intentos = intentos - 1|?
\end{think}

\needspace{42\baselineskip}
\step{La versión de Python}
\begin{typethis}
import random

otra = "s"
while otra == "s":
    print("Niveles: 1 fácil, 2 medio, 3 difícil")
    nivel = input("Elige un nivel: ")
    if nivel == "1":
        maximo = 10
        intentos = 4
    elif nivel == "2":
        maximo = 50
        intentos = 6
    else:
        maximo = 100
        intentos = 7
    secreto = random.randint(1, maximo)
    print("Pensé un número del 1 al", maximo)

    ganaste = False
    while intentos > 0 and not ganaste:
        print("Intentos que quedan:", intentos)
        numero = int(input("Tu número: "))
        intentos = intentos - 1
        if numero == secreto:
            ganaste = True
        elif numero < secreto:
            print("¡Muy bajo!")
        else:
            print("¡Muy alto!")

    if ganaste:
        print("¡Lo adivinaste! Era", secreto)
    else:
        print("¡Se acabaron los intentos! Era", secreto)
    otra = input("¿Jugar otra vez? (s/n) ").lower()

print("¡Gracias por jugar!")
\end{typethis}

\needspace{14\baselineskip}
\step{Cómo funciona}
\begin{longtable}{>{\centering\arraybackslash}p{1.1cm} P{5.3cm} P{9.1cm}}
\toprule
\textbf{Línea} & \textbf{Bloque de Scratch} & \textbf{Qué significa} \\
\midrule
3--4 & \chip[control]{repetir hasta que}\newline\chip[operators]{respuesta = n} & El bucle grande: una vuelta = un juego completo. Sigue mientras \py|otra| sea \py|"s"| (de «sí»). \\
7--15 & \chip[control]{si entonces si no} & El nivel fija dos variables: el número más grande y cuántos intentos tienes. Cualquier cosa que no sea 1 o 2 da el nivel difícil. \\
19 & \chip[variables]{fijar ganaste a no} & \py|ganaste| es una variable de verdadero o falso. Empieza en \py|False| (falso). \\
20 & \chip[control]{repetir hasta que}\newline\chip[operators]{intentos = 0 o ganaste} & El bucle pequeño: una vuelta = un intento. Sigue mientras todavía tengas intentos \textbf{y} todavía \textbf{no} hayas ganado. \\
22--23 & \chip[sensing]{preguntar y esperar}\newline\chip[variables]{cambiar intentos en -1} & Pide un número, conviértelo en número y gasta un intento. \\
24--25 & \chip[variables]{fijar ganaste a sí} & ¡Le atinaste! \py|ganaste = True| (verdadero) hace que el bucle pequeño se detenga. \\
31--34 & \chip[control]{si entonces si no} & Después del bucle pequeño: ¿ganaste o se acabaron los intentos? \py|if ganaste:| es la forma corta de \py|if ganaste == True:|. \\
35 & \chip[sensing]{preguntar ¿Jugar otra vez?} & \py|.lower()| convierte \texttt{S} en \texttt{s}, así también funcionan las mayúsculas. Escribe solo \texttt{s}, sin acento. \\
37 & \chip[looks]{decir} & Esta línea está otra vez pegada al margen, así que se ejecuta una vez, cuando el bucle grande ya terminó. \\
\bottomrule
\end{longtable}

\begin{think}[¿Es justo?]
En difícil tienes 7 intentos para 100 números. ¿Alcanza? ¡Sí, \textbf{si} eres listo! Siempre di el número de en medio de los que quedan (50, luego 25 o 75, y así). Cada intento corta los números a la mitad, y 7 mitades te llevan de 100 hasta 1. Por eso fácil tiene 4 intentos y medio tiene 6. ¡Pruébalo!
\end{think}

\enlargethispage{\baselineskip}
\begin{challenge}[Sube de nivel]\small
\tick Agrega un nivel 4, «superdifícil»: números hasta 1000. ¿Cuántos intentos son justos? (Parte 1000 a la mitad hasta llegar a 1.)\par
\tick Cuenta los juegos que ganaste e imprímelo al final.\par
\tick Avisa cuando el número esté a 5 o menos del secreto: imprime \texttt{¡Caliente!}\par
\tick Ignora los números demasiado grandes: si \py|numero > maximo|, imprime un aviso y no gastes un intento.
\end{challenge}

% =====================================================================
\mission{Elige tu propia historia}{funciones como escenas, una función que revisa la respuesta y \texttt{return}}

Una aventura en la que \textbf{tú} decides qué pasa. Cada elección lleva a un lugar diferente. Hay cuatro finales. ¿Puedes encontrarlos todos?

\step{La versión de Scratch}
En Scratch, cada escena podría ser un \textbf{mensaje}. El programa principal \textbf{envía} el mensaje, y otro programa lo \textbf{recibe}:

\begin{center}
\begin{minipage}[t]{0.52\linewidth}
\begin{scratch}[scale=0.8]
\blockinit{al presionar \greenflag}
\blocksensing{preguntar \ovalnum{¿Nombre del héroe?} y esperar}
\blocklook{decir \ovalnum{... está en un bosque.} durante \ovalnum{1} segundos}
\blocksensing{preguntar \ovalnum{¿Izquierda o derecha?} y esperar}
\blockifelse{si \booloperator{\ovalsensing{respuesta} = \ovalnum{izquierda}} entonces}{
  \blockevent{enviar \selectmenu{cueva} y esperar}
}{
  \blockevent{enviar \selectmenu{río} y esperar}
}
\blocklook{decir \ovalnum{FIN}}
\end{scratch}
\end{minipage}\hfill
\begin{minipage}[t]{0.44\linewidth}
\begin{scratch}[scale=0.73]
\blockinit{al recibir \selectmenu{cueva}}
\blocklook{decir \ovalnum{La cueva está oscura...} durante \ovalnum{1} segundos}
\blocksensing{preguntar \ovalnum{¿Agarrar o escapar?} y esperar}
\blockif{si \booloperator{\ovalsensing{respuesta} = \ovalnum{agarrar}} entonces}{
  \blocklook{decir \ovalnum{¡Un dragón amigable!}}
}
\end{scratch}
\end{minipage}
\end{center}

En Python, cada escena es una \textbf{función}. Llamar a la función, como \py|cueva()|, es lo mismo que \chip[events]{enviar cueva y esperar}.

\begin{think}[Primero predice]
Cuenta las escenas. ¿Cuántos finales diferentes puede tener la historia? Dibuja un pequeño mapa: una caja para cada escena y una flecha para cada elección.
\end{think}

\needspace{50\baselineskip}
\step{La versión de Python}
\begin{typethis}
import time
def elegir(pregunta, a, b):
    mensaje = pregunta + " (" + a + " o " + b + ") "
    respuesta = ""
    while respuesta != a and respuesta != b:
        respuesta = input(mensaje).lower()
    return respuesta

def decir(palabras):
    print(palabras)
    time.sleep(1)

def cueva():
    decir("La cueva está oscura. Algo brilla...")
    decir("...¡y algo gruñe!")
    opcion = elegir("¿Qué haces?", "agarrar", "escapar")
    if opcion == "agarrar":
        decir("¡Un dragón amigable te comparte su oro!")
    else:
        decir("Sales de puntitas, sano y salvo.")

def rio():
    decir("Un río ancho. Hay un bote que se tambalea.")
    opcion = elegir("¿Qué haces?", "remar", "nadar")
    if opcion == "remar":
        decir("¡Llegas a una isla de loros que hablan!")
    else:
        decir("¡Brrr! Muy fría. ¡Hora del chocolate!")

heroe = input("¿Cómo se llama tu héroe? ")
decir(heroe + " está a la orilla de un bosque oscuro.")
lado = elegir("¿Hacia dónde?", "izquierda", "derecha")
if lado == "izquierda":
    cueva()
else:
    rio()
decir("FIN")
\end{typethis}

\needspace{14\baselineskip}
\step{Cómo funciona}
\begin{longtable}{>{\centering\arraybackslash}p{1.1cm} P{5.3cm} P{9.1cm}}
\toprule
\textbf{Línea} & \textbf{Bloque de Scratch} & \textbf{Qué significa} \\
\midrule
2--7 & \chip[myblocks]{definir elegir} & Un bloque ayudante. Hace la pregunta y \textbf{sigue preguntando} hasta que escribas una de las dos opciones. Luego \py|return| te devuelve tu elección. \\
5 & \chip[control]{repetir hasta que} & Sigue mientras la respuesta no sea \py|a| \textbf{y} no sea \py|b|. \py|!=| quiere decir «no es igual a». \\
9--11 & \chip[looks]{decir durante 1 segundos} & Nuestro propio bloque de «decir durante 1 segundo»: imprime y luego espera. \\
13--20 & \chip[events]{al recibir cueva} & La escena de la cueva. Definirla todavía no la ejecuta. Espera hasta que la llamen, igual que un programa «al recibir» espera su mensaje. \\
16 & \chip[sensing]{preguntar y esperar} & \py|elegir()| devuelve \py|"agarrar"| o \py|"escapar"|, y lo guardamos en \py|opcion|. Escribe las opciones tal cual, sin acentos. \\
22--28 & \chip[events]{al recibir río} & La escena del río. \\
30--37 & \chip[events]{al presionar bandera verde} & La historia empieza aquí, pegada al margen. \\
34, 36 & \chip[events]{enviar cueva y esperar} & Llamar a una escena. La historia salta dentro de la función, la ejecuta, luego regresa y sigue en la línea 37. \\
\bottomrule
\end{longtable}

\begin{challenge}[Sube de nivel]
\tick Cambia las palabras para que la historia trate de algo que te encante: el espacio, el futbol, los dinosaurios.\par
\tick Agrega una escena nueva \py|def puente():| y dale al río una tercera opción que lleve a ella. (Pista: puedes llamar a \py|elegir()| desde dentro de una escena, y puedes llamar a una escena desde dentro de otra.)\par
\tick Cuenta el tesoro: dale a cada final un número diferente de monedas de oro e imprímelo.\par
\tick Primero haz un mapa de la historia en papel: cajas para las escenas, flechas para las elecciones. Luego escribe una función para cada caja.
\end{challenge}

% =====================================================================
\mission{Mascota virtual}{varias variables que cambian en cada turno, un bucle por siempre con \texttt{break} y mantener números dentro de un rango}

Cuida a tu propia mascota. Dale de comer y juega con ella. Si tiene demasiada hambre o está demasiado triste, ¡se escapa!

\step{La versión de Scratch}
\begin{center}
\begin{scratch}[scale=0.8]
\blockinit{al presionar \greenflag}
\blockvariable{fijar \selectmenu{hambre} a \ovalnum{3}}
\blockvariable{fijar \selectmenu{alegría} a \ovalnum{7}}
\blockinfloop{por siempre}{
  \blocksensing{preguntar \ovalnum{¿comer, jugar o salir?} y esperar}
  \blockif{si \booloperator{\ovalsensing{respuesta} = \ovalnum{comer}} entonces}{
    \blockvariable{cambiar \selectmenu{hambre} en \ovalnum{-3}}
  }
  \blockif{si \booloperator{\ovalsensing{respuesta} = \ovalnum{jugar}} entonces}{
    \blockvariable{cambiar \selectmenu{alegría} en \ovalnum{2}}
  }
  \blockvariable{cambiar \selectmenu{hambre} en \ovalnum{1}}
  \blockif{si \booloperator{\ovalvariable{hambre} > \ovalnum{9}} entonces}{
    \blocklook{decir \ovalnum{¡Me escapé!}}
    \blockstop{detener \selectmenu{todos}}
  }
}
\end{scratch}
\end{center}

\begin{think}[Primero predice]
¿Qué crees que pasa si solo escribes \texttt{jugar} y nunca \texttt{comer}? ¿Cuántos días va a aguantar tu mascota? Ejecútalo y mira si tenías razón.
\end{think}

\needspace{42\baselineskip}
\step{La versión de Python}
\begin{typethis}
import random

nombre = input("Ponle nombre a tu mascota: ")
hambre = 3
alegria = 7
dia = 0

while True:
    dia = dia + 1
    print("Día", dia, "|", nombre, "| hambre", hambre,
          "| alegría", alegria)
    accion = input("¿comer, jugar o salir? ").lower()
    if accion == "comer":
        hambre = hambre - 3
        print(nombre, "come feliz. ¡Ñam!")
    elif accion == "jugar":
        alegria = alegria + 2
        print(nombre, "¡persigue la pelota!")
    elif accion == "salir":
        print("¡Adiós!", nombre, "te va a extrañar.")
        break
    else:
        print("Prueba comer, jugar o salir.")

    hambre = hambre + 1
    alegria = alegria - random.randint(0, 1)
    if hambre < 0:
        hambre = 0
    if alegria > 10:
        alegria = 10
    if hambre >= 10 or alegria <= 0:
        print("¡Oh, no!", nombre, "¡se escapó!")
        print("Cuida mejor a tu próxima mascota.")
        break
\end{typethis}

\needspace{14\baselineskip}
\step{Cómo funciona}
\begin{longtable}{>{\centering\arraybackslash}p{1.1cm} P{5.3cm} P{9.1cm}}
\toprule
\textbf{Línea} & \textbf{Bloque de Scratch} & \textbf{Qué significa} \\
\midrule
3--6 & \chip[variables]{fijar hambre a 3} & Las «estadísticas» de la mascota. El hambre va de 0 (llena) a 10 (muerta de hambre); la alegría va de 0 (triste) a 10 (muy feliz). En el código escribimos \py|alegria| sin acento. \\
8 & \chip[control]{por siempre} & \py|while True:| repite por siempre. Aquí no hay peligro, porque \py|input()| te espera en cada turno, y \py|break| nos saca. \\
10--11 & \chip[looks]{decir} & Muestra las estadísticas. Los signos \texttt{|} son solo una cerca entre los números. \\
13--23 & \chip[control]{si entonces si no} & Una respuesta, cuatro acciones posibles. \py|elif| revisa la siguiente opción solo si las anteriores no coincidieron. Escribe las palabras sin acentos. \\
21 & \chip[control]{detener este programa} & \py|break| salta fuera del bucle en ese mismo momento. \\
25--26 & \chip[variables]{cambiar hambre en 1} & Cada día, la mascota tiene más hambre y quizá se pone un poco más triste (pierde 0 o 1 punto de alegría al azar). \\
27--30 & \chip[control]{si entonces} & Mantén los números dentro del rango: el hambre no puede bajar de 0, y la alegría no puede pasar de 10. \\
31--34 & \chip[control]{detener todos} & \py|or| quiere decir que basta con una de las dos: demasiada hambre \textbf{o} demasiada tristeza, y la mascota se escapa. \\
\bottomrule
\end{longtable}

\needspace{8\baselineskip}
\begin{outputbox}[Podrías ver]
Ponle nombre a tu mascota: Galleta\\
Día 1 | Galleta | hambre 3 | alegría 7\\
¿comer, jugar o salir? comer\\
Galleta come feliz. ¡Ñam!\\
Día 2 | Galleta | hambre 1 | alegría 7
\end{outputbox}

\begin{challenge}[Sube de nivel]
\tick Agrega una tercera estadística, \py|energia|, y una opción nueva \texttt{dormir} que la llene. Jugar debería gastar algo de energía.\par
\tick Dale a la mascota un mensaje de humor: si \py|alegria| es 8 o más, imprime \texttt{(ronronea)}.\par
\tick Haz que la mascota gane: si la cuidas durante 15 días, imprime un diploma y usa \py|break|.\par
\tick ¡Dibuja a tu mascota! Después de las tres líneas de inicio, dibuja una cara cuyo color dependa de \py|alegria| (por ejemplo \py|"gold"| cuando está feliz, \py|"gray"| cuando está triste), y vuelve a dibujarla en cada turno.
\end{challenge}

% =====================================================================
\mission{Juego de persecución}{tres objetos, \texttt{towards()} para apuntar a un objeto y \texttt{distance()} para «tocar»}

Eres una tortuga. Junta las estrellas doradas, pero ¡cuidado!: un gato naranja te persigue. En cada turno escribes \texttt{w}, \texttt{a}, \texttt{s} o \texttt{d} y presionas Enter, y luego el gato da un paso hacia ti.

\step{La versión de Scratch}
En Scratch el jugador, el gato y la estrella tienen cada uno sus propios programas:

\begin{center}
\begin{minipage}[t]{0.32\linewidth}
\textbf{Jugador}\par
\begin{scratch}[scale=0.72]
\blockinit{al presionar tecla \selectmenu{w}}
\blockmove{apuntar en dirección \ovalnum{0}}
\blockmove{mover \ovalnum{40} pasos}
\end{scratch}
\end{minipage}\hfill
\begin{minipage}[t]{0.32\linewidth}
\textbf{Gato}\par
\begin{scratch}[scale=0.66]
\blockinit{al presionar \greenflag}
\blockinfloop{por siempre}{
  \blockmove{apuntar hacia \selectmenu{Jugador}}
  \blockmove{mover \ovalnum{2} pasos}
  \blockif{si \boolsensing{¿tocando \selectmenu{Jugador}?} entonces}{
    \blockstop{detener \selectmenu{todos}}
  }
}
\end{scratch}
\end{minipage}\hfill
\begin{minipage}[t]{0.32\linewidth}
\textbf{Estrella}\par
\begin{scratch}[scale=0.66]
\blockinit{al presionar \greenflag}
\blockinfloop{por siempre}{
  \blockif{si \boolsensing{¿tocando \selectmenu{Jugador}?} entonces}{
    \blockvariable{cambiar \selectmenu{puntos} en \ovalnum{1}}
    \blockmove{ir a \selectmenu{posición aleatoria}}
  }
}
\end{scratch}
\end{minipage}
\end{center}

En Python hay \textbf{un solo} bucle de juego. Cada vuelta del bucle es un turno: el jugador se mueve, revisamos la estrella, el gato se mueve, revisamos al gato.

\begin{think}[Primero predice]
Encuentra los tres objetos en el Python. Dos de ellos tienen nombre (\py|gato| y \py|estrella|). ¿Qué objeto se mueve con las órdenes sencillas como \py|forward(40)|, sin ningún nombre delante?
\end{think}

\needspace{30\baselineskip}
\step{La versión de Python}
La primera parte crea los tres objetos. Las órdenes sencillas son \textbf{tú}; \py|gato| y \py|estrella| son objetos nuevos.

\begin{typethis}[Escribe esto en una celda de código (parte 1)]
from ColabTurtlePlus.Turtle import *
clearscreen()
initializeTurtle((480, 360))
import random

speed(13)
shape("turtle")
color("green")
penup()

gato = Turtle()
gato.speed(13)
gato.shape("circle")
gato.color("orange")
gato.penup()
gato.goto(200, 140)

estrella = Turtle()
estrella.speed(13)
estrella.shape("triangle")
estrella.color("gold")
estrella.penup()
estrella.goto(120, -80)
\end{typethis}

Ejecútalo ya: deberías ver una tortuga verde en el centro, un círculo naranja (el gato) arriba a la derecha y un triángulo dorado (la estrella). Luego escribe la parte 2 \textbf{en la misma celda}, debajo, y ejecuta otra vez toda la celda.

\lstset{firstnumber=25}
\begin{typethis}[Escribe esto en una celda de código (parte 2)]
puntos = 0
for turno in range(20):
    jugada = input("Muévete con w, a, s o d: ").lower()
    if jugada == "w":
        setheading(90)
    elif jugada == "s":
        setheading(270)
    elif jugada == "a":
        setheading(180)
    elif jugada == "d":
        setheading(0)
    forward(40)
    if xcor() > 220 or xcor() < -220:
        back(40)
    if ycor() > 160 or ycor() < -160:
        back(40)

    if estrella.distance(pos()) < 20:
        puntos = puntos + 1
        print("¡Atrapaste la estrella! Puntos:", puntos)
        x = random.randint(-5, 5) * 40
        y = random.randint(-4, 4) * 40
        estrella.goto(x, y)

    gato.setheading(gato.towards(pos()))
    gato.forward(25)
    if gato.distance(pos()) < 30:
        print("¡El gato te atrapó!")
        break

print("¡Fin del juego! Estrellas:", puntos)
\end{typethis}
\lstset{firstnumber=1}

\needspace{14\baselineskip}
\step{Cómo funciona}
\begin{longtable}{>{\centering\arraybackslash}p{1.1cm} P{5.3cm} P{9.1cm}}
\toprule
\textbf{Línea} & \textbf{Bloque de Scratch} & \textbf{Qué significa} \\
\midrule
6--9 & \chip[looks]{cambiar disfraz a} & Eres una tortuga verde con el lápiz arriba. \\
11--16 & (un objeto nuevo) & \py|gato = Turtle()| crea un objeto nuevo llamado \py|gato|. Cada orden para él empieza con \py|gato.| \\
18--23 & (un objeto nuevo) & La estrella, de la misma manera. \\
26 & \chip[control]{por siempre} & 20 turnos. Cada vuelta del bucle es un turno. \\
27 & \chip[events]{al presionar tecla} & Pide una tecla. Colab no puede sentir cuando presionas teclas, así que escribes la tecla y presionas Enter. Las teclas \texttt{w}, \texttt{a}, \texttt{s}, \texttt{d} están juntas en el teclado, como flechas. \\
28--35 & \chip[motion]{apuntar en dirección} & Mira hacia donde elegiste. Direcciones de Python: 90 es arriba, 0 es derecha (recuerda: Python = 90 $-$ Scratch). Cualquier otra tecla: sigue en la misma dirección. \\
36--40 & \chip[motion]{mover 40 pasos}\newline\chip[motion]{si toca un borde, rebotar} & Muévete. Si eso te sacó del borde del escenario, da un paso atrás. \\
42 & \chip[sensing]{¿tocando Estrella?} & \py|pos()| da tu posición como una pareja \py|(x, y)|. \py|estrella.distance(pos())| es qué tan lejos está la estrella de ti. Menos de 20 cuenta como tocarla. \\
45--47 & \chip[motion]{ir a posición aleatoria} & Elige un lugar nuevo al azar. El \py|* 40| pone la estrella en lugares a los que puedes llegar, porque te mueves de 40 en 40 pasos. \\
49 & \chip[motion]{apuntar hacia Jugador} & \py|gato.towards(pos())| calcula qué dirección apunta del gato hacia ti, y el gato gira hacia allá. \\
50--53 & \chip[sensing]{¿tocando Jugador?}\newline\chip[control]{detener todos} & El gato da 25 pasos hacia ti. Si está a menos de 30, te atrapó y \py|break| termina el juego. \\
\bottomrule
\end{longtable}

\begin{think}[¿Por qué el gato es más lento que tú?]
Tú te mueves 40 pasos por turno y el gato 25, así que siempre puedes escapar, ¡pero solo hasta que llegues al borde! Las esquinas son peligrosas. Cambia el 25 del gato por 35 y mira cuánto más difícil se vuelve.
\end{think}

\enlargethispage{\baselineskip}
\begin{challenge}[Sube de nivel]\small
\tick Haz el juego más largo (más turnos) o más difícil (un gato más rápido).\par
\tick Acelera al gato: cada estrella que atrapes suma 5 a su paso (usa \py|paso_gato|).\par
\tick Agrega un segundo gato que empiece abajo a la izquierda.\par
\tick Deja un rastro: usa \py|pendown()| en vez de \py|penup()| y un \py|pencolor| bonito.
\end{challenge}

% =====================================================================
\mission{Juego de atrapar}{un objeto que cae, un bucle de juego dentro de otro, \texttt{jumpto()} y \texttt{distance()}}

Caen manzanas del cielo, y el viento las empuja de lado. Mira la primera mitad de la caída, luego calcula dónde va a caer la manzana y mueve tu canasta ahí. ¿Puedes atrapar las cinco?

\step{La versión de Scratch}
\begin{center}
\begin{minipage}[t]{0.6\linewidth}
\textbf{Manzana}\par
\begin{scratch}[scale=0.9]
\blockinit{al presionar \greenflag}
\blockrepeat{repetir \ovalnum{5}}{
  \blockmove{ir a x: \ovaloperator{elegir número al azar entre \ovalnum{-80} y \ovalnum{80}} y: \ovalnum{160}}
  \blockrepeat{repetir hasta que \booloperator{\ovalmove{posición en y} < \ovalnum{-150}}}{
    \blockmove{cambiar y en \ovalnum{-5}}
  }
  \blockif{si \boolsensing{¿tocando \selectmenu{Canasta}?} entonces}{
    \blockvariable{cambiar \selectmenu{puntos} en \ovalnum{1}}
  }
}
\end{scratch}
\end{minipage}\hfill
\begin{minipage}[t]{0.37\linewidth}
\textbf{Canasta}\par
\begin{scratch}[scale=0.82]
\blockinit{al presionar tecla \selectmenu{flecha izquierda}}
\blockmove{cambiar x en \ovalnum{-10}}
\end{scratch}
\end{minipage}
\end{center}

\begin{think}[Primero predice]
La canasta espera a que escribas un número. ¿Qué crees que pasa si escribes una palabra, como \texttt{izquierda}, en vez de un número? (¡Pruébalo! Luego lee la última línea del mensaje de error como un detective. Python siempre escribe sus errores en inglés; la palabra importante es el nombre del error, aquí \texttt{ValueError}: un valor que no sirve.)
\end{think}

\needspace{26\baselineskip}
\step{La versión de Python}
\begin{typethis}[Escribe esto en una celda de código (parte 1)]
from ColabTurtlePlus.Turtle import *
clearscreen()
initializeTurtle((480, 360))
import random

bgcolor("skyblue")
hideturtle()

canasta = Turtle()
canasta.shape("square")
canasta.color("brown")
canasta.penup()
canasta.goto(0, -150)

manzana = Turtle()
manzana.shape("circle")
manzana.color("red")
manzana.penup()
manzana.speed(13)
\end{typethis}

Ejecútalo: un cuadrado café (la canasta) abajo y un círculo rojo (la manzana). Ahora agrega la parte 2 debajo, en la misma celda.

\lstset{firstnumber=21}
\begin{typethis}[Escribe esto en una celda de código (parte 2)]
puntos = 0
for ronda in range(5):
    x = random.choice([-80, 0, 80])
    viento = random.choice([-20, 0, 20])
    manzana.jumpto(x, 160)
    for paso in range(4):
        manzana.goto(manzana.xcor() + viento,
                     manzana.ycor() - 40)

    lugar = int(input("Canasta x (-240 a 240): "))
    canasta.goto(lugar, -150)

    for paso in range(4):
        manzana.goto(manzana.xcor() + viento,
                     manzana.ycor() - 40)
    if canasta.distance(manzana.pos()) < 50:
        puntos = puntos + 1
        print("¡Atrapada!")
    else:
        print("¡Plaf! Cayó en x =",
              round(manzana.xcor()))

print("Atrapaste", puntos, "de 5 manzanas")
\end{typethis}
\lstset{firstnumber=1}

\needspace{14\baselineskip}
\step{Cómo funciona}
\begin{longtable}{>{\centering\arraybackslash}p{1.1cm} P{5.3cm} P{9.1cm}}
\toprule
\textbf{Línea} & \textbf{Bloque de Scratch} & \textbf{Qué significa} \\
\midrule
6--7 & \chip[looks]{cambiar fondo a} & Un escenario azul cielo (\py|"skyblue"|). Escondemos la tortuga sencilla porque este juego solo necesita la canasta y la manzana. Los colores van en inglés: \py|"brown"| es café y \py|"red"| es rojo. \\
9--19 & (dos objetos nuevos) & La canasta y la manzana, cada una con su propio nombre. \\
22 & \chip[control]{repetir 5} & El bucle grande: una vuelta = una manzana. \\
23--24 & \chip[operators]{elegir número al azar} & \py|random.choice()| elige un elemento de una lista. La manzana empieza en $-80$, 0 u 80, y el viento la empuja $-20$, 0 o 20 pasos de lado en cada paso hacia abajo. \\
25 & \chip[motion]{ir a x: y:} & \py|jumpto()| mueve al instante, sin deslizarse, así la manzana aparece de golpe arriba. \\
26--28 & \chip[control]{repetir hasta que}\newline\chip[motion]{cambiar y en -40} & El bucle pequeño: cuatro pasos hacia abajo (y de lado). Esta es la primera mitad de la caída. \\
30--31 & \chip[events]{al presionar tecla}\newline\chip[motion]{fijar x a} & Escribes una posición x y la canasta se desliza hasta ahí. \\
33--35 & \chip[motion]{cambiar y en -40} & Cuatro pasos más: la segunda mitad de la caída. \\
36 & \chip[sensing]{¿tocando Canasta?} & \py|manzana.pos()| es la posición de la manzana. Si la canasta está a menos de 50 pasos de ella, cuenta como atrapada. \\
\bottomrule
\end{longtable}

\begin{think}[El secreto para atrapar]
La manzana cae 8 pasos y se mueve de lado lo mismo en cada paso. Cuando se detiene a la mitad, ya dio 4 pasos. Así que en la segunda mitad se moverá de lado \textbf{exactamente lo mismo otra vez}. Si empezó en 0 y ahora está en 80, ¿dónde va a caer?
\end{think}

\begin{challenge}[Sube de nivel]
\tick Deja que la manzana empiece en cualquier lugar de $-160$ a 160 (pero cuidado: un viento fuerte podría sacarla del escenario).\par
\tick Hazlo más difícil: solo cuenta si la canasta está a menos de 30.\par
\tick ¡Haz un jugador computadora! Borra la línea del \py|input| y deja que la canasta calcule sola dónde va a caer la manzana.\par
\tick Deja una marca: después de que caiga cada manzana, agrega \py|manzana.stamp()| para que veas dónde cayeron las cinco manzanas. Si tienes \py|pitido()| de la Parte 7, ¡haz un bip cuando atrapes una!
\end{challenge}

% =====================================================================
\banner{pen}{TE TOCA}{Tu propio proyecto}

Ya construiste ocho grandes proyectos. Ahora viene el mejor: \textbf{el tuyo}. La mejor manera de empezar es \textbf{hacer un remix de uno de tus propios proyectos de Scratch}. Ábrelo en Scratch, mira sus programas y planea cómo convertirlo en Python.

\subsection*{¿Es bueno para Colab?}
Casi todos los proyectos de Scratch se pueden convertir en Python. Así se cambian los bloques difíciles:

\begin{center}
\begin{tabular}{P{5.6cm} P{10.6cm}}
\toprule
\textbf{En Scratch} & \textbf{En Python (Colab)} \\
\midrule
\chip[events]{al presionar tecla} & un turno: \py|jugada = input("¿w, a, s o d? ")| \\
\chip[sensing]{¿tocando Objeto2?} & \py|gato.distance(pos()) < 20| \\
\chip[motion]{apuntar hacia Objeto2} & \py|gato.setheading(gato.towards(pos()))| \\
\chip[control]{por siempre} & \py|for turno in range(20):| o \py|while True:| con \py|input()| y \py|break| \\
\chip[events]{enviar y esperar} & llamar a tu propia función: \py|cueva()| \\
\chip[looks]{decir ¡Hola!} & \py|print("¡Hola!")| o \py|write("¡Hola!")| \\
\chip[sound]{iniciar sonido} & \py|pitido(262)| (de la Parte 7) \\
un segundo objeto & \py|gato = Turtle()| \\
disfraces & formas (\py|"turtle"|, \py|"circle"|, \ldots) y colores \\
\bottomrule
\end{tabular}
\end{center}

\needspace{12\baselineskip}
\subsection*{Mi plan}
\begin{tcolorbox}[colback=white, colframe=codeframe, arc=5pt, title=\textbf{Mi proyecto}, coltitle=ink, colbacktitle=codebg]
\textbf{Nombre de mi proyecto:}\par\vspace{9pt}\rule{\linewidth}{0.4pt}\par\vspace{6pt}
\textbf{Es un remix de mi proyecto de Scratch:}\par\vspace{9pt}\rule{\linewidth}{0.4pt}\par\vspace{6pt}
\textbf{Qué pasa (en tres oraciones):}\par\vspace{4pt}
\foreach \n in {1,...,3}{\rule{\linewidth}{0.4pt}\par\vspace{9pt}}
\end{tcolorbox}

\needspace{14\baselineskip}
\begin{tcolorbox}[colback=white, colframe=codeframe, arc=5pt, title=\textbf{Mis objetos}, coltitle=ink, colbacktitle=codebg]
\renewcommand{\arraystretch}{2.1}
\begin{tabular}{@{}p{3.6cm}p{3.4cm}p{8.2cm}@{}}
\textbf{Objeto en Scratch} & \textbf{Nombre en Python} & \textbf{Qué hace} \\ \hline
& & \\ \hline
& & \\ \hline
& & \\ \hline
\end{tabular}
\end{tcolorbox}

\needspace{10\baselineskip}
\subsection*{Constrúyelo paso a paso}
\begin{enumerate}[label=\tick]
  \item Empieza un cuaderno nuevo de Colab. Pon el nombre de tu proyecto en una celda de texto arriba de todo.
  \item Escribe las tres líneas de inicio (si dibuja) y crea tus objetos. Ejecútalo.
  \item Agrega \textbf{una} cosa a la vez: el movimiento, luego los toques, luego los puntos. Ejecuta después de cada una.
  \item Cuando algo salga mal, lee el mensaje de error como un detective: primero la última línea, luego el número de línea.
  \item Muéstrale tu proyecto a alguien en casa. Pídele que lo juegue, y fíjate en lo que hace.
  \item Agrega una idea más que te haya dado. ¡Eso también es un remix!
\end{enumerate}


\needspace{20\baselineskip}
\begin{tcolorbox}[colback=white, colframe=codeframe, arc=5pt, title=\textbf{Bloques que voy a necesitar, y su Python}, coltitle=ink, colbacktitle=codebg]
\renewcommand{\arraystretch}{2.1}
\begin{tabular}{@{}p{7.4cm}p{7.9cm}@{}}
\textbf{Bloque de Scratch} & \textbf{Línea de Python} \\ \hline
& \\ \hline
& \\ \hline
& \\ \hline
& \\ \hline
& \\ \hline
& \\ \hline
\end{tabular}
\end{tcolorbox}

\begin{tcolorbox}[colback=white, colframe=codeframe, arc=5pt, title=\textbf{Cosas que arreglar o agregar la próxima vez}, coltitle=ink, colbacktitle=codebg]
\vspace{4pt}
\foreach \n in {1,...,4}{\rule{\linewidth}{0.4pt}\par\vspace{9pt}}
\end{tcolorbox}

\begin{center}
\begin{tcolorbox}[colback=ink, colframe=ink, arc=10pt, width=13cm, halign=center]
{\color{white}\Large\bfseries ¡Ya eres un verdadero programador de Python!}\par\vspace{4pt}
{\color{white}Puedes convertir tus propias ideas en programas, con bloques y con Python. ¡Sigue creando cosas!}
\end{tcolorbox}
\end{center}

% =====================================================================
% Capítulo 09: El Zoológico de Bichos (errores y depuración)
% Cada mensaje de error de abajo salió de ejecutar el fragmento con
% python3.13 -I (Python 3.13.5), la misma versión de Python que usa Colab.
% Los mensajes de error se dejan en inglés, tal como los escribe Python.
% =====================================================================

% ---------- Cajas que solo se usan en este capítulo ----------
\newtcblisting{bzerror}{listing only, enhanced, colback=errbg, colframe=errframe,
  boxrule=0.6pt, arc=4pt, left=6pt, right=4pt, top=0pt, bottom=0pt, before skip=2pt, after skip=3pt,
  listing options={basicstyle=\ttfamily\footnotesize, columns=fixed, keepspaces=true,
    upquote=true, showstringspaces=false, numbers=none, breaklines=false}}
\newtcblisting{bzerrorsmall}{listing only, enhanced, colback=errbg, colframe=errframe,
  boxrule=0.6pt, arc=4pt, left=6pt, right=4pt, top=0pt, bottom=0pt, before skip=2pt, after skip=3pt,
  listing options={basicstyle=\ttfamily\scriptsize, columns=fixed, keepspaces=true,
    upquote=true, showstringspaces=false, numbers=none, breaklines=false}}
\newtcblisting{bzcolab}{listing only, enhanced, colback=errbg, colframe=errframe,
  boxrule=0.8pt, arc=5pt, left=6pt, right=4pt, top=2pt, bottom=2pt,
  title={\textbf{Lo que Colab muestra debajo de la celda}}, coltitle=white, colbacktitle=errframe,
  fonttitle=\small,
  listing options={basicstyle=\ttfamily\scriptsize, columns=fixed, keepspaces=true,
    upquote=true, showstringspaces=false, numbers=none, breaklines=false}}
\newtcolorbox{bzcard}[3]{enhanced, colback=white, colframe=errframe, boxrule=1pt, arc=8pt,
  left=8pt, right=8pt, top=3pt, bottom=4pt, before skip=6pt, after skip=8pt,
  title={\textbf{Bicho #1: #2}\hfill{\itshape #3}}, coltitle=ink, colbacktitle=errhead,
  fonttitle=\normalsize, fontupper=\small, before upper={\setlength{\parskip}{3pt}}}
\newtcblisting{bzsnip}{listing only, enhanced, colback=codebg, colframe=codeframe,
  boxrule=0.6pt, arc=4pt, left=6pt, right=4pt, top=1pt, bottom=1pt, before skip=0pt, after skip=0pt,
  listing options={style=py, numbers=none, basicstyle=\ttfamily\footnotesize}}
\newcommand{\bzlabel}[1]{{\bfseries\color{errframe}#1}\ }
\newcommand{\bzruled}{\vspace{4pt}\rule{\linewidth}{0.4pt}\par\vspace{6pt}\rule{\linewidth}{0.4pt}\par}

\partopener{errhead}{CUANDO ALGO SALE MAL}{El Zoológico de Bichos}{Conoce los doce bichos que todo programador atrapa y aprende a domarlos como un detective.}

A un error en el código se le llama \textbf{bicho} (en inglés, \textit{bug}), y encontrarlo y arreglarlo se llama \textbf{depurar} (o cazar bichos). Todos los programadores del mundo hacen bichos todos los días, incluso las personas que inventaron Python. En este zoológico vas a conocer cada tipo de bicho y vas a aprender a atraparlo.

\subsection*{¿Por qué Python se queja si Scratch nunca lo hacía?}

Scratch fue \textbf{diseñado} para no tener mensajes de error. Sus bloques son como LEGO: solo encajan de maneras que tienen sentido, así que nunca puedes poner un bloque en el lugar equivocado. Y cuando pasa algo raro, Scratch hace en silencio algo que no hace daño y sigue adelante. Prueba esto en Scratch:

\begin{tcbraster}[raster columns=2, raster equal height=rows, raster column skip=12pt, raster valign=center]
\begin{tcolorbox}[colback=white, colframe=white, boxrule=0pt, valign=center]
\begin{scratch}
\blockinit{al presionar \greenflag}
\blocklook{decir \ovaloperator{\ovalnum{10} / \ovalnum{0}}}
\end{scratch}
\par\smallskip\small El gato solo dice \textbf{Infinity} y sigue como si nada.
\end{tcolorbox}
\begin{code}
print(10 / 0)
\end{code}
\end{tcbraster}

Python no esconde los problemas debajo de la alfombra. Cuando llega a algo que no puede hacer, se \textbf{detiene} y te deja una nota:

\begin{bzerror}
ZeroDivisionError: division by zero
\end{bzerror}

Piensa en Python como el ayudante de un detective. Cuando algo sale mal, congela la escena y te escribe una \textbf{nota con pistas}. La nota te dice \textbf{dónde} pasó (el número de línea, como un cuarto de una casa), \textbf{qué tipo} de problema fue (el nombre del error, como el nombre de un sospechoso) y te da una \textbf{pista}. Las letras rojas no quieren decir «fallaste». Quieren decir «¡llegó una pista!».

\begin{think}[Dos familias de bichos]
\textbf{Bichos de gramática} (\texttt{SyntaxError} e \texttt{IndentationError}): Python lee primero toda la celda, encuentra una línea que no entiende y no ejecuta \textbf{nada}, ni siquiera la línea 1.

\textbf{Bichos de ejecución} (todos los demás): Python ejecuta las líneas una por una y se detiene en la mala. Todo lo que se imprimió \textbf{antes} de esa línea sí aparece, y eso también es una pista: «llegó hasta aquí».
\end{think}

% ---------------------------------------------------------------------
\needspace{14\baselineskip}
\subsection*{Cómo leer un error en Colab}

Aquí hay una celda con un bicho. ¿Puedes encontrarlo antes de seguir leyendo?

\begin{code}
puntos = 10
print(puntso)
\end{code}

Cuando la ejecutas, Colab muestra algo así debajo de la celda (las primeras líneas pueden verse un poco distintas en tu pantalla, pero la última línea será la misma):

\begin{bzcolab}
---------------------------------------------------------------------
NameError                         Traceback (most recent call last)
/tmp/ipykernel_1604/434711600.py in <cell line: 0>()
      1 puntos = 10
----> 2 print(puntso)

NameError: name 'puntso' is not defined
\end{bzcolab}

\begin{watchout}[Python siempre escribe sus errores en inglés]
Python habla inglés: sus palabras (\texttt{for}, \texttt{print}, \texttt{forward}...) siempre van en inglés, y sus mensajes de error también. Aunque tu Colab esté en español, la nota roja saldrá en inglés. ¡No te preocupes! La palabra más importante es el \textbf{nombre del error} (\texttt{NameError}, \texttt{SyntaxError}...), y en este capítulo te explicamos en español qué quiere decir cada uno. Por ejemplo, \texttt{name 'puntso' is not defined} quiere decir «el nombre \texttt{puntso} no está definido»: Python no conoce esa palabra.
\end{watchout}

\needspace{10\baselineskip}
Los bichos de gramática se ven un poco distintos. Aquí Python muestra la línea y señala el lugar exacto con unas flechitas, \texttt{\textasciicircum}:

\begin{bzcolab}
  File "/tmp/ipython-input-2961531340.py", line 1
    if 5 > 3
            ^
SyntaxError: expected ':'
\end{bzcolab}

\needspace{12\baselineskip}
\begin{tcolorbox}[enhanced, colback=errbg, colframe=errframe, boxrule=0.8pt, arc=6pt,
  title={\textbf{La rutina del detective}}, coltitle=white, colbacktitle=errframe]
\begin{enumerate}[leftmargin=*]
  \item \textbf{Lee primero la ÚLTIMA línea.} La palabra antes de los dos puntos es el \textbf{tipo} de bicho (qué animal del zoológico es). Las palabras de después son la pista.
  \item \textbf{Busca el número de línea.} \texttt{line 2} (línea 2), o la flecha \texttt{---->} junto a una línea, te dice en qué cuarto pasó el crimen. Activa los números de línea en Colab: \textbf{Herramientas} (\textit{Tools}), luego \textbf{Configuración} (\textit{Settings}), luego \textbf{Editor}, luego \textbf{Mostrar números de línea} (\textit{Show line numbers}).
  \item \textbf{Mira las marcas.} En los bichos de gramática, Python pone marcas \texttt{\textasciicircum} debajo de la parte exacta de la línea que no le gustó. En los bichos de ejecución, Colab señala la línea con una flecha \texttt{---->}.
  \item \textbf{Lee la pista.} Algunas notas traen una pista, como \texttt{expected ':'} («esperaba \texttt{:}»). Si la nota dice que un nombre \textbf{is not defined} (no está definido), revisa con mucho cuidado cómo está escrito y las mayúsculas: Colab no va a adivinar la palabra correcta por ti.
\end{enumerate}
\end{tcolorbox}

\begin{watchout}[Una pista a la vez]
Python se detiene en el \textbf{primer} problema que encuentra. Si tu código tiene tres bichos, los vas a conocer uno por uno: arreglas uno, ejecutas otra vez, lees la siguiente pista. Eso es completamente normal.
\end{watchout}

Debajo del error, Colab también muestra un botón \textbf{Explicar error} (\textit{Explain error}). Primero sé tú el detective: lee la última línea, haz tu suposición y solo después presiónalo para comprobar.

\needspace{16\baselineskip}
\subsection*{El decodificador de pistas}

\small
\renewcommand{\arraystretch}{1.3}
\begin{longtable}{P{4.4cm} P{5cm} P{6.1cm}}
\toprule
\textbf{La última línea dice} & \textbf{Quiere decir} & \textbf{Prueba esto} \\
\midrule
\texttt{SyntaxError} & Python no puede leer la línea. & Busca si faltan dos puntos \texttt{:}, un paréntesis \texttt{)} o unas comillas \texttt{"}. \\
\texttt{IndentationError} & Una línea no está bien alineada. & Las líneas dentro de un bucle, un \texttt{if} o un \texttt{def} necesitan cuatro espacios delante. \\
\texttt{NameError} & Python no conoce esa palabra. & Revisa cómo está escrita y las mayúsculas. ¿Creaste primero la variable? ¿Ejecutaste las celdas de arriba? \\
\texttt{TypeError} & Revolviste texto y números. & Usa \py|int( )| para hacer un número, o \py|str( )| para hacer texto. \\
\texttt{ValueError} & Ese valor no tiene sentido aquí. & ¿Escribiste una palabra donde hacía falta un número? \\
\texttt{IndexError} & Ese número de elemento no existe. & Python cuenta desde 0. Compruébalo con \py|len( )|. \\
\texttt{ZeroDivisionError} & Dividiste entre cero. & Revisa que el número no sea 0 antes de dividir. \\
\texttt{AttributeError} & Ese objeto no tiene ese comando. & Revisa cómo está escrito lo que va después del punto, como en \py|gato.forward|. \\
\texttt{ModuleNotFoundError} & Las herramientas de la tortuga no están cargadas. & Vuelve a ejecutar la celda «En el principio». \\
\texttt{KeyboardInterrupt} & Presionaste detener. & ¡No hay nada roto! Python solo dice «me detuviste». \\
\bottomrule
\end{longtable}
\normalsize

% ---------------------------------------------------------------------
\needspace{20\baselineskip}
\subsection*{¡Bienvenido al Zoológico de Bichos!}

Aquí están los doce bichos que vas a encontrar más seguido. Cada tarjeta muestra la costumbre de Scratch que deja entrar al bicho, un programita roto, la pista \textbf{real} que da Python (copiada de Python 3.13, la versión que usa Colab) y cómo arreglarlo. ¡Intenta escribir cada programa roto en Colab para ver al bicho en libertad!

\begin{bzcard}{1}{Faltan los dos puntos}{la Cobra de los Dos Puntos}
\noindent\begin{minipage}[c]{0.43\linewidth}
\bzlabel{Costumbre de Scratch:} el bloque en forma de C \chip[control]{si} abre su propia boca. En Python, los dos puntos \texttt{:} son la boca que se abre, y tienes que escribirlos.
\end{minipage}\hfill
\begin{minipage}[c]{0.54\linewidth}
\begin{bzsnip}
if 5 > 3
    print("sí")
\end{bzsnip}
\end{minipage}\par
\begin{bzerror}
    if 5 > 3
            ^
SyntaxError: expected ':'
\end{bzerror}
\bzlabel{Arréglalo:} pon \texttt{:} al final de cada línea con \texttt{if}, \texttt{for}, \texttt{while} y \texttt{def}: \py|if 5 > 3:|. La pista \texttt{expected ':'} quiere decir «esperaba dos puntos».
\end{bzcard}

\begin{bzcard}{2}{Faltan espacios}{la Torre Tambaleante}
\noindent\begin{minipage}[c]{0.43\linewidth}
\bzlabel{Costumbre de Scratch:} los bloques se meten solitos \textbf{dentro} del bloque en C. En Python, cuatro espacios al principio de una línea la ponen adentro.
\end{minipage}\hfill
\begin{minipage}[c]{0.54\linewidth}
\begin{bzsnip}
for i in range(3):
print(i)
\end{bzsnip}
\end{minipage}\par
\begin{bzerrorsmall}
    print(i)
    ^^^^^
IndentationError: expected an indented block after 'for' statement on line 1
\end{bzerrorsmall}
\bzlabel{Arréglalo:} pon cuatro espacios (o presiona Tab) antes de \py|print(i)|. La pista dice «esperaba un bloque con sangría después del \texttt{for} de la línea 1»: Python hasta te dice qué línea abrió la boca.
\end{bzcard}

\begin{bzcard}{3}{Espacios que sobran}{el Vagabundo Despistado}
\noindent\begin{minipage}[c]{0.43\linewidth}
\bzlabel{Costumbre de Scratch:} puedes arrastrar un bloque a cualquier parte de la pantalla de Scratch. En Python, una línea con espacios delante tiene que estar dentro de un bloque en C.
\end{minipage}\hfill
\begin{minipage}[c]{0.54\linewidth}
\begin{bzsnip}
print("Hola")
    print("amigo")
\end{bzsnip}
\end{minipage}\par
\begin{bzerror}
    print("amigo")
IndentationError: unexpected indent
\end{bzerror}
\bzlabel{Arréglalo:} la línea 2 no está dentro de nada, así que bórrale los espacios. Esta vez no hay flechitas: la pista es la palabra \textbf{unexpected} («inesperada»): una sangría que nadie esperaba.
\end{bzcard}

\begin{bzcard}{4}{Letras cambiadas}{el Mono Mal Escrito}
\noindent\begin{minipage}[c]{0.43\linewidth}
\bzlabel{Costumbre de Scratch:} eliges las variables de la paleta, así que nunca puedes escribir mal una. En Python escribes cada nombre.
\end{minipage}\hfill
\begin{minipage}[c]{0.54\linewidth}
\begin{bzsnip}
puntos = 10
print(puntso)
\end{bzsnip}
\end{minipage}\par
\begin{bzerror}
    print(puntso)
          ^^^^^^
NameError: name 'puntso' is not defined
\end{bzerror}
\bzlabel{Arréglalo:} escríbelo exactamente igual. Los comandos de la tortuga también son nombres: \py|forwad(100)| da \texttt{NameError: name 'forwad' is not defined}
\end{bzcard}

\begin{bzcard}{5}{Mayúsculas}{el Copión Mayúsculo}
\noindent\begin{minipage}[c]{0.43\linewidth}
\bzlabel{Costumbre de Scratch:} las palabras de los bloques ya vienen escritas. Para Python, \texttt{Print} y \texttt{print} son tan distintos como \texttt{gato} y \texttt{perro}.
\end{minipage}\hfill
\begin{minipage}[c]{0.54\linewidth}
\begin{bzsnip}
Print("hola")
\end{bzsnip}
\end{minipage}\par
\begin{bzerror}
    Print("hola")
    ^^^^^
NameError: name 'Print' is not defined
\end{bzerror}
\bzlabel{Arréglalo:} las palabras de Python van en minúsculas: \py|print|, \py|if|, \py|for|. Un \texttt{If} con mayúscula da solo \texttt{SyntaxError: invalid syntax} («gramática no válida»). ¡Y \texttt{Nombre} y \texttt{nombre} son dos cajas distintas!
\end{bzcard}

\begin{bzcard}{6}{Faltan las comillas}{el Ladrón de Comillas}
\noindent\begin{minipage}[c]{0.43\linewidth}
\bzlabel{Costumbre de Scratch:} escribes el texto en una burbuja blanca, sin comillas. En Python, las comillas dicen «esto es solo texto».
\end{minipage}\hfill
\begin{minipage}[c]{0.54\linewidth}
\begin{bzsnip}
print(hola)
\end{bzsnip}
\end{minipage}\par
\begin{bzerror}
    print(hola)
          ^^^^
NameError: name 'hola' is not defined
\end{bzerror}
\bzlabel{Arréglalo:} \py|print("hola")|. Sin comillas, Python cree que \texttt{hola} es el nombre de una variable. Si olvidas solo las comillas \textbf{de cierre}, \py|print("hola)|, te sale \texttt{SyntaxError: unterminated string literal}: «empezaste un texto y nunca lo terminaste».
\end{bzcard}

\begin{bzcard}{7}{Falta un paréntesis}{el Murciélago sin Paréntesis}
\noindent\begin{minipage}[c]{0.43\linewidth}
\bzlabel{Costumbre de Scratch:} los huecos de un bloque ya vienen dibujados. En Python, los huecos son paréntesis \texttt{( )}, y cada \texttt{(} necesita su \texttt{)}.
\end{minipage}\hfill
\begin{minipage}[c]{0.54\linewidth}
\begin{bzsnip}
print("hola"
\end{bzsnip}
\end{minipage}\par
\begin{bzerror}
    print("hola"
         ^
SyntaxError: '(' was never closed
\end{bzerror}
\bzlabel{Arréglalo:} la pista dice «el \texttt{(} nunca se cerró». Pon el paréntesis de cierre: \py|print("hola")|. Su primo, \py|print "hola"|, da \texttt{SyntaxError: Missing parentheses in call to 'print'. Did you mean print(...)?} («faltan los paréntesis al llamar a \texttt{print}»).
\end{bzcard}

\begin{bzcard}{8}{Un igual en vez de dos}{la Anguila del Igual}
\noindent\begin{minipage}[c]{0.43\linewidth}
\bzlabel{Costumbre de Scratch:} el bloque verde \chip[operators]{( ) = ( )} hace una pregunta, y \chip[variables]{fijar} es otro bloque. Python usa \texttt{=} para fijar y \texttt{==} para preguntar.
\end{minipage}\hfill
\begin{minipage}[c]{0.54\linewidth}
\begin{bzsnip}
x = 5
if x = 5:
    print("cinco")
\end{bzsnip}
\end{minipage}\par
\begin{bzerror}
    if x = 5:
       ^^^^^
SyntaxError: invalid syntax. Maybe you meant '==' or ':=' instead of '='?
\end{bzerror}
\bzlabel{Arréglalo:} \py|if x == 5:|. La pista pregunta «¿quizás querías \texttt{==}?». No hagas caso de la parte de \texttt{:=}; es una herramienta para mucho más adelante.
\end{bzcard}

\begin{bzcard}{9}{Texto más un número}{el Duende del Pegamento}
\noindent\begin{minipage}[c]{0.43\linewidth}
\bzlabel{Costumbre de Scratch:} \chip[operators]{unir} pega cualquier cosa con cualquier cosa. Python no pega texto con un número usando \texttt{+}.
\end{minipage}\hfill
\begin{minipage}[c]{0.54\linewidth}
\begin{bzsnip}
puntos = 10
print("Puntos: " + puntos)
\end{bzsnip}
\end{minipage}\par
\begin{bzerror}
    print("Puntos: " + puntos)
          ~~~~~~~~~~~^~~~~~~~
TypeError: can only concatenate str (not "int") to str
\end{bzerror}
\bzlabel{Arréglalo:} «concatenate» quiere decir pegar, \texttt{str} quiere decir texto e \texttt{int} quiere decir número entero. Usa una coma, \py|print("Puntos:", puntos)|, o convierte el número en texto: \py|"Puntos: " + str(puntos)|. El mismo duende muerde en \py|input() + 1|, porque \py|input()| siempre da texto.
\end{bzcard}

\begin{bzcard}{10}{Una palabra que no es número}{el Número Falso}
\noindent\begin{minipage}[c]{0.43\linewidth}
\bzlabel{Costumbre de Scratch:} Scratch trata una palabra como 0 y sigue adelante. El \py|int()| de Python sabe leer \texttt{"5"} pero no \texttt{"cinco"}.
\end{minipage}\hfill
\begin{minipage}[c]{0.54\linewidth}
\begin{bzsnip}
edad = int("cinco")
\end{bzsnip}
\end{minipage}\par
\begin{bzerror}
    edad = int("cinco")
ValueError: invalid literal for int() with base 10: 'cinco'
\end{bzerror}
\bzlabel{Arréglalo:} no hagas caso de «literal» ni de «base 10»; la pista es \texttt{'cinco'}. Te sale la misma pista si alguien escribe \texttt{cinco} en \py|int(input("¿Cuántos años tienes? "))|. Escribe cifras: \texttt{5}.
\end{bzcard}

\begin{bzcard}{11}{Un elemento que no existe}{el Búho del Uno de Más}
\noindent\begin{minipage}[c]{0.43\linewidth}
\bzlabel{Costumbre de Scratch:} las listas empiezan en el elemento 1, y \chip[lists]{elemento 99 de lista} solo da un espacio vacío. Python cuenta desde \textbf{0}.
\end{minipage}\hfill
\begin{minipage}[c]{0.54\linewidth}
\begin{bzsnip}
frutas = ["manzana", "plátano", "cereza"]
print(frutas[3])
\end{bzsnip}
\end{minipage}\par
\begin{bzerror}
    print(frutas[3])
          ~~~~~~^^^
IndexError: list index out of range
\end{bzerror}
\bzlabel{Arréglalo:} «index» quiere decir número de elemento, y la pista dice «ese número se sale de la lista». Las tres frutas son \py|frutas[0]|, \py|frutas[1]| y \py|frutas[2]|. La última también es \py|frutas[-1]|. Comprueba cuántas hay con \py|len(frutas)|.
\end{bzcard}

\begin{bzcard}{12}{Dividir entre cero}{el Zombi del Cero}
\noindent\begin{minipage}[c]{0.43\linewidth}
\bzlabel{Costumbre de Scratch:} Scratch dice \textbf{Infinity} y sigue adelante. Repartir 12 dulces entre 0 amigos no tiene sentido, así que Python se detiene.
\end{minipage}\hfill
\begin{minipage}[c]{0.54\linewidth}
\begin{bzsnip}
dulces = 12
amigos = 0
print(dulces / amigos)
\end{bzsnip}
\end{minipage}\par
\begin{bzerror}
    print(dulces / amigos)
          ~~~~~~~^~~~~~~~
ZeroDivisionError: division by zero
\end{bzerror}
\bzlabel{Arréglalo:} comprueba primero con \py|if amigos != 0:| (\texttt{!=} quiere decir «no es igual a»), y divide solo dentro del \texttt{if}.
\end{bzcard}

% ---------------------------------------------------------------------
\clearpage
\subsection*{Bichos sigilosos: ¡sin letras rojas!}

Los animales más sigilosos del zoológico no hacen ningún ruido. El programa se ejecuta, no hay ningún error, pero hace \textbf{lo que no debe}. Son igualitos a los bichos de Scratch, así que ya sabes cómo cazarlos: ejecuta el programa y revisa lo que \textbf{de verdad} hizo.

\newcommand{\bzsneaky}[2]{\needspace{8\baselineskip}\par\medskip{\bfseries\color{errframe}#1}\quad{\itshape #2}\par}

\bzsneaky{El Hurra Perdido}{sangría en el nivel equivocado}
\noindent\begin{minipage}[c]{0.55\linewidth}
\begin{code}
for i in range(3):
    print("Hip, hip")
    print("¡Hurra!")
\end{code}
\end{minipage}\hfill
\begin{minipage}[c]{0.42\linewidth}\small\raggedright
Esto imprime «¡Hurra!» tres veces. Si solo querías uno al final, quítale los espacios a la línea 3. Los espacios deciden qué está dentro del bucle, y Python nunca te avisa si eliges el nivel equivocado.
\end{minipage}

\bzsneaky{El Loro Quisquilloso}{«Si» no es «si»}
\noindent\begin{minipage}[c]{0.55\linewidth}
\begin{code}
respuesta = input("¿Te gusta la pizza? ")
if respuesta == "si":
    print("¡A mí también!")
\end{code}
\end{minipage}\hfill
\begin{minipage}[c]{0.42\linewidth}\small\raggedright
Escribe \texttt{Si} con S mayúscula y no se imprime nada. Scratch piensa que \texttt{Si} y \texttt{si} son iguales; Python no. (Escribimos \texttt{si} sin acento para que sea fácil de teclear.) Arreglo: \py|if respuesta.lower() == "si":| (\py|.lower()| primero convierte la respuesta en minúsculas).
\end{minipage}

\bzsneaky{El Héroe Cero}{contar desde 0}
\noindent\begin{minipage}[c]{0.55\linewidth}
\begin{code}
frutas = ["manzana", "plátano"]
print(frutas[1])
\end{code}
\end{minipage}\hfill
\begin{minipage}[c]{0.42\linewidth}\small\raggedright
Querías la primera fruta, pero imprime \texttt{plátano}. El \chip[lists]{elemento 1} de Scratch es el \py|[0]| de Python.
\end{minipage}

\bzsneaky{El Fantasma Infinito}{un bucle que nunca termina}
\noindent\begin{minipage}[c]{0.55\linewidth}
\begin{code}
contador = 0
while contador < 3:
    print("hola")
\end{code}
\end{minipage}\hfill
\begin{minipage}[c]{0.42\linewidth}\small\raggedright
Imprime \texttt{hola} una y otra vez y nunca se detiene, porque \texttt{contador} se queda en 0. Presiona el botón de detener \stopbtn\ de la celda (aparece en lugar del botón de reproducir mientras la celda se ejecuta). Luego agrega \py|contador = contador + 1| dentro del bucle: el \chip[variables]{cambiar contador en 1} de Scratch.
\end{minipage}

\bzsneaky{La Receta Silenciosa}{olvidar los paréntesis}
\noindent\begin{minipage}[c]{0.55\linewidth}
\begin{code}
def festejar():
    print("¡Hurra!")

festejar
\end{code}
\end{minipage}\hfill
\begin{minipage}[c]{0.42\linewidth}\small\raggedright
No sale ningún «¡Hurra!». Sin \texttt{()}, solo \textbf{nombraste} la receta; no la cocinaste. Colab puede mostrar una línea rara que empieza con \texttt{<function}. Arreglo: \py|festejar()|.
\end{minipage}

\bzsneaky{El Acertijo del Redondeo}{round(2.5) da 2}
\noindent\begin{minipage}[c]{0.55\linewidth}
\begin{code}
print(round(2.5))
print(round(3.5))
\end{code}
\end{minipage}\hfill
\begin{minipage}[c]{0.42\linewidth}\small\raggedright
¡Esto imprime \texttt{2} y luego \texttt{4}! Scratch redondea 2.5 hacia arriba, a 3. Python redondea una mitad exacta al número \textbf{par} más cercano. En realidad no es un bicho de Python, pero puede causar un bicho en tu programa, así que ahora ya lo sabes.
\end{minipage}

% ---------------------------------------------------------------------
\needspace{18\baselineskip}
\subsection*{Trucos de detective}

Los programadores de Scratch que son buenos arreglando bichos usan los mismos trucos una y otra vez. Los investigadores Karen Brennan y Mitchel Resnick les preguntaron a jóvenes de Scratch cómo cazan bichos, y estos son los trucos que nombraron. Funcionan igual de bien en Python.

\needspace{10\baselineskip}
\textbf{1. Haz que diga lo que sabe.} En Scratch puedes agregar un bloque \chip[looks]{decir} para ver qué hay dentro de una variable. En Python, agrega un \py|print|:

\begin{tcbraster}[raster columns=2, raster equal height=rows, raster column skip=12pt, raster valign=center]
\begin{tcolorbox}[colback=white, colframe=white, boxrule=0pt, valign=center]
\begin{scratch}
\blocklook{decir \ovalvariable{puntos} durante \ovalnum{2} segundos}
\end{scratch}
\end{tcolorbox}
\begin{code}
print("puntos vale", puntos)
\end{code}
\end{tcbraster}

Pon uno antes y otro después de la línea sospechosa. Cuando el número no es el que esperabas, encontraste la escena del crimen. Borra los \py|print| de más cuando termines.

\textbf{2. Encuentra dónde está el problema.} Lee el número de línea, o ejecuta el programa en pedacitos. Divide una celda grande en celdas más pequeñas y ejecútalas una por una, como cuando haces clic en un solo bloque de Scratch para probarlo.

\textbf{3. Léelo en voz alta.} Explica tu código, línea por línea, a un adulto, a un amigo o a un patito de hule. ¡Hay programadores que de verdad tienen un patito de juguete en su escritorio para esto! Muchas veces tú mismo oyes el bicho a la mitad de la explicación.

\textbf{4. Cambia una sola cosa a la vez.} Haz un cambio, ejecuta, mira. Si cambias cinco cosas al mismo tiempo, no vas a saber cuál ayudó (ni cuál hizo un bicho nuevo).

\textbf{5. Compara con un ejemplo que funcione.} Busca en este libro un programa que funcione y que haga algo parecido. Ponlos lado a lado y encuentra la diferencia, como en un juego de encontrar las diferencias.

\textbf{6. Toma un descanso.} Toma agua, estírate, mira por la ventana. El cerebro sigue trabajando en los acertijos sin que te des cuenta, y muchas veces el bicho es obvio cuando regresas.

\textbf{7. Cuéntaselo a alguien.} Pedir ayuda es una habilidad de programador de verdad. Di qué querías, qué pasó en su lugar y qué dice la última línea del error.


% ---------------------------------------------------------------------
\clearpage
\subsection*{¡Caza de bichos! Atrápalos a todos}

Cada programa de abajo tiene un bicho: algunos son ruidosos (letras rojas), otros son sigilosos (sin error, pero con un resultado equivocado). Escribe cada uno en Colab, ejecútalo y lee la pista. Luego arréglalo y escribe en las líneas qué estaba mal. Antes de ejecutarlo, ¿puedes \textbf{adivinar} qué animal del zoológico va a aparecer? (Las respuestas están al final del libro, ¡pero inténtalo primero!)

\newcommand{\bzhunt}[2]{\par\needspace{7\baselineskip}\smallskip\textbf{Caza de bichos #1}\quad #2\par\nopagebreak}

\bzhunt{1}{Debería decir Miau 5 veces.}
\begin{code}
contador = 1
while contador <= 5:
print("Miau número", contador)
    contador = contador + 1
\end{code}
\bzruled

\bzhunt{2}{Debería mostrar tus puntos.}
\begin{code}
puntos = 0
puntos = puntos + 10
Print("Tus puntos son", puntos)
\end{code}
\bzruled

\bzhunt{3}{Escribe \texttt{9} cuando te pregunte.}
\begin{code}
edad = input("¿Cuántos años tienes? ")
print("El año que viene tendrás", edad + 1)
\end{code}
\bzruled

\bzhunt{4}{Escribe \texttt{gato123} cuando te pregunte.}
\begin{code}
clave = input("¿Contraseña? ")
if clave = "gato123":
    print("¡Adelante!")
\end{code}
\bzruled

\bzhunt{5}{Debería imprimir la primera y la última fruta. (¡Dos bichos!)}
\begin{code}
frutas = ["manzana", "plátano", "cereza"]
print("Primera:", frutas[1])
print("Última:", frutas[3])
\end{code}
\bzruled

\bzhunt{6}{Debería repartir los dulces en partes iguales, o decir que no hay con quién repartirlos.}
\begin{code}
dulces = 20
amigos = 0
print("Cada amigo recibe", dulces / amigos)
\end{code}
\bzruled

\bzhunt{7}{Debería festejar.}
\begin{code}
def festejar():
    print("¡Hip, hip, hurra!")

festejar
\end{code}
\bzruled

\bzhunt{8}{Debería contar 5, 4, 3, 2, 1 y despegar. (¡Ten listo el botón de detener!)}
\begin{code}
n = 5
while n > 0:
    print(n)
print("¡Despegue!")
\end{code}
\bzruled

\bzhunt{9}{Debería dibujar un cuadrado. (Pon las tres líneas de inicio arriba.)}
\begin{code}
for i in range(4):
    forward(50)
right(90)
\end{code}
\bzruled

\bzhunt{10}{Escribe tu nombre cuando te pregunte.}
\begin{code}
nombre = input("¿Cómo te llamas? ")
print("Hola, " + nombre)
print("¡Mucho gusto!)
\end{code}
\bzruled

% ---------------------------------------------------------------------
\needspace{10\baselineskip}
\subsection*{Lista del Zoológico de Bichos}

Marca cada animal cuando lo hayas atrapado en tu propio código. ¡Cada bicho que atrapas te hace mejor detective!

\begin{tcbraster}[raster columns=3, raster column skip=8pt, raster equal height]
\begin{tcolorbox}[colback=errbg, colframe=errframe, boxrule=0.5pt, arc=4pt, fontupper=\small]
\tick Cobra de los Dos Puntos\par \tick Torre Tambaleante\par \tick Vagabundo Despistado\par \tick Mono Mal Escrito\par \tick Copión Mayúsculo\par \tick Ladrón de Comillas
\end{tcolorbox}
\begin{tcolorbox}[colback=errbg, colframe=errframe, boxrule=0.5pt, arc=4pt, fontupper=\small]
\tick Murciélago sin Paréntesis\par \tick Anguila del Igual\par \tick Duende del Pegamento\par \tick Número Falso\par \tick Búho del Uno de Más\par \tick Zombi del Cero
\end{tcolorbox}
\begin{tcolorbox}[colback=errbg, colframe=errframe, boxrule=0.5pt, arc=4pt, fontupper=\small]
\tick Hurra Perdido\par \tick Loro Quisquilloso\par \tick Héroe Cero\par \tick Fantasma Infinito\par \tick Receta Silenciosa\par \tick Acertijo del Redondeo
\end{tcolorbox}
\end{tcbraster}

\begin{itemize}[label={}, leftmargin=0pt]
  \item \tick Ya sé leer primero la \textbf{última línea} de un error.
  \item \tick Puedo encontrar el número de línea y las marcas \texttt{\textasciicircum\textasciicircum\textasciicircum}.
  \item \tick Puedo detener un bucle infinito con el botón de detener.
  \item \tick Ya sé usar \py|print| para que mi programa diga lo que sabe.
  \item \tick Cambio una sola cosa a la vez, y tomo un descanso cuando me atoro.
\end{itemize}

% =====================================================================
% 10 EL DICCIONARIO DE BLOQUES y MI HOJA DE AYUDA DE PYTHON
% =====================================================================
% Un bloque de la extensión Música (del mismo color que Lápiz, pero sin el ícono del lápiz)
{\catcode`\_=11 \gdef\blockmusic{\_penfalse\scr_normalblock{scrpen}}}

\partopener{looks}{BÚSCALO AQUÍ}{El diccionario de bloques}{¿Conoces el bloque de Scratch pero no el Python? Encuéntralo aquí.}

Todos los bloques de Scratch 3.0 están en este diccionario, en el mismo orden que en la paleta de Scratch: \chip[motion]{Movimiento}, \chip[looks]{Apariencia}, \chip[sound]{Sonido}, \chip[events]{Eventos}, \chip[control]{Control}, \chip[sensing]{Sensores}, \chip[operators]{Operadores}, \chip[variables]{Variables}, \chip[lists]{Listas}, \chip[myblocks]{Mis Bloques}, \chip[pen]{Lápiz} y un poquito de \chip[pen]{Música}.

\begin{itemize}
  \item La columna de Python usa comandos sencillos, como \py|forward(10)|. Mueven a \textbf{la tortuga}, que hace el papel de tu objeto. Primero pon las tres líneas de inicio arriba de tu celda.
  \item \texttt{gato} es un segundo objeto que hiciste con \py|gato = Turtle()|. Entonces \py|gato.forward(10)| mueve al gato.
  \item Algunos bloques necesitan una línea hasta arriba de tu celda: \py|import random|, \py|import time| o \py|import math|.
  \item Algunos bloques no pueden funcionar en Colab. La columna Consejo siempre te dice qué hacer en su lugar.
  \item Python habla inglés: sus palabras (\py|for|, \py|print|, \py|forward|...) siempre van en inglés. Los nombres que tú inventas (\texttt{gato}, \texttt{puntos}, \texttt{comida}) pueden ir en español, sin acentos ni ñ.
  \item Los colores se escriben en inglés: \texttt{"red"} es rojo, \texttt{"blue"} azul, \texttt{"gold"} dorado, \texttt{"skyblue"} celeste.
\end{itemize}

\setscratch{scale=0.8}
\begingroup
\small
\renewcommand{\arraystretch}{1.0}
\begin{longtable}{P{5.8cm} P{4.9cm} P{4.9cm}}
\toprule
\textbf{Scratch} & \textbf{Python} & \textbf{Consejo} \\
\midrule
\endhead
% ---------------------------------------------------------------- MOVIMIENTO
\cat{motion}{Movimiento}
\begin{scratch}\blockmove{mover \ovalnum{10} pasos}\end{scratch} & \py|forward(10)| & Hacia atrás: \py|back(10)|. La tortuga dibuja mientras se mueve, a menos que uses \py|penup()|. \dictrow
\begin{scratch}\blockmove{girar \turnright{} \ovalnum{15} grados}\end{scratch} & \py|right(15)| & Girar funciona igualito que en Scratch. \dictrow
\begin{scratch}\blockmove{girar \turnleft{} \ovalnum{15} grados}\end{scratch} & \py|left(15)| & \dictrow
\begin{scratch}\blockmove{ir a \selectmenu{posición aleatoria}}\end{scratch} & \py|x = random.randint(|\newline\py|    -240, 240)|\newline\py|y = random.randint(|\newline\py|    -180, 180)|\newline\py|goto(x, y)| & Pon \py|import random| hasta arriba. Una línea puede seguir dentro de un paréntesis. \dictrow
\begin{scratch}\blockmove{ir a \selectmenu{puntero del ratón}}\end{scratch} & \py|x = int(input("¿x? "))|\newline\py|y = int(input("¿y? "))|\newline\py|goto(x, y)| & No se puede en Colab: no tiene ratón. Mejor escribe el lugar. \dictrow
\begin{scratch}\blockmove{ir a \selectmenu{Objeto2}}\end{scratch} & \py|goto(gato.pos())| & \py|pos()| te da la x y la y de un objeto, juntas. \dictrow
\begin{scratch}\blockmove{ir a x: \ovalnum{0} y: \ovalnum{0}}\end{scratch} & \py|goto(0, 0)| & Los mismos números que en Scratch: el escenario es del mismo tamaño. \py|home()| va a (0, 0) y apunta a la derecha. \dictrow
\begin{scratch}[scale=0.68]\blockmove{desplazar en \ovalnum{1} segs a \selectmenu{posición aleatoria}}\end{scratch} & \py|speed(2)|\newline y luego las líneas de \textbf{ir a posición aleatoria} & Cualquier \textbf{ir a} se vuelve un desplazamiento cuando la velocidad es lenta. \dictrow
\begin{scratch}\blockmove{desplazar en \ovalnum{1} segs a x: \ovalnum{0} y: \ovalnum{0}}\end{scratch} & \py|speed(1)|\newline\py|goto(0, 0)| & La velocidad va de 1 (lenta) a 13 (la más rápida). No está en segundos. \dictrow
\begin{scratch}\blockmove{apuntar en dirección \ovalnum{90}}\end{scratch} & \py|setheading(0)| & \textbf{¡La regla del 90 menos!} Usa \py|setheading(90 - d)|: arriba 0 es 90, derecha 90 es 0, abajo 180 es $-90$, izquierda $-90$ es 180. \dictrow
\begin{scratch}\blockmove{apuntar hacia \selectmenu{Objeto2}}\end{scratch} & \py|p = gato.pos()|\newline\py|d = towards(p)|\newline\py|setheading(d)| & Aquí no hace falta la regla: \py|towards| ya usa la brújula de Python. Puntero del ratón: no en Colab. \dictrow
\begin{scratch}\blockmove{cambiar x en \ovalnum{10}}\end{scratch} & \py|setx(xcor() + 10)| & Toma la x, súmale 10 y luego fíjala. \dictrow
\begin{scratch}\blockmove{fijar x a \ovalnum{0}}\end{scratch} & \py|setx(0)| & \dictrow
\begin{scratch}\blockmove{cambiar y en \ovalnum{10}}\end{scratch} & \py|sety(ycor() + 10)| & \dictrow
\begin{scratch}\blockmove{fijar y a \ovalnum{0}}\end{scratch} & \py|sety(0)| & \dictrow
\begin{scratch}\blockmove{si toca un borde, rebotar}\end{scratch} & \py|h = heading()|\newline\py|if abs(xcor()) > 240:|\newline\py|    setheading(180-h)|\newline\py|if abs(ycor()) > 180:|\newline\py|    setheading(-h)| & Python no tiene paredes, así que tú revisas los bordes. Pon estas líneas dentro de tu bucle, después de \py|forward|. \dictrow
\begin{scratch}\blockmove{fijar estilo de rotación a \selectmenu{no rotar}}\end{scratch} & \py|shape("circle")| & Una tortuga siempre apunta hacia donde va. Una forma redonda se ve igual hacia todos lados, así que parece que no gira. \dictrow
\ovalmove{posición en x}\ \ovalmove{posición en y} & \py|xcor()|\newline\py|ycor()| & Para imprimir un número redondito: \py|round(xcor())| \dictrow
\ovalmove{dirección} & \py|90 - heading()| & La regla del 90 menos también funciona al revés. ¿Te da menos de $-180$? Súmale 360. \dictrow
% ---------------------------------------------------------------- APARIENCIA
\cat{looks}{Apariencia}
\begin{scratch}\blocklook{decir \ovalnum{¡Hola!} durante \ovalnum{2} segundos}\end{scratch} & \py|print("¡Hola!")|\newline\py|time.sleep(2)| & Pon \py|import time| hasta arriba. \dictrow
\begin{scratch}\blocklook{decir \ovalnum{¡Hola!}}\end{scratch} & \py|print("¡Hola!")| & Las palabras aparecen debajo de la celda. Para escribir en el escenario: \py|write("¡Hola!")| \dictrow
\begin{scratch}\blocklook{pensar \ovalnum{Umm...} durante \ovalnum{2} segundos}\end{scratch} & \py|print("(Umm...)")|\newline\py|time.sleep(2)| & En Python no hay globitos. Los paréntesis pueden mostrar que es un pensamiento. \dictrow
\begin{scratch}\blocklook{pensar \ovalnum{Umm...}}\end{scratch} & \py|print("(Umm...)")| & \dictrow
\begin{scratch}\blocklook{cambiar disfraz a \selectmenu{tortuga}}\end{scratch} & \py|shape("turtle")| & Los disfraces son formas, con nombre en inglés: \texttt{"turtle"} (tortuga), \texttt{"arrow"} (flecha), \texttt{"circle"} (círculo), \texttt{"square"} (cuadrado), \texttt{"triangle"} (triángulo), \texttt{"classic"} (clásica). En Colab no hay dibujos. \dictrow
\begin{scratch}\blocklook{siguiente disfraz}\end{scratch} & \py|n = (n + 1) % 4|\newline\py|shape(disfraces[n])| & Primero haz una lista de 4 formas llamada \texttt{disfraces}, y \py|n = 0|. El \py|% 4| regresa al principio después de la última. \dictrow
\begin{scratch}\blocklook{cambiar fondo a \selectmenu{cielo azul}}\end{scratch} & \py|bgcolor("skyblue")| & Un color liso, no un dibujo. Los nombres de los colores van en minúsculas. \dictrow
\begin{scratch}\blocklook{siguiente fondo}\end{scratch} & \py|b = (b + 1) % 3|\newline\py|bgcolor(cielos[b])| & Igual que siguiente disfraz, con una lista de 3 colores llamada \texttt{cielos}. \dictrow
\begin{scratch}\blocklook{cambiar tamaño en \ovalnum{10}}\end{scratch} & \py|medida = medida + 10|\newline\py|shapesize(medida / 100)| & Guarda el tamaño en tu propia caja, empezando con \py|medida = 100|. \dictrow
\begin{scratch}\blocklook{fijar tamaño a \ovalnum{200} \%}\end{scratch} & \py|shapesize(2)| & 1 quiere decir 100\%, así que 2 quiere decir 200\%. Nunca uses 0: eso es un error. \dictrow
\begin{scratch}\blocklook{cambiar el efecto \selectmenu{color} en \ovalnum{25}}\end{scratch} & \py|fillcolor("gold")| & En Python no hay efectos, así que escoge un color nuevo. \py|color("gold")| pinta la tortuga y el lápiz. \dictrow
\begin{scratch}\blocklook{fijar efecto \selectmenu{desvanecer} a \ovalnum{50}}\end{scratch} & \py|fillcolor("pink")| & No se puede en Colab: no hay tortugas transparentes. Prueba un color más pálido, o \py|hideturtle()| para 100. Ojo de pez, remolino, pixelar y mosaico tampoco funcionan. \dictrow
\begin{scratch}\blocklook{quitar efectos gráficos}\end{scratch} & \py|color("black")| & Una tortuga nueva es negra. \dictrow
\begin{scratch}\blocklook{mostrar}\end{scratch} & \py|showturtle()| & \dictrow
\begin{scratch}\blocklook{esconder}\end{scratch} & \py|hideturtle()| & Una tortuga escondida puede seguir dibujando. \dictrow
\begin{scratch}\blocklook{ir a la capa \selectmenu{de adelante}}\end{scratch} & Haz este objeto \textbf{al final} & No hay comando de capas. El objeto que haces al final se dibuja encima. \dictrow
\begin{scratch}\blocklook{ir \ovalnum{1} capas hacia \selectmenu{adelante}}\end{scratch} & Cambia el orden en que haces los objetos & Las líneas y lo que escribes siempre quedan debajo de las tortugas. \dictrow
\ovallook{disfraz \selectmenu{número}} & \py|n + 1|\newline\py|shape()| & Python cuenta desde 0, así que suma 1. \py|shape()| con el paréntesis vacío te da el nombre del disfraz. \dictrow
\ovallook{fondo \selectmenu{número}} & \py|b + 1|\newline\py|bgcolor()| & \py|bgcolor()| con el paréntesis vacío te da el color. \dictrow
\ovallook{tamaño} & \py|medida| & Tu propia caja de \textbf{cambiar tamaño en}. \dictrow
% ---------------------------------------------------------------- SONIDO
\cat{sound}{Sonido}
\begin{scratch}[scale=0.76]\blocksound{tocar sonido \selectmenu{Miau} hasta que termine}\end{scratch} & \py|pitido(262)|\newline\py|time.sleep(0.5)| & \py|pitido| es la ayuda de la Parte 7. Colab no espera a que termine un sonido, así que esperas tú. No hay Miau: toca notas. \dictrow
\begin{scratch}\blocksound{iniciar sonido \selectmenu{Miau}}\end{scratch} & \py|pitido(262)| & Python sigue de inmediato. \dictrow
\begin{scratch}\blocksound{detener todos los sonidos}\end{scratch} & El botón de detener \stopbtn & Los pitidos son cortos, así que se detienen solos. \dictrow
\begin{scratch}\blocksound{cambiar efecto \selectmenu{tono} en \ovalnum{10}}\end{scratch} & \py|frec = frec * 1.06|\newline\py|pitido(frec)| & Guarda la nota en una caja llamada \texttt{frec}. Por 1.06 es un «10» de Scratch. \dictrow
\begin{scratch}\blocksound{fijar efecto \selectmenu{tono} a \ovalnum{120}}\end{scratch} & \py|frec = 523| & 262 es el do central; 523 es el do de arriba (tono 120 en Scratch). \dictrow
\begin{scratch}\blocksound{quitar efectos de sonido}\end{scratch} & \py|frec = 262| & De regreso al do central. \dictrow
\begin{scratch}\blocksound{cambiar volumen en \ovalnum{-10}}\blocksound{fijar volumen a \ovalnum{100} \%}\end{scratch}\ovalsound{volumen} & Usa el volumen de tu computadora & \py|pitido| no tiene perilla de volumen, así que sube o baja tus bocinas. \dictrow
% ---------------------------------------------------------------- EVENTOS
\cat{events}{Eventos}
\begin{scratch}\blockinit{al presionar \greenflag}\end{scratch} & El botón de reproducir \playbtn & O presiona Mayús y Enter. Python empieza en la línea 1 de la celda. \dictrow
\begin{scratch}\blockinit{al presionar tecla \selectmenu{espacio}}\end{scratch} & \py|k = input("¿Tecla? ")|\newline\py|if k == "s":|\newline\py|    saltar()| & No se puede en Colab: prueba \py|input()| y luego presiona Enter. Las teclas de verdad funcionan en Thonny con \py|onkeypress| (no \py|onkey|). \dictrow
\begin{scratch}\blockinit{al hacer clic en este objeto}\end{scratch} & \py|c = input("¿Clic? ")| & En Colab no hay ratón, así que pregunta con \py|input()|. Thonny puede usar clics con \py|onclick|. \dictrow
\begin{scratch}\blockinit{cuando el fondo cambie a \selectmenu{noche}}\end{scratch} & \py|def noche():| & Haz tu propio bloque y llama a \py|noche()| justo después de cambiar el \py|bgcolor|. \dictrow
\begin{scratch}\blockinit{cuando \selectmenu{volumen del sonido} > \ovalnum{10}}\end{scratch} & \py|input("¡Aplaude! ")| & No se puede: Colab no te oye. Presionar Enter puede ser tu aplauso de mentiritas. Para \textbf{cronómetro}, mira la fila del cronómetro en Sensores. \dictrow
\begin{scratch}\blockinit{al recibir \selectmenu{empezar}}\end{scratch} & \py|def empezar():| & Un mensaje es como tu propio bloque. Las líneas de adentro llevan cuatro espacios. \dictrow
\begin{scratch}\blockevent{enviar \selectmenu{empezar}}\end{scratch} & \py|empezar()| & Python ejecuta todo el bloque y luego sigue. \dictrow
\begin{scratch}\blockevent{enviar \selectmenu{empezar} y esperar}\end{scratch} & \py|empezar()| & ¡Es lo mismo! Python siempre espera a que un bloque termine. \dictrow
% ---------------------------------------------------------------- CONTROL
\cat{control}{Control}
\begin{scratch}\blockcontrol{esperar \ovalnum{1} segundos}\end{scratch} & \py|time.sleep(1)| & Pon \py|import time| hasta arriba. \dictrow
\begin{scratch}\blockrepeat{repetir \ovalnum{10}}{\blockspace}\end{scratch} & \py|for i in range(10):|\newline\py|    forward(10)| & Los cuatro espacios son la boca de la C. Usa números enteros: \py|range(2.5)| es un error. \dictrow
\begin{scratch}\blockinfloop{por siempre}{\blockspace}\end{scratch} & \py|while True:|\newline\py|    forward(10)| & Nunca se detiene solo, así que presiona \stopbtn{}. En Colab es más seguro un bucle \py|for| grande, como \py|range(200)|. \dictrow
\begin{scratch}\blockif{si \booloperator{} entonces}{\blockspace}\end{scratch} & \py|if puntos > 5:|\newline\py|    print("¡Guau!")| & Los dos puntos abren la boca. \dictrow
\begin{scratch}\blockifelse{si \booloperator{} entonces}{\blockspace}{\blockspace}\end{scratch} & \py|if puntos > 5:|\newline\py|    print("grande")|\newline\py|else:|\newline\py|    print("chico")| & \py|else:| se alinea con el \py|if|. Python también tiene \py|elif| para «si no, si...», que Scratch no tiene. \dictrow
\begin{scratch}\blockcontrol{esperar hasta que \booloperator{}}\end{scratch} & \py|a = ""|\newline\py|while a != "si":|\newline\py|    a = input("¿Vamos? ")| & Sigue preguntando hasta que escribas \texttt{si} (sin acento, para que sea fácil de escribir). Python hace una cosa a la vez, así que algo dentro del bucle tiene que cambiar. \dictrow
\begin{scratch}\blockrepeat{repetir hasta que \booloperator{}}{\blockspace}\end{scratch} & \py|while puntos != 5:| & ¡Al revés! Repetir \textbf{hasta que} sea 5 quiere decir seguir \textbf{mientras} no sea 5. \dictrow
\begin{scratch}\blockstop{detener \selectmenu{todos}}\end{scratch} & \py|sys.exit()| & Pon \py|import sys| hasta arriba. Colab muestra un pequeño aviso: no pasa nada. \dictrow
\begin{scratch}\blockstop{detener \selectmenu{este programa}}\end{scratch} & \py|break|\newline\py|return| & \py|break| salta fuera de un bucle. \py|return| sale de tu propio bloque. \dictrow
\begin{scratch}[scale=0.75]\blockcontrol{detener \selectmenu{otros programas en el objeto}}\end{scratch} & No hace falta & Python solo ejecuta un programa a la vez. \dictrow
\begin{scratch}\blockinitclone{al comenzar como clon}\end{scratch} & \py|s = Turtle()|\newline\py|s.shape("circle")|\newline\py|s.goto(0, 100)| & Las líneas justo después de \py|Turtle()| preparan cada objeto nuevo. \dictrow
\begin{scratch}\blockcontrol{crear clon de \selectmenu{mí mismo}}\end{scratch} & \py|estrella = Turtle()|\newline o, para muchos:\newline\py|estrellas = []|\newline\py|s = Turtle()|\newline\py|estrellas.append(s)| & Esto hace un objeto nuevo y le pone un nombre. Si son muchos clones, guárdalos en una lista. \dictrow
\begin{scratch}\blockcontrol{eliminar este clon}\end{scratch} & \py|s.hideturtle()|\newline\py|estrellas.remove(s)| & Escóndelo y luego sácalo de la lista. \dictrow
% ---------------------------------------------------------------- SENSORES
\cat{sensing}{Sensores}
\boolsensing{¿tocando \selectmenu{Objeto2}?} & \py|p = gato.pos()|\newline\py|distance(p) < 20| & Pregunta: ¿está el gato a menos de 20 pasos? El borde: \py|abs(xcor()) > 240| \dictrow
\boolsensing{¿tocando el color \selectmenu{rojo}?} & \py|ycor() < -100| & No se puede: las tortugas no ven colores. Mejor recuerda dónde están las cosas: «la lava está debajo de $y = -100$». \dictrow
\boolsensing{¿color \selectmenu{rojo} está tocando \selectmenu{azul}?} & Usa lugares, no colores & El mismo truco que en \textbf{¿tocando el color?}. \dictrow
\ovalsensing{distancia a \selectmenu{Objeto2}} & \py|distance(gato.pos())| & Puntero del ratón: no en Colab. \dictrow
\begin{scratch}\blocksensing{preguntar \ovalnum{¿Cómo te llamas?} y esperar}\end{scratch} & \py|nombre = input("¿Quién? ")| & Tú escoges el nombre de la caja. \dictrow
\ovalsensing{respuesta} & \py|nombre| & \py|input| siempre te da \textbf{texto}. Para un número: \py|int(nombre)| \dictrow
\boolsensing{¿tecla \selectmenu{espacio} presionada?} & \py|k = input("¿Tecla? ")|\newline\py|k == " "| & Colab no siente las teclas. Escribe y luego presiona Enter. Thonny sí puede: \py|onkeypress|. \dictrow
\boolsensing{¿ratón presionado?} & \py|input("¿Clic? ")| & No se puede en Colab: no hay ratón. Mejor pregunta. \dictrow
\ovalsensing{posición x del ratón}\newline\ovalsensing{posición y del ratón} & \py|x = int(input("¿x? "))| & Mejor escribe los números. \dictrow
\begin{scratch}\blocksensing{fijar modo de arrastre a \selectmenu{arrastrable}}\end{scratch} & No en Colab & Pruébalo en Thonny: \py|ondrag(goto)| \dictrow
\ovalsensing{volumen del sonido} & \py|input("¿Qué tan fuerte? ")| & No se puede: Colab no puede usar tu micrófono. Mejor pide un número. \dictrow
\ovalsensing{cronómetro} & \py|time.time() - inicio| & Primero escribe \py|inicio = time.time()|. Es tu propio cronómetro. \dictrow
\begin{scratch}\blocksensing{reiniciar cronómetro}\end{scratch} & \py|inicio = time.time()| & Anota una nueva hora de inicio. \dictrow
\ovalsensing{\selectmenu{posición en x} de \selectmenu{Objeto2}} & \py|gato.xcor()| & También \py|gato.ycor()| y \py|gato.heading()|. \dictrow
\ovalsensing{\selectmenu{año} actual} & \py|d = time.localtime()|\newline\py|d.tm_year| & También \py|d.tm_mon| (mes), \py|d.tm_mday| (día), \py|d.tm_hour| (hora), \py|d.tm_min| (minuto). El reloj de Colab puede estar en la hora de Londres. \dictrow
\ovalsensing{días desde el 2000} & \py|s = time.time()|\newline\py|s / 86400 - 10957| & Los segundos desde 1970 se vuelven días; luego resta los días de 1970 a 2000. \dictrow
\ovalsensing{nombre de usuario} & \py|nombre = input("¿Quién? ")| & Colab no sabe cómo te llamas, ¡así que pregunta! \dictrow
% ---------------------------------------------------------------- OPERADORES
\cat{operators}{Operadores}
\ovaloperator{\ovalnum{7} + \ovalnum{2}}\ \ovaloperator{\ovalnum{7} - \ovalnum{2}}\newline\ovaloperator{\ovalnum{7} * \ovalnum{2}}\ \ovaloperator{\ovalnum{7} / \ovalnum{2}} & \py|7 + 2|\quad\py|7 - 2|\newline\py|7 * 2|\quad\py|7 / 2| & \py|/| siempre da un decimal: \py|6 / 3| es \texttt{2.0}. Dividir entre 0 es un error. \dictrow
\ovaloperator{elegir número al azar entre \ovalnum{1} y \ovalnum{10}} & \py|random.randint(1, 10)| & Pon \py|import random| hasta arriba. Pueden salir el 1 y el 10. \dictrow
\booloperator{\ovalvariable{a} > \ovalnum{50}} & \py|a > 50| & «Mayor o igual» es \py|a >= 50|. \dictrow
\booloperator{\ovalvariable{a} < \ovalnum{50}} & \py|a < 50| & \dictrow
\booloperator{\ovalvariable{a} = \ovalnum{50}} & \py|a == 50| & ¡Dos signos de igual para comparar! Las mayúsculas cuentan: \py|"Gato" == "gato"| es \py|False|. \dictrow
\booloperator{\booloperator{} y \booloperator{}} & \py|a > 1 and a < 9| & En inglés y en minúsculas: \py|and| \dictrow
\booloperator{\booloperator{} o \booloperator{}} & \py|a < 1 or a > 9| & Aquí es \py|or|. \dictrow
\booloperator{no \booloperator{}} & \py|not a == 5| & Igual que \py|a != 5| \dictrow
\ovaloperator{unir \ovalnum{manzana } \ovalnum{banana}} & \py|"manzana " + "banana"| & ¿Vas a unir un número? Usa \py|str(5)|. \dictrow
\ovaloperator{letra \ovalnum{1} de \ovalnum{manzana}} & \py|"manzana"[0]| & ¡Python cuenta desde 0! La letra \texttt{n} es \py|[n - 1]|. \dictrow
\ovaloperator{longitud de \ovalnum{manzana}} & \py|len("manzana")| & \dictrow
\booloperator{¿\ovalnum{manzana} contiene \ovalnum{a}?} & \py|"a" in "manzana"| & El orden se voltea. Las mayúsculas cuentan. \dictrow
\ovaloperator{\ovalnum{7} módulo \ovalnum{3}} & \py|7 % 3| & El residuo. Números pares: \py|n % 2 == 0| \dictrow
\ovaloperator{redondear \ovalnum{2.5}} & \py|round(2.5)|\newline\py|math.floor(x + 0.5)| & ¡Cuidado! \py|round(2.5)| da 2, pero Scratch da 3. La segunda línea redondea como Scratch. \dictrow
\ovaloperator{\selectmenu{valor absoluto} de \ovalnum{-5}} & \py|abs(-5)| & No hace falta importar nada. \dictrow
\ovaloperator{\selectmenu{piso} de \ovalnum{2.7}}\newline\ovaloperator{\selectmenu{techo} de \ovalnum{2.1}} & \py|math.floor(2.7)|\newline\py|math.ceil(2.1)| & Pon \py|import math| hasta arriba. Piso es \py|floor| y techo es \py|ceil| (de \emph{ceiling}, techo en inglés). \dictrow
\ovaloperator{\selectmenu{raíz cuadrada} de \ovalnum{16}} & \py|math.sqrt(16)| & \dictrow
\ovaloperator{\selectmenu{sen} de \ovalnum{30}}\newline\ovaloperator{\selectmenu{cos}}\ \ovaloperator{\selectmenu{tan}} & \py|r = math.radians(30)|\newline\py|math.sin(r)| & Python mide los ángulos en radianes, así que primero cambia los grados. Lo mismo para \py|cos| y \py|tan|. Ojo: seno es \py|sin|. \dictrow
\ovaloperator{\selectmenu{arcsen} de \ovalnum{0.5}}\newline\ovaloperator{\selectmenu{arccos}}\ \ovaloperator{\selectmenu{arctan}} & \py|a = math.asin(0.5)|\newline\py|math.degrees(a)| & Vuelve a cambiar la respuesta a grados. \dictrow
\ovaloperator{\selectmenu{ln} de \ovalvariable{x}}\ \ovaloperator{\selectmenu{log} de \ovalvariable{x}} & \py|math.log(x)|\newline\py|math.log10(x)| & El \py|log| de Python es el \texttt{ln} de Scratch. \dictrow
\ovaloperator{\selectmenu{e \^{}} de \ovalvariable{x}}\ \ovaloperator{\selectmenu{10 \^{}} de \ovalvariable{x}} & \py|math.exp(x)|\newline\py|10 ** x| & La potencia es \py|**|. El signo \texttt{\^{}} quiere decir otra cosa en Python. \dictrow
% ---------------------------------------------------------------- VARIABLES
\cat{variables}{Variables}
\chip[variables]{Crear una variable} & \py|puntos = 0| & No hay botón: el primer \py|=| hace la caja. \dictrow
\ovalvariable{puntos} & \py|puntos| & Usa el nombre en cualquier lugar donde va un número. \dictrow
\begin{scratch}\blockvariable{fijar \selectmenu{puntos} a \ovalnum{0}}\end{scratch} & \py|puntos = 0|\newline\py|nombre = "Ana"| & \py|=| mete un valor en la caja. El texto necesita comillas. \dictrow
\begin{scratch}\blockvariable{cambiar \selectmenu{puntos} en \ovalnum{1}}\end{scratch} & \py|puntos = puntos + 1|\newline\py|puntos += 1| & Las dos líneas hacen lo mismo. Dentro de tu propio bloque, escribe primero \py|global puntos|. \dictrow
\begin{scratch}\blockvariable{mostrar variable \selectmenu{puntos}}\end{scratch} & \py|print(puntos)| & Muestra el valor una vez. Imprime otra vez para ver un valor nuevo. \dictrow
\begin{scratch}\blockvariable{esconder variable \selectmenu{puntos}}\end{scratch} & Simplemente no lo imprimas & \dictrow
% ---------------------------------------------------------------- LISTAS
\cat{lists}{Listas}
\chip[lists]{Crear una lista} & \py|comida = []| & Los corchetes hacen una lista vacía. \dictrow
\ovallist{comida} & \py|" ".join(comida)| & Todos los elementos, con espacios entre ellos. \dictrow
\begin{scratch}\blocklist{añadir \ovalnum{manzana} a \selectmenu{comida}}\end{scratch} & \py|comida.append("manzana")| & Lo agrega al final. \dictrow
\begin{scratch}\blocklist{eliminar \ovalnum{1} de \selectmenu{comida}}\end{scratch} & \py|del comida[0]| & ¡Contando desde 0! El elemento \texttt{n} es \py|[n - 1]|. El último: \py|comida.pop()| \dictrow
\begin{scratch}\blocklist{Eliminar todos de \selectmenu{comida}}\end{scratch} & \py|comida.clear()| & \dictrow
\begin{scratch}\blocklist{insertar \ovalnum{higo} en \ovalnum{1} de \selectmenu{comida}}\end{scratch} & \py|comida.insert(0, "higo")| & \dictrow
\begin{scratch}[scale=0.62]\blocklist{reemplazar elemento \ovalnum{1} de \selectmenu{comida} con \ovalnum{kiwi}}\end{scratch} & \py|comida[0] = "kiwi"| & \dictrow
\ovallist{elemento \ovalnum{1} de \selectmenu{comida}} & \py|comida[0]| & ¡El primer elemento es el número 0! El último: \py|comida[-1]|. Al azar: \py|random.choice(comida)| \dictrow
\ovallist{\# de elemento de \ovalnum{higo} en \selectmenu{comida}} & \py|comida.index("higo") + 1| & Suma 1, porque Python cuenta desde 0. Revisa primero con \py|in|: un elemento que no está es un error. \dictrow
\ovallist{longitud de \selectmenu{comida}} & \py|len(comida)| & \dictrow
\boollist{¿\ovalnum{manzana} está en \selectmenu{comida}?} & \py|"manzana" in comida| & ¡Aquí el orden es igual que en Scratch! Las mayúsculas cuentan: \texttt{"Manzana"} es distinta. \dictrow
\begin{scratch}\blocklist{mostrar lista \selectmenu{comida}}\end{scratch} & \py|print(comida)| & \dictrow
\begin{scratch}\blocklist{esconder lista \selectmenu{comida}}\end{scratch} & Simplemente no la imprimas & \dictrow
% ---------------------------------------------------------------- MIS BLOQUES
\cat{myblocks}{Mis Bloques}
\begin{scratch}\initmoreblocks{definir \namemoreblocks{saltar}}\end{scratch} & \py|def saltar():|\newline\py|    forward(20)| & Las líneas de adentro llevan cuatro espacios. Ejecuta esta celda antes de usar el bloque. \dictrow
\begin{scratch}\blockmoreblocks{saltar}\end{scratch} & \py|saltar()| & El paréntesis hace que se ejecute. \dictrow
\begin{scratch}\initmoreblocks{definir \namemoreblocks{saltar \ovalmoreblocks{altura}}}\end{scratch} & \py|def saltar(altura):|\newline\py|    forward(altura)| & La entrada es una caja que solo funciona dentro del bloque. \dictrow
\begin{scratch}\blockmoreblocks{saltar \ovalnum{30}}\end{scratch} & \py|saltar(30)| & \dictrow
\begin{scratch}[scale=0.78]\initmoreblocks{definir \namemoreblocks{hola \ovalmoreblocks{nombre} ¿fuerte? \boolmoreblocks{fuerte}}}\end{scratch} & \py|def hola(nombre, fuerte):|\newline\py|    if fuerte:|\newline\py|        print("¡HOLA!")| & Llámalo con \py|hola("Ana", True)|. \py|True| y \py|False| empiezan con mayúscula. \dictrow
\chip[myblocks]{Ejecutar sin actualizar}\newline\chip[myblocks]{la pantalla} & \py|speed(13)| & La velocidad más rápida que todavía muestra el dibujo. \dictrow
Un bloque que te devuelve una respuesta (¡Scratch no puede hacer esto!) & \py|def doble(n):|\newline\py|    return n * 2| & \py|print(doble(4))| muestra 8. Como un bloque reportero redondo hecho por ti. \dictrow
% ---------------------------------------------------------------- LÁPIZ
\cat{pen}{Lápiz}
\begin{scratch}\blockpen{borrar todo}\end{scratch} & \py|clear()| & Borra el dibujo de esta tortuga. \py|clearscreen()| borra todo, también los objetos. \dictrow
\begin{scratch}\blockpen{sellar}\end{scratch} & \py|stamp()| & Sella la forma de la tortuga. \py|clearstamps()| borra todos los sellos: ¡Scratch no puede hacer eso! \dictrow
\begin{scratch}\blockpen{bajar lápiz}\end{scratch} & \py|pendown()| & La tortuga empieza con el lápiz ya bajado. \dictrow
\begin{scratch}\blockpen{subir lápiz}\end{scratch} & \py|penup()| & Úsalo para moverte sin dibujar. \dictrow
\begin{scratch}\blockpen{fijar color de lápiz a \ovalnum{rojo}}\end{scratch} & \py|pencolor("red")| & Nombres de colores en inglés y en minúsculas. \py|color("red")| pinta la tortuga y el lápiz. \dictrow
\begin{scratch}[scale=0.77]\blockpen{cambiar \selectmenu{color} de lápiz en \ovalnum{10}}\end{scratch} & \py|c = (c + 1) % 3|\newline\py|pencolor(arcoiris[c])| & Primero haz una lista de 3 colores llamada \texttt{arcoiris}, y \py|c = 0|. \dictrow
\begin{scratch}\blockpen{fijar \selectmenu{color} de lápiz a \ovalnum{50}}\end{scratch} & \py|pencolor("cyan")| & Números de Scratch: 0 rojo, 8 naranja, 16 amarillo, 33 verde, 50 cian, 66 azul, 75 morado. \dictrow
\begin{scratch}\blockpen{fijar \selectmenu{brillo} de lápiz a \ovalnum{50}}\end{scratch} & \py|pencolor("darkred")| & Saturación y brillo: escoge un nombre más pálido o más oscuro, como \texttt{"pink"} (rosa) o \texttt{"darkred"} (rojo oscuro). \dictrow
\begin{scratch}[scale=0.71]\blockpen{fijar \selectmenu{transparencia} de lápiz a \ovalnum{50}}\end{scratch} & \py|pencolor("pink")| & No se puede con el lápiz. Un color más pálido se ve un poco transparente. \dictrow
\begin{scratch}\blockpen{cambiar tamaño de lápiz en \ovalnum{1}}\end{scratch} & \py|g = pensize() + 1|\newline\py|pensize(g)| & \py|pensize()| con el paréntesis vacío te da el tamaño de ahora. \dictrow
\begin{scratch}\blockpen{fijar tamaño de lápiz a \ovalnum{5}}\end{scratch} & \py|pensize(5)| & \dictrow
Sin bloque de Scratch: ¡un regalo extra! & \py|fillcolor("gold")|\newline\py|begin_fill()|\newline\py|circle(50)|\newline\py|end_fill()| & Rellena una figura de color. \dictrow
% ---------------------------------------------------------------- MÚSICA
\cat{pen}{Música}
\begin{scratch}\blockmusic{tocar nota \ovalnum{60} durante \ovalnum{0.5} tiempos}\end{scratch} & \py|pitido(262)|\newline\py|time.sleep(0.5)| & La nota 60 es el do central (262); mi es 330. Nota \texttt{n}:\newline\py|f = 2 ** ((n-69)/12)|\newline\py|pitido(440 * f)| \dictrow
\begin{scratch}[scale=0.65]\blockmusic{tocar tambor \selectmenu{(1) Caja} durante \ovalnum{0.25} tiempos}\end{scratch} & \py|pitido(100, 0.1)|\newline\py|time.sleep(0.25)| & Una nota muy grave y corta suena un poco como un tambor. \dictrow
\begin{scratch}\blockmusic{fijar instrumento a \selectmenu{(1) Piano}}\end{scratch} & No se puede con \py|pitido| & \py|pitido| tiene un solo sonido. Mejor prueba notas cortas y largas. \dictrow
\begin{scratch}\blockmusic{fijar tempo a \ovalnum{60}}\blockmusic{silenciar durante \ovalnum{1} tiempos}\end{scratch} & \py|tiempo = 60 / 60|\newline\py|time.sleep(tiempo)| & Tempo 60 quiere decir un tiempo cada segundo; un silencio es una espera. Un \texttt{tiempo} más chico es más rápido. \dictrow
Una melodía & \py|melodia = [262, 330]|\newline\py|for f in melodia:|\newline\py|    pitido(f)|\newline\py|    time.sleep(0.5)| & ¡Una lista de notas y un bucle hacen una melodía! Agrega más notas a la lista. \\
\bottomrule
\end{longtable}
\endgroup
\setscratch{scale=1}

% =====================================================================
\banner{operators}{TEN ESTA PÁGINA CERCA}{Mi hoja de ayuda de Python}

\begin{tcbraster}[raster columns=2, raster column skip=10pt, raster row skip=8pt, raster equal height=rows,
  raster every box/.style={arc=6pt, coltitle=ink, fontupper=\small, left=6pt, right=4pt, top=3pt, bottom=3pt}]
\begin{tcolorbox}[raster multicolumn=2, colback=white, colframe=pen, colbacktitle=pen!40, title=\textbf{Empieza un dibujo: las tres líneas de inicio}]
\py|from ColabTurtlePlus.Turtle import *|\\
\py|clearscreen()|\\
\py|initializeTurtle((480, 360))|\\[2pt]
{\color{muted}Una vez por sesión, primero: \py|!pip install ColabTurtlePlus|}
\end{tcolorbox}
\begin{tcolorbox}[colback=white, colframe=motion, colbacktitle=motion!40, title=\textbf{Moverse}]
\py|forward(100)| \quad \py|back(50)|\\
\py|right(90)| \quad \py|left(90)|\\
\py|goto(0, 0)|\\
\py|setheading(90 - d)| {\color{muted}(la regla del 90 menos)}\\
\py|speed(13)| {\color{muted}(1 lento a 13 rápido)}
\end{tcolorbox}
\begin{tcolorbox}[colback=white, colframe=pen, colbacktitle=pen!40, title=\textbf{Lápiz}]
\py|penup()| \quad \py|pendown()|\\
\py|pencolor("red")|\\
\py|pensize(5)|\\
\py|clear()| \quad \py|stamp()|\\
\py|bgcolor("skyblue")|
\end{tcolorbox}
\begin{tcolorbox}[colback=white, colframe=looks, colbacktitle=looks!40, title=\textbf{Decir y preguntar}]
\py|print("Hola")|\\
\py|write("Hola")| {\color{muted}(en el escenario)}\\
\py|nombre = input("¿Quién? ")|\\
\py|edad = int(input("¿Edad? "))|\\
{\color{muted}¡\py|input| siempre te da texto!}
\end{tcolorbox}
\begin{tcolorbox}[colback=white, colframe=control, colbacktitle=control!40, title=\textbf{Bucles}]
\py|for i in range(4):|\\
\py|    forward(100)|\\[3pt]
\py|while puntos != 5:|\\
\py|    puntos = puntos + 1|
\end{tcolorbox}
\begin{tcolorbox}[colback=white, colframe=control, colbacktitle=control!40, title=\textbf{Si (if)}]
\py|if intento < secreto:|\\
\py|    print("Muy bajo")|\\
\py|elif intento > secreto:|\\
\py|    print("Muy alto")|\\
\py|else:|\\
\py|    print("¡Sí!")|
\end{tcolorbox}
\begin{tcolorbox}[colback=white, colframe=variables, colbacktitle=variables!40, title=\textbf{Variables}]
\py|puntos = 0|\\
\py|puntos = puntos + 1|\\
\py|puntos += 1| {\color{muted}(lo mismo)}\\
\py|print("Puntos: " + str(puntos))|
\end{tcolorbox}
\begin{tcolorbox}[colback=white, colframe=lists, colbacktitle=lists!40, title=\textbf{Listas}]
\py|comida = ["manzana", "higo"]|\\
\py|comida.append("kiwi")|\\
\py|comida[0]| {\color{muted}(¡el primer elemento!)}\\
\py|len(comida)| \quad \py|"higo" in comida|\\
\py|for f in comida:|\\
\py|    print(f)|
\end{tcolorbox}
\begin{tcolorbox}[colback=white, colframe=operators, colbacktitle=operators!40, title=\textbf{Operadores}]
\py|+  -  *  /| \quad \py|7 % 3| {\color{muted}(residuo)}\\
\py|==  !=  <  >  <=  >=|\\
\py|and  or  not|\\
\py|import random|\\
\py|random.randint(1, 10)|
\end{tcolorbox}
\begin{tcolorbox}[colback=white, colframe=myblocks, colbacktitle=myblocks!40, title=\textbf{Mi propio bloque}]
\py|def cuadrado(lado):|\\
\py|    for i in range(4):|\\
\py|        forward(lado)|\\
\py|        right(90)|\\
\py|cuadrado(50)|\\[3pt]
\py|def doble(n):|\\
\py|    return n * 2|\\
\py|print(doble(4))| {\color{muted}(muestra 8)}
\end{tcolorbox}
\begin{tcolorbox}[colback=white, colframe=motion, colbacktitle=motion!40, title=\textbf{Más objetos}]
\py|gato = Turtle()|\\
\py|gato.shape("turtle")|\\
\py|gato.goto(-100, 0)|\\
\py|gato.forward(50)|\\
\py|gato.distance(pos()) < 20| {\color{muted}(¿tocando?)}
\end{tcolorbox}
\begin{tcolorbox}[colback=white, colframe=sound, colbacktitle=sound!40, title=\textbf{Esperar y sonido}]
\py|import time|\\
\py|time.sleep(1)| {\color{muted}(espera 1 segundo)}\\
\py|pitido(262)| {\color{muted}(do central, ayuda de la Parte 7)}\\
\py|time.sleep(0.5)|
\end{tcolorbox}
\end{tcbraster}

\vspace{8pt}
\begin{tcolorbox}[colback=warnbg, colframe=warnframe, arc=6pt, title=\textbf{Recuerda}, coltitle=ink, colbacktitle=warnframe!35]
\begin{tabular}{@{}P{7.7cm} P{7.7cm}@{}}
\textbf{:} \ Los dos puntos abren un bucle, un \texttt{if} o un \texttt{def}. & \textbf{Cuatro espacios} meten una línea adentro, como la boca de un bloque en forma de C. \\
\textbf{=} \ mete en una caja. \ \textbf{==} \ compara. \ \textbf{!=} \ quiere decir «no es igual». & \textbf{"Las palabras"} necesitan comillas. Los números no. \\
Python cuenta desde \textbf{0}: el elemento 1 es \py|[0]|. & Las mayúsculas importan: \py|print|, no \py|Print|, y \py|"Gato" == "gato"| es \py|False|. \\
\py|input()| da texto. Usa \py|int( )| para un número. & \py|round(2.5)| es 2 en Python, pero 3 en Scratch. \\
Direcciones: \py|setheading(90 - d)|. & Nombres de colores en inglés y en minúsculas: \py|"red"|. \\
Ejecuta la celda «En el principio» en cada sesión. & Lee primero la \textbf{última línea} de un error (Python la escribe en inglés). ¡Es una pista!
\end{tabular}
\end{tcolorbox}

\partopener{myblocks}{PARA LOS ADULTOS}{Cómo ayudar a tu programador}{No necesitas saber Python para ayudar. Estas páginas te dicen cómo.}

Este libro está hecho para dos personas: tu hijo o hija frente al teclado y tú a su lado. No necesitas saber Python. Tu trabajo es ser entrenador y público: hacer preguntas, leer juntos en voz alta los mensajes de error y asombrarte cuando algo funciona.

Algunos de los consejos que siguen vienen de investigaciones sobre niños que pasan de programar con bloques a programar con texto. Cuando un consejo es solo nuestra propia experiencia y no un resultado de investigación, lo marcamos como \textit{(criterio propio)}. Los artículos aparecen al final de esta sección.

\subsection*{Cómo funciona el libro}

El libro tiene siete Partes. \textbf{No} siguen el orden de los colores de los bloques en Scratch. Siguen el orden que la investigación sugiere para pasar de los bloques al texto:

\begin{center}
\renewcommand{\arraystretch}{1.35}
\begin{tabular}{l P{4.6cm} P{9.4cm}}
\toprule
 & \textbf{Parte} & \textbf{Por qué va aquí} \\
\midrule
1 & \chip[motion]{Movimiento y Lápiz} & Dibujar da una imagen de inmediato y coincide con bloques que tu hijo ya conoce. Ejecutar una celda hace el papel de la bandera verde. \\
2 & \chip[control]{Bucles} & Las figuras necesitan \chip[control]{repetir} desde el principio. Los bucles son más difíciles en texto que en bloques, así que se practican pronto. \\
3 & \chip[variables]{Variables y Listas} & Puntos, nombres y tamaños, en cuanto un proyecto los necesita. \\
4 & \chip[operators]{Operadores} & Matemáticas, números al azar y comparaciones, que la Parte siguiente necesita. \\
5 & \chip[sensing]{Sensores y Decisiones} & \texttt{input()} e \texttt{if}: cuestionarios y juegos de adivinar. \\
6 & \chip[myblocks]{Mis Bloques} & Los superbloques propios de tu hijo, escritos con \texttt{def}. \\
7 & \chip[looks]{Apariencia, Sonido}\newline\chip[looks]{y Eventos} & Al final, a propósito (mira más abajo). \\
\bottomrule
\end{tabular}
\end{center}

\textbf{Por qué los eventos van al final.} En Scratch, bloques como \chip[events]{al presionar bandera verde}, \chip[events]{al presionar tecla espacio}, y tener muchos programas funcionando a la vez, no cuestan nada. En código de texto son lo más difícil. Cuando Dann y sus colegas (2012) enseñaron a estudiantes a pasar del lenguaje de bloques Alice a Java, los eventos y los programas que funcionan al mismo tiempo fueron lo más difícil de trasladar. En Python, presionar una tecla necesita una \emph{función} que se ejecute cuando se presiona la tecla, así que las funciones (Parte 6) tienen que ir antes.

Después de las siete Partes vienen tres secciones que tu hijo puede usar en cualquier momento:
\begin{itemize}
  \item \textbf{Grandes Proyectos}: programas más largos que usan todo, para cuando terminen las Partes.
  \item \textbf{Zoológico de Bichos}: los mensajes de error comunes de Python, qué significan y cómo arreglarlos. Tenlo abierto desde el primer día.
  \item \textbf{Diccionario de Bloques}: busca un bloque de Scratch y encuentra su Python.
\end{itemize}

Todas las Partes tienen la misma forma: páginas cortas que ponen cada bloque de Scratch junto a su Python, luego de dos a cuatro \textbf{misiones}, después \textbf{Tradúcelo tú}, un \textbf{Acertijo de dos líneas} y una lista «Ya sé\ldots». Las respuestas están al final del libro.

\textbf{Para familias de habla hispana.} Python habla inglés: sus palabras clave (\texttt{for}, \texttt{print}, \texttt{forward}\ldots) y sus mensajes de error siempre están en inglés, y no se pueden cambiar. Scratch, en cambio, se puede poner en español con el ícono del globo terráqueo de la barra superior: elige «Español (Latinoamérica)», que es la versión cuyos nombres de bloques usa este libro. Los menús de Colab aparecen en el idioma de tu cuenta de Google; aquí los escribimos en español con el inglés entre paréntesis. Y esta edición usa nombres en español para todo lo que tu hijo inventa: variables, funciones, objetos y textos.

\subsection*{Cómo llevar una sesión}

\begin{itemize}
  \item \textbf{Que sea corta: de 20 a 30 minutos.} Terminen cuando algo funcione y se le pueda mostrar a alguien. Paren antes de que crezca la frustración. Una misión suele llevar dos sesiones: una para que funcione y otra para los retos de «Hazlo tuyo». \textit{(criterio propio: la investigación no da duraciones de sesión)}
  \item \textbf{Primero Scratch, siempre.} El Paso 1 de cada misión es construirla en Scratch. Así lo único nuevo es escribir, y el salto se siente pequeño. \textit{(criterio propio)}
  \item \textbf{Digan las conexiones en voz alta.} Este es el paso que más importa. Weintrop y Wilensky (2019) siguieron a estudiantes que habían empezado con bloques: su ventaja inicial se había desvanecido al pasar a un lenguaje de texto, en una clase donde nadie conectaba los bloques con el código. El método con mejor evidencia es poner los bloques y el código lado a lado, señalar desde cada bloque la línea que hace el mismo trabajo, y luego traducir un programa que no han visto antes (Dann et al., 2012). Las páginas lado a lado, las tablas de «Léelo línea por línea» y la página \textbf{Tradúcelo tú} de cada Parte están hechas exactamente para esto. Di cosas como «\texttt{for i in range(4):} es el bloque repetir~4, y los espacios son la boca de la C».
  \item \textbf{Predecir, luego ejecutar, luego cambiar.} Empezar con código que funciona es más fácil que empezar con una celda vacía. Dos métodos de enseñanza muy conocidos usan este orden. PRIMM (Sentance, Waite y Kallia, 2019) tiene cinco pasos (sus letras vienen del inglés): \textbf{P}redecir qué hará el código, ejecutarlo (\textbf{R}un), \textbf{I}nvestigar («¿qué hace esta línea?»), \textbf{M}odificar (cambiar un número, un color o una palabra) y crear algo nuevo (\textbf{M}ake). Usar-Modificar-Crear (Lee et al., 2011) es la misma idea en tres pasos. Antes de que tu hijo ejecute algo, pregúntale: «¿Qué crees que va a pasar?».
  \item \textbf{Deja que escriba.} Copiar código letra por letra es como los niños aprenden los dos puntos, los paréntesis y los espacios. Intenta no quitarle el teclado ni corregirle los errores de escritura. \textit{(criterio propio)}
  \item \textbf{Los errores son normales.} Scratch nunca muestra un mensaje de error, así que probablemente tu hijo nunca ha visto uno. El texto rojo es una pista, no un fracaso. Lean juntos en voz alta la \textbf{última línea}, busquen el número de línea y busquen el error en el Zoológico de Bichos. «Tomar un descanso» es una estrategia real para depurar que los niños mencionaron a Brennan y Resnick (2012), así que puedes enseñarla como tal.
  \item \textbf{Celebra los momentos de «programador de verdad».} Los estudiantes que pasaron al código de texto dijeron que se sentían programadores de verdad, y esa sensación motiva (Weintrop y Wilensky, 2017).
\end{itemize}

\subsection*{La revisión de cuatro preguntas al final de una Parte}

En este libro no hay exámenes. En su lugar, cuando tu hijo termine una Parte, elige uno de sus programas y hazle cuatro preguntas. Vienen de las entrevistas que Brennan y Resnick (2012) usaron para ver lo que los niños entendían de verdad:

\begin{tcolorbox}[enhanced, colback=myblocks!10, colframe=myblocks!80!black, arc=6pt, boxrule=0.8pt, left=8pt]
\begin{enumerate}[leftmargin=*]
  \item \textbf{Explicar:} «¿Qué hace tu programa? Muéstrame la línea que hace\ldots».
  \item \textbf{Ampliar:} «¿Cómo podrías hacerlo más grande o mejor?».
  \item \textbf{Arreglar un bicho:} cambia una cosa en secreto (quita unos dos puntos, un espacio o un paréntesis) y pídele que lo encuentre.
  \item \textbf{Remezclar:} «Agrega una cosa nueva que sea tuya».
\end{enumerate}
\end{tcolorbox}

Si una de las cuatro cuesta trabajo, dediquen otra sesión a esa Parte antes de seguir. Saber qué \emph{es} un bucle importa menos que poder usarlo, cambiarlo y arreglarlo.

\needspace{19\baselineskip}
\subsection*{Un ritmo sugerido}

Esto es solo una guía \textit{(criterio propio)}. Supone dos o tres sesiones de 20 a 30 minutos por semana. Vayan más despacio si una Parte cuesta trabajo; es tiempo bien empleado.

\begin{center}
\renewcommand{\arraystretch}{1.1}
\begin{tabular}{l P{7.4cm} c c}
\toprule
 & \textbf{Parte} & \textbf{Sesiones} & \textbf{Semanas} \\
\midrule
 & Preparándonos (Colab, las líneas de inicio) & 1--2 & 1 \\
1 & Movimiento y Lápiz & 6--8 & 2--3 \\
2 & Bucles & 5--7 & 2--3 \\
3 & Variables y Listas & 6--8 & 2--3 \\
4 & Operadores & 5--7 & 2--3 \\
5 & Sensores y Decisiones & 6--8 & 2--3 \\
6 & Mis Bloques (funciones) & 6--8 & 2--3 \\
7 & Apariencia, Sonido y Eventos & 6--8 & 2--3 \\
 & Grandes Proyectos & 6--10 & 3--4 \\
\midrule
 & \textbf{Total} & \textbf{unas 50--65} & \textbf{unos 5 meses} \\
\bottomrule
\end{tabular}
\end{center}

Las Partes 2 y 6 suelen necesitar más tiempo. En estudios con estudiantes mayores, las preguntas sobre bucles y sobre llamar funciones se respondieron mejor en bloques que en texto (Weintrop y Wilensky, 2017).

\subsection*{Consejos para Colab}

\begin{itemize}
  \item \textbf{Cuenta de Google.} Colab necesita una. Las cuentas para niños (Family Link o algunas cuentas escolares) pueden tener Colab bloqueado. Si es así, inicia sesión con tu propia cuenta.
  \item \textbf{Guardar.} Los cuadernos se guardan solos en tu Google Drive, en una carpeta llamada \textbf{Colab Notebooks}. \textbf{Archivo} (\textit{File}) $\rightarrow$ \textbf{Guardar una copia en Drive} (\textit{Save a copy in Drive}) hace una copia, algo útil antes de un cambio grande.
  \item \textbf{Cuando Colab se reinicia.} Después de un rato sin usarlo, Colab olvida todo, incluida la biblioteca instalada. Si ves \texttt{No module named 'ColabTurtlePlus'} o un \texttt{NameError} con algo que antes funcionaba, abre el menú \textbf{Entorno de ejecución} (\textit{Runtime}) y elige \textbf{Ejecutar todas} (\textit{Run all}). No es un error de tu hijo.
  \item \textbf{Números de línea.} Actívalos para que coincidan con el libro: \textbf{Herramientas} (\textit{Tools}) $\rightarrow$ \textbf{Configuración} (\textit{Settings}) $\rightarrow$ \textbf{Editor} $\rightarrow$ \textbf{Mostrar números de línea} (\textit{Show line numbers}).
  \item \textbf{La biblioteca.} La tortuga de este libro viene de \textbf{ColabTurtlePlus}. Probamos todo en Colab el 7 de noviembre de 2026 con Python 3.13 y ColabTurtlePlus 2.2.0. La celda «En el principio» la instala una vez por sesión, y todos los programas de dibujo empiezan con las mismas tres líneas:
\end{itemize}

\begin{code}
!pip install ColabTurtlePlus
\end{code}
\begin{code}
from ColabTurtlePlus.Turtle import *
clearscreen()
initializeTurtle((480, 360))
\end{code}

\texttt{initializeTurtle((480, 360))} crea un escenario exactamente como el de Scratch: (0,~0) en el centro, x de $-240$ a 240, y de $-180$ a 180.

\begin{watchout}[Cuatro cosas que salen mal]
\begin{itemize}[leftmargin=*]
  \item \textbf{Instala la biblioteca correcta.} Hay una biblioteca más antigua con un nombre muy parecido, \textbf{ColabTurtle} (sin «Plus»). Es una biblioteca distinta: no tiene \texttt{Turtle()}, dibuja en una ventana negra de 800 por 500 y su (0,~0) está en la esquina superior izquierda. Si tu dibujo se ve así, revisa que la primera celda diga \texttt{ColabTurtlePlus}.
  \item \textbf{No uses} \py|mode("logo")|. Algunos tutoriales en internet lo usan para que las direcciones coincidan con Scratch. En ColabTurtlePlus tiene un bicho: comprobamos en Colab que \py|setheading(90)| entonces apunta \emph{hacia abajo}. El libro se queda en el modo normal y usa la regla «dirección en Python = 90 menos dirección en Scratch».
  \item \textbf{Usa} \py|speed(13)| \textbf{para dibujar rápido, no} \py|speed(0)|. Con 0, no aparece nada hasta que el programa llama a \py|done()|.
  \item \textbf{Los nombres de colores van en minúsculas y en inglés:} \py|"red"|, \py|"skyblue"|, no \py|"Red"|.
\end{itemize}
\end{watchout}

\textbf{Lo que Colab no puede hacer.} ColabTurtlePlus dibuja una imagen dentro del cuaderno. No puede escuchar teclas, clics del ratón ni temporizadores, y no tiene sonido ni disfraces con imágenes. Por eso, en Colab el libro usa \texttt{input()} (escribir una respuesta y presionar Enter) para interactuar. Los juegos de verdad con teclas son para Thonny, más abajo. Por favor, díselo a tu hijo con claridad, para que no piense que hizo algo mal.

\subsection*{Sonido en Colab}

Colab funciona en las computadoras de Google, así que Python no puede llegar directamente a tus bocinas. Pero sí puede enviar un sonido a la página web. La Parte~7 usa esta función de ayuda. Ejecútala una vez en una celda (no necesita ColabTurtlePlus):

\begin{code}
import numpy as np
from IPython.display import Audio, display

def pitido(frec=440, segs=0.5):
    tiempos = np.linspace(0, segs, int(22050 * segs))
    sonido = np.sin(2 * np.pi * frec * tiempos)
    display(Audio(sonido, rate=22050, autoplay=True))
\end{code}

Después, \py|pitido(262)| toca el do central, y \py|pitido(523, 1)| toca el do de arriba durante un segundo. Algunas notas para probar: do 262, re 294, mi 330, fa 349, sol 392.

\begin{itemize}
  \item \textbf{Reproducción automática.} El sonido empieza solo porque ejecutar la celda cuenta como un clic. Algunos navegadores (a menudo Safari), o una pestaña silenciada, lo bloquean. Entonces aparece un pequeño reproductor: presiona su botón de reproducir.
  \item \textbf{Notas encimadas.} Python termina la celda antes de que el navegador empiece a tocar, así que varios \texttt{pitido} en una misma celda pueden sonar encimados. Pon \py|time.sleep(0.6)| después de cada \texttt{pitido} (con \py|import time| al principio), o pon cada nota en su propia celda.
\end{itemize}

\subsection*{Más allá de Colab}

Cuando tu hijo quiera juegos de verdad con teclas, o si no tienes cuenta de Google, estos son los siguientes lugares a donde ir.

\textbf{Thonny} (\url{thonny.org}) es un editor de Python gratuito para Windows, Mac y Linux, hecho para principiantes. No necesita cuenta, funciona sin internet y nada sale de tu computadora. Incluye el \texttt{turtle} estándar de Python, que abre su propia ventana y \emph{sí} puede reaccionar al teclado. Los comandos de la tortuga son los mismos que en Colab, y la regla de «90 menos» para las direcciones sigue funcionando. Aquí tienes un pequeño juego de monedas: las flechas del teclado dirigen la tortuga, y tocar la moneda dorada suma un punto.

En Thonny, primero haz clic en la ventana de la tortuga y luego presiona las flechas. Los puntos aparecen en la barra de título de la ventana. Tres cosas para notar en el programa de abajo:
\begin{itemize}
  \item \py|screen.onkeypress(avanzar, "Up")| es \chip[events]{al presionar tecla flecha arriba}. Usa \texttt{onkeypress}, no \texttt{onkey}: \texttt{onkey} espera hasta que se suelta la tecla. Los nombres de las teclas (\texttt{"Up"}, \texttt{"Left"}, \texttt{"Right"}) van en inglés.
  \item Escribe \texttt{avanzar} sin paréntesis. Le estás diciendo a Python qué función ejecutar \emph{después}, no la estás ejecutando ahora.
  \item \py|screen.listen()| enciende el teclado, y \py|turtle.mainloop()| mantiene la ventana esperando teclas. Tiene que ser la última línea.
\end{itemize}
% caja de código partible (mismo aspecto que code) para no dejar media página vacía
\begin{tcblisting}{listing only, enhanced, breakable, colback=codebg, colframe=codeframe, boxrule=0.6pt, arc=5pt,
  left=22pt, right=6pt, top=2pt, bottom=2pt, listing options={style=py}}
import turtle
import random

screen = turtle.Screen()
screen.setup(480, 360)
jugador = turtle.Turtle()
jugador.shape("turtle")
jugador.penup()
moneda = turtle.Turtle()
moneda.shape("circle")
moneda.color("gold")
moneda.penup()
moneda.goto(100, 80)
puntos = 0

def revisar_moneda():
    global puntos
    if jugador.distance(moneda) < 20:
        puntos = puntos + 1
        screen.title("Puntos: " + str(puntos))
        x = random.randint(-200, 200)
        y = random.randint(-150, 150)
        moneda.goto(x, y)

def avanzar():
    jugador.forward(15)
    revisar_moneda()

def girar_izquierda():
    jugador.left(30)

def girar_derecha():
    jugador.right(30)

screen.onkeypress(avanzar, "Up")
screen.onkeypress(girar_izquierda, "Left")
screen.onkeypress(girar_derecha, "Right")
screen.listen()
turtle.mainloop()
\end{tcblisting}

Este ejemplo usa solo cosas de las Partes 1 a 6 más \texttt{global} (Parte 6), y por eso los eventos se dejan para el final.

\textbf{Pygame Zero, más adelante.} Pygame Zero es una biblioteca de juegos cuyas ideas se parecen mucho a Scratch: un \texttt{Actor} es un objeto, \texttt{update()} es un bucle por siempre, y \texttt{on\_key\_down} es un bloque de sombrero «al presionar tecla». Thonny tiene un menú \textbf{Ejecutar} (\textit{Run}) $\rightarrow$ \textbf{Modo Pygame Zero} (\textit{Pygame Zero mode}) y lo instala desde \textbf{Herramientas} (\textit{Tools}) $\rightarrow$ \textbf{Administrar paquetes} (\textit{Manage packages}). Una advertencia: se ha reportado un problema (Thonny issue \#3158) en el que el modo Pygame Zero falla con versiones nuevas de Python. Pruébalo en tu propia computadora antes de prometerle un juego a tu hijo.

\textbf{Sin computadora donde instalar (iPad, Chromebook):} el \textbf{editor de Python de micro:bit} (\url{python.microbit.org}) funciona en el navegador, no necesita cuenta y tiene un simulador, así que no necesitas una placa micro:bit. Tiene botones, sonido y bucles por siempre. Dibuja en una cuadrícula de 5 por 5 luces, no con una tortuga. \textbf{MakeCode Arcade} (\url{arcade.makecode.com}) crea juegos retro y deja que tu hijo cambie entre bloques y una versión de Python, lo que es un puente suave.

\textbf{Herramientas de las que quizá leas pero que conviene evitar:} Trinket cerró en agosto de 2026, y el editor Mu está archivado y ya no se actualiza, según sus propios avisos.

\subsection*{Mostrarle a tu hijo su propio proyecto de Scratch como Python}

Este es un momento maravilloso: «Mira, este es \emph{tu} proyecto, en Python».

Un proyecto de Scratch guardado con \textbf{Archivo}, \textbf{Guardar en tu computador} es un archivo \texttt{.sb3}. Un archivo \texttt{.sb3} es en realidad un archivo zip. Dentro están las imágenes de los disfraces, los sonidos y un archivo llamado \texttt{project.json}, que guarda todos los programas como texto. Cada bloque se guarda con un nombre (en inglés) como \texttt{motion\_movesteps} (mover pasos) o \texttt{control\_repeat} (repetir). Scratch no tiene un botón de «exportar como Python», pero puedes asomarte adentro con el propio Python. Sube el \texttt{.sb3} a Colab (el ícono de la carpeta a la izquierda y luego el botón de subir), cambia el nombre del archivo en la línea 3 y ejecuta:

\begin{code}
import zipfile, json

with zipfile.ZipFile("mi-proyecto.sb3") as z:
    proyecto = json.loads(z.read("project.json"))

for objeto in proyecto["targets"]:
    print("Objeto:", objeto["name"])
    for bloque in objeto["blocks"].values():
        if isinstance(bloque, dict):
            print("   ", bloque["opcode"])
\end{code}

Para un proyecto que dibuja un cuadrado, imprime \texttt{event\_whenflagclicked}, \texttt{control\_repeat}, \texttt{motion\_movesteps} y \texttt{motion\_turnright}. Puede que los bloques no salgan en el orden del programa, pero tu hijo los reconocerá. (El escenario también aparece en la lista, con el nombre \texttt{Stage}.)

\textbf{¿Qué proyectos se traducen bien?} Los proyectos sencillos de dibujo (bandera verde, repetir, mover, girar, lápiz, variables, preguntar, Mis Bloques) se traducen casi línea por línea al Python de este libro, y el Diccionario de Bloques lo convierte en una buena actividad para hacer juntos a mano. El escenario de Scratch y la tortuga comparten las mismas coordenadas, así que el dibujo sale igual. Los juegos basados en disfraces con imágenes, objetos que se tocan, clones o sonido no se traducen de forma sencilla. Para esos, hablen mejor de la estructura: un bucle es un \texttt{for}, un bloque de Mis Bloques es un \texttt{def}. Reescribir uno de sus propios proyectos con el Diccionario de Bloques abierto es el verdadero logro de este libro.

\subsection*{Referencias}

{\footnotesize
\begin{itemize}[leftmargin=*, itemsep=1pt]
  \item Brennan, K., \& Resnick, M. (2012). New frameworks for studying and assessing the development of computational thinking. Paper presented at the Annual Meeting of the American Educational Research Association, Vancouver, Canada. \url{https://scratched.gse.harvard.edu/ct/files/AERA2012.pdf}
  \item Dann, W., Cosgrove, D., Slater, D., Culyba, D., \& Cooper, S. (2012). Mediated transfer: Alice 3 to Java. \emph{Proceedings of the 43rd ACM Technical Symposium on Computer Science Education (SIGCSE '12)}, 141--146.
  \item Lee, I., Martin, F., Denner, J., Coulter, B., Allan, W., Erickson, J., Malyn-Smith, J., \& Werner, L. (2011). Computational thinking for youth in practice. \emph{ACM Inroads, 2}(1), 32--37. \url{https://doi.org/10.1145/1929887.1929902}
  \item Sentance, S., Waite, J., \& Kallia, M. (2019). Teaching computer programming with PRIMM: A sociocultural perspective. \emph{Computer Science Education, 29}(2--3), 136--176. \url{https://doi.org/10.1080/08993408.2019.1608781}
  \item Weintrop, D., \& Wilensky, U. (2017). Comparing block-based and text-based programming in high school computer science classrooms. \emph{ACM Transactions on Computing Education, 18}(1), Article 3. \url{https://doi.org/10.1145/3089799}
  \item Weintrop, D., \& Wilensky, U. (2019). Transitioning from introductory block-based and text-based environments to professional programming languages in high school computer science classrooms. \emph{Computers \& Education, 142}, 103646. \url{https://doi.org/10.1016/j.compedu.2019.103646}
\end{itemize}
}

\vspace{10pt}
\begin{tcolorbox}[colback=white, colframe=muted!50, boxrule=0.5pt, arc=5pt]
{\small\color{muted}Este es un libro de ejercicios independiente. No está hecho, respaldado ni relacionado con la Scratch Foundation, el MIT ni Google. Scratch es un proyecto de la Scratch Foundation. Google Colab es un producto de Google. ColabTurtlePlus, Thonny, Pygame Zero, el editor de Python de micro:bit y MakeCode Arcade pertenecen a sus respectivos creadores.}
\end{tcolorbox}


% =====================================================================
\banner{chalbg}{¡NO MIRES HASTA HABERLO INTENTADO!}{Respuestas}
\subsection*{Parte 1: Apuntar en una dirección (la tabla)}

\begin{center}
\renewcommand{\arraystretch}{1.3}
\begin{tabular}{P{2.8cm} P{6.0cm} P{3.0cm} P{2.8cm}}
\toprule
\textbf{Dirección en Scratch} & \textbf{Código de Python} & \textbf{Dirección en Python} & \textbf{¿Hacia dónde?} \\
\midrule
90 & \py|setheading(90 - 90)| & 0 & derecha \\
0 & \py|setheading(90 - 0)| & 90 & arriba \\
-90 & \py|setheading(90 - -90)| & 180 & izquierda \\
45 & \py|setheading(90 - 45)| & 45 & arriba y a la derecha \\
135 & \py|setheading(90 - 135)| & $-45$ & abajo y a la derecha \\
\bottomrule
\end{tabular}
\end{center}

Restar un número negativo es lo mismo que sumar, así que $90 - (-90) = 180$. Python también entiende $-45$: es la misma dirección que 315.

\subsection*{Parte 1: Tradúcelo tú}

\begin{code}
from ColabTurtlePlus.Turtle import *
clearscreen()
initializeTurtle((480, 360))
penup()
goto(-100, 0)
pendown()
pencolor("red")
pensize(4)
setheading(90 - 0)
forward(80)
right(90)
forward(80)
stamp()
\end{code}

La tortuga dibuja una línea roja hacia arriba desde $(-100, 0)$, gira a la derecha, dibuja una línea de lado a lado y se sella en $(-20, 80)$. Escribir \py|setheading(90)| en lugar de \py|setheading(90 - 0)| también funciona. Si agregaste \py|shape("turtle")|, ¡muy bien!: tu sello se ve como una tortuga.

\subsection*{Parte 1: Acertijo de dos líneas}

\begin{code}
setheading(90 - 180)
forward(100)
\end{code}

\py|setheading(-90)| también funciona, y también \py|right(90)| en lugar de \py|setheading|, porque la tortuga empieza mirando a la derecha. Un cuarto de vuelta en el sentido de las agujas del reloj la deja mirando hacia abajo.

\subsection*{Parte 2, Dibuja figuras por el camino corto: la tabla de figuras}
Pentágono: $360 \div 5 = 72$. \quad Hexágono: $360 \div 6 = 60$. \quad Octágono: $360 \div 8 = 45$. \quad Decágono: $360 \div 10 = 36$. \quad Casi un círculo: $360 \div 36 = 10$.

La figura de 36 lados es la que más se parece a un círculo. Con \py|pencolor("red")| dentro del bucle, todo el triángulo sale rojo: la línea 4 pone el lápiz azul, pero la línea del rojo se ejecuta antes de cada \py|forward|, así que el azul nunca se usa. Con \py|print("¡Tarán!")| pegado al borde se escribe una vez; con cuatro espacios queda dentro del bucle y se escribe 3 veces.

\subsection*{Parte 2, Estrellas: la tabla de estrellas}
\begin{center}
\renewcommand{\arraystretch}{1.3}
\begin{tabular}{P{2.4cm} P{3.4cm} P{2.4cm}}
\toprule
\textbf{Puntas} & \textbf{$180 \div$ puntas} & \textbf{Giro} \\
\midrule
9 & 20 & 160 \\
15 & 12 & 168 \\
45 & 4 & 176 \\
\bottomrule
\end{tabular}
\end{center}

\subsection*{Parte 2, Espiral cuadrada: la línea que falta}
En la primera vuelta, \texttt{i} vale 0, así que \py|forward(i * 5)| es \py|forward(0)|. La tortuga no se mueve, así que la primera «línea» no mide nada. Con \py|range(60)| y \py|forward(i * 3)|, la línea más larga mide $59 \times 3 = 177$ pasos, cerca de los $39 \times 5 = 195$ pasos de la primera espiral, así que el tamaño queda más o menos igual mientras las líneas quedan más juntas.

\subsection*{Parte 2, Flor de cuadrados}
36 cuadrados $\times$ 10 grados $= 360$, y 12 cuadrados $\times$ 30 grados $= 360$: los giros entre los cuadrados tienen que sumar un círculo completo. \py|hideturtle()| no lleva \textbf{ningún} espacio, así que se ejecuta una vez, después de toda la flor.

\subsection*{Parte 2, Tradúcelo tú}
Tu código puede verse un poco distinto y aun así estar bien. Si la tortuga dibuja el mismo dibujo, ¡lo lograste!
\begin{codenonum}
from ColabTurtlePlus.Turtle import *
clearscreen()
initializeTurtle((480, 360))
pensize(4)
pencolor("purple")
for i in range(8):
    for j in range(6):
        forward(40)
        right(60)
    right(45)
penup()
print("¡Listo!")
\end{codenonum}
\py|forward(40)| y \py|right(60)| llevan 8 espacios (dentro de los dos bucles). \py|right(45)| lleva 4 espacios. \py|penup()| y \py|print("¡Listo!")| no llevan ninguno, porque van después de los bucles. Te salen 8 hexágonos en un anillo.

\subsection*{Parte 2, Acertijo de dos líneas}
\begin{codenonum}
for i in range(5):
    print(i + 1)
\end{codenonum}
\texttt{i} va 0, 1, 2, 3, 4, así que \py|i + 1| va 1, 2, 3, 4, 5.

\subsection*{Parte 3: Examen rápido, ¿caja o fracaso?}
Python acepta \texttt{vidas}, \texttt{jugador\_1} y \texttt{mejorTiempo}. \texttt{mejor marca} tiene un espacio, \texttt{1er\_lugar} empieza con un dígito, e \texttt{import} es una palabra que Python ya usa.

\subsection*{Parte 3, Misión: Contador de puntos}
\py|puntos += 10| hace exactamente lo mismo que \py|puntos = puntos + 10|, así que nada cambia. Para contar hacia atrás, empieza con \py|puntos = 50| y usa \py|puntos -= 10|: la última línea dice \texttt{Puntos finales: 0}. Con \py|vidas -= 1| dentro del bucle, las vidas van 2, 1, 0, luego $-1$ y $-2$. Python no se detiene en 0 a menos que tú se lo digas (aprenderás cómo con \py|if| en la Parte 5).

\textbf{Superreto:}
\begin{codenonum}
import time
puntos = 0
premio = 10
print("Puntos:", puntos)
for i in range(5):
    time.sleep(1)
    puntos = puntos + premio
    premio = premio + 10
    print("Puntos:", puntos)
print("¡Fin del juego! Puntos finales:", puntos)
\end{codenonum}
Los puntos van 10, 30, 60, 100, 150. Los puntos finales son \texttt{150}.

\subsection*{Parte 3, Misión: Espiral que crece}
Con \texttt{+ 8} la última línea mide 317 pasos, y la espiral apenas cabe en el escenario. Para el lápiz más grueso, pon \py|grueso = 1| junto a \py|medida = 5|, y agrega \py|pensize(grueso)| y \py|grueso = grueso + 0.2| dentro del bucle.

\textbf{Superreto} (después de las tres líneas de inicio):
\begin{codenonum}
speed(13)
pencolor("purple")
penup()
goto(-100, 100)
pendown()
medida = 200
for i in range(40):
    forward(medida)
    right(90)
    medida = medida - 5
\end{codenonum}

\subsection*{Parte 3, Misión: Arcoíris con una lista de colores}
\py|colores.append("pink")| añade el rosa al final de la lista, así que el bucle da 7 vueltas y el rosa es la franja de hasta abajo. Las franjas roja, blanca y azul, de arriba abajo, forman la bandera de los Países Bajos. Prueba \py|["white", "blue", "red"]| para Rusia.

\textbf{Superreto:} las líneas 8 a 16 quedan así
\begin{codenonum}
y = 60
largo = 300
for color in colores:
    pencolor(color)
    penup()
    goto(-150, y)
    pendown()
    forward(largo)
    y = y - 15
    largo = largo - 30
hideturtle()
\end{codenonum}

\subsection*{Parte 3, Misión: Inventor de historias al azar}
Para los lugares, agrega \py|lugares = ["en la tina", "en el camión"]| antes del bucle, \py|lugar = random.choice(lugares)| dentro de él, y cambia el print por \py|print(heroe, accion, cosa, lugar)|. Después del bucle, \py|len(cosas)| es \texttt{4}, y \texttt{"mi tarea"} es el último elemento, porque \py|append| siempre añade al final:
\begin{outputbox}
4\\
{[}'un plátano', 'la luna', 'un calcetín', 'mi tarea']
\end{outputbox}

\subsection*{Parte 3: Tradúcelo tú}
\begin{codenonum}
mascotas = []
mascotas.append("gato")
mascotas.append("perro")
mascotas.append("pez")
print(len(mascotas))
print(mascotas[1])
\end{codenonum}
Imprime \texttt{3} y luego \texttt{perro}. El elemento 2 de Scratch es el \texttt{[1]} de Python. \py|mascotas = []| hace el trabajo de \chip[lists]{Eliminar todos de mascotas}, porque empieza una lista vacía nuevecita.

\subsection*{Parte 3: Acertijo de dos líneas}
\begin{codenonum}
for palabra in palabras:
    cuenta += 1
\end{codenonum}
El bucle da una vuelta por cada elemento, y en cada vuelta le suma 1 a \texttt{cuenta}. ¡Así es exactamente como cuenta \py|len|!

\subsection*{Parte 4: ¡Adivina qué sale!}
\begin{center}
\renewcommand{\arraystretch}{1.3}
\begin{tabular}{l l P{7.5cm}}
\toprule
\textbf{Código} & \textbf{Muestra} & \textbf{Por qué} \\
\midrule
\py|print(9 / 3)| & \texttt{3.0} & Dividir siempre da un decimal. \\
\py|print("4" + "4")| & \texttt{44} & Son strings, así que \texttt{+} los une. \\
\py|print("Perro" == "perro")| & \texttt{False} & En Python, las mayúsculas cuentan. \\
\py|print("pizza"[1])| & \texttt{i} & Python cuenta desde 0, así que \texttt{[1]} es la segunda letra. \\
\py|print(round(4.5))| & \texttt{4} & Un .5 va al número par más cercano. \\
\py|print(17 % 5)| & \texttt{2} & $17 = 3 \times 5 + 2$, así que sobran 2. \\
\bottomrule
\end{tabular}
\end{center}

\subsection*{Parte 4, Misión 1: Lanzador de dados}
Para 100 tiradas, cambia la línea 3 por \py|for i in range(100):| y la línea 8 por \py|str(total / 100)|. Con dos dados, el total más grande es $6 + 6 = 12$.

\subsection*{Parte 4, Misión 2: Impresora de tablas}
Para la tabla del 9, cambia la línea 1 por \py|tabla = 9|. Para llegar hasta 20, cambia la línea 3 por \py|range(20)|. Para los números al cuadrado, la línea 4 queda \py|print(f"{n} x {n} = {n * n}")|. El \texttt{True} o \texttt{False} del final te dice si el resultado es par.

\subsection*{Parte 4, Misión 3: Arte de puntos al azar}
Con 7 colores, la línea 15 necesita \py|random.randint(0, 6)|. Para la mitad de arriba, la línea 12 queda \py|y = random.randint(0, 180)|.

\subsection*{Parte 4, Misión 4: Máquina de nombres}
Dice \py|i + 1| porque \py|i| empieza en 0, pero las personas cuentan las letras desde 1.

\subsection*{Parte 4: Tradúcelo tú}
Tu código puede verse un poco diferente y aun así estar bien. Si imprime tres totales, ¡lo lograste!
\begin{codenonum}
import random
for i in range(3):
    a = random.randint(1, 10)
    b = random.randint(1, 10)
    print("Total: " + str(a + b))
\end{codenonum}
Una ejecución real imprimió \texttt{Total: 9}, \texttt{Total: 19} y \texttt{Total: 8}. Un f-string también funciona: \py|print(f"Total: {a + b}")|.

\subsection*{Parte 4: Acertijo de dos líneas}
\begin{codenonum}
palabra = "Tortuga"
print(palabra[0] + str(len(palabra)))
\end{codenonum}
\py|palabra[0]| es la primera letra, \texttt{T}. \py|len(palabra)| es 7, y \py|str( )| lo convierte en string para poder unirlo.

\subsection*{Parte 5, Misión «Di hola»: retos}
La pregunta de la edad. \py|str( )| vuelve a convertir el número en texto para poder unirlo:
\begin{codenonum}
nombre = input("¿Cómo te llamas? ")
edad = int(input("¿Cuántos años tienes? "))
print("En 10 años tendrás " + str(edad + 10))
\end{codenonum}
\textbf{Súper reto:} sin \py|int( )|, \py|edad + 10| intenta sumar un número a un texto. La última línea dice \texttt{TypeError: can only concatenate str (not "int") to str}. Un \texttt{TypeError} quiere decir que se mezclaron textos y números.

\subsection*{Parte 5, Misión «Adivina mi número»: solo 5 intentos}
\begin{codenonum}
import random
secreto = random.randint(1, 10)
intento = 0
veces = 0
while intento != secreto and veces < 5:
    intento = int(input("Adivina del 1 al 10: "))
    veces += 1
    if intento < secreto:
        print("Muy bajo")
    elif intento > secreto:
        print("Muy alto")
if intento == secreto:
    print("¡Lo adivinaste en " + str(veces) + " intentos!")
else:
    print("¡Se acabaron los intentos! Era " + str(secreto))
\end{codenonum}
Si escribes \texttt{gato} sale \texttt{ValueError: invalid literal for int() with base 10: 'gato'}, porque \py|int( )| no puede convertir la palabra \texttt{gato} en un número.

\subsection*{Parte 5, Misión «Concurso con puntos»: el mensaje final}
Pon esto al final de todo, después de escribir los puntos:
\begin{codenonum}
if puntos == 3:
    print("¡Campeón del concurso!")
elif puntos >= 1:
    print("¡Bien hecho!")
else:
    print("¡Más suerte la próxima vez!")
\end{codenonum}
\textbf{Súper reto:} pon \py|import time| e \py|inicio = time.time()| al principio, y \py|print(round(time.time() - inicio))| al final.

\subsection*{Parte 5, Misión «Prueba de reflejos»: el mejor de tres}
\begin{codenonum}
import time
import random
mejor = 100
for i in range(3):
    input("Presiona Enter para empezar ")
    print("Prepárate...")
    time.sleep(random.randint(2, 5))
    inicio = time.time()
    input("¡YA! ¡Presiona Enter! ")
    tiempo = time.time() - inicio
    if tiempo < mejor:
        mejor = tiempo
print("Tu mejor tiempo: " + str(round(mejor, 2)) + " segundos")
\end{codenonum}
\py|mejor| empieza en 100 segundos, así que tu primer intento seguro lo mejora.

\subsection*{Parte 5: Tradúcelo tú}
Tu código puede verse un poco diferente y aun así estar bien. Si hace lo mismo, ¡lo lograste!
\begin{codenonum}
respuesta = input("¿Cuál es la contraseña? ")
while respuesta != "banana":
    print("¡Mal! Intenta otra vez.")
    respuesta = input("¿Cuál es la contraseña? ")
print("¡Adelante!")
\end{codenonum}
Usar \texttt{while not respuesta == "banana":} en la línea 2 también está bien.

\subsection*{Parte 5: Acertijo de dos líneas}
\begin{codenonum}
edad = int(input("¿Cuántos años tienes? "))
print("El próximo año tendrás " + str(edad + 1))
\end{codenonum}

\subsection*{Parte 6, Misión: Haz tu propio bloque}
\textbf{Triángulo:}
\begin{codenonum}
def triangulo(lado):
    for i in range(3):
        forward(lado)
        right(120)
\end{codenonum}
\textbf{Entrada de color:}
\begin{codenonum}
def cuadrado(lado, tinta):
    pencolor(tinta)
    for i in range(4):
        forward(lado)
        right(90)
\end{codenonum}
\textbf{Sin paréntesis:} una línea que solo dice \py|cuadrado| no hace nada, y no sale ningún error. Sin los paréntesis solo dices el nombre del bloque. Tienes que llamarlo, por ejemplo \py|cuadrado(60)|.

\subsection*{Parte 6, Misión: Máquina de figuras}
\py|flor(8, 4, 60)| dibuja 8 pétalos, y cada pétalo es un cuadrado. Con \py|animationOff()| la tortuga salta en vez de caminar, así que el dibujo aparece más rápido. Entrada de color:
\begin{codenonum}
def flor(petalos, lados, medida, tinta):
    pencolor(tinta)
    for i in range(petalos):
        poligono(lados, medida)
        right(360 / petalos)
\end{codenonum}
Llámalo con cuatro entradas, por ejemplo \py|flor(12, 6, 40, "red")|.

\subsection*{Parte 6, Misión: Máquina de respuestas}
\begin{codenonum}
def mitad(n):
    return n / 2

def perimetro(ancho, alto):
    return 2 * ancho + 2 * alto

def es_par(n):
    return n % 2 == 0
\end{codenonum}
\py|print(mitad(9))| muestra \texttt{4.5}. \py|es_par(7)| es \py|False|, así que la parte del \py|if| se salta.

Con \py|print(n * 2)| en la línea 2, \py|doble| solo \textbf{muestra} el número y devuelve \texttt{None}. Así que la línea 7 muestra \texttt{8} y luego \texttt{None}. Después, la línea 11 intenta hacer \py|area(None, 10)|, y Python se detiene con un \texttt{TypeError}, porque no puede multiplicar «nada» por 10.

\subsection*{Parte 6, Misión: Ejército de objetos}
\begin{codenonum}
def girar_todos(angulo):
    for s in ejercito:
        s.left(angulo)
\end{codenonum}
Llama \py|girar_todos(90)| antes de marchar, y el ejército marcha hacia arriba. Para la carrera, pon \py|import random| hasta arriba y usa \py|s.forward(random.randint(10, 80))| en \texttt{marchar}. \py|ejercito[0]| es el primer objeto que se creó: el rojo de arriba (las listas cuentan desde 0).

\subsection*{Parte 6, Tradúcelo tú}
\begin{codenonum}
from ColabTurtlePlus.Turtle import *
clearscreen()
initializeTurtle((480, 360))

def estrella(medida):
    for i in range(5):
        forward(medida)
        right(144)

pencolor("gold")
estrella(100)
estrella(50)
\end{codenonum}
Deberías ver una estrella dorada grande con una pequeña adentro, que comparten una punta.

\subsection*{Parte 6, Acertijo de dos líneas}
\begin{codenonum}
def triple(n):
    return n * 3
\end{codenonum}
\py|triple(2)| es 6, y \py|triple(6)| es 18.

\subsection*{Parte 7, Desfile de disfraces: cinco disfraces}
Si agregas \py|"arrow"| a la lista de disfraces, la lista de colores también necesita un quinto color, por ejemplo \py|colores = ["green", "red", "blue", "orange", "purple"]|. Si no, cuando \py|n| vale 4, \py|colores[4]| no existe y Python se detiene con un \texttt{IndexError}. La línea 17 no necesita cambios: \py|len(disfraces)| ahora vale 5 solita.

\subsection*{Parte 7, Hacer música: ¡Feliz cumpleaños!}
\begin{codenonum}
Sol = 392
La = 440
Si = 494
Do2 = 523
melodia = [Sol, Sol, La, Sol, Do2, Si]
tiempos = [0.75, 0.25, 1, 1, 1, 2]
\end{codenonum}
Pon esto en lugar de las líneas 2 a 10 y deja el bucle como está. Si suena demasiado rápido, cambia \py|0.4| en la línea 12 por \py|0.5|.

\subsection*{Parte 7, Carrera de tortugas: tres corredores}
\begin{codenonum}
verde = crear_corredor("green", 0)
\end{codenonum}
va después de la línea 17. El bucle y lo que tiene dentro quedan así:
\begin{codenonum}
while (rojo.xcor() < 200 and azul.xcor() < 200
       and verde.xcor() < 200):
    rojo.forward(random.randint(1, 20))
    azul.forward(random.randint(1, 20))
    verde.forward(random.randint(1, 20))
\end{codenonum}
Los paréntesis dejan que la condición ocupe dos líneas.

\subsection*{Parte 7, Tradúcelo tú}
\begin{codenonum}
from ColabTurtlePlus.Turtle import *
clearscreen()
initializeTurtle((480, 360))
import time

shape("turtle")
shapesize(3)
for i in range(3):
    hideturtle()
    time.sleep(0.5)
    showturtle()
    time.sleep(0.5)
print("¡Bu!")
time.sleep(2)
\end{codenonum}
El tamaño 300\,\% es \py|shapesize(3)|. «Decir durante 2 segundos» es un \py|print| y luego una espera de 2 segundos. Si usaste una tortuga globo para escribir \texttt{¡Bu!} en el escenario y la borraste con \py|clear()| después de 2 segundos, ¡eso se parece todavía más a Scratch!

\subsection*{Parte 7, Acertijo de dos líneas}
\begin{codenonum}
def festejo():
    print("¡Hurra!")
\end{codenonum}
La tercera línea, la que envía el mensaje, es \py|festejo()|. Cuando ejecutas las tres, ves \texttt{¡Hurra!}.

% =====================================================================
% 08-projects-answers.tex : pistas para los retos «Sube de nivel» de Grandes proyectos
% =====================================================================

\subsection*{Grandes proyectos: pistas para «Sube de nivel»}
Para subir de nivel no hay una sola respuesta correcta. Si tu programa hace lo que tú querías, ¡está bien! Estas pistas te ayudan si te atoras.

\textbf{Anima tu nombre.} Cambia \py|range(8)| por \py|range(20)| en la línea 20 y el número de \py|time.sleep()| en la línea 29. Para las letras que crecen, cambia el tipo de letra de la línea 27 por \py|font=("Arial", 20 + cuadro * 5, "bold")|. Para el arcoíris, reemplaza \py|random.choice(colores)| en la línea 24 por \py|colores[cuadro % 7]| (hay 7 colores en la lista).

\textbf{Flor de lápiz.} \py|petalo(150)| hace pétalos más grandes que pueden salirse del escenario. Para una flor dentro de otra flor, copia las líneas 18--21 y usa \py|petalo(50)| en la copia. Para el jardín, esto funciona (reemplaza las líneas 17--31):
\begin{codenonum}
def flor(x, y):
    penup()
    goto(x, y)
    pendown()
    for p in range(petalos):
        color(colores[p % 6])
        petalo(50)
        left(360 / petalos)

petalos = int(input("¿Cuántos pétalos? (3 a 20) "))
flor(-150, 0)
flor(0, 0)
flor(150, 0)
hideturtle()
\end{codenonum}

\textbf{Concurso de tablas.} Para 10 preguntas cambia \py|range(5)| en la línea 16 y el \py|5| de los mensajes (líneas 22 y 24). Para un concurso de la tabla del siete, la línea 17 se vuelve \py|a = 7|. Para los sonidos: \py|pitido(523)| va después de \py|print("¡Correcto!")| y \py|pitido(196)| después de la línea del \texttt{¡Uy!}

\textbf{Adivina el número de lujo.} Partiendo 1000 a la mitad: 500, 250, 125, 63, 32, 16, 8, 4, 2, 1. Son 10 intentos. Agrega \py|elif nivel == "4":| con \py|maximo = 1000| e \py|intentos = 10| antes del \py|else:|. Para contar los juegos ganados, pon \py|ganados = 0| antes del bucle grande, agrega \py|ganados = ganados + 1| dentro de \py|if ganaste:| e imprímelo después del bucle. Para «¡Caliente!»: \py|if numero != secreto and numero - secreto < 6 and secreto - numero < 6:|. Para «demasiado grande»: revisa \py|if numero > maximo:| justo después de la línea del \py|input|, imprime un mensaje y súmale 1 otra vez a \py|intentos|.

\textbf{Elige tu propia historia.} Una tercera opción necesita otro ayudante, o puedes hacer dos preguntas seguidas: primero \py|elegir("¿Qué haces?", "remar", "nadar")|, y dentro de la parte de \texttt{nadar}, \py|elegir("¿Nadar o cruzar el puente?", "nadar", "puente")|, que llama a \py|puente()|. Pon \py|def puente():| cerca del principio con las otras escenas, para que Python ya lo haya leído antes de que empiece la historia.

\textbf{Mascota virtual.} Agrega \py|energia = 7| al principio, un \py|elif accion == "dormir":| nuevo que haga \py|energia = 10|, y \py|energia = energia - 2| en la parte de \texttt{jugar}. Un mensaje de humor: \py|if alegria >= 8:| y luego \py|print("(ronronea)")|. Para ganar: \py|if dia == 15:| imprime tu diploma y luego \py|break|.

\textbf{Juego de persecución.} Pon \py|paso_gato = 25| junto a \py|puntos = 0|, usa \py|gato.forward(paso_gato)| y agrega \py|paso_gato = paso_gato + 5| dentro del \py|if| de la estrella. Un gato rápido puede saltar por encima de ti, así que cambia también la prueba de atrapar por \py|gato.distance(pos()) < paso_gato + 5|. Un segundo gato: copia las líneas 11--16 con el nombre \py|gato2| y \py|gato2.goto(-200, -140)|, y luego copia las líneas 49--53 para \py|gato2|.

\textbf{Juego de atrapar.} Para empezar más a lo ancho, la línea 23 se vuelve \py|x = random.randint(-160, 160)|; para que la manzana no se salga del escenario, usa \py|viento = random.choice([-10, 0, 10])|. Para el jugador computadora, reemplaza la línea 30 por
\begin{codenonum}
    lugar = manzana.xcor() + 4 * viento
\end{codenonum}
A la manzana le faltan 4 pasos por caer, y se mueve \py|viento| de lado en cada paso. ¡La computadora atrapa todas las manzanas! Para las marcas, pon \py|manzana.stamp()| después de la línea 35.

\subsection*{El Zoológico de Bichos: caza de bichos}
Cada error de abajo es la última línea real de Python 3.13 (Python siempre escribe sus errores en inglés). Tu arreglo puede verse un poco distinto y aun así estar bien: si ahora el programa hace lo que debe, ¡atrapaste al bicho!

\small
\renewcommand{\arraystretch}{1.3}
\begin{longtable}{P{1.1cm} P{6.4cm} P{8.1cm}}
\toprule
\textbf{Bicho} & \textbf{La pista} & \textbf{El arreglo} \\
\midrule
1 & \texttt{IndentationError: expected an indented block after 'while' statement on line 2} (la Torre Tambaleante) & La línea 3 necesita cuatro espacios delante, para que las líneas 3 y 4 queden las dos dentro del bucle. \\
2 & \texttt{NameError: name 'Print' is not defined} (el Copión Mayúsculo) & Usa una p minúscula: \py|print("Tus puntos son", puntos)| \\
3 & \texttt{TypeError: can only concatenate str (not "int") to str} (el Duende del Pegamento) & \py|input()| siempre da texto. Conviértelo en número: \py|edad = int(input("¿Cuántos años tienes? "))| \\
4 & \texttt{SyntaxError: invalid syntax. Maybe you meant '==' or ':=' instead of '='?} (la Anguila del Igual) & Para hacer una pregunta, usa dos iguales: \py|if clave == "gato123":| \\
5 & No hay error en la línea 2, pero imprime \texttt{Primera: plátano} (el Héroe Cero). Luego \texttt{IndexError: list index out of range} (el Búho del Uno de Más). & Python cuenta desde 0: usa \py|frutas[0]| para la primera y \py|frutas[2]| (o \py|frutas[-1]|) para la última. \\
6 & \texttt{ZeroDivisionError: division by zero} (el Zombi del Cero) & Comprueba primero:\newline \py|if amigos == 0:|\newline \py|    print("¡No hay con quién repartir!")|\newline \py|else:|\newline \py|    print("Cada uno recibe", dulces / amigos)| \\
7 & No hay error, pero no se imprime nada (la Receta Silenciosa). & Pon los paréntesis para cocinar la receta: \py|festejar()| \\
8 & No hay error, pero imprime \texttt{5} para siempre (el Fantasma Infinito). Presiona el botón de detener; entonces ves \texttt{KeyboardInterrupt}. & \texttt{n} nunca cambia. Agrega \py|n = n - 1| dentro del bucle, después de \py|print(n)|. \\
9 & No hay error, pero la tortuga dibuja una sola línea larga de 200 pasos en vez de un cuadrado (el primo del Hurra Perdido). & \py|right(90)| no tiene espacios, así que está fuera del bucle y pasa una sola vez. Ponle cuatro espacios delante. \\
10 & \texttt{SyntaxError: unterminated string literal (detected at line 3)} (el Ladrón de Comillas) & Pon las comillas de cierre: \py|print("¡Mucho gusto!")|. Como es un bicho de gramática, no se ejecuta nada, ni siquiera la pregunta de la línea 1. \\
\bottomrule
\end{longtable}
\normalsize


\end{document}
